% !TeX program = xelatex
% ============================================================
%  S. J. Blundell, Magnetism in Condensed Matter (OUP 2001)
%  中译重排版：正文译为中文，公式/编号/交叉引用照原书，插图用原书扫描
% ============================================================
\documentclass[11pt,a4paper,twoside,openright]{ctexbook}

\usepackage[textwidth=13.2cm,textheight=23.5cm,marginparwidth=3.1cm,marginparsep=0.4cm]{geometry}
\usepackage{amsmath,amssymb,bm}
\usepackage{graphicx}
\graphicspath{{figs/}}
\setmainfont{Times New Roman}
\usepackage{fancyhdr}
\usepackage{setspace}
\usepackage{enumitem}
\usepackage{booktabs}
\usepackage{array}
\setlength{\tabcolsep}{4.5pt}

\newcommand{\ve}[1]{\bm{#1}}
\newcommand{\dd}{\,\mathrm{d}}

\ctexset{
  chapter = {
    name = {第,章},
    number = \arabic{chapter},
    format = {\centering\zihao{2}\heiti},
    beforeskip = 1.5em,
    afterskip  = 1.8em,
  },
  section = {
    format = {\zihao{4}\heiti},
    beforeskip = 1.6em,
    afterskip  = 0.8em,
  },
  subsection = {
    format = {\zihao{-4}\heiti},
    beforeskip = 1.2em,
    afterskip  = 0.6em,
  },
}

% 页眉页脚
\pagestyle{fancy}
\fancyhf{}
\fancyhead[LE]{\small\kaishu 凝聚态中的磁学}
\fancyhead[RO]{\small\kaishu\nouppercase{\leftmark}}
\fancyfoot[C]{\small\thepage}
\renewcommand{\headrulewidth}{0.4pt}
\renewcommand{\chaptermark}[1]{\markboth{第\arabic{chapter}章\quad #1}{}}

\setstretch{1.18}

% 防止段首/段末孤行、显示公式前断页
\clubpenalty=10000
\widowpenalty=10000
\predisplaypenalty=10000

% 重排 \cleardoublepage：空白页不带页眉页脚
\makeatletter
\def\cleardoublepage{\clearpage\if@twoside\ifodd\c@page\else
\hbox{}\newpage\if@twocolumn\hbox{}\newpage\fi\fi\fi}
\makeatother

% 例题（编号照原书手写，不用计数器）
\newenvironment{example}[1]{\par\medskip\noindent\rule{\textwidth}{0.8pt}\par\nopagebreak\noindent{\heiti 例 #1}\quad}{\par\nopagebreak\rule{\textwidth}{0.8pt}\par\medskip}

% 原书编号边注 → 同号脚注（编号不依赖出现顺序，防漏注错位）
\newcommand{\fn}[2]{\begingroup\setcounter{footnote}{\numexpr#1-1\relax}\footnote{#2}\endgroup}

% 无编号边注与人物小传 → 边栏（模拟原书页边注）
\newcommand{\mnote}[1]{\marginpar{\raggedright\footnotesize\kaishu #1}}
\newcommand{\bio}[1]{\marginpar{\raggedright\scriptsize\itshape #1\par}}

% 习题条目（编号照原书）
\newcommand{\exer}[1]{\par\medskip\noindent{\heiti (#1)}\enspace}

\begin{document}
% ---------------- 扉页 ----------------
\begin{titlepage}
\thispagestyle{empty}
\centering
\vspace*{3.2cm}
{\zihao{0}\heiti 凝聚态中的磁学}\\[1.2em]
{\zihao{3}\kaishu Magnetism in Condensed Matter}\\[3.5em]
{\zihao{4}S. J. 布伦德尔（S. J. Blundell）\quad 著}\\[1.0em]
{\zihao{-4}\kaishu 牛津大学出版社 2001 年版}\\[2.5em]
{\zihao{-4}\kaishu 中译重排版}\\[1em]
{\zihao{5}\kaishu 公式与编号照原书\quad 插图取自原书扫描}
\vfill
\end{titlepage}

% 扉页背面留空：保证正文第 1 页起于奇数页（正面）
\thispagestyle{empty}\mbox{}\clearpage

\frontmatter
\chapter*{译者说明}
\addcontentsline{toc}{chapter}{译者说明}
本书为 Stephen J. Blundell 所著 \emph{Magnetism in Condensed Matter}
（Oxford University Press, 2001）的中文译本，由扫描版整理重排而成，
供学习研究之用。整理与翻译遵循以下约定：
\begin{enumerate}[itemsep=0.2em]
  \item 正文全部译为中文，力求信达；专业名词在首次出现或易生歧义处
        以圆括号附注英文原文。
  \item 全部公式、公式编号（如 (2.15)）、图表编号与正文交叉引用
        均照录原书，未作改动；书中引用的文献题录不予翻译。
  \item 插图取自原书扫描件，图注译为中文并保留英文原图注要点；
        人名不译（首次出现附原文）。
  \item 例题（Example）与各章末习题（Exercises）照原书翻译编号；
        书末附部分习题答案与提示。
  \item 原书按 SI 单位制行文，附录 A 给出与 cgs 制的换算关系。
\end{enumerate}

% ---------------- 原书前言（Preface，牛津 2001） ----------------
\chapter*{前言}
\addcontentsline{toc}{chapter}{前言}

\begin{quote}
\centering \kaishu “……万有也靠他而立。”\quad（《歌罗西书》一章十七节）
\end{quote}

磁学是一门被研究了近三千年的学问。磁石（lodestone）——一种铁矿石——最早引起了希腊
学者与哲学家的注意，而导航用磁性罗盘则是这项研究结出的第一件技术成果。罗盘固然至迟在公元十二世纪已为西欧所知，但直到约 1600 年，才出现堪称近代水平的对罗盘原理的
解释。最近两个世纪以来，进展大大加快，并形成了两项把磁学与其他物理现象联系起来的
重大成果。其一，磁与电密不可分，二者共同构成了光——所谓电磁波。其二，这一联系的
根源在于相对论，因此磁学可以说是一种纯粹的相对论性效应，它源于观察者与导线中、
乃至铁原子中运动电荷的相对运动。然而，时至今日真正引人入胜的，仍是凝聚态体系——
包括铁磁体、自旋玻璃与低维体系——中的磁性。宏观体系所表现出的磁性质，与原子、
分子有着根本的不同，尽管二者由完全相同的基本组分构成。其原因在于，磁是一种集体
现象，需要极大量粒子的彼此协同；在这个意义上，它与超导、超流、乃至固态现象本身
颇为相似。对这类基本问题的探究，与寻找新型永磁材料、传感器与记录应用材料的技术
驱动，始终并行不悖。

本书源于作者在牛津大学为选择了凝聚态物理方向的三年级、四年级本科生开设的一门
课程。当时显而易见，我们需要一本既讲述基础、又提供背景材料与课堂上无法涵盖的
补充专题的教材。本书的目标，是把这一领域作为一个连贯的整体呈现出来，提供有用而
有趣的原始素材，并且读来令人愉悦。本书同时也是《牛津凝聚态物理 master 系列》
（Oxford Master Series in Condensed Matter Physics）的一种；该系列其他各卷分别
讨论电子性质、光学性质、超导电性、结构与软凝聚态物质。

阅读本书的前提是掌握基本的量子力学与电磁学，并熟悉原子物理的若干结果。
这些内容已总结于书末附录，既便于读者查阅，也统一了全书的记号。

凝聚态体系中那些有趣的磁效应有两个关键要素：其一，原子必须具有磁矩；其二，
这些磁矩之间必须以某种方式发生相互作用。这两个主题分别在第二章与第四章中讨论。
第二章回答“原子为什么具有磁矩？”这一问题，并说明磁矩彼此并无相互作用时它们的行为
及研究方法。\mnote{本课程的一些可能安排：（1）短课程（假定已知第 1 章内容）：
第 2 章（略去 2.6--2.8）；第 3 章（略去 3.2）；第 4 章（略去 4.2.5、4.2.6）；
第 5 章（略去 5.4--5.7）；第 6 章（略去 6.4--6.5）；第 7 章（略去 7.1--7.2）。
（2）较长课程：第 1--4 章；第 5 章（5.6、5.7 作背景阅读）；第 6 章；第 7 章
（7.5--7.9 作背景阅读）；第 8 章选讲专题。}第三章描述晶体内的局域环境如何影响这些磁矩，以及研究这种影响的各种
技术。第四章则回答“不同原子上的磁矩如何彼此相互作用？”。有了这两个要素，
磁有序便可能出现，这是第五、六章的主题。第五章描述固体中出现的各种类型的磁有序。
第六章再次讨论有序问题，但从对称性破缺的基本观念出发，描述相变、激发与磁畴。
本章特别强调磁有序与超导电性等其他类型的对称性破缺基态之间的联系。第七章专门
讨论金属的磁性质——在金属中，磁性往往与巡游的传导电子相关。第八章描述当存在
相互竞争的磁相互作用、或体系维度降低时出现的一些精微而复杂效应。这些专题至今
仍是研究热点，尚有许多悬而未决的问题。全书始终贯穿对材料性质与应用的讨论，
以展示这些思想对真实材料的意义——包括铁氧体、永磁体，以及近年来具有巨大技术
价值的各种磁光效应与磁电阻效应背后的物理。这是一本为物理学家而写的书，因此
重点放在支撑整个领域的清晰的物理原理上：量子力学、对称性与电磁学。但它并不只是
一本“理论书”，而是力求把这一学科与实验物理学家当前实际使用的测量手段与实验
技术联系起来，在本科基础物理的原理与当前研究热点之间架起一座桥梁。

第一章至第七章的末尾附有延伸阅读书目与习题。习题难度不一，用以深化正文中
讨论的问题。第八章未设习题（该章所述皆为当前研究课题），但提供了详尽的延伸
阅读书目。

本书写作过程中得到了许多人的帮助，在此谨致谢忱。感谢牛津大学出版社的
Sönke Adlung 及其团队的支持，也感谢本 master 系列的其他作者。牛津大学
Mansfield 学院与牛津大学物理系为我提供了令人振奋的工作环境。我要感谢我的
学生们——他们常常促使我对自己原以为已经理解的东西作更深入的思考。在本书
各个方面的准备过程中，我与 Hideo Aoki、Arzhang Ardavan、Deepto Chakrabarty、
Amalia Coldea、Radu Coldea、Roger Cowley、Steve Cox、Gillian Gehring、
Matthias Gester、John Gregg、Martin Greven、Mohamedally Kurmoo、Steve Lee、
Wilson Poon、Francis Pratt、John Singleton 与 Candadi Sukumar 的讨论使我
受益良多。我尤其要感谢那些阅读了本书各轮草稿的挚友与同行，他们严苛的批评与
富有洞见的问题使最终成品大为改进：Katherine Blundell、Richard Blundell、
Andrew Boothroyd、Geoffrey Brooker、Bill Hayes、Brendon Lovett、
Lesley Parry-Jones 与 Peter Riedi。本书付印后若发现错误，将公布于本书网页：
\begin{center}
\texttt{http://users.ox.ac.uk/\textasciitilde sjb/magnetism/}
\end{center}

最重要的是，我要感谢 Katherine——我亲爱的妻子与灵魂伴侣。她给予我的启迪、
忠告、友谊与爱，远胜他人。本书谨献给她。

\begin{flushright}
牛津\quad 2001 年 5 月\\
S.~J.~B.
\end{flushright}

\clearpage


\tableofcontents

\mainmatter

% 第1章 引言（Introduction）
\chapter{引言}

本书讨论的是磁性在凝聚态（condensed matter）中的种种表现。固体中含有大量磁矩
（magnetic moments），它们能够协同行动，产生与彼此孤立时截然不同的行为；再加上
可以存在的磁相互作用类型多种多样，这就使得真实体系的磁性质丰富得惊人。本书的
思路是放缓节奏、逐块搭建这幅图景。在这一导论性的章节里，我们先回顾初等经典物理
和量子物理中关于磁矩的一些事实；随后各章依次讨论：第二章讨论当大量磁矩被置入
固体但彼此隔离、也与环境隔离时的行为；第三章考虑它们的近邻环境所产生的影响；
接着第四章讨论磁矩之间可能存在的各种磁相互作用。到第五章我们就可以讨论长程序
（long range order）的出现，第六章讨论它与对称性破缺（broken symmetry）概念的
联系。最后两章把这一概念的含义贯彻到各种不同情形之中。全书采用 SI 单位制
（cgs 单位制的描述与换算表见附录 A）。

\section{磁矩}

磁学中最基本的对象是磁矩。在经典电磁学中，可以把磁矩等同于一个电流环。若有电流
$I$ 绕一个元（即无穷小的）有向环路流动，环面积为 $|d\mathbf{S}|$（见图~1.1(a)），
则磁矩 $d\boldsymbol{\mu}$ 为
\begin{equation*}
\mathbf{d}\boldsymbol{\mu} = I\,\mathbf{dS},\tag{1.1}
\end{equation*}
磁矩的单位是 A\,m$^2$。矢量 $\mathbf{dS}$ 的长度等于环的面积，其方向垂直于环面，
指向由环绕元电流的方向按右手定则确定。

这一对象也等价于一个磁偶极子（magnetic dipole）；之所以这样称呼，是因为它的行为
类似于电偶极子（相隔一小段距离的一正一负两个电荷）。因此可以设想：磁偶极子是这样
一种对象——两个带相反磁荷的磁单极子（magnetic monopoles），沿 $\mathbf{dS}$ 方向
相隔一小段距离（关于电磁学的背景知识见附录 B）。

\begin{figure}
\centering
\includegraphics[width=0.62\textwidth]{fig1_01.pdf}
\caption{（a）元磁矩：由元电流环产生的 $\mathrm{d}\boldsymbol{\mu}=I\,\mathrm{d}\mathbf{S}$。（b）一个有限回路（俯视图）上的电流 $I$ 的磁矩 $\boldsymbol{\mu}=I\int\mathrm{d}\mathbf{S}$，可由把许多无穷小电流环的磁矩求和得到。}
\end{figure}

磁矩 $d\boldsymbol{\mu}$ 垂直于电流环平面，因而既可以与环上电荷的角动量矢量平行，
也可以反平行。对于有限大小的回路，可以把其磁矩 $\boldsymbol{\mu}$ 表示为分布在
回路面积上的许多相同的无穷小电流环的磁矩之和（见图~1.1(b)）：相邻无穷小环的电流
彼此抵消，只留下沿回路周界流动的电流。于是
\begin{equation*}
\boldsymbol{\mu} = \int \mathrm{d}\boldsymbol{\mu} = I\int \mathrm{d}\mathbf{S}.\tag{1.2}
\end{equation*}

\subsection{磁矩与角动量}

电流环源于电荷的运动。本书考虑的所有电荷都与具有质量的粒子相联系，因此书中所有
电流环里质量也在作轨道运动——磁矩总是与角动量相伴随。

在原子中，沿轨道运动的电子所具有的磁矩 $\boldsymbol{\mu}$ 与该电子的角动量
$\mathbf{L}$ 方向一致，且与它成正比。我们写成
\begin{equation*}
\boldsymbol{\mu} = \gamma\,\mathbf{L},\tag{1.3}
\end{equation*}
其中 $\gamma$ 是一个常数，称为旋磁比（gyromagnetic ratio）。磁矩与角动量之间的
这一关系由 1915 年发现的爱因斯坦--德哈斯效应（Einstein--de Haas effect）所证实：
一根铁磁棒沿其轴向用细丝竖直悬挂（见图~1.2），起初静止且未磁化，随后沿棒长方向
加竖直磁场使其磁化。这一竖直磁化来自原子磁矩的排列，对应一定的角动量；为守恒
总角动量，铁棒将开始沿相反方向绕轴转动。测出棒的角动量，便可推知原子磁矩所携带
的角动量，从而得到旋磁比。爱因斯坦--德哈斯效应是磁化诱导的转动；还存在逆效应，
即巴尼特效应（Barnett effect）——转动诱导磁化。两种现象都表明磁矩与角动量相联系。
\bio{Albert Einstein（1879--1955）\\Wander J. de Haas（1878--1960）}

\begin{figure}
\centering
\includegraphics[width=0.5\textwidth]{fig1_02.pdf}
\caption{爱因斯坦--德哈斯效应。铁磁棒悬挂于细丝之下。线圈提供磁场使铁磁体磁化并产生
转动。实验可以共振方式实现：周期性地反转线圈中的电流（从而反转铁磁体中的磁化），
观察角响应随频率的变化。}
\end{figure}

\subsection{进动}

现在考虑处于磁场 $\mathbf{B}$ 中的磁矩 $\boldsymbol{\mu}$（见图~1.3）。磁矩的能量为
\begin{equation*}
E = -\boldsymbol{\mu}\cdot\mathbf{B},\tag{1.4}
\end{equation*}
（见附录 B）因此磁矩沿磁场方向时能量最低。磁矩还会受到力矩（torque）
\begin{equation*}
\mathbf{G} = \boldsymbol{\mu}\times\mathbf{B},\tag{1.5}
\end{equation*}
（见附录 B）——假如磁矩不携带任何角动量，这个力矩将把磁矩转向磁场方向\fn{1}{对于处于电场 $\boldsymbol{\mathcal{E}}$ 中的电偶极子 $\mathbf{p}$，能量为 $E=-\mathbf{p}\cdot\boldsymbol{\mathcal{E}}$，力矩为 $\mathbf{G}=\mathbf{p}\times\boldsymbol{\mathcal{E}}$。静止的电偶极矩只是两个相隔一定距离的静止电荷，不携带任何角动量，因此若 $\boldsymbol{\mathcal{E}}$ 与 $\mathbf{p}$ 不共线，力矩 $\mathbf{G}$ 将把 $\mathbf{p}$ 转向 $\boldsymbol{\mathcal{E}}$。而静止的磁矩则与角动量相联系，故行为不同。}。

然而，由于磁矩通过式~(1.3) 与角动量 $\mathbf{L}$ 相联系，而力矩等于角动量的时间
变化率，式~(1.5) 可改写为
\begin{equation*}
\frac{\mathrm{d}\boldsymbol{\mu}}{\mathrm{d}t} = \gamma\,\boldsymbol{\mu}\times\mathbf{B}.\tag{1.6}
\end{equation*}
这意味着 $\boldsymbol{\mu}$ 的变化同时垂直于 $\boldsymbol{\mu}$ 和 $\mathbf{B}$。
磁场并不把 $\boldsymbol{\mu}$ 拉向 $\mathbf{B}$，而是使 $\boldsymbol{\mu}$ 的方向
绕 $\mathbf{B}$ 进动（precession）。式~(1.6) 还表明 $|\boldsymbol{\mu}|$ 不随时间
改变。这一情形与陀螺的旋转完全类似。\fn{2}{设想一个自转轴相对竖直方向倾斜的陀螺。向下的重力对陀螺施加一个（水平的）力矩。若陀螺不自转，它只会倒下；但正因为它在自转，它具有平行于自转轴的角动量，力矩使自转轴在水平面内平行于力矩方向运动——陀螺进动了。}
\bio{Joseph Larmor（1857--1942）}

下面的例子将针对一个具体情形详细求解式~(1.6)。

\begin{example}{1.1}
设 $\mathbf{B}$ 沿 $z$ 方向，$\boldsymbol{\mu}$ 初始时与 $\mathbf{B}$ 成 $\theta$ 角
并位于 $xz$ 平面内（见图~1.4）。则
\begin{align*}
\dot{\mu}_x &= \gamma B\mu_y\tag{1.7}\\
\dot{\mu}_y &= -\gamma B\mu_x\tag{1.8}\\
\dot{\mu}_z &= 0,\tag{1.9}
\end{align*}
所以 $\mu_z$ 不随时间变化，而 $\mu_x$、$\mu_y$ 都在振荡。解这组微分方程得
\begin{align*}
\mu_x(t) &= |\boldsymbol{\mu}|\sin\theta\cos(\omega_{\mathrm{L}}t)\tag{1.10}\\
\mu_y(t) &= |\boldsymbol{\mu}|\sin\theta\sin(\omega_{\mathrm{L}}t)\tag{1.11}\\
\mu_z(t) &= |\boldsymbol{\mu}|\cos\theta,\tag{1.12}
\end{align*}
其中
\begin{equation*}
\omega_{\mathrm{L}} = \gamma B\tag{1.13}
\end{equation*}
称为拉莫尔进动频率（Larmor precession frequency）。
\end{example}

\begin{figure}
\centering
\includegraphics[width=0.32\textwidth]{fig1_03.pdf}
\caption{处于磁场 $\mathbf{B}$ 中的磁矩 $\boldsymbol{\mu}$，能量为
$-\boldsymbol{\mu}\cdot\mathbf{B}=-\mu B\cos\theta$。}
\end{figure}

\begin{figure}
\centering
\includegraphics[width=0.4\textwidth]{fig1_04.pdf}
\caption{磁场 $\mathbf{B}$ 中的磁矩 $\boldsymbol{\mu}$ 以拉莫尔进动频率 $\gamma B$ 绕磁场
进动，$\gamma$ 为旋磁比。$\mathbf{B}$ 沿 $z$ 轴，磁矩初始时位于 $xz$ 平面内、与
$\mathbf{B}$ 成 $\theta$ 角。磁矩绕半角为 $\theta$ 的圆锥进动。}
\end{figure}

注意旋磁比 $\gamma$ 是把角动量与磁矩（通过式~1.3）、又把进动频率与磁场（通过式~1.13）
联系起来的比例常数。进动现象暗示了后面内容的精微之处：磁场不仅使磁矩排列整齐，
还能诱发各种动力学效应。

\subsection{玻尔磁子}

在继续之前，值得快速估算一下原子磁矩的大小，从而推知旋磁比的量级。考虑一个电子
（电荷 $-e$，质量 $m_{\mathrm{e}}$）绕氢原子核作圆周轨道运动（见图~1.5）。绕核
电流为 $I=-e/\tau$，其中 $\tau=2\pi r/v$ 是轨道周期，$v=|\mathbf{v}|$ 是速率，
$r$ 是圆轨道半径。基态中电子角动量的大小 $m_{\mathrm{e}}vr$ 必等于 $\hbar$，于是
电子的磁矩为
\begin{equation*}
\mu = \pi r^{2} I = -\frac{e\hbar}{2m_{\mathrm{e}}} \equiv -\mu_{\mathrm{B}}\tag{1.14}
\end{equation*}
其中 $\mu_{\mathrm{B}}$ 是玻尔磁子（Bohr magneton），
\begin{equation*}
\mu_{\mathrm{B}} = \frac{e\hbar}{2m_{\mathrm{e}}}.\tag{1.15}
\end{equation*}
这是描述原子磁矩大小的方便单位，数值为 $9.274\times10^{-24}$~A\,m$^2$。注意式~(1.14)
中磁矩为负：由于电子带负电，其磁矩与角动量反平行。电子的旋磁比为
$\gamma=-e/2m_{\mathrm{e}}$，拉莫尔频率则为 $\omega_{\mathrm{L}}=|\gamma|B=eB/2m_{\mathrm{e}}$。

\begin{figure}
\centering
\includegraphics[width=0.42\textwidth]{fig1_05.pdf}
\caption{氢原子中绕核（单个质子）以速度 $\mathbf{v}$ 作轨道运动的电子。}
\end{figure}
\bio{Niels Bohr（1885--1962）}

\subsection{磁化强度与磁场}

磁性固体由大量带磁矩的原子组成。磁化强度（magnetization）$\mathbf{M}$ 定义为单位
体积的磁矩。通常在“连续介质近似”下考虑这个矢量：即在足够大的空间尺度上观察，
看不到单个原子磁矩带来的颗粒感。这样 $\mathbf{M}$ 可视为一个光滑的矢量场，只在
磁性固体边界处不连续。

自由空间（真空）中没有磁化。磁场可用线性相关的两个矢量场 $\mathbf{B}$ 与
$\mathbf{H}$ 描述：
\begin{equation*}
\mathbf{B} = \mu_{0}\mathbf{H},\tag{1.16}
\end{equation*}
其中 $\mu_{0}=4\pi\times10^{-7}$~H\,m$^{-1}$ 是真空磁导率。$\mathbf{B}$ 与 $\mathbf{H}$
只是彼此的标度版本：前者以特斯拉（T）度量，后者以 A\,m$^{-1}$ 度量。

在磁性固体中，$\mathbf{B}$ 与 $\mathbf{H}$ 的关系更为复杂，两个矢量场的大小和方向
都可以相差很大。一般的矢量关系是
\begin{equation*}
\mathbf{B} = \mu_{0}(\mathbf{H} + \mathbf{M}).\tag{1.17}
\end{equation*}
\mnote{由于历史原因，标准约定把 $\mathbf{B}$ 称为磁感应强度（magnetic induction）或磁通密度（magnetic flux density），把 $\mathbf{H}$ 称为磁场强度（magnetic field strength）。然而这些名称既繁琐又易误导。遵照通常的用法，我们把二者都简称为磁场，由字母 B 和 H 表明所指为何。}

当磁化强度 $\mathbf{M}$ 与磁场 $\mathbf{H}$ 呈线性关系时，该固体称为线性材料，
写作
\begin{equation*}
\mathbf{M} = \chi\mathbf{H},\tag{1.18}
\end{equation*}
$\chi$ 是无量纲量，称为磁化率（magnetic susceptibility）。此时 $\mathbf{B}$ 与
$\mathbf{H}$ 仍呈线性关系：
\begin{equation*}
\mathbf{B} = \mu_{0}(1+\chi)\mathbf{H} = \mu_{0}\mu_{\mathrm{r}}\mathbf{H},\tag{1.19}
\end{equation*}
$\mu_{\mathrm{r}}=1+\chi$ 是材料的相对磁导率。

下面是一段警示。它源于我们必须在可磁化介质中非常小心地定义场。考虑自由空间中
施加磁场 $\mathbf{B}_{\mathrm{a}}$、$\mathbf{H}_{\mathrm{a}}$ 的某个区域，二者满足
$\mathbf{B}_{\mathrm{a}}=\mu_{0}\mathbf{H}_{\mathrm{a}}$。至此一切都很简单。现在把
一块磁性固体插入该区域。固体内部的场 $\mathbf{B}_{\mathrm{i}}$、$\mathbf{H}_{\mathrm{i}}$
可能与 $\mathbf{B}_{\mathrm{a}}$、$\mathbf{H}_{\mathrm{a}}$ 大不相同，原因是固体中
所有磁矩产生的磁场。事实上 $\mathbf{B}_{\mathrm{i}}$ 与 $\mathbf{H}_{\mathrm{i}}$
都可以依赖于测量点在固体内的位置——\fn{3}{被磁化的样品也会影响其外部的磁场，正如它在内部（这里所考虑的）的影响一样；玩过条形磁铁与铁屑的人都清楚这一点。}除非是
椭球形的特殊样品（见图~1.6）。若磁场沿椭球某个主轴施加，则在样品内处处有
\begin{equation*}
\mathbf{H}_{\mathrm{i}} = \mathbf{H}_{\mathrm{a}} - N\mathbf{M},\tag{1.20}
\end{equation*}
其中 $N$ 是相应的退磁因子（demagnetizing factor，见附录 D）。为从 $\mathbf{H}_{\mathrm{a}}$
得到 $\mathbf{H}_{\mathrm{i}}$ 而需加的“修正项” $\mathbf{H}_{\mathrm{d}}=-N\mathbf{M}$
称为退磁场（demagnetizing field）。类似地
\begin{equation*}
\mathbf{B}_{\mathrm{i}} = \mu_{0}(\mathbf{H}_{\mathrm{i}} + \mathbf{M}) = \mathbf{B}_{\mathrm{a}} + \mu_{0}(1-N)\mathbf{M}.\tag{1.21}
\end{equation*}

\begin{figure}
\centering
\includegraphics[width=0.55\textwidth]{fig1_06.pdf}
\caption{主轴为 $a$、$b$、$c$ 的椭球形磁化固体样品。球（$a=b=c$）与薄板（$a,b\to\infty$，
$c=0$）是其特例。}
\end{figure}

\begin{example}{1.2}
对球形样品，$N=\frac{1}{3}$，球内场为
\begin{equation*}
\mathbf{H}_{\mathrm{i}} = \mathbf{H}_{\mathrm{a}} - \frac{\mathbf{M}}{3},\tag{1.22}
\end{equation*}
\begin{equation*}
\mathbf{B}_{\mathrm{i}} = \mathbf{B}_{\mathrm{a}} + \frac{2\mu_{0}\mathbf{M}}{3}.\tag{1.23}
\end{equation*}
当磁化强度与（插入样品前测得的）外加场 $|H_{\mathrm{a}}|=|B_{\mathrm{a}}|/\mu_0$
相比不够小时，这些退磁修正必须认真对待。不过，对弱磁性的特殊情形可以绕开这些
复杂因素：对 $\chi\ll1$ 的线性材料，有 $M\ll H$、$H_{\mathrm{i}}\approx H_{\mathrm{a}}$
及 $B_{\mathrm{i}}\approx\mu_0 H_{\mathrm{i}}$，于是可以放心地假定材料中的磁场就是
我们施加的磁场。第二、三章讨论抗磁性的较弱效应时将采用这一近似。\fn{4}{即使对这些材料，在精确的实验工作中退磁场仍然必须予以考虑。}在铁磁体中，退磁效应则总是重要的。
\end{example}

\begin{example}{1.3}
材料的内禀磁化率为
\begin{equation*}
\chi_{\text{intrinsic}} = \frac{M}{H_{\mathrm{i}}}.\tag{1.24}
\end{equation*}
但实验上测得的并不是这个内禀量：实验测的是磁化强度 $M$ 对外加场 $H_{\mathrm{a}}$
的响应，即
\begin{equation*}
\chi_{\text{experimental}} = \frac{M}{H_{\mathrm{a}}}.\tag{1.25}
\end{equation*}
两个量之间有关系
\begin{equation*}
\chi_{\text{experimental}} = \frac{M}{H_{\mathrm{i}} + NM} = \frac{M/H_{\mathrm{i}}}{1 + NM/H_{\mathrm{i}}} = \frac{\chi_{\text{intrinsic}}}{1 + N\chi_{\text{intrinsic}}}.\tag{1.26}
\end{equation*}
当 $\chi_{\text{intrinsic}}\ll1$ 时二者的区分无关紧要；当 $\chi_{\text{intrinsic}}$
接近或超过 1 时，区分可能极为重要。例如，从上方逼近居里温度的铁磁体（见第四章）
有 $\chi_{\text{intrinsic}}\to\infty$，而 $\chi_{\text{experimental}}\to1/N$。
\end{example}

此处还要提出最后一点提醒：铁磁材料可能没有净磁矩，因为它由磁畴（domains）构成。
\fn{5}{关于磁畴的更多内容见第 6.7 节。}每个畴内有均匀的磁化，但相邻畴的磁化方向
不同。因此样品看起来可能没有磁化，尽管在足够小的尺度上所有磁矩都局域地排列一致。

本章余下部分将讨论磁矩与经典力学（1.2 节）和量子力学（1.3 节）相关的若干进一步
的内容。

\section{经典力学与磁矩}

本节用纯经典的论证描述外磁场对电荷体系的影响。先考虑单个电荷，再用该结果计算
电荷体系的磁化强度。电磁学重要结果摘要见附录 B。

\subsection{正则动量}

经典力学中，电荷为 $q$、速度为 $\mathbf{v}$ 的粒子在电场 $\boldsymbol{\varepsilon}$
和磁场 $\mathbf{B}$ 中所受力为
\begin{equation*}
\mathbf{F} = q(\boldsymbol{\varepsilon} + \mathbf{v}\times\mathbf{B})\tag{1.27}
\end{equation*}
称为洛伦兹力（Lorentz force）。用这个熟悉的方程可以说明磁场如何修正带电粒子的
动量。利用 $\mathbf{F}=m\,\mathrm{d}\mathbf{v}/\mathrm{d}t$、$\mathbf{B}=\nabla\times\mathbf{A}$
和 $\boldsymbol{\varepsilon}=-\nabla V-\partial\mathbf{A}/\partial t$（$V$ 为电势，
$\mathbf{A}$ 为磁矢势，$m$ 为粒子质量），式~(1.27) 可改写为
\begin{equation*}
m\frac{\mathrm{d}\mathbf{v}}{\mathrm{d}t} = -q\nabla V - q\frac{\partial\mathbf{A}}{\partial t} + q\,\mathbf{v}\times(\nabla\times\mathbf{A}).\tag{1.28}
\end{equation*}
\mnote{矢量恒等式列表见附录 G。另注意 $\mathbf{v}$ 不随位置变化。}
利用矢量恒等式
\begin{equation*}
\mathbf{v}\times(\nabla\times\mathbf{A}) = \nabla(\mathbf{v}\cdot\mathbf{A}) - (\mathbf{v}\cdot\nabla)\mathbf{A}\tag{1.29}
\end{equation*}
化简式~(1.28)，得
\begin{equation*}
m\frac{\mathrm{d}\mathbf{v}}{\mathrm{d}t} + q\left(\frac{\partial\mathbf{A}}{\partial t} + (\mathbf{v}\cdot\nabla)\mathbf{A}\right) = -q\nabla(V - \mathbf{v}\cdot\mathbf{A}).\tag{1.30}
\end{equation*}
注意 $m\,\mathrm{d}\mathbf{v}/\mathrm{d}t$ 是在与粒子一起运动的坐标系中测得的力；
偏导数 $\partial\mathbf{A}/\partial t$ 度量空间固定点上 $\mathbf{A}$ 的变化率。
于是式~(1.30) 可写作
\begin{equation*}
\frac{\mathrm{d}}{\mathrm{d}t}(m\mathbf{v} + q\mathbf{A}) = -q\nabla(V - \mathbf{v}\cdot\mathbf{A})\tag{1.31}
\end{equation*}
其中 $\mathrm{d}\mathbf{A}/\mathrm{d}t$ 是 $\mathbf{A}$ 的对流导数（convective derivative）
\begin{equation*}
\frac{\mathrm{d}\mathbf{A}}{\mathrm{d}t} = \frac{\partial\mathbf{A}}{\partial t} + (\mathbf{v}\cdot\nabla)\mathbf{A},\tag{1.32}
\end{equation*}
它度量运动粒子所在处 $\mathbf{A}$ 的变化率。式~(1.31) 具有牛顿第二定律的形式
（即“某个像动量的量的变化率等于某个像势能的量的梯度”），由此自然地引导我们定义
正则动量（canonical momentum）
\begin{equation*}
\mathbf{p} = m\mathbf{v} + q\mathbf{A}\tag{1.33}
\end{equation*}
以及带电粒子感受到的一个依赖于速度的有效势能 $q(V-\mathbf{v}\cdot\mathbf{A})$。
无磁场（$\mathbf{A}=0$）时正则动量退化为熟悉的动量 $m\mathbf{v}$。动能仍等于
$\frac{1}{2}mv^{2}$，用正则动量可写成 $(\mathbf{p}-q\mathbf{A})^{2}/2m$。这一结果
下文就会用到，书中后面在磁场中与动能相联系的量子力学算符也写作
$(-\mathrm{i}\hbar\nabla-q\mathbf{A})^{2}/2m$。

\subsection{玻尔--范列文定理}

下一步计算固体中电子系统的净磁矩，即求磁场诱导的磁化强度。由式~(1.4)，磁化强度
正比于系统能量对外磁场的导数。\fn{6}{另见附录 E。}式~(1.27) 表明磁场对带电粒子的
力总是垂直于其速度，因此不做功，系统的能量不可能依赖于外磁场。若系统能量与外磁场
无关，则不会有磁化强度。
\bio{Niels Bohr（1885--1962）\\Hendrika J. van Leeuwen（1887--1974）\\Ludwig Boltzmann（1844--1906）}

这一思想凝结为玻尔--范列文定理（Bohr--van Leeuwen theorem）：经典体系中不存在
热平衡磁化强度。可概述其证明如下：在经典统计力学中，$N$ 个电荷均为 $q$ 的粒子
的配分函数 $Z$ 正比于
\begin{equation*}
\int\!\!\int\!\cdots\!\int \exp(-\beta E(\{\mathbf{r}_i,\mathbf{p}_i\}))\,\mathrm{d}\mathbf{r}_1\cdots\mathrm{d}\mathbf{r}_N\,\mathrm{d}\mathbf{p}_1\cdots\mathrm{d}\mathbf{p}_N,\tag{1.34}
\end{equation*}
其中 $\beta=1/k_{\mathrm{B}}T$，$k_{\mathrm{B}}$ 是玻尔兹曼常数，$T$ 是温度，
$i=1,\ldots,N$。这里 $E(\{\mathbf{r}_i,\mathbf{p}_i\})$ 是 $N$ 个带电粒子处于位置
$\mathbf{r}_1,\mathbf{r}_2,\ldots,\mathbf{r}_N$、动量 $\mathbf{p}_1,\mathbf{p}_2,\ldots,\mathbf{p}_N$
时的能量。积分于是展布在 $6N$ 维相空间（$3N$ 个位置坐标与 $3N$ 个动量坐标）上。
如上节所示，磁场的效果是把每个粒子的动量平移 $q\mathbf{A}$，因此须把 $\mathbf{p}_i$
换成 $\mathbf{p}_i-q\mathbf{A}$。动量积分限从 $-\infty$ 到 $\infty$，这个平移可以
通过移动动量积分的原点而吸收掉。所以配分函数不是磁场的函数，自由能
$F=-k_{\mathrm{B}}T\log Z$ 亦然（见附录 E）。故经典体系中磁化强度必为零。

这个结论初看令人意外。没有外磁场时电子走直线；加了磁场，路径弯曲（实为螺旋线）
并作回旋轨道（cyclotron orbits）。人们难免想说：这些方向一致的弯曲回旋轨道必然
贡献净磁矩，因而外磁场应当影响能量。但这一论证的谬误可借图~1.7 理解：图中画出
经典体系中电子在外磁场下的轨道。电子确实作回旋轨道，也确实对应净磁矩；把这些轨道
求和，得到系统边缘处净的逆时针环流（如图~1.1）。然而表面附近的电子无法完成完整的
圆环，它们与表面反复发生弹性碰撞，沿样品周界作所谓的跳跃轨道（skipping orbits）。
体内电子的逆时针电流与在表面反射、散射的电子的顺时针跳跃轨道电流恰好抵消。

因此玻尔--范列文定理似乎是对的，却与实验相抵触：许多含电子的真实体系确有净磁化
强度。那么定理的前提必定有问题——而这些前提正是经典力学！于是我们得出结论：
经典力学不足以解释磁性材料这一最基本的性质，要说明真实材料的磁性，量子理论
不可避免。下一节将较详细地讨论电子的量子力学。

\begin{figure}
\centering
\includegraphics[width=0.5\textwidth]{fig1_07.pdf}
\caption{经典体系中的电子在外磁场下在体内作回旋轨道（图中逆时针进动），贡献净的逆时针
轨道电流（见图~1.1）。这一净电流与在表面散射、沿样品顺时针进动的电子的跳跃轨道电流
恰好抵消。}
\end{figure}

\section{自旋的量子力学}

本节简要回顾与电子自旋量子力学有关的一些结果。关于角动量量子力学的更完整论述
可参见章末延伸阅读；与量子物理和原子物理相关的一些结果也列于附录 C。

\subsection{轨道角动量与自旋角动量}

1.1 节讨论的电子角动量与电子绕核的轨道运动相联系，称为轨道角动量（orbital angular
momentum）。在真实原子中它依赖于电子占据的电子态。按通常方式定义量子数 $l$ 与
$m_l$（见附录 C），沿固定轴（取 $z$ 轴）的轨道角动量分量为 $m_l\hbar$，轨道角动量
的大小\fn{7}{严格地说，well-defined（有确切定义）的量是角动量的平方和磁偶极矩的平方：算符 $\hat{\mathbf{L}}^2$ 的本征值为 $l(l+1)\hbar^2$，$\hat{L}_z$ 的本征值为 $m_l\hbar$；类似地，算符 $\hat{\boldsymbol{\mu}}^2$ 的本征值为 $l(l+1)\mu_{\mathrm{B}}^2$，$\hat{\mu}_z$ 的本征值为 $-m_l\mu_{\mathrm{B}}$。}为
$\sqrt{l(l+1)}\hbar$。于是沿 $z$ 轴的磁矩分量为 $-m_l\mu_{\mathrm{B}}$，总磁偶极矩的
大小为 $\sqrt{l(l+1)}\mu_{\mathrm{B}}$。

情况还因下述事实而更复杂：电子还具有与内禀角动量相联系的内禀磁矩。电子的内禀
角动量称为自旋（spin）。这个名字源于人们曾以为电子绕自身轴进动；但电子是点粒子，
这实在难以想象。\mnote{这一概念已经改变，但名称沿用至今。这倒也不是坏事：电子自旋的行为如此违反直觉，很难找到任何一个词能完全恰当地描述它！}

电子自旋由自旋量子数 $s$ 刻画，对电子 $s=\frac{1}{2}$。角动量任一分量只能取
$2s+1$ 个可能值：$s\hbar,(s-1)\hbar,\ldots,-s\hbar$。自旋角动量分量写作 $m_s\hbar$。
对 $s=\frac{1}{2}$ 的电子，只有两个可能值 $m_s=\pm\frac{1}{2}$，沿特定轴的角动量
分量为 $\hbar/2$ 或 $-\hbar/2$，分别称为“向上”和“向下”。电子自旋角动量的大小\fn{8}{严格地，算符 $\hat{\mathbf{S}}^2$ 的本征值是 $s(s+1)\hbar^2$。}为
$\sqrt{s(s+1)}\hbar=\sqrt{3}\hbar/2$。

自旋角动量联系着磁矩，其沿特定轴的分量为 $-g\mu_{\mathrm{B}}m_s$，大小为
$\sqrt{s(s+1)}\,g\mu_{\mathrm{B}}=\sqrt{3}\,g\mu_{\mathrm{B}}/2$。式中 $g$ 是常数，
称为 $g$ 因子（g-factor）。\mnote{$g$ 因子将在附录 C.6 中更详细地讨论。}$g$ 因子
近似为 2，因此尽管自旋是半整数的，电子内禀磁矩沿 $z$ 轴的分量\fn{9}{$\mp$ 号这样上下颠倒，是因为磁矩与角动量反平行，其根源是电子带负电：$m_s=+\frac{1}{2}$ 时磁矩为 $-\mu_{\mathrm{B}}$；$m_s=-\frac{1}{2}$ 时磁矩为 $+\mu_{\mathrm{B}}$。}为 $\mp\mu_{\mathrm{B}}$。
磁场 $\mathbf{B}$ 中电子的能量因此为
\begin{equation*}
E = g\mu_{\mathrm{B}}m_s B.\tag{1.35}
\end{equation*}
于是电子的能级在磁场中分裂 $g\mu_{\mathrm{B}}B$，称为塞曼分裂（Zeeman splitting）。
\bio{Pieter Zeeman（1865--1943）}

一般地，原子中的电子可以既有轨道角动量又有自旋角动量，二者组合起来。$g$ 因子在
真实原子中因此可取不同数值，依赖于自旋与轨道角动量的相对贡献。这一点下一章
再谈。

电子的角动量总是 $\hbar$ 的整数或半整数倍。为方便起见，在角动量算符的表达式中
略去 $\hbar$ 因子，即约定这些算符以 $\hbar$ 为单位度量角动量。本书今后定义的角动量
算符（如 $\hat{\mathbf{L}}$）都使角动量为 $\hbar\hat{\mathbf{L}}$。这可以简化后文
的表达式。
\bio{Wolfgang Pauli（1900--1958）}

\subsection{泡利自旋矩阵与旋量}

电子自旋的行为联系着一套相当奇特的代数，其基础是三个泡利自旋矩阵（Pauli spin
matrices）：
\begin{equation*}
\hat{\sigma}_x = \begin{pmatrix} 0 & 1 \\ 1 & 0 \end{pmatrix},\qquad
\hat{\sigma}_y = \begin{pmatrix} 0 & -\mathrm{i} \\ \mathrm{i} & 0 \end{pmatrix},\qquad
\hat{\sigma}_z = \begin{pmatrix} 1 & 0 \\ 0 & -1 \end{pmatrix}.\tag{1.36}
\end{equation*}
把它们看作“矩阵构成的矢量”会很方便：
\begin{equation*}
\boldsymbol{\sigma} = (\hat{\sigma}_x, \hat{\sigma}_y, \hat{\sigma}_z).\tag{1.37}
\end{equation*}

在继续之前，回顾几个可以直接代入验证的结果。设
\begin{equation*}
\mathbf{a} = \begin{pmatrix} a_1 \\ a_2 \\ a_3 \end{pmatrix}\tag{1.38}
\end{equation*}
为三分量矢量，则 $\boldsymbol{\sigma}\cdot\mathbf{a}$ 是矩阵
\begin{equation*}
\boldsymbol{\sigma}\cdot\mathbf{a} = \begin{pmatrix} a_3 & a_1-\mathrm{i}a_2 \\ a_1+\mathrm{i}a_2 & -a_3 \end{pmatrix}.\tag{1.39}
\end{equation*}
这样的矩阵可以相乘，例如
\begin{equation*}
(\boldsymbol{\sigma}\cdot\mathbf{a})(\boldsymbol{\sigma}\cdot\mathbf{b}) = \mathbf{a}\cdot\mathbf{b} + \mathrm{i}\,\boldsymbol{\sigma}\cdot(\mathbf{a}\times\mathbf{b})\tag{1.40}
\end{equation*}
以及
\begin{equation*}
(\boldsymbol{\sigma}\cdot\mathbf{a})^{2} = |\mathbf{a}|^{2}.\tag{1.41}
\end{equation*}
\mnote{注意式 1.40 与 1.41 中的每一项都是矩阵。$\mathbf{a}\cdot\mathbf{b}$ 是 $\mathbf{a}\cdot\mathbf{b}\,\mathbf{I}$ 的简写，$\mathbf{I}=\begin{pmatrix}1&0\\0&1\end{pmatrix}$ 是单位矩阵；类似地 $|\mathbf{a}|^2$ 是 $|\mathbf{a}|^2\mathbf{I}$ 的简写。}

现在定义自旋角动量算符
\begin{equation*}
\hat{\mathbf{S}} = \frac{1}{2}\hat{\boldsymbol{\sigma}}\tag{1.42}
\end{equation*}
于是
\begin{equation*}
\hat{S}_x = \frac{1}{2}\begin{pmatrix} 0 & 1 \\ 1 & 0 \end{pmatrix},\qquad
\hat{S}_y = \frac{1}{2}\begin{pmatrix} 0 & -\mathrm{i} \\ \mathrm{i} & 0 \end{pmatrix},\qquad
\hat{S}_z = \frac{1}{2}\begin{pmatrix} 1 & 0 \\ 0 & -1 \end{pmatrix}.\tag{1.43}
\end{equation*}
再次注意我们采用角动量以 $\hbar$ 为单位的约定，与电子相联系的角动量实际上是
$\hbar\hat{\mathbf{S}}$。（有些书把 $\hat{\mathbf{S}}$ 定义为 $\hbar\hat{\boldsymbol{\sigma}}/2$。）

只有 $\hat{S}_z$ 是对角的，因此当电子自旋沿 $z$ 方向时表象特别简单。$\hat{S}_z$
的本征值记为 $m_s$，取值 $m_s=\pm\frac{1}{2}$，相应本征态为 $|\!\uparrow_z\rangle$
与 $|\!\downarrow_z\rangle$：
\begin{equation*}
|\!\uparrow_z\rangle = \begin{pmatrix}1\\0\end{pmatrix}\tag{1.44}
\end{equation*}
\begin{equation*}
|\!\downarrow_z\rangle = \begin{pmatrix}0\\1\end{pmatrix}\tag{1.45}
\end{equation*}
分别对应自旋平行与反平行于 $z$ 轴。（“左右矢”记号即形如 $|\psi\rangle$ 的态写法，
见附录 C。）于是
\begin{equation*}
\hat{S}_z|\!\uparrow_z\rangle = \frac{1}{2}|\!\uparrow_z\rangle,\qquad
\hat{S}_z|\!\downarrow_z\rangle = -\frac{1}{2}|\!\downarrow_z\rangle.\tag{1.46}
\end{equation*}
自旋平行或反平行于 $x$、$y$ 轴的本征态是
\begin{equation*}
|\!\uparrow_x\rangle = \frac{1}{\sqrt{2}}\begin{pmatrix}1\\1\end{pmatrix},\qquad
|\!\downarrow_x\rangle = \frac{1}{\sqrt{2}}\begin{pmatrix}1\\-1\end{pmatrix}\tag{1.47}
\end{equation*}
\begin{equation*}
|\!\uparrow_y\rangle = \frac{1}{\sqrt{2}}\begin{pmatrix}1\\i\end{pmatrix},\qquad
|\!\downarrow_y\rangle = \frac{1}{\sqrt{2}}\begin{pmatrix}1\\-i\end{pmatrix}.\tag{1.48}
\end{equation*}
自旋波函数的这种双分量表示称为旋量（spinor）表示，这些态称为旋量。一般态可写为
\begin{equation*}
|\psi\rangle = \begin{pmatrix}a\\b\end{pmatrix} = a|\!\uparrow_z\rangle + b|\!\downarrow_z\rangle\tag{1.49}
\end{equation*}
$a$、$b$ 为复数，\fn{10}{当然，在一般情形下它们可以是位置的函数。}按惯例归一化为
\begin{equation*}
|a|^{2} + |b|^{2} = 1.\tag{1.50}
\end{equation*}

\begin{figure}
\centering
\includegraphics[width=0.5\textwidth]{fig1_08.pdf}
\caption{黎曼球面表示自旋 $\frac{1}{2}$ 粒子的自旋态。自旋矢量 $\mathbf{S}$ 位于
单位球面上。从球面南极到 $\mathbf{S}$ 的连线与水平赤道面（阴影）交于
$q=x+\mathrm{i}y$，该水平面被视为复平面。图中标出六种情形（$\mathbf{S}$ 平行或
反平行于 $x$、$y$、$z$ 轴）下复数 $q$ 的数值。}
\end{figure}
\bio{George F. B. Riemann（1826--1866）}

总自旋角动量算符 $\hat{\mathbf{S}}$ 定义为
\begin{equation*}
\hat{\mathbf{S}} = \begin{pmatrix}\hat{S}_x\\ \hat{S}_y\\ \hat{S}_z\end{pmatrix} = \mathbf{i}\hat{S}_x + \mathbf{j}\hat{S}_y + \mathbf{k}\hat{S}_z,\tag{1.51}
\end{equation*}
$\mathbf{i}$、$\mathbf{j}$、$\mathbf{k}$ 为笛卡尔单位矢量。算符 $\hat{\mathbf{S}}^{2}$
则为
\begin{equation*}
\hat{\mathbf{S}}^{2} = \hat{S}_x^{2} + \hat{S}_y^{2} + \hat{S}_z^{2}.\tag{1.52}
\end{equation*}
由于 $\hat{S}_x^2$、$\hat{S}_y^2$、$\hat{S}_z^2$ 的本征值总是
$\frac{1}{4}=(\pm\frac{1}{2})^{2}$，对任何自旋态 $|\psi\rangle$ 都有
\begin{equation*}
\hat{\mathbf{S}}^{2}|\psi\rangle = (\hat{S}_x^{2}+\hat{S}_y^{2}+\hat{S}_z^{2})|\psi\rangle = \left(\frac{1}{4}+\frac{1}{4}+\frac{1}{4}\right)|\psi\rangle = \frac{3}{4}|\psi\rangle.\tag{1.53}
\end{equation*}

其中许多结果可以推广到自旋量子数 $s>\frac{1}{2}$ 的粒子：最重要的结论是
$\hat{\mathbf{S}}^{2}$ 的本征值变为 $s(s+1)$。对本章考虑的 $s=\frac{1}{2}$，
$s(s+1)=\frac{3}{4}$，与式~(1.53) 一致。自旋算符间的对易关系为
\begin{equation*}
[\hat{S}_x, \hat{S}_y] = i\hat{S}_z\tag{1.54}
\end{equation*}
及其循环置换。用式~(1.40) 与 (1.42) 极易证明。这些算符都与 $\hat{\mathbf{S}}^{2}$
对易，故
\begin{equation*}
[\hat{\mathbf{S}}^{2}, \hat{S}_z] = 0.\tag{1.55}
\end{equation*}
因此可以同时确定总自旋与其一个分量，但不能同时确定多于一个的分量。

图~1.8 给出一个有助于思考自旋的几何构造。自旋矢量 $\mathbf{S}$ 指向三维空间中的
一点；由于量子态已归一化，$\mathbf{S}$ 位于单位球面上。从矢量 $\mathbf{S}$ 的端点
向球面南极引一条线，观察它与水平面（图~1.8 中阴影）的交点 $q$。把该水平面当作
阿甘德图（Argand diagram），$x$ 轴为实轴、$y$ 轴为虚轴，则 $q=x+\mathrm{i}y$ 是
复数。于是 $\mathbf{S}$ 的旋量表示为 $\begin{pmatrix}1\\q\end{pmatrix}$，归一化后为
\begin{equation*}
\mathbf{S} = \frac{1}{\sqrt{1+|q|^{2}}}\begin{pmatrix}1\\q\end{pmatrix}.\tag{1.56}
\end{equation*}
这个表示中的球面称为黎曼球面（Riemann sphere）。

\subsection{升降算符}

升降算符（raising and lowering operators）$\hat{S}_{+}$ 与 $\hat{S}_{-}$ 定义为
\begin{equation*}
\begin{aligned}
\hat{S}_{+} &= \hat{S}_x + i\hat{S}_y\\
\hat{S}_{-} &= \hat{S}_x - i\hat{S}_y.
\end{aligned}\tag{1.57}
\end{equation*}

算符 $\hat{A}$ 厄米（Hermitian）要求 $\hat{A}^{\dagger}=\hat{A}$，$\dagger$ 表示
伴随操作（对矩阵即“先转置再对每个元素取复共轭”）。升降算符不是厄米的（因为
$\hat{S}_{+}^{\dagger}=\hat{S}_{-}$，$\hat{S}_{-}^{\dagger}=\hat{S}_{+}$），因此不
对应可观测的量，但它们非常有用。直接运用式~(1.54) 与 (1.57) 得到对易关系：
\begin{equation*}
[\hat{S}_{+}, \hat{S}_{-}] = 2\hat{S}_z,\tag{1.58}
\end{equation*}
\begin{equation*}
[\hat{S}_z, \hat{S}_{\pm}] = \pm\hat{S}_{\pm},\tag{1.59}
\end{equation*}
以及
\begin{equation*}
[\hat{\mathbf{S}}^{2}, \hat{\mathbf{S}}_{\pm}] = 0.\tag{1.60}
\end{equation*}
另一个直接代入即可证明的有用关系是
\begin{equation*}
\hat{S}_{+}\hat{S}_{-} + \hat{S}_{-}\hat{S}_{+} = 2(\hat{S}_x^{2}+\hat{S}_y^{2})\tag{1.61}
\end{equation*}
它给出 $\hat{\mathbf{S}}^{2}$ 的一个方便表示：
\begin{align*}
\hat{\mathbf{S}}^{2} &= \hat{S}_x^{2}+\hat{S}_y^{2}+\hat{S}_z^{2}\tag{1.62}\\
&= \frac{1}{2}\left(\hat{S}_{+}\hat{S}_{-}+\hat{S}_{-}\hat{S}_{+}\right)+\hat{S}_z^{2}.\tag{1.63}
\end{align*}

写成矩阵，升降算符为
\begin{equation*}
\hat{S}_{+} = \begin{pmatrix}0 & 1\\ 0 & 0\end{pmatrix}\tag{1.64}
\end{equation*}
\begin{equation*}
\hat{S}_{-} = \begin{pmatrix}0 & 0\\ 1 & 0\end{pmatrix}.\tag{1.65}
\end{equation*}
\mnote{升降算符得名于它们对自旋态的作用。可以直接证明：$\hat{S}_+|\!\uparrow_z\rangle=0$，$\hat{S}_+|\!\downarrow_z\rangle=|\!\uparrow_z\rangle$，$\hat{S}_-|\!\uparrow_z\rangle=|\!\downarrow_z\rangle$，$\hat{S}_-|\!\downarrow_z\rangle=0$。因此升算符把自旋角动量的 $z$ 分量升高 $\hbar$，降算符把它降低 $\hbar$；若 $z$ 分量已处于最大（最小）值，$\hat{S}_+$（$\hat{S}_-$）将直接湮灭该态。}
利用式~(1.43)、(1.63)、(1.64) 与 (1.65) 即得
\begin{equation*}
\hat{\mathbf{S}}^{2} = \frac{3}{4}\begin{pmatrix}1 & 0\\ 0 & 1\end{pmatrix},\tag{1.66}
\end{equation*}
与式~(1.53) 一致。

\subsection{两个自旋的耦合}

现在考虑两个自旋 $\frac{1}{2}$ 粒子，由哈密顿量\fn{11}{式 1.67 这类相互作用在本书中将极为重要：超精细相互作用（见第二章）与海森伯交换相互作用（见第四章）都取这种形式。}
\begin{equation*}
\hat{\mathcal{H}} = A\,\hat{\mathbf{S}}^{a}\cdot\hat{\mathbf{S}}^{b},\tag{1.67}
\end{equation*}
描述的相互作用耦合，其中 $\hat{\mathbf{S}}^{a}$ 与 $\hat{\mathbf{S}}^{b}$ 是两个
粒子自旋的算符。把二者当作一个联合体时，总自旋也可用一个算符表示：
\begin{equation*}
\hat{\mathbf{S}}^{\mathrm{tot}} = \hat{\mathbf{S}}^{a}+\hat{\mathbf{S}}^{b}\tag{1.68}
\end{equation*}
于是
\begin{equation*}
(\hat{\mathbf{S}}^{\mathrm{tot}})^{2} = (\hat{\mathbf{S}}^{a})^{2}+(\hat{\mathbf{S}}^{b})^{2}+2\hat{\mathbf{S}}^{a}\cdot\hat{\mathbf{S}}^{b}.\tag{1.69}
\end{equation*}

两个自旋 $\frac{1}{2}$ 粒子的组合体的自旋量子数为 $s=0$ 或 1。
$(\hat{\mathbf{S}}^{\mathrm{tot}})^{2}$ 的本征值为 $s(s+1)$，故 $s=0$ 或 1 时分别为
0 或 2。由式~(1.53)，$(\hat{\mathbf{S}}^{a})^{2}$ 与 $(\hat{\mathbf{S}}^{b})^{2}$ 的
本征值都是 $\frac{3}{4}$。于是由式~(1.69)
\begin{equation*}
\hat{\mathbf{S}}^{a}\cdot\hat{\mathbf{S}}^{b} = \begin{cases} \frac{1}{4} & \text{若 } s=1\\[2pt] -\frac{3}{4} & \text{若 } s=0. \end{cases}\tag{1.70}
\end{equation*}
因为哈密顿量为 $\hat{\mathcal{H}}=A\hat{\mathbf{S}}^{a}\cdot\hat{\mathbf{S}}^{b}$，
系统对 $s=0$、$1$ 有两个能级：
\begin{equation*}
E = \begin{cases} \frac{A}{4} & \text{若 } s=1\\[2pt] -\frac{3A}{4} & \text{若 } s=0. \end{cases}\tag{1.71}
\end{equation*}
每个态的简并度为 $2s+1$，所以 $s=0$ 的态是单态（singlet），$s=1$ 的态是三重态
（triplet）。该态自旋的 $z$ 分量 $m_s$：单态取 0，三重态取 $-1$、$0$、$1$ 三个值
之一。

式~(1.70) 列出了 $\hat{\mathbf{S}}^{a}\cdot\hat{\mathbf{S}}^{b}$ 的本征值，描述其
本征态也有用处。先考虑如下基底：
\begin{equation*}
|\!\uparrow\uparrow\rangle,\quad |\!\uparrow\downarrow\rangle,\quad |\!\downarrow\uparrow\rangle,\quad |\!\downarrow\downarrow\rangle.\tag{1.72}
\end{equation*}
在该表示中第一个箭头指标记 $a$ 的自旋的 $z$ 分量，第二个箭头指标记 $b$ 的自旋的
$z$ 分量。$\hat{\mathbf{S}}^{a}\cdot\hat{\mathbf{S}}^{b}$ 的本征态是这些基底态的
线性组合，列于表~1.1；其计算见习题 1.9。注意 $m_s$ 等于各单个自旋 $z$ 分量之和。
而且由于本征态是原基底下态的混合，一般不可能同时知道原来两个自旋各自的 $z$ 分量
和联合体的总自旋——这是一般性的特征，在更复杂的情形中会变得更加重要。

\begin{table}
\centering
\caption{$\hat{\mathbf{S}}^{a}\cdot\hat{\mathbf{S}}^{b}$ 的本征态及相应的 $m_s$、$s$ 与 $\hat{\mathbf{S}}^{a}\cdot\hat{\mathbf{S}}^{b}$ 的本征值。}
\begin{tabular}{cccc}
\toprule
本征态 & $m_s$ & $s$ & $\hat{\mathbf{S}}^{a}\cdot\hat{\mathbf{S}}^{b}$ \\
\midrule
$|\!\uparrow\uparrow\rangle$ & $1$ & $1$ & $\frac{1}{4}$ \\[2pt]
$\dfrac{|\!\uparrow\downarrow\rangle+|\!\downarrow\uparrow\rangle}{\sqrt{2}}$ & $0$ & $1$ & $\frac{1}{4}$ \\[8pt]
$|\!\downarrow\downarrow\rangle$ & $-1$ & $1$ & $\frac{1}{4}$ \\[2pt]
$\dfrac{|\!\uparrow\downarrow\rangle-|\!\downarrow\uparrow\rangle}{\sqrt{2}}$ & $0$ & $0$ & $-\frac{3}{4}$ \\
\bottomrule
\end{tabular}
\end{table}

我们的基底式~(1.72) 从另一个角度看也不令人满意：波函数对两个电子的交换必须是
反对称的。波函数是空间函数 $\psi_{\text{space}}(\mathbf{r}_1,\mathbf{r}_2)$ 与自旋
函数 $\chi$ 的乘积，$\chi$ 是式~(1.72) 所列态的线性组合。空间波函数对电子交换可以
是对称或反对称的。例如空间波函数
\begin{equation*}
\psi_{\text{space}}(\mathbf{r}_1,\mathbf{r}_2) = \frac{\phi(\mathbf{r}_1)\xi(\mathbf{r}_2)\pm\phi(\mathbf{r}_2)\xi(\mathbf{r}_1)}{\sqrt{2}}\tag{1.73}
\end{equation*}
依 $\pm$ 号而对电子交换对称（$+$）或反对称（$-$）。这类对称性称为交换对称性
（exchange symmetry）。式~(1.73) 中 $\phi(\mathbf{r}_i)$ 与 $\xi(\mathbf{r}_i)$ 是
第 $i$ 个电子的单粒子波函数。无论空间波函数的交换对称性如何，自旋波函数 $\chi$
必须具有相反的交换对称性：空间波函数对称时 $\chi$ 必须反对称，反之亦然——这样才能
保证乘积 $\psi_{\text{space}}(\mathbf{r}_1,\mathbf{r}_2)\times\chi$ 整体反对称。

$|\!\uparrow\uparrow\rangle$ 与 $|\!\downarrow\downarrow\rangle$ 这样的态在电子
交换下显然对称；但交换 $|\!\uparrow\downarrow\rangle$ 中的两个电子得到
$|\!\downarrow\uparrow\rangle$，它并不等于 $|\!\uparrow\downarrow\rangle$ 的倍数。
因此 $|\!\uparrow\downarrow\rangle$（同理 $|\!\downarrow\uparrow\rangle$）在两电子
交换下既非对称也非反对称。所以毫不奇怪，我们需要取这两个态的线性组合作为本征态。
组合列于表~1.1：$(|\!\uparrow\downarrow\rangle+|\!\downarrow\uparrow\rangle)/\sqrt{2}$
在电子交换下对称（与其他两个 $s=1$ 态一样），而
$(|\!\uparrow\downarrow\rangle-|\!\downarrow\uparrow\rangle)/\sqrt{2}$ 是反对称的。

对交换的这种不对称性的另一后果是泡利不相容原理（Pauli exclusion principle）：
两个电子不能处于同一量子态。若两个电子处于完全相同的空间与自旋量子态（比如都处于
空间态 $\phi(\mathbf{r})$ 且都自旋向上），则其自旋波函数必对电子交换对称，空间波
函数就必为反对称，于是
\begin{equation*}
\psi_{\text{space}}(\mathbf{r}_1,\mathbf{r}_2) = \frac{\phi(\mathbf{r}_1)\phi(\mathbf{r}_2)-\phi(\mathbf{r}_2)\phi(\mathbf{r}_1)}{\sqrt{2}} = 0.\tag{1.74}
\end{equation*}
态消失——这正说明两个电子不能处于同一量子态。

\begin{figure}
\centering
\includegraphics[width=0.45\textwidth]{fig1_09.pdf}
\caption{两个电子以 $A\mathbf{S}^{a}\cdot\mathbf{S}^{b}$ 型相互作用耦合，给出三重态
（$s=1$）与单态（$s=0$）。若 $A>0$，单态在下、三重态在上；三重态可被磁场 $\mathbf{B}$
分裂为三个分量。}
\end{figure}

我们经常会遇到两个自旋通过 $A\hat{\mathbf{S}}^{a}\cdot\hat{\mathbf{S}}^{b}$ 型
相互作用耦合的情形（$A$ 为常数）。若 $A>0$，较低的能级是单态（能量 $-3A/4$），
其上是三重态（能量 $A/4$），比单态高出 $A$（见图~1.9）。磁场可把三重态分裂成
$m_s$ 不同的三个态。若 $A<0$，则三重态是最低能级。

\section*{延伸阅读}

\begin{itemize}[itemsep=0.2em]
\item B.~I.~Bleaney 与 B.~Bleaney，《Electricity and Magnetism》，OUP 1989，
对电磁学有全面的论述（另见附录 B）。
\item A.~I.~Rae，《Introduction to Quantum Mechanics》，IOP Publishing 1992，
是量子力学入门层次的清晰阐述。
\item 量子角动量的精彩论述见《Feynman Lectures in Physics》第三卷第 1--3 章，
R.~P.~Feynman，Addison-Wesley 1975。
\item J.~J.~Sakurai，《Modern Quantum Mechanics》，第 2 版，Addison-Wesley 1994，
是极好的量子力学参考书。
\end{itemize}

\section*{习题}

\exer{1.1} 计算电子（取 $g=2$）的磁矩。在磁通密度 0.3~T 的磁场中，该电子的拉莫尔
进动频率是多少？若电子自旋分别平行、反平行于磁场，能量差是多少？把这个能量差换算
成频率。

\exer{1.2} 利用式 1.43 中自旋算符的定义，证明式 1.53 及对易关系式 1.54 与 1.55。

\exer{1.3} 利用式 1.57 中升降算符的定义，证明式 1.58 与 1.61。

\exer{1.4} 利用自旋的对易关系 $[\hat{S}_x,\hat{S}_y]=\mathrm{i}\hat{S}_z$（及其循环
置换），证明
\begin{equation*}
[\hat{\mathbf{S}}\cdot\mathbf{X}, \hat{\mathbf{S}}] = \mathrm{i}\hat{\mathbf{S}}\times\mathbf{X},\tag{1.75}
\end{equation*}
其中 $\mathbf{X}$ 为矢量。

\exer{1.5} 利用式 1.58 与 1.61，证明
\begin{equation*}
\hat{S}_{\pm}|S, S_z\rangle = \sqrt{S(S+1)-S_z(S_z\pm1)}\,|S, S_z\pm1\rangle,\tag{1.76}
\end{equation*}
其中 $|S,S_z\rangle$ 表示总自旋角动量为 $S(S+1)\hbar^{2}$、自旋角动量 $z$ 分量为
$S_z\hbar$ 的态。据此证明式 1.76 的如下特例：
\begin{equation*}
\hat{S}_{-}|S, S\rangle = \sqrt{2}S\,|S, S-1\rangle\tag{1.77}
\end{equation*}
\begin{equation*}
\hat{S}_{+}|S, S-1\rangle = \sqrt{2}S\,|S, S\rangle\tag{1.78}
\end{equation*}

\exer{1.6} 若磁场 $\mathbf{B}$ 在空间中均匀，证明这与取 $\mathbf{A}=\frac{1}{2}(\mathbf{B}\times\mathbf{r})$ 相一致，并证明 $\nabla\cdot\mathbf{A}=0$。是否还有其他 $\mathbf{A}$ 的选取能给出同样的 $\mathbf{B}$？

\exer{1.7} 电子的动能算符为 $\hat{p}^{2}/2m$。用式 1.41 证明它可以改写为
\begin{equation*}
\frac{(\boldsymbol{\sigma}\cdot\hat{\mathbf{p}})^{2}}{2m_{\mathrm{e}}}.\tag{1.79}
\end{equation*}
施加磁场时须把 $\hat{\mathbf{p}}$ 换成 $\hat{\mathbf{p}}+e\mathbf{A}$。借助式 1.40
证明：把这一替换代入式 1.79 得到如下形式的动能
\begin{equation*}
\frac{(\hat{\mathbf{p}}+e\mathbf{A})^{2}}{2m_{\mathrm{e}}} + g\mu_{\mathrm{B}}\mathbf{B}\cdot\mathbf{S}\tag{1.80}
\end{equation*}
此处 $g$ 因子为 $g=2$。（注意本题中须小心运用式 1.40 与 1.41，因为 $\hat{\mathbf{p}}$
是算符，与 $\mathbf{A}$ 不对易。）

\exer{1.8} 某原子的轨道角动量为零，自旋量子数为 $\frac{1}{2}$，处于 $|\!\uparrow_z\rangle$
态。沿与 $z$ 轴成 $\theta$ 角的方向测量其角动量。证明测得角动量平行于新轴的概率为
$\cos^{2}(\theta/2)$。

\exer{1.9} 以式 1.72 的基底可以构造诸如 $\hat{\mathbf{S}}^{a}\cdot\hat{\mathbf{S}}^{b}$
等算符的矩阵表示；注意例如 $\hat{S}_z^{a}$ 这类算符只作用于波函数中与第一个自旋
有关的部分。于是
\begin{equation*}
\hat{S}_z^{a} = \frac{1}{2}\begin{pmatrix} 1 & 0 & 0 & 0\\ 0 & 1 & 0 & 0\\ 0 & 0 & -1 & 0\\ 0 & 0 & 0 & -1 \end{pmatrix}\tag{1.81}
\end{equation*}
\begin{equation*}
\hat{S}_z^{b} = \frac{1}{2}\begin{pmatrix} 1 & 0 & 0 & 0\\ 0 & -1 & 0 & 0\\ 0 & 0 & 1 & 0\\ 0 & 0 & 0 & -1 \end{pmatrix}\tag{1.82}
\end{equation*}
构造 $\hat{S}_x^{a}$、$\hat{S}_x^{b}$、$\hat{S}_y^{a}$、$\hat{S}_y^{b}$ 的类似表示，
进而证明
\begin{equation*}
\hat{\mathbf{S}}^{a}\cdot\hat{\mathbf{S}}^{b} = \frac{1}{4}\begin{pmatrix} 1 & 0 & 0 & 0\\ 0 & -1 & 2 & 0\\ 0 & 2 & -1 & 0\\ 0 & 0 & 0 & 1 \end{pmatrix}.\tag{1.83}
\end{equation*}
求该算符的本征值与本征矢，并验证与表 1.1 一致。

\exer{1.10} 对球形样品（a）水与（b）$\mathrm{MnSO_4\cdot4H_2O}$ 施加 0.5~T 的磁场。
分别求样品内部的 H 与 B 场与自由空间值偏离的比例。（水与 $\mathrm{MnSO_4\cdot4H_2O}$
的磁化率见表 2.1。）你会发现修正确实非常小。

\exer{1.11} 证明算符
\begin{equation*}
\hat{S}_{\theta,\phi} = \sin\theta\cos\phi\,\hat{S}_x + \sin\theta\sin\phi\,\hat{S}_y + \cos\theta\,\hat{S}_z,\tag{1.84}
\end{equation*}
（它表示自旋沿由球极角 $\theta$、$\phi$ 所确定方向的分量算符）本征值为
$\pm\frac{1}{2}$，本征态形如
\begin{align*}
|\!\uparrow\rangle &= \begin{pmatrix}\cos(\theta/2)\\ \sin(\theta/2)e^{i\phi}\end{pmatrix}\tag{1.85}\\
|\!\downarrow\rangle &= \begin{pmatrix}\sin(\theta/2)\\ -\cos(\theta/2)e^{i\phi}\end{pmatrix}.\tag{1.86}
\end{align*}
验证这些结果与图 1.8 的黎曼球面表示一致。再证明
\begin{equation*}
\hat{S}_{\theta,\phi}^{2} = \frac{1}{4}\begin{pmatrix}1 & 0\\ 0 & 1\end{pmatrix}.\tag{1.87}
\end{equation*}

\exer{1.12} 磁场沿 $z$ 方向时电子的哈密顿量（能量）算符为
\begin{equation*}
\hat{\mathcal{H}} = g\mu_{\mathrm{B}}\mathbf{B}\cdot\hat{\mathbf{S}} = g\mu_{\mathrm{B}}B\hat{S}_z.\tag{1.88}
\end{equation*}
含时薛定谔方程为
\begin{equation*}
\hat{\mathcal{H}}\psi(t) = \mathrm{i}\hbar\frac{\mathrm{d}\psi(t)}{\mathrm{d}t}\tag{1.89}
\end{equation*}
于是
\begin{equation*}
\psi(t) = \exp(-\mathrm{i}\hat{\mathcal{H}}t/\hbar)\,\psi(0).\tag{1.90}
\end{equation*}
利用式 1.41 证明
\begin{equation*}
\exp(\mathrm{i}\alpha\sigma_m) = I\cos\alpha + \mathrm{i}\sigma_m\sin\alpha\tag{1.91}
\end{equation*}
其中 $I$ 是单位矩阵，$\sigma_m$ 是某个泡利自旋矩阵，$\alpha$ 为实数。进而证明：
若把 $\psi(t)$ 写成旋量，
\begin{equation*}
\psi(t) = \begin{pmatrix} \exp(-\mathrm{i}g\mu_{\mathrm{B}}Bt/2\hbar) & 0\\ 0 & \exp(\mathrm{i}g\mu_{\mathrm{B}}Bt/2\hbar) \end{pmatrix}\psi(0),\tag{1.92}
\end{equation*}
并利用上一题的结果证明：自旋态如此演化，使 $\theta$ 的期望值不变而 $\phi$ 以
角频率 $geB/2m$ 旋转。这表明拉莫尔进动现象也可以从量子力学处理中导出。

\exer{1.13} 这是推导自旋进动的另一途径。从式 1.88 出发，用式 C.7 证明
\begin{align*}
\frac{\mathrm{d}}{\mathrm{d}t}\langle\hat{\mathbf{S}}\rangle &= \frac{1}{\mathrm{i}\hbar}\langle[\hat{\mathbf{S}}, \hat{\mathcal{H}}]\rangle\tag{1.93}\\
&= -\frac{g\mu_{\mathrm{B}}}{\hbar}\langle\hat{\mathbf{S}}\rangle\times\mathbf{B},\tag{1.94}
\end{align*}
它与式 1.6 类似，且
\begin{equation*}
\gamma = -\frac{g\mu_{\mathrm{B}}}{\hbar} = -\frac{ge}{2m_{\mathrm{e}}}.\tag{1.95}
\end{equation*}
负号来自电子带负电。

\exer{1.14} 本题讨论电偶极子的相应情形。（a）把电偶极矩为 $\mathbf{p}$、转动惯量
为 $I$ 的电偶极子置于电场 $\boldsymbol{\mathcal{E}}$ 中。经典地证明：$\mathbf{p}$ 与
$\boldsymbol{\mathcal{E}}$ 之间的夹角 $\theta$ 满足微分方程
\begin{equation*}
\ddot{\theta} = -\frac{p\mathcal{E}\sin\theta}{I}.\tag{1.96}
\end{equation*}
并证明 $\theta$ 很小时该方程给出简谐运动。

（b）用量子力学重做此题。考虑哈密顿量
\begin{equation*}
\hat{\mathcal{H}} = -\frac{\hbar^{2}}{2I}\frac{\partial^{2}}{\partial\theta^{2}} - p\mathcal{E}\cos\theta,\tag{1.97}
\end{equation*}
并说明为什么它是合适的哈密顿量。用式 C.7 证明
\begin{equation*}
\frac{\mathrm{d}}{\mathrm{d}t}\langle\hat{\theta}\rangle = \frac{\langle\hat{L}\rangle}{I}\tag{1.98}
\end{equation*}
其中 $\hat{L}=-\mathrm{i}\hbar\,\partial/\partial\theta$，且
\begin{equation*}
\frac{\mathrm{d}}{\mathrm{d}t}\langle\hat{L}\rangle = -p\mathcal{E}\langle\sin\theta\rangle.\tag{1.99}
\end{equation*}
进而推出
\begin{equation*}
\frac{\mathrm{d}^{2}}{\mathrm{d}t^{2}}\langle\hat{\theta}\rangle = -\frac{p\mathcal{E}}{I}\langle\sin\theta\rangle,\tag{1.100}
\end{equation*}
在适当的极限下它退化为经典表达式。把这与磁偶极子的情形作比较：二者为何不同？
为何电情形不出现自旋进动？

我们已经看到：电场中的电偶极子在电场平面内来回振荡，而磁场中的磁偶极子绕磁场
进动。对每一种情形，我们熟悉的“施加场（电场或磁场），偶极子（电的或磁的）就会
沿场排列”的观念有什么问题？

% 第2章 孤立磁矩（Isolated magnetic moments）
\chapter{孤立磁矩}

本章考察孤立磁矩的性质。在此阶段，我们忽略不同原子的磁矩之间的相互作用，也忽略
磁矩与其近邻环境之间的相互作用，剩下的就只是孤立原子的物理，以及它们与外加磁场
的相互作用。这当然不意味着问题变简单了，但复杂性的来源是单个原子内电子的组合
方式，而不是凝聚态中原子数目巨大这一事实。经过这一简化，大量的原子只是使磁化率
这类性质多出一个因子 $n$——单位体积内的原子数。

\section{磁场中的原子}

1.1 节（式 1.35）已经证明：沿 $z$ 轴的磁场中，电子自旋的能量为
\begin{equation*}
E = g\mu_{\mathrm{B}}Bm_s\tag{2.1}
\end{equation*}
其中 $g\approx2$，$m_s=\pm\frac{1}{2}$，故 $E\approx\pm\mu_{\mathrm{B}}B$。除自旋
角动量外，原子中的电子还具有轨道角动量。若原子中第 $i$ 个电子的位置为 $\mathbf{r}_i$、
动量为 $\mathbf{p}_i$，则总角动量 $\hbar\mathbf{L}$ 为
\begin{equation*}
\hbar\mathbf{L} = \sum_i \mathbf{r}_i\times\mathbf{p}_i\tag{2.2}
\end{equation*}
求和遍及原子中所有电子。考虑哈密顿量为 $\hat{\mathcal{H}}_0$ 的原子：
\begin{equation*}
\hat{\mathcal{H}}_0 = \sum_{i=1}^{Z}\left(\frac{p_i^{2}}{2m} + V_i\right)\tag{2.3}
\end{equation*}
它是对原子中 $Z$ 个电子求和的电子动能（第 $i$ 个电子的 $p_i^{2}/2m_{\mathrm{e}}$）
与势能（第 $i$ 个电子的 $V_i$）之和。设 $\hat{\mathcal{H}}_0$ 的本征态与本征值
均为已知。

现在加入磁场
\begin{equation*}
\mathbf{B} = \nabla\times\mathbf{A}\tag{2.4}
\end{equation*}
其中 $\mathbf{A}$ 是磁矢势。我们选取规范\fn{1}{式 2.4 把 B 与 A 联系起来。但对给定的磁场 B，磁矢势 A 并非唯一确定：可以给 A 加上某个标量势的梯度而仍得到同样的 B。我们对 A 所作的这种选择称为规范选择。}使得
\begin{equation*}
\mathbf{A}(\mathbf{r}) = \frac{\mathbf{B}\times\mathbf{r}}{2}.\tag{2.5}
\end{equation*}
于是动能须按 1.2 节的方案修改。由于电子电荷为 $-e$，动能为
$[\mathbf{p}_i+e\mathbf{A}(\mathbf{r}_i)]^{2}/2m_{\mathrm{e}}$，受扰动的哈密顿量
应写成
\begin{align*}
\hat{\mathcal{H}} &= \sum_{i=1}^{Z}\left(\frac{[\mathbf{p}_i+e\mathbf{A}(\mathbf{r}_i)]^{2}}{2m_{\mathrm{e}}} + V_i\right) + g\mu_{\mathrm{B}}\mathbf{B}\cdot\mathbf{S}\tag{2.6}\\
&= \sum_i\left(\frac{p_i^{2}}{2m_{\mathrm{e}}} + V_i\right) + \mu_{\mathrm{B}}(\mathbf{L}+g\mathbf{S})\cdot\mathbf{B} + \frac{e^{2}}{8m_{\mathrm{e}}}\sum_i(\mathbf{B}\times\mathbf{r}_i)^{2}\tag{2.7}\\
&= \hat{\mathcal{H}}_0 + \mu_{\mathrm{B}}(\mathbf{L}+g\mathbf{S})\cdot\mathbf{B} + \frac{e^{2}}{8m_{\mathrm{e}}}\sum_i(\mathbf{B}\times\mathbf{r}_i)^{2}\tag{2.8}
\end{align*}
对原哈密顿量 $\hat{\mathcal{H}}_0$ 的主要微扰通常是 $\mu_{\mathrm{B}}(\mathbf{L}+g\mathbf{S})\cdot\mathbf{B}$
项，但正如将要看到的，它有时会消失。这是原子自身磁矩的效应，称为顺磁项（paramagnetic
term）。第三项 $(e^{2}/8m_{\mathrm{e}})\sum_i(\mathbf{B}\times\mathbf{r}_i)^{2}$ 来自
抗磁矩。这两项贡献将在 2.3 节（抗磁性）与 2.4 节（顺磁性）中详细讨论。下一节概述
需要解释的效应。

\section{磁化率}

如 1.1.4 节所示，对线性材料 $M=\chi H$，$M$ 是单位体积的磁矩（磁化强度），$\chi$
是磁化率（无量纲）。注意 $M$ 的定义意味着 $\chi$ 表示单位体积中由磁场 H 诱导的
磁矩。磁化率常以摩尔磁化率（molar magnetic susceptibility）$\chi_{\mathrm{m}}$ 列表：
\begin{equation*}
\chi_{\mathrm{m}} = \chi V_{\mathrm{m}}.\tag{2.9}
\end{equation*}
式中 $V_{\mathrm{m}}$ 是摩尔体积，即 1 摩尔（$6.022\times10^{23}$ 个化学式单位）
物质占据的体积。摩尔体积（单位 m$^3$）等于物质的相对原子质量\fn{2}{相对原子质量是 1 摩尔的质量。注意相对原子质量通常以克为单位列表。}（单位 kg）除以密度
$\rho$（单位 kg\,m$^{-3}$）。质量磁化率 $\chi_g$ 定义为
\begin{equation*}
\chi_g = \frac{\chi}{\rho},\tag{2.10}
\end{equation*}
单位为 m$^3$\,kg$^{-1}$。若干物质的磁化率数值列于表 2.1。磁化率为负者以抗磁性为主，
为正者以顺磁性为主。

\begin{table}
\centering
\caption{若干物质在 298 K 下的磁化率 $\chi$ 与摩尔磁化率 $\chi_{\mathrm{m}}$。水、苯和 NaCl
是弱抗磁体（磁化率为负）；$\mathrm{CuSO_4\cdot5H_2O}$、$\mathrm{MnSO_4\cdot4H_2O}$、Al
和 Na 是顺磁体（磁化率为正）。}
\begin{tabular}{lcc}
\toprule
 & $\chi/10^{-6}$ & $\chi_{\mathrm{m}}/10^{-10}$（m$^3$\,mol$^{-1}$）\\
\midrule
水 & $-90$ & $-16.0$\\
苯 & $-7.2$ & $-6.4$\\
NaCl & $-13.9$ & $-3.75$\\
石墨（$\parallel$） & $-260$ & $-31$\\
石墨（$\perp$） & $-3.8$ & $-4.6$\\
Cu & $-1.1$ & $-0.078$\\
Ag & $-2.4$ & $-0.25$\\
\midrule
$\mathrm{CuSO_4\cdot5H_2O}$ & $176$ & $192$\\
$\mathrm{MnSO_4\cdot4H_2O}$ & $2640$ & $2.79\times10^{3}$\\
Al & $22$ & $2.2$\\
Na & $7.3$ & $1.7$\\
\bottomrule
\end{tabular}
\end{table}

周期表前 60 个元素的磁化率绘于图 2.1。其中一些为负，表明 2.3 节将讨论的抗磁性
占主导；另一些为正，表明存在顺磁性，这将在 2.4 节讨论。

\begin{figure}
\centering
\includegraphics[width=0.75\textwidth]{fig2_01.pdf}
\caption{周期表前 60 个元素在室温下的质量磁化率，按原子序数排列。Fe、Co、Ni 是铁磁体，
在无外加磁场时也有自发磁化。}
\end{figure}

\section{抗磁性}

所有材料都表现出一定程度的抗磁性（diamagnetism）\fn{3}{前缀 dia- 意为“相对”“横跨”（英语的 diagonal（对角）、diameter（直径）等词即由此而来）。}——一种弱的、
负的磁化率。对抗磁物质，磁场诱导出的磁矩与产生它的外加磁场方向相反。

这一效应常被经典地讨论：磁场对电子轨道运动的作用产生反电动势\fn{4}{electromotive force（电动势）。}，由楞次定律，它反抗产生自己的磁场。然而，
上一章的玻尔--范列文定理提醒我们警惕这类论证——它们试图证明经典体系中磁场可以
诱导磁矩。\fn{5}{延伸阅读中有进一步讨论。}抗磁性是完全量子力学的现象，必须按量子
力学处理。

用量子力学方法很容易说明这一效应。考虑没有未满电子壳层的原子，此时式 2.8 中的
顺磁项可以忽略。若 $\mathbf{B}$ 平行于 $z$ 轴，则 $\mathbf{B}\times\mathbf{r}_i = B(-y_i, x_i, 0)$，且
\begin{equation*}
(\mathbf{B}\times\mathbf{r}_i)^{2} = B^{2}(x_i^{2}+y_i^{2})\tag{2.11}
\end{equation*}
于是抗磁项引起的基态能量一阶移动为
\begin{equation*}
\Delta E_0 = \frac{e^{2}B^{2}}{8m_{\mathrm{e}}}\sum_{i=1}^{Z}\langle 0|(x_i^{2}+y_i^{2})|0\rangle,\tag{2.12}
\end{equation*}
$|0\rangle$ 是基态波函数。若假定原子球对称\fn{6}{若总角动量 J 为零，这是一个良好的近似。}，则 $\langle x_i^{2}\rangle=\langle y_i^{2}\rangle=\frac{1}{3}\langle r_i^{2}\rangle$，可得
\begin{equation*}
\Delta E_0 = \frac{e^{2}B^{2}}{12m_{\mathrm{e}}}\sum_{i=1}^{Z}\langle 0|r_i^{2}|0\rangle.\tag{2.13}
\end{equation*}

考虑体积 $V$ 内由 $N$ 个离子（每个有 $Z$ 个质量为 $m$ 的电子）组成、所有壳层均填满
的固体。求 $T=0$ 时的磁化强度可以循附录 E 的做法：
\begin{equation*}
M = -\frac{\partial F}{\partial B} = -\frac{N}{V}\frac{\partial\Delta E_0}{\partial B} = -\frac{Ne^{2}B}{6m_{\mathrm{e}}}\sum_{i=1}^{Z}\langle r_i^{2}\rangle,\tag{2.14}
\end{equation*}
$F$ 是亥姆霍兹自由能。由此可提取抗磁磁化率 $\chi=M/H\approx\mu_0M/B$（设 $\chi\ll1$）：
\begin{equation*}
\chi = -\frac{N}{V}\frac{e^{2}\mu_0}{6m_{\mathrm{e}}}\sum_{i=1}^{Z}\langle r_i^{2}\rangle.\tag{2.15}
\end{equation*}
此式使用了一阶微扰论（二阶项将在 2.4.4 节考虑）。温度自零上升后，基态以上的态对
抗磁磁化率的贡献逐渐变大，但这是边缘效应。抗磁磁化率通常基本不依赖温度。

把实验测得的多种离子的抗磁摩尔磁化率对 $Z_{\mathrm{eff}}r^{2}$ 作图，可以相当粗糙地
检验上式，其中 $Z_{\mathrm{eff}}$ 是离子外壳层中的电子数\fn{7}{对离子而言这一数值不同于原子序数 Z，故用符号 $Z_{\mathrm{eff}}$ 表示“有效”原子序数；我们忽略了内壳层电子。}，$r$ 是实测离子半径。此近似假定离子外壳层的所有电子的
$\langle r_i\rangle^{2}$ 大致相同，即
\begin{equation*}
\sum_{i=1}^{Z_{\mathrm{eff}}}\langle r_i^{2}\rangle \approx Z_{\mathrm{eff}}r^{2}.\tag{2.16}
\end{equation*}

若干离子的抗磁磁化率示于图 2.2。实验值由比较一系列离子盐（NaF、NaCl、NaBr、KCl、
KBr、…）的抗磁磁化率推得。这一做法并不精确——离子中并非所有电子的平均半径平方都
相同（故式 2.16 绝非严格），但符合程度仍相当可观。之所以选离子，是因为例如 Na 和
Cl 原子有未成对电子，而 $\mathrm{Na^{+}}$ 与 $\mathrm{Cl^{-}}$ 离子都是闭壳层结构，
与 Ne、Ar 类似（可参照下文图 2.13 的周期表）。于是本会主导原子磁响应的顺磁效应
在离子中可以忽略。

\begin{figure}
\centering
\includegraphics[width=0.6\textwidth]{fig2_02.pdf}
\caption{若干离子的抗磁摩尔磁化率实验值 $\chi_m$ 对 $Z_{\mathrm{eff}}r^{2}$ 作图，
$Z_{\mathrm{eff}}$ 为离子外层电子数，$r$ 为实测离子半径。}
\end{figure}

在具有离域 $\pi$ 电子的分子（如萘和石墨）中，可以观察到较大且各向异性的抗磁磁化率。
萘由两个苯环沿一边并合而成（图 2.3(a)）。$\pi$ 电子非常活跃，诱导电流可以沿环边缘
流动，产生大的抗磁磁化率；当磁场垂直于环平面时该效应最大。

\begin{figure}
\centering
\includegraphics[width=0.45\textwidth]{fig2_03.pdf}
\caption{（a）萘由两个稠合的苯环组成。（b）石墨由六角层堆叠而成，碳原子画成黑点。
碳原子在隔层（而非相邻层）之间对齐（见竖直虚线）。}
\end{figure}

有效环直径是原子直径的数倍，故效应很大。石墨同理：它由结合疏松的六角层构成
（图 2.3(b)）。磁场垂直于层面时的抗磁磁化率远大于平行方向。

抗磁性存在于一切材料之中，但它是弱效应：要么可以忽略，要么只是对更强效应的小
修正。

\section{顺磁性}

顺磁性（paramagnetism）\fn{8}{前缀 para- 意为“伴随”“沿着”（英语的 parallel（平行）等词即由此而来）。}对应正的磁化率：外加磁场诱导的磁化强度与磁场平行。
上一节考虑的是不含未成对电子的材料，其原子或分子在未加场时没有磁矩。本节关心
的是因未成对电子而具有非零磁矩的原子。无外加磁场时，这些磁矩指向随机——相邻原子
的磁矩之间相互作用极弱，可视为彼此独立。加上磁场后磁矩被排齐，排齐的程度（因而
诱导的磁化强度）取决于磁场强度。

原子磁矩与原子的总角动量 $\mathbf{J}$ 相联系，$\mathbf{J}$ 是轨道角动量 $\mathbf{L}$
与自旋角动量 $\mathbf{S}$ 之和：
\begin{equation*}
\mathbf{J} = \mathbf{L} + \mathbf{S}.\tag{2.17}
\end{equation*}
与全书一致，这些量都以 $\hbar$ 为单位度量。自旋与轨道部分如何组合将在后面几节
详述；本节只假定每个原子具有大小为 $\mu$ 的磁矩。

增强磁场会使自旋排齐，升温则使它们随机化，因此可以预期顺磁材料的磁化强度依赖于
比值 $B/T$。顺磁效应一般远强于抗磁效应，尽管抗磁性总是作为弱的负贡献存在。

\subsection{顺磁性的半经典处理}

我们先半经典地处理顺磁性（下面会看到这对应 $J=\infty$），忽略量子化导致的“磁矩
只能指向某些方向”这一事实。考虑与外场 $\mathbf{B}$（不失一般性取沿 $z$ 方向）成
$\theta$ 到 $\theta+\mathrm{d}\theta$ 角的磁矩，其能量为 $-\mu B\cos\theta$，沿
$\mathbf{B}$ 的净磁矩为 $\mu\cos\theta$。若磁矩可以完全随机地选择方向，则取角介于
$\theta$ 与 $\theta+\mathrm{d}\theta$ 之间的比例正比于图 2.4 中环带的面积——单位
半径球面上为 $2\pi\sin\theta\,\mathrm{d}\theta$。单位球总面积为 $4\pi$，故该比例
为 $\frac{1}{2}\sin\theta\,\mathrm{d}\theta$。温度 $T$ 下磁矩取角介于 $\theta$ 与
$\theta+\mathrm{d}\theta$ 的概率正比于统计因子 $\frac{1}{2}\sin\theta\,\mathrm{d}\theta$
与玻尔兹曼因子 $\exp(\mu B\cos\theta/k_{\mathrm{B}}T)$ 之积（$k_{\mathrm{B}}$ 为玻尔
兹曼常数）。于是沿 $\mathbf{B}$ 的平均磁矩为
\bio{Paul Langevin\\（1872--1946）}
\begin{align*}
\langle\mu_z\rangle &= \frac{\int_0^{\pi}\mu\cos\theta\,\exp(\mu B\cos\theta/k_{\mathrm{B}}T)\,\frac{1}{2}\sin\theta\,\mathrm{d}\theta}{\int_0^{\pi}\exp(\mu B\cos\theta/k_{\mathrm{B}}T)\,\frac{1}{2}\sin\theta\,\mathrm{d}\theta}\tag{2.18}\\
&= \mu\,\frac{\int_{-1}^{1}x e^{yx}\,\mathrm{d}x}{\int_{-1}^{1}e^{yx}\,\mathrm{d}x},\tag{2.19}
\end{align*}
其中定义 $y=\mu B/k_{\mathrm{B}}T$，$x=\cos\theta$。由此得
\begin{equation*}
\frac{\langle\mu_z\rangle}{\mu} = \coth y - \frac{1}{y} \equiv L(y)\tag{2.20}
\end{equation*}
$L(y)=\coth y-1/y$ 称为朗之万函数（Langevin function），绘于图 2.5。$y$ 小时
\begin{equation*}
\coth(y) = \frac{1}{y} + \frac{y}{3} + O(y^{3})\tag{2.21}
\end{equation*}
于是
\begin{equation*}
L(y) = \frac{y}{3} + O(y^{3}).\tag{2.22}
\end{equation*}

\begin{figure}
\centering
\includegraphics[width=0.42\textwidth]{fig2_04.pdf}
\caption{计算顺磁材料的平均磁矩：考虑磁矩与 $z$ 轴夹角介于 $\theta$ 与
$\theta+\mathrm{d}\theta$ 的概率，它正比于单位球上阴影环带的面积
$2\pi\sin\theta\,\mathrm{d}\theta$。}
\end{figure}

\begin{figure}
\centering
\includegraphics[width=0.5\textwidth]{fig2_05.pdf}
\caption{经典顺磁体的磁化强度由朗之万函数 $L(y)=\coth y-\frac{1}{y}$ 描述。$y$ 小时
$L(y)\approx y/3$，如图中在原点附近与曲线相切的直线所示。磁场增强或温度降低时，
磁化强度的数值增大。}
\end{figure}

用 $n$ 表示单位体积的磁矩数。饱和磁化强度 $M_{\mathrm{s}}$ 是所有磁矩完全排齐时
能得到的最大磁化强度，$M_{\mathrm{s}}=n\mu$。实际得到的磁化强度为
$M=n\langle\mu_z\rangle$，磁化强度与饱和磁化强度之比是有用的量：
\begin{equation*}
\frac{M}{M_{\mathrm{s}}} = \frac{\langle\mu_z\rangle}{\mu} \approx \frac{y}{3} = \frac{\mu B}{3k_{\mathrm{B}}T}\tag{2.23}
\end{equation*}
再利用小场下成立的 $\chi=M/H\approx\mu_0M/B$，\fn{9}{小场时 $\chi\ll1$，故 $B\approx\mu_0H$。}得
\begin{equation*}
\chi = \frac{n\mu_0\mu^{2}}{3k_{\mathrm{B}}T}.\tag{2.24}
\end{equation*}
这表明磁化率与温度成反比，即居里定律（Curie's law，以其发现者 Pierre Curie 命名）。
\fn{10}{人们常写 $\chi=C_{\text{Curie}}/T$，$C_{\text{Curie}}$ 为居里常数。}

\subsection{$J=\frac{1}{2}$ 的顺磁性}

现在用量子系统重做上述计算：经典磁矩换成 $J=\frac{1}{2}$ 的量子自旋。磁矩 $z$
分量只有两个可能值 $m_J=\pm\frac{1}{2}$：或平行、或反平行于 $\mathbf{B}$。磁矩为
$\mp\mu_{\mathrm{B}}$（取 $g=2$），相应能量 $\pm\mu_{\mathrm{B}}B$（两个解示于
图 2.6）。于是
\begin{align*}
\langle g\mu_{\mathrm{B}}m_J\rangle &= \frac{-\mu_{\mathrm{B}}e^{\mu_{\mathrm{B}}B/k_{\mathrm{B}}T}+\mu_{\mathrm{B}}e^{-\mu_{\mathrm{B}}B/k_{\mathrm{B}}T}}{e^{\mu_{\mathrm{B}}B/k_{\mathrm{B}}T}+e^{-\mu_{\mathrm{B}}B/k_{\mathrm{B}}T}}\tag{2.25}\\
&= \mu_{\mathrm{B}}\tanh\left(\frac{\mu_{\mathrm{B}}B}{k_{\mathrm{B}}T}\right)\tag{2.26}
\end{align*}
写 $y=\mu_{\mathrm{B}}B/k_{\mathrm{B}}T=g\mu_{\mathrm{B}}JB/k_{\mathrm{B}}T$（此处
$J=\frac{1}{2}$，$g=2$），便有
\begin{equation*}
\frac{M}{M_{\mathrm{s}}} = \frac{\langle m_J\rangle}{J} = \tanh y.\tag{2.27}
\end{equation*}
这一函数不同于朗之万函数，但形状相当接近（见图 2.7）。小场下
$\tanh(\mu_{\mathrm{B}}B/k_{\mathrm{B}}T)\approx\mu_{\mathrm{B}}B/k_{\mathrm{B}}T$，
且
\begin{equation*}
\chi = \frac{n\mu_0\mu_{\mathrm{B}}^{2}}{k_{\mathrm{B}}T}.\tag{2.28}
\end{equation*}

式 2.27 也可以用另一种方法非常高效地导出。配分函数 $Z$ 是按简并度加权的玻尔兹曼
概率之和。单个自旋的配分函数为
\begin{equation*}
Z = \mathrm{e}^{\mu_{\mathrm{B}}B/k_{\mathrm{B}}T} + \mathrm{e}^{-\mu_{\mathrm{B}}B/k_{\mathrm{B}}T} = 2\cosh\left(\frac{\mu_{\mathrm{B}}B}{k_{\mathrm{B}}T}\right),\tag{2.29}
\end{equation*}
\mnote{关于 Z、F 以及 $M=-(\partial F/\partial B)_T$ 这类表达式，详见附录 E。}
由 $F=-k_{\mathrm{B}}T\ln Z$ 可得亥姆霍兹自由能；对单位体积 $n$ 个自旋，
\begin{equation*}
F = nk_{\mathrm{B}}T\ln\left[2\cosh\left(\frac{\mu_{\mathrm{B}}B}{k_{\mathrm{B}}T}\right)\right].\tag{2.30}
\end{equation*}
磁化强度由 $M=-(\partial F/\partial B)_T$ 给出，同样得到
\begin{equation*}
\frac{M}{M_{\mathrm{s}}} = \tanh\left(\frac{\mu_{\mathrm{B}}B}{k_{\mathrm{B}}T}\right),\tag{2.31}
\end{equation*}
与式 2.27 一致。

\begin{figure}
\centering
\includegraphics[width=0.42\textwidth]{fig2_06.pdf}
\caption{自旋 $\frac{1}{2}$ 磁矩的能量随磁场的变化。}
\end{figure}

\begin{figure}
\centering
\includegraphics[width=0.5\textwidth]{fig2_07.pdf}
\caption{自旋 $\frac{1}{2}$ 顺磁体的磁化强度遵循 $\tanh y$；$y$ 小时
$\tanh y\approx y$，如图中在原点附近与曲线相切的直线所示。}
\end{figure}

这一方法还可用于导出该模型的其他热力学量（见习题 2.4），结果按 $k_{\mathrm{B}}T/\mu_{\mathrm{B}}B$
绘于图 2.8。图 2.8(a) 与图 2.7 信息相同，只是横轴反向——因为要理解材料的热性质，
我们真正关心的是固定磁场下升温的效果。样品升温时磁矩随机化、磁化强度下降，但同时
能量密度 $E=-M_{\mathrm{s}}B$ 上升（见图 2.8(b)）。$T\to\infty$ 时能量为零：此时
磁矩相对外场完全随机，能量得失相抵。降温对应能量降低（这一点在 2.6 节还会提到）。

热容 $C=(\partial E/\partial T)_B$ 在 $k_{\mathrm{B}}T\sim\mu_{\mathrm{B}}B$ 附近
有一个宽峰，称为肖特基反常（Schottky anomaly，见图 2.8(c)）。其成因是：在这个
温度附近，热运动恰好能够激发系统两个态之间的跃迁。极低温下系统难以改变能量——
没有足够能量激发离开基态的跃迁，所有自旋被“冻住”，都与磁场排齐；极高温下系统也
难以改变能量——两个态被同等占据。两者之间出现极大。因此热容峰可以有效地提示
“有有趣的事情正在发生”。但要注意：肖特基反常并不是相变对应的那种尖峰，而是一个
平滑、宽缓的极大。

熵 $S=-(\partial F/\partial T)_B$ 随温度上升而增大（见图 2.8(d)），符合预期——
它反映自旋的无序程度。反之，降温对应有序化、熵减少。这一事实对磁制冷技术非常有用，
将在 2.6 节描述。

下一节考虑总角动量量子数为任意 $J$ 的一般顺磁体：前面经典与量子两种情形都作为
特例包含在内。

\begin{figure}
\centering
\includegraphics[width=0.42\textwidth]{fig2_08.pdf}
\caption{单位体积含 $n$ 个无相互作用自旋 $\frac{1}{2}$ 离子的顺磁盐的（a）磁化强度
$M$（以饱和磁化强度归一）、（b）能量 $E$、（c）热容 $C$（恒定外磁场下）与
（d）熵 $S$ 随 $k_{\mathrm{B}}T/\mu_{\mathrm{B}}B$ 的变化。$E$、$C$、$S$ 按单位体积
顺磁盐给出。}
\end{figure}

\subsection{布里渊函数}

现在推导 $J$ 取任意整数或半整数的一般情形。前面各情形（$J=\frac{1}{2}$ 与
$J=\infty$）的许多特征在这里依然成立，例如磁场使磁矩趋于排齐、温度使它们趋于
无序。配分函数为
\begin{equation*}
Z = \sum_{m_J=-J}^{J}\exp(m_J g_J\mu_{\mathrm{B}}B/k_{\mathrm{B}}T).\tag{2.32}
\end{equation*}
写 $x=g_J\mu_{\mathrm{B}}B/k_{\mathrm{B}}T$，有
\begin{equation*}
\langle m_J\rangle = \frac{\sum_{m_J=-J}^{J}m_J\,\mathrm{e}^{m_Jx}}{\sum_{m_J=-J}^{J}\mathrm{e}^{m_Jx}} = \frac{1}{Z}\frac{\partial Z}{\partial x},\tag{2.33}
\end{equation*}
于是
\begin{equation*}
M = ng_J\mu_{\mathrm{B}}\langle m_J\rangle = \frac{ng_J\mu_{\mathrm{B}}}{Z}\frac{\partial Z}{\partial B}\frac{\partial B}{\partial x} = nk_{\mathrm{B}}T\frac{\partial\ln Z}{\partial B}.\tag{2.34}
\end{equation*}
配分函数 $Z$ 是首项 $a=\mathrm{e}^{-Jx}$、公比 $r=\mathrm{e}^{x}$ 的几何级数，
可用熟知的公式求和：
\begin{equation*}
a+ar+ar^{2}+\cdots+ar^{M-1} = \sum_{j=1}^{M}ar^{j-1} = \frac{a(1-r^{M})}{1-r},\tag{2.35}
\end{equation*}
$M$ 是级数项数，此处 $M=2J+1$。稍加整理得
\begin{equation*}
Z = \frac{\sinh[(2J+1)\frac{x}{2}]}{\sinh[\frac{x}{2}]}\tag{2.36}
\end{equation*}
作代换
\begin{equation*}
y = xJ = g_J\mu_{\mathrm{B}}JB/k_{\mathrm{B}}T,\tag{2.37}
\end{equation*}
得
\begin{equation*}
M = M_{\mathrm{s}}\,\mathrm{B}_J(y)\tag{2.38}
\end{equation*}
其中饱和磁化强度
\begin{equation*}
M_{\mathrm{s}} = ng_J\mu_{\mathrm{B}}J\tag{2.39}
\end{equation*}
而 $\mathrm{B}_J(y)$ 是布里渊函数（Brillouin function）：
\bio{Léon Brillouin\\（1889--1969）}
\begin{equation*}
\mathrm{B}_J(y) = \frac{2J+1}{2J}\coth\left(\frac{2J+1}{2J}y\right) - \frac{1}{2J}\coth\frac{y}{2J}.\tag{2.40}
\end{equation*}
该函数对不同 $J$ 绘于图 2.9。布里渊函数具有恰当的极限：$J\to\infty$ 时退化为朗之万
函数
\begin{equation*}
\mathrm{B}_{\infty}(y) = L(y),\tag{2.41}
\end{equation*}
$J=\frac{1}{2}$ 时退化为 $\tanh$ 函数
\begin{equation*}
\mathrm{B}_{1/2}(y) = \tanh(y).\tag{2.42}
\end{equation*}
即回到前几节讨论过的两种情形。

$y$ 的典型量级可估计如下：取 $J=\frac{1}{2}$、$g_J=2$、$B=1$~T，室温下
$y\sim2\times10^{-3}$。因此除极低温和/或极强磁场外，实验情形都对应 $y\ll1$（因而
$\chi\ll1$）。$y$ 小时利用 $\coth y$ 的麦克劳林展开可得
\begin{equation*}
\mathrm{B}_J(y) = \frac{(J+1)y}{3J} + O(y^{3}).\tag{2.43}
\end{equation*}
低场下磁化率为
\begin{equation*}
\chi = \frac{M}{H}\approx\frac{\mu_0M}{B} = \frac{n\mu_0\mu_{\mathrm{eff}}^{2}}{3k_{\mathrm{B}}T}\tag{2.44}
\end{equation*}
形如经典的居里定律。\fn{11}{即 $\chi=C_{\text{Curie}}/T$，其中 $C_{\text{Curie}}=n\mu_0g_J^{2}J(J+1)/3k_{\mathrm{B}}$ 为居里常数。}于是测量 $\chi$ 便可推定有效磁矩
$\mu_{\mathrm{eff}}$：
\begin{equation*}
\mu_{\mathrm{eff}} = g_J\mu_{\mathrm{B}}\sqrt{J(J+1)}\tag{2.45}
\end{equation*}
其中
\begin{equation*}
g_J = \frac{3}{2} + \frac{S(S+1)-L(L+1)}{2J(J+1)}\tag{2.46}
\end{equation*}
常数 $g_J$ 称为 Landé $g$ 值（见附录 C）。
\bio{Alfred Landé\\（1888--1975）}

\begin{figure}
\centering
\includegraphics[width=0.55\textwidth]{fig2_09.pdf}
\caption{磁矩量子数为 $J$ 的顺磁体的磁化强度遵循布里渊函数 $\mathrm{B}_J(y)$，
图中对不同 $J$ 绘出：$J=\frac{1}{2},1,\frac{3}{2},2,\frac{5}{2},\ldots$ 及 $J=\infty$。}
\end{figure}

居里定律的 $\chi\propto1/T$ 意味着 $1/\chi$ 对 $T$ 作图是直线、$\chi T$ 对 $T$ 作图
为常数（见图 2.10）。后续章节考虑相互作用的作用时，记住这几点会有用。还须注意：
磁化率是在外磁场趋于零的极限下取值，由式 2.44 以
$\mu_{\mathrm{eff}}=g_J\mu_{\mathrm{B}}\sqrt{J(J+1)}$ 给出；而在强场下磁化强度饱和
为 $M_{\mathrm{s}}$，由式 2.39 相当于每离子磁矩 $g_J\mu_{\mathrm{B}}J$。这两个数值
$\mu_{\mathrm{eff}}=g_J\mu_{\mathrm{B}}\sqrt{J(J+1)}$ 与 $M_{\mathrm{s}}/n=g_J\mu_{\mathrm{B}}J$
一般不等，仅当 $J\to\infty$（经典极限）时才相等。

\begin{figure}
\centering
\includegraphics[width=0.32\textwidth]{fig2_10.pdf}
\caption{居里定律 $\chi\propto1/T$（a）；$1/\chi$ 对 $T$ 作图为直线（b）；$\chi T$
对 $T$ 作图为常数（c）。}
\end{figure}

\subsection{范弗莱克顺磁性}

若基态 $|0\rangle$ 的 $J=0$，则没有顺磁效应，因为
\begin{equation*}
\langle 0|\hat{\boldsymbol{\mu}}|0\rangle = g_J\mu_{\mathrm{B}}\langle 0|\hat{\mathbf{J}}|0\rangle = 0.\tag{2.47}
\end{equation*}
这意味着加磁场不改变系统基态能量，因而没有顺磁磁化率。但这一结论只在一阶微扰论
中正确。二阶微扰论预言基态能量 $E_0$ 会变——因为它计入了 $J\neq0$ 的激发态混入。
$J=0$ 离子的基态能量变化为
\begin{equation*}
\Delta E_0 = \sum_n\frac{|\langle 0|(\mathbf{L}+g\mathbf{S})\cdot\mathbf{B}|n\rangle|^{2}}{E_0-E_n} + \frac{e^{2}}{8m_{\mathrm{e}}}\sum_i(\mathbf{B}\times\mathbf{r}_i)^{2},\tag{2.48}
\end{equation*}
第二项来自抗磁性，第一项求和遍及系统所有激发态。磁化率于是为
\begin{equation*}
\chi = \frac{N}{V}\left(2\mu_{\mathrm{B}}^{2}\sum_n\frac{|\langle 0|(L_z+gS_z)|n\rangle|^{2}}{E_n-E_0} - \frac{e^{2}\mu_0}{6m_{\mathrm{e}}}\sum_{i=1}^{Z}\langle r_i^{2}\rangle\right),\tag{2.49}
\end{equation*}
第一项为正（因 $E_n>E_0$），称为范弗莱克顺磁性（Van Vleck paramagnetism）；
第二项为负，即前面讨论过的常规抗磁磁化率（式 2.15）。与抗磁性一样，范弗莱克
顺磁性既小、又不依赖温度。
\bio{John H. van Vleck\\（1899--1980）}

\section{离子的基态与洪特规则}

典型的原子含有很多电子，其中大部分位于没有净角动量的填满壳层中。但可能存在未满
壳层，其中电子可以组合出非零的自旋与轨道角动量。1.3.4 节已经看到两个自旋
$\frac{1}{2}$ 电子如何组合成自旋为 0 或 1 的联合体。在原子中，未满壳层所有电子的
自旋角动量可以全部组合在一起，轨道角动量亦然。于是原子具有总轨道角动量
$\hbar\mathbf{L}$ 与总自旋角动量 $\hbar\mathbf{S}$。轨道与自旋角动量可以按
\begin{equation*}
(2L+1)(2S+1)\tag{2.50}
\end{equation*}
种方式组合：这是 $\mathbf{L}$ 的 $z$ 分量的取值数（级数 $-L,-L+1,\ldots,L-1,L$
共 $(2L+1)$ 项）乘以 $\mathbf{S}$ 的 $z$ 分量的取值数（同理为 $(2S+1)$）。这些以
不同方式组合未满壳层电子角动量（自旋与轨道）的组态具有不同的能量。原因在于：自旋
角动量的选择影响波函数的空间部分，轨道角动量影响电子绕核的运动方式——两者都影响
电子相互回避的好坏，从而影响静电排斥能。下面说明如何找到能量最低的组态。

\subsection{精细结构}

迄今我们一直把自旋与轨道角动量分开处理，因为它们彼此独立。然而它们确实通过
自旋--轨道相互作用（spin--orbit interaction，见附录 C）弱耦合，它对具有确定 L、S
的态是一个微扰。因此 L 与 S 不再分别守恒，但总角动量 $\mathbf{J}=\mathbf{L}+\mathbf{S}$
守恒。若把相对论效应作为微扰处理（通常可以），则 $\mathbf{L}^{2}=L(L+1)$ 与
$\mathbf{S}^{2}=S(S+1)$ 可视为守恒。于是具有给定 L、S 的态分裂成若干不同 J 的能级，
称为精细结构（fine structure）。J 取 $|L-S|$ 到 $L+S$ 之间的值。由 J 的定义
（式 2.17），
\begin{equation*}
\mathbf{J}^{2} = \mathbf{L}^{2}+\mathbf{S}^{2}+2\mathbf{L}\cdot\mathbf{S},\tag{2.51}
\end{equation*}
自旋--轨道相互作用取 $\lambda\mathbf{L}\cdot\mathbf{S}$ 的形式（见附录 C），$\lambda$
为常数，其能量期望值为
\begin{equation*}
\langle\lambda\mathbf{L}\cdot\mathbf{S}\rangle = \frac{\lambda}{2}[J(J+1)-L(L+1)-S(S+1)].\tag{2.52}
\end{equation*}
原子的能量主要由静电考虑通过 S、L 的取值决定，因此能量本征态可以用 S 和 L 标记。
在无自旋--轨道相互作用时，J 在 $|L-S|$ 到 $L+S$ 中的具体取值并不影响能量。每个
能级是 $(2S+1)(2L+1)$ 个态组成的多重态（multiplet）。把自旋--轨道相互作用作为
微扰加入后，多重态分裂成以 J 标记的不同精细结构能级；每个能级自身又有 $2J+1$ 的
简并，故加磁场后可按不同 $m_J$ 进一步分裂。精细结构能级的分裂遵循朗德间隔定则
（Land\'e interval rule）：给定多重态中相邻能级 $E(J)$ 与 $E(J-1)$ 的间隔为
\begin{align*}
E(J)-E(J-1) &= \frac{\lambda}{2}[J(J+1)-L(L+1)-S(S+1)]\\
&\quad -\frac{\lambda}{2}[(J-1)J-L(L+1)-S(S+1)]\\
&= \lambda J.\tag{2.53}
\end{align*}
即（考虑 $J-1$ 与 $J$ 两能级的间隔）分裂正比于 $J$。

\begin{figure}
\centering
\includegraphics[width=0.5\textwidth]{fig2_11.pdf}
\caption{角动量组合 $L=3$、$S=\frac{3}{2}$ 给出从 $\frac{3}{2}$ 到 $\frac{9}{2}$ 的
各个 $J$ 值。}
\end{figure}

\begin{example}{2.1}
图 2.11 以 $L=3$、$S=\frac{3}{2}$ 为例演示这种角动量组合。此时
$(2S+1)(2L+1)=28$；由图可见，若考虑 J 从 $|L-S|=\frac{3}{2}$ 到 $L+S=\frac{9}{2}$
的各个态，由于每个 J 态简并度为 $2J+1$，总共也是 28 个态。这形象地说明了
\begin{equation*}
\sum_{J=|L-S|}^{L+S}(2J+1) = (2L+1)(2S+1).\tag{2.54}
\end{equation*}
\end{example}

角动量量子数的组合显然非常多；对一个特定离子，哪一个是基态？

\subsection{洪特规则}

使能量取极小的角动量量子数组合可以用洪特规则（Hund's rules）估计。\fn{12}{洪特规则只适用于基态组态，对最低能级之上各能级的次序不作任何断言；它还假定只有一个未满支壳层。}\bio{Friederich Hund\\（1896--1997）}这三条经验规则按重要性递减排列：先满足第一条，在此前提下再尽量满足第二条，依此类推。

（1）安排电子波函数使 S 最大。由于泡利不相容原理阻止自旋平行的电子占据同一位置，
库仑能得以减小，从而降低电子间的静电排斥。

（2）在第一条已定的波函数基础上使 L 最大。这同样降低能量：同向旋转轨道上的电子
能更有效地相互回避，从而减小库仑排斥。

（3）最后用 $J=|L-S|$（壳层填充不足半满）或 $J=|L+S|$（超过半满）定出 J。第三条
规则源于使自旋--轨道能量取极小的尝试。注意它只在一定条件下适用：如下一章所示，
在许多体系（过渡金属离子是很好的例子）中，自旋--轨道能量不如晶体场等其他能量项
重要，此时洪特第三条规则被违反；而对稀土离子，如下所示，第三条规则非常好用。

求出 S、L、J 之后，基态可用形如 ${}^{2S+1}L_J$ 的谱项符号（term symbol）概括。
其中 L 不写成数字，而按下列序列用字母表示：
\begin{center}
\begin{tabular}{lcccccccc}
\toprule
$L$ & 0 & 1 & 2 & 3 & 4 & 5 & 6 & $\cdots$\\
\midrule
 & S & P & D & F & G & H & I & $\cdots$\\
\bottomrule
\end{tabular}
\end{center}
$2S+1$ 是自旋多重度（spin multiplicity）。

\begin{example}{2.2}
以稀土离子 $\mathrm{Dy^{3+}}$（外壳层 $4f^{9}$）为例：f 电子 $l=3$，为满足洪特
第一条规则，$2l+1=7$ 个自旋向上，剩下 2 个自旋向下（见图 2.12）。于是
$S=7\times\frac{1}{2}-2\times\frac{1}{2}=\frac{5}{2}$（自旋简并 $2S+1=6$）。自旋
向上的电子对轨道角动量没有净贡献，因此只需使其余 2 个自旋向下的电子的轨道贡献最大：
$L=3+2=5$，对应字母 H。壳层超过半满，故 $J=|5+\frac{5}{2}|=\frac{15}{2}$。谱项
符号为 ${}^{6}\mathrm{H}_{15/2}$。
\end{example}

\begin{figure}
\centering
\includegraphics[width=0.35\textwidth]{fig2_12.pdf}
\caption{$\mathrm{Dy^{3+}}$ 的基态。}
\end{figure}

洪特规则预言基态，但对激发态及其与基态的接近程度不置一词。因此它使我们能在
“只有基态被占据”的假定下估计离子的磁矩。

本章和下一章都要把这些预言与实验比较，我们将集中于含 3d 与 4f 离子的化合物——
它们在许多磁性体系中都很重要。3d 元素是第一行过渡金属（Sc--Zn），4f 元素即镧系
元素或稀土元素（La--Lu），见图 2.13 的周期表。

\begin{figure}
\centering
\includegraphics[width=0.85\textwidth]{fig2_13.pdf}
\caption{周期表。3d 元素（Sc--Zn）与 4f 元素（La--Lu）以阴影标出。每个元素旁的数字
是原子序数 Z（核内质子数）。}
\end{figure}

洪特规则用于 3d 与 4f 离子的结果示于图 2.14（注意：如下一章详述，4f 离子的洪特
规则预言比 3d 离子更符合实验）。\fn{13}{本章只讨论自由原子或离子。把原子放入晶体环境中情况就会改变；对 3d 离子变化相当大，见第三章。}在每个组中段 S 升至极大；L 与
J 的极大大致出现在四分之一与四分之三填充处——J 的两个极大不对称，反映了壳层不足
半满与超过半满时规则的不同。

\begin{figure}
\centering
\includegraphics[width=0.48\textwidth]{fig2_14.pdf}
\caption{按洪特规则得到的 3d 与 4f 离子的 S、L、J。图中 $n$ 是支壳层（3d 或 4f）
中的电子数。}
\end{figure}

由式 2.44，测量磁化率可以推定有效磁矩。以玻尔磁子为单位：
\begin{equation*}
p = \mu_{\mathrm{eff}}/\mu_{\mathrm{B}},\tag{2.55}
\end{equation*}
洪特规则的预言（用式 2.45）是 $p=g_J[J(J+1)]^{\frac{1}{2}}$。如表 2.2 所示，对
固体中的 4f 离子，这一预言与实测 $p=\mu_{\mathrm{eff}}/\mu_{\mathrm{B}}$ 通常符合
极好。Sm 与 Eu 出现偏差，原因是它们存在 J 不同于基态的低激发态，这些态
因距基态很近而也有显著布居，使 $p$ 偏离基态值。4f 离子总体符合良好，但下一章会
看到 3d 离子的符合程度差得多，这是局域晶体环境影响所致。这一效应对 4f 离子不
重要，因为部分填充的 4f 壳层深藏于离子内部、处于填满的 5s 与 5p 壳层之下。

\begin{table}
\centering
\caption{用洪特规则得到的 4f 离子磁基态。对每个离子列出壳层组态及预言的基态
S、L、J；并给出由洪特规则计算的 $p=\mu_{\mathrm{eff}}/\mu_{\mathrm{B}}=g_J[J(J+1)]^{1/2}$。
下一列是实验值 $p_{\mathrm{exp}}$：除 Sm、Eu 外符合极好。实验值来自顺磁盐在
$k_{\mathrm{B}}T\gg E_{\mathrm{CEF}}$ 温度下的磁化率测量，$E_{\mathrm{CEF}}$ 为
晶体场能量。}
\small
\begin{tabular}{lccccccc}
\toprule
离子 & 壳层 & $S$ & $L$ & $J$ & 谱项 & $p$ & $p_{\mathrm{exp}}$\\
\midrule
$\mathrm{Ce^{3+}}$ & $4f^{1}$ & $\frac{1}{2}$ & 3 & $\frac{5}{2}$ & ${}^{2}F_{5/2}$ & 2.54 & 2.51\\
$\mathrm{Pr^{3+}}$ & $4f^{2}$ & 1 & 5 & 4 & ${}^{3}H_{4}$ & 3.58 & 3.56\\
$\mathrm{Nd^{3+}}$ & $4f^{3}$ & $\frac{3}{2}$ & 6 & $\frac{9}{2}$ & ${}^{4}I_{9/2}$ & 3.62 & 3.3--3.7\\
$\mathrm{Pm^{3+}}$ & $4f^{4}$ & 2 & 6 & 4 & ${}^{5}I_{4}$ & 2.68 & --\\
$\mathrm{Sm^{3+}}$ & $4f^{5}$ & $\frac{5}{2}$ & 5 & $\frac{5}{2}$ & ${}^{6}I_{5/2}$ & 0.85 & 1.74\\
$\mathrm{Eu^{3+}}$ & $4f^{6}$ & 3 & 3 & 0 & ${}^{7}F_{0}$ & 0.0 & 3.4\\
$\mathrm{Gd^{3+}}$ & $4f^{7}$ & $\frac{7}{2}$ & 0 & $\frac{7}{2}$ & ${}^{8}S_{7/2}$ & 7.94 & 7.98\\
$\mathrm{Tb^{3+}}$ & $4f^{8}$ & 3 & 3 & 6 & ${}^{7}F_{6}$ & 9.72 & 9.77\\
$\mathrm{Dy^{3+}}$ & $4f^{9}$ & $\frac{5}{2}$ & 5 & $\frac{15}{2}$ & ${}^{6}H_{15/2}$ & 10.63 & 10.63\\
$\mathrm{Ho^{3+}}$ & $4f^{10}$ & 2 & 6 & 8 & ${}^{5}I_{8}$ & 10.60 & 10.4\\
$\mathrm{Er^{3+}}$ & $4f^{11}$ & $\frac{3}{2}$ & 6 & $\frac{15}{2}$ & ${}^{4}I_{15/2}$ & 9.59 & 9.5\\
$\mathrm{Tm^{3+}}$ & $4f^{12}$ & 1 & 5 & 6 & ${}^{3}H_{6}$ & 7.57 & 7.61\\
$\mathrm{Yb^{3+}}$ & $4f^{13}$ & $\frac{1}{2}$ & 3 & $\frac{7}{2}$ & ${}^{2}F_{7/2}$ & 4.53 & 4.5\\
$\mathrm{Lu^{3+}}$ & $4f^{14}$ & 0 & 0 & 0 & ${}^{1}S_{0}$ & 0 & 0\\
\bottomrule
\end{tabular}
\end{table}

本章我们把每个原子的磁矩与局域在该原子上的电子相联系。金属中的传导电子并不局域
于特定原子，而是处于能带态、遍布整个材料。传导电子既表现抗磁效应也表现顺磁效应，
这将在第七章讨论。

\subsection{L--S 耦合与 j--j 耦合}

以上原子模型假定 L--S 耦合（L--S coupling，又称 Russell--Saunders 耦合）：自旋--
轨道相互作用是弱微扰，主要能量项由决定 L 与 S 取值的静电相互作用给出——即先把
各电子的轨道与自旋角动量分别组合起来；直到知道整个原子的总轨道与总自旋角动量
之后，才把自旋--轨道相互作用作为弱微扰加入，把每个谱项分裂成以 J 标记的精细结构
能级。

对高原子序数 Z 的原子这不再适用，因为自旋--轨道相互作用能正比于
$Z^{4}$\fn{14}{$Z^{4}$ 律只对类氢原子成立；对中性原子没有这么简单，但接近 $Z^{2}$。}（见附录 C），在这些原子中不能当作小微扰。更好的方案称为 j--j 耦合（j--j coupling）：
此时自旋--轨道相互作用是主导能量，先把每个电子自身的自旋与轨道角动量耦合起来，
随后较弱的静电效应再耦合各电子的总角动量。

\begin{example}{2.3}
以碳（C）与铅（Pb）为例。二者外壳层都有两个 p 电子，但 Pb 在周期表中比 C 靠后得
多，其 $Z^{4}$ 大近五个数量级。碳（组态 $2p^{2}$）可用 L--S 耦合处理：由洪特规则，
S=1（规则一“S 最大”，两电子 $\frac{1}{2}+\frac{1}{2}$）、L=1（规则二“L 最大”，
p 电子一个可取 $m_l=1$，另一个只能取 $m_l=0$）、J=0（规则三，不足半满，
$J=|1-1|$），故 C 基态谱项符号为 ${}^{3}P_{0}$。

铅（组态 $6p^{2}$）中自旋--轨道相互作用占主导，更适合 j--j 耦合。每个电子
$s=\frac{1}{2}$、$l=1$，组合 j 可为 $\frac{1}{2}$ 或 $\frac{3}{2}$，取决于自旋--
轨道能的符号；实际给出 $j=\frac{1}{2}$。两个电子都有 $j=\frac{1}{2}$，经较弱的
静电效应耦合成可为 0 或 1 的总 J。由泡利不相容原理基态为 J=0（两个电子若量子数
全同就必须处于反对称态而分开），故基态的好量子数为 $j_1=\frac{1}{2}$、
$j_2=\frac{1}{2}$、J=0。注意此时 L 与 S 不是好量子数。（好量子数的定义见附录 C。）
\end{example}

\section{绝热退磁}

本节描述一种把样品冷却到极低温度的技术，它要用到我们关于顺磁材料的一些想法，
并相当清楚地带出熵的概念。考虑含 $N$ 个独立磁矩的顺磁盐样品。不加磁场时磁矩方向
随机（我们假定它们彼此无相互作用），系统没有净磁化强度。加磁场则使磁矩趋于排齐、
产生磁化强度。升温减小磁化强度，加场增大磁化强度——这体现于磁化强度是 $B/T$ 的
函数（见式 2.37 与 2.38）。

极高温下所有磁矩方向完全随机，净磁化强度为零。热能 $k_{\mathrm{B}}T$ 太大，所有
态被同等占据，无论其能量是否有利。若磁矩 $J=\frac{1}{2}$，只能平行或反平行于磁场：
排列上下磁矩的方式数 $W=2^{N}$。于是磁矩对熵 S 的贡献为
\begin{equation*}
S = k_{\mathrm{B}}\ln W = Nk_{\mathrm{B}}\ln 2.\tag{2.56}
\end{equation*}
一般 $J>\frac{1}{2}$ 时 $W=(2J+1)^{N}$，熵为
\begin{equation*}
S = Nk_{\mathrm{B}}\ln(2J+1).\tag{2.57}
\end{equation*}

计算较低温度下的熵需要换个方法。离子总角动量 $z$ 分量取 $m_J$ 的概率
$p(m_J)$ 正比于玻尔兹曼因子 $\exp(-m_Jg\mu_{\mathrm{B}}B/k_{\mathrm{B}}T)$：
\begin{equation*}
p(m_J) = \frac{1}{Z}\exp(m_J g\mu_{\mathrm{B}}B/k_{\mathrm{B}}T)\tag{2.58}
\end{equation*}
$Z$ 是配分函数（式 2.32）。系统的有序程度由熵 S 描述；既然现在概率化地处理系统，
就采用表达式
\begin{equation*}
S = -Nk_{\mathrm{B}}\sum_{m_J}p(m_J)\ln p(m_J).\tag{2.59}
\end{equation*}
熵的表达式也可以这样得到：先由 $F=-Nk_{\mathrm{B}}T\ln Z$ 算出亥姆霍兹自由能 F，
再用 $S=-(\partial F/\partial T)_B$。

考察式 2.59 的后果。无外场或高温下系统完全无序，各 $m_J$ 等概率
$p(m_J)=1/(2J+1)$，熵退化为
\begin{equation*}
S = Nk_{\mathrm{B}}\ln(2J+1)\tag{2.60}
\end{equation*}
与式 2.57 一致。温度降低时低能量态越来越可几；磁矩沿外场排列的程度（磁化强度）
增加，熵下降。极低温下所有磁矩都与磁场排齐以降低能量——此时系统只有一种排列方式
（全部自旋排齐），$W=1$，$S=0$。

磁制冷样品的原理如下。先用液氦把顺磁体冷却到一个低的起始温度，然后分两步进行
磁制冷（见图 2.15）。

\begin{figure}
\centering
\includegraphics[width=0.6\textwidth]{fig2_15.pdf}
\caption{顺磁盐的熵随温度的函数关系，对应零与某个最大值 $B_b$ 之间的若干外加磁场。
把顺磁盐从温度 $T_{\mathrm{i}}$ 磁制冷到 $T_{\mathrm{f}}$ 分两步：先等温磁化，
在恒温 $T_{\mathrm{i}}$ 下把磁场从 0 增至 $B_b$（$a\to b$）；再绝热退磁
（$b\to c$）。$S(T)$ 曲线的计算取 $J=\frac{1}{2}$（见式 2.76），并加了一项
$\propto T^{3}$ 以模拟晶格振动的熵。$B=0$ 的曲线实际对应小而非零的 B，以模拟小的
剩余场。}
\end{figure}

第一步是等温磁化。顺磁体的能量因磁矩沿场排列而降低，故给定温度下增强外场可以增加
排齐程度。这一步等温进行（图 2.15 中 $a\to b$）：让样品与液氦池保持热接触（氦在
大气压下的沸点为 4.2 K，减压还可低于 4.2 K）。样品温度不变，氦池吸收样品能量与熵
下降时释放的热。热接触通常由样品室中的低压氦气提供——它在样品与室壁之间导热，
室本身浸在氦池中。（这种气体常称“交换”气体，因为它让样品与热池交换热量。）

第二步是使样品与氦池热隔离（抽走交换气体），然后缓慢把磁场降至零——慢到过程是
准静态的、熵保持不变。这一步称为绝热退磁（adiabatic demagnetization，图 2.15 中
$b\to c$），它降低系统温度。绝热退磁期间样品的总熵不变：磁矩的熵增加（场减弱时
磁矩随机化），恰好被声子（晶格振动）的熵减少（样品降温）所平衡。熵就这样在声子
与自旋之间交换。

这种制冷方法有极限吗？初看之下，$B=0$ 时熵在一切 $T>0$ 似乎都是
$S=Nk_{\mathrm{B}}\ln(2J+1)$，只在绝对零度才降为零，那么绝热退磁似乎可以一路制冷
到绝对零度。然而真实系统中总存在磁矩间相互作用导致的微小剩余内场，使熵在略高于
绝对零度处就提前跌落（见图 2.15）。这个场的大小限制了盐能冷却到的最低温度。某些
剩余内场极小的盐可以达到几个毫开尔文。

\section{核自旋}

原子中具有磁矩的不只是电子。原子核常因核角动量而有非零自旋。每个核有一个称为核
自旋量子数的量子数 $I$，代表以 $\hbar$ 为单位的核总角动量。但核磁矩非常小，因为
其大小与所涉粒子的质量成反比：核磁矩典型地比电子磁矩小三个量级（大多介于
$10^{-3}$ 与 $10^{-4}\mu_{\mathrm{B}}$ 之间）。磁矩之小，加上相邻原子核之间不存在
强相互作用，使核自旋体系在常温实验室条件下无法有序。

\begin{table}
\centering
\caption{一些常见核自旋的性质。首先列出中子（n）、质子（p）、氘核（d）与氚核（t）。
Z 是原子序数（质子数），N 是中子数，核 X 记作 ${}^{A}X$，$A=Z+N$ 为质量数。I 是核
自旋量子数，$\mu$ 是以核磁子（$\mu_{\mathrm{N}}$）量度的核磁矩，$g_I$ 是核 $g$
因子。最后一列是这些磁矩在 1 T 磁场中的进动频率 $\nu$。}
\small
\begin{tabular}{lcccccc}
\toprule
核 & $Z$ & $N$ & $I$ & $\mu/\mu_{\mathrm{N}}$ & $g_I$ & $\nu$（B=1 T，MHz）\\
\midrule
n & 0 & 1 & $\frac{1}{2}$ & $-1.913$ & $-3.826$ & 29.17\\
p=${}^{1}$H & 1 & 0 & $\frac{1}{2}$ & 2.793 & 5.586 & 42.58\\
d=${}^{2}$H & 1 & 1 & 1 & 0.857 & 0.857 & 6.536\\
t=${}^{3}$H & 1 & 2 & $\frac{1}{2}$ & $-2.128$ & 4.255 & 32.43\\
${}^{12}$C & 6 & 6 & 0 & 0 & 0 & 0\\
${}^{13}$C & 6 & 7 & $\frac{1}{2}$ & 0.702 & 1.404 & 10.71\\
${}^{14}$N & 7 & 7 & 1 & 0.404 & 0.404 & 3.076\\
${}^{16}$O & 8 & 8 & 0 & 0 & 0 & 0\\
${}^{17}$O & 8 & 9 & $\frac{5}{2}$ & $-1.893$ & $-0.757$ & 5.772\\
${}^{19}$F & 9 & 10 & $\frac{1}{2}$ & 2.628 & 5.257 & 40.05\\
${}^{31}$P & 15 & 16 & $\frac{1}{2}$ & 1.132 & 2.263 & 17.24\\
${}^{33}$S & 16 & 17 & $\frac{3}{2}$ & 0.643 & 0.429 & 3.266\\
\bottomrule
\end{tabular}
\end{table}

核磁性的单位是核磁子（nuclear magneton）$\mu_{\mathrm{N}}$：
\begin{equation*}
\mu_{\mathrm{N}} = \frac{e\hbar}{2m_{\mathrm{p}}} = 5.0508\times10^{-27}\ \mathrm{A\,m^{2}},\tag{2.61}
\end{equation*}
$m_p$ 是质子质量。（这远小于 1s 电子的磁矩——玻尔磁子
$\mu_{\mathrm{B}}=e\hbar/2m_{\mathrm{e}}=9.27\times10^{-24}$ A\,m$^2$，$m_{\mathrm{e}}$
为电子质量，见式 1.15。）核自旋量子数 I 取下列数值之一：0、$\frac{1}{2}$、1、
$\frac{3}{2}$、2、……角动量沿 $z$ 方向的分量为 $m_I$（以 $\hbar$ 为单位），$m_I$
可取
\begin{equation*}
-I, -I+1, \dots, I-1, I.\tag{2.62}
\end{equation*}
把核磁矩投影到 $z$ 方向、取 $m_I$ 最大的自旋态，其值为
\begin{equation*}
\mu = g_I\mu_{\mathrm{N}}I\tag{2.63}
\end{equation*}
$g_I$ 是核 $g$ 因子，为 1 量级的数，反映核的细致结构。若干核的 I、$g_I$ 与 $\mu$
列于表 2.3。中子与质子的磁矩都不是 $\mu_{\mathrm{N}}$ 的简单倍数——这表明每个粒子
内部结构复杂：中子与质子都由三个夸克组成。质子 $I=\frac{1}{2}$、$g_I=5.586$，磁矩
$\mu_p=1.410\times10^{-26}$ A\,m$^2$。

通常我们忽略核磁性：水中质子的磁矩并不能让水吸附在磁极上（事实上由于水分子中
电子云，水主要是弱抗磁体）。不过，借助非常灵敏的共振技术实验可以探测核磁矩。这一
实验方法称为核磁共振（nuclear magnetic resonance），将在 3.2.1 节讨论。

\begin{example}{2.4}
许多常见核完全没有磁矩，例如最常见的碳与氧——${}^{12}\mathrm{C}$ 与
${}^{16}\mathrm{O}$。这是因为这两种核的质子数与中子数都是偶数；核中全同粒子对
倾向于配成单态。有未成对核子的核才有非零磁矩：${}^{13}\mathrm{C}$ 有一个未成对
中子，$I=\frac{1}{2}$；${}^{14}\mathrm{N}$ 有一个未成对中子和一个未成对质子，
组合给出 I=1。虽然碳与氧最丰的同位素没有核磁矩，其次丰的同位素却有（天然碳约
1.1\% 是 ${}^{13}\mathrm{C}$，天然氧约 0.04\% 是 ${}^{17}\mathrm{O}$；二者都有
非零磁矩，见表 2.3）。
\end{example}

\section{超精细结构}

原子内部，核磁矩可以与电子磁矩发生磁相互作用，但极其微弱。由此产生的能级劈裂
比 2.5.1 节的精细结构还要小，故称超精细结构（hyperfine structure）。核与电子之间
的弱磁相互作用称为超精细相互作用（hyperfine interactions）。其起源相当复杂，
但基本原理可以这样理解：考虑核磁矩 $\boldsymbol{\mu}$ 处于磁场
$\mathbf{B}_{\text{electrons}}$ 中（由所有电子的运动与自旋产生），产生能量项
$-\boldsymbol{\mu}\cdot\mathbf{B}_{\text{electrons}}$。$\mathbf{B}_{\text{electrons}}$
预期正比于所有电子的总角动量 $\mathbf{J}$，故超精细相互作用的哈密顿量可写为
\begin{equation*}
\hat{\mathcal{H}}_{\mathrm{hf}} = A\,\mathbf{I}\cdot\mathbf{J}.\tag{2.64}
\end{equation*}
I 是核角动量，A 是可由实验确定的参数，量度超精细相互作用的强度。

对一般原子，超精细相互作用的严格形式很复杂，但对与点核相互作用的单电子原子可以
详细算出（见习题 2.8）。这里只作概述，以阐明相互作用的基本机制。

重要的磁相互作用有两类。第一类是磁偶极相互作用（4.1 节详述）。核磁矩
$\boldsymbol{\mu}_I=g_I\mu_I\mathbf{I}$ 与电子磁矩
$\boldsymbol{\mu}_{\mathrm{e}}=g_e\mu_{\mathrm{B}}\mathbf{S}$ 之间的偶极相互作用为
$\hat{\mathcal{H}}_{\text{dipole}}$：
\begin{align*}
\hat{\mathcal{H}}_{\text{dipole}} &= \frac{\mu_0}{4\pi r^{3}}\left[\boldsymbol{\mu}_{\mathrm{e}}\cdot\boldsymbol{\mu}_I - \frac{3}{r^{2}}(\boldsymbol{\mu}_{\mathrm{e}}\cdot\mathbf{r})(\boldsymbol{\mu}_I\cdot\mathbf{r})\right]\tag{2.65}\\
&= \frac{\mu_0 g_e g_I\mu_{\mathrm{B}}\mu_I}{4\pi r^{3}}\left[\mathbf{I}\cdot\mathbf{S} - \frac{3}{r^{2}}(\mathbf{S}\cdot\mathbf{r})(\mathbf{I}\cdot\mathbf{r})\right].\tag{2.66}
\end{align*}
若两个磁矩都沿 $z$ 轴排列，则 $\hat{\mathcal{H}}_{\text{dipole}}\propto(1-3\cos^{2}\theta)I_zS_z/r^{3}$，
相互作用能正比于
\begin{equation*}
\int\mathrm{d}^{3}\mathbf{r}\,|\psi(\mathbf{r})|^{2}(1-3\cos^{2}\theta)m_sm_I/r^{3},\tag{2.67}
\end{equation*}
对 s 轨道该积分为零——本质原因是偶极场在球面上平均为零，可借图 2.16 理解。

球对称的 s 轨道波函数把它平均为零。但对 $l>0$ 轨道（如 p 轨道）中的未成对电子，
积分不为零。上面的处理也只考虑了电子磁矩的自旋部分；电子的轨道磁矩同样在核处产生
磁场，给出正比于 $\mathbf{I}\cdot\mathbf{L}$ 的项。（所有这些项在习题 2.8 的处理中
都会自然出现。）

\begin{figure}
\centering
\includegraphics[width=0.5\textwidth]{fig2_16.pdf}
\caption{磁性核（设为球形，图中以阴影表示）内外磁场分布的示意图。磁场分布由绕核赤道
流动的电流（图中阴影）产生。球外的场看似来自球心处的点偶极子，对球对称体积求平均
为零——场向上与向下的机会均等。核球内部的场平均向上、对球对称体积平均不为零。
这正是费米接触相互作用的起源。}
\end{figure}

第二类机制是费米接触相互作用（Fermi contact interaction）——与上一机制相反，
它对所有非 s 轨道为零，只对 s 轨道不为零。这里考虑式 2.66 在 $r\to0$ 时的情形。
核当然不是点偶极子，而有有限体积，记为 $V$。我们进一步假设\fn{15}{这一假设只是为了让计算简单。它实际上给出正确答案，但更好的处理见习题 2.8。}核是磁化强度均匀的球
$\mathbf{M}=\boldsymbol{\mu}_I/V$。球内的磁通密度为 $\mu_0\mathbf{M}$，但须先扣除
退磁场——对球为 $\mu_0\mathbf{M}/3$（见附录 D）——剩下 $2\mu_0\mathbf{M}/3$。
于是闯入核球内的电子将感受到磁通密度
\begin{equation*}
\mathbf{B} = \frac{2\mu_0\boldsymbol{\mu}_I}{3V}.\tag{2.68}
\end{equation*}
能量代价由它乘以 $\mu_{\mathrm{e}}|\psi(0)|^{2}V$——处于核内的电子磁矩量（注意
$|\psi(0)|^{2}V$ 是在核内找到电子的概率）——得到。故接触能 $E_{\text{contact}}$ 为
\begin{equation*}
E_{\text{contact}} = \left(\frac{2\mu_0}{3}\right)2\mu_{\mathrm{B}}g_I\mu_{\mathrm{N}}\,\mathbf{I}\cdot\mathbf{S},\tag{2.69}
\end{equation*}
其中用了 $\boldsymbol{\mu}_{\mathrm{e}}=-2\mu_{\mathrm{B}}\mathbf{S}$ 与
$\boldsymbol{\mu}_I=g_I\mu_I\mathbf{I}$。电子在核内感受到的净场也可借图 2.16 理解。
只有 s 电子在 $r=0$ 处有振幅，因而只有它们表现出这一效应。费米接触相互作用引起的
s 电子超精细劈裂，通常远大于 $l>0$ 电子的磁偶极超精细劈裂。

如习题 2.8 所示，同时包含偶极与费米接触项的超精细相互作用形式相当复杂，但其要点
是：对 s 电子（$l=0$）为常数倍的 $\mathbf{I}\cdot\mathbf{S}$，对 $l>0$ 电子为
$\mathbf{I}\cdot\mathbf{N}$，其中 $\mathbf{N}=\mathbf{L}-\mathbf{S}+3(\mathbf{S}\cdot\mathbf{r})\mathbf{r}/r^{2}$。
对 s 电子 $\mathbf{J}=\mathbf{S}$，相互作用与式 2.64 一致。可以证明 $\mathbf{N}$
能投影到 $\mathbf{J}$ 上，使 $\mathbf{I}\cdot\mathbf{N}$ 的期望值正比于
$\mathbf{I}\cdot\mathbf{J}$（对单电子原子比例常数为 $L(L+1)/J(J+1)$），故式 2.64
对 $l>0$ 的电子也适用。

\begin{figure}
\centering
\includegraphics[width=0.42\textwidth]{fig2_17.pdf}
\caption{（a）$I=\frac{1}{2}$、$J=\frac{1}{2}$ 与（b）$I=1$、$J=2$ 时超精细劈裂的
示意。量子数 $F$ 可取 $|J-I|,\ldots,J+I-1,J+I$；每个态的能量为
$\frac{A}{2}[F(F+1)-I(I+1)-J(J+1)]$。每个以 $F$ 标记的超精细态有 $2F+1$ 简并，
加磁场 $\mathbf{B}$ 后可按 $m_F$ 分裂。注意朗德间隔定则成立。磁场中劈裂的大小依赖
$F$、$I$、$J$，经由量子数 $F$ 的 $g$ 因子，近似为
$g_F=g_J[F(F+1)+J(J+1)-I(I+1)]/2F(F+1)$。}
\end{figure}

原子（核与电子的组合体）的总角动量 $\mathbf{F}$ 为
\begin{equation*}
\mathbf{F} = \mathbf{J} + \mathbf{I}.\tag{2.70}
\end{equation*}
量子数 $F$ 可取 $|J-I|,\ldots,J+I-1,J+I$。下面几行与自旋--轨道相互作用的处理
几乎完全平行。用式 2.70 可证
\begin{equation*}
\mathbf{F}^{2} = \mathbf{J}^{2}+\mathbf{I}^{2}+2\mathbf{I}\cdot\mathbf{J},\tag{2.71}
\end{equation*}
超精细相互作用取 $A\mathbf{I}\cdot\mathbf{J}$ 形式，其能量期望值为
\begin{equation*}
\langle A\mathbf{I}\cdot\mathbf{J}\rangle = \frac{A}{2}[F(F+1)-I(I+1)-J(J+1)].\tag{2.72}
\end{equation*}
超精细相互作用可作为微扰加入，它把电子能级分裂成以 F 标记的不同超精细结构能级；
每个能级有 $2F+1$ 简并，加磁场可按 $m_F$ 分裂。与式 2.53 直接类比，不同超精细
结构能级的分裂也遵循朗德间隔定则：相邻能级 $E(F)$ 与 $E(F-1)$ 的间隔为
\begin{equation*}
E(F)-E(F-1) = AF.\tag{2.73}
\end{equation*}
即分裂正比于 F——所考虑的相邻两能级中较大的量子数。下面的例子将演示这一点。

\begin{example}{2.5}
图 2.17 对两个例子演示了这一效应。基态原子氢含一个质子（$I=\frac{1}{2}$）和一个
电子（$J=S=\frac{1}{2}$），故 F=0 或 1（图 2.17(a)）。这两态间的能量劈裂极小，
比典型电子能量小约 6 个量级，对应约 1420 MHz 的电磁频率，或约 21 cm 的波长。
这一跃迁是原子氢微波激射器（maser）的基础，在天文学中也很重要：天文学家称之为
H I 发射。当宇宙中原子氢的浓度不足以大量形成 $\mathrm{H_2}$ 分子时，21 cm 辐射
可以热激发产生（所需能量低于宇宙微波背景的温度，见图 3.8）。若气体相对地球运动，
该辐射经历多普勒频移，于是可以由地面测量遥远星系不同区域内气体的转动；它也可用于
我们自己的银河系。

图 2.17(b) 给出更复杂的例子：I=1、J=2 的原子。图中可见朗德间隔定则成立——相邻
能级的劈裂大小随上能级的 F 值按比例变化。
\end{example}

\section*{延伸阅读}

\begin{itemize}[itemsep=0.2em]
\item B.~H.~Bransden 与 C.~J.~Joachain，《Physics of atoms and molecules》，
Longman 1983，提供了关于孤立原子的大量信息。
\item 有用的背景知识还可见 P.~W.~Atkins，《Molecular quantum mechanics》，
OUP 1983。
\item A.~Abragam 与 B.~Bleaney 的全面专著《Electron paramagnetic resonance of
transition ions》，Dover 1986，也很有用。
\item D.~J.~Griffiths，《Introduction to electromagnetism》，Prentice Hall 1989，
对物质中的静磁场给出了易读的论述。
\item 关于抗磁性的经典推导与量子推导优劣的讨论，见 S.~L.~O'Dell 与 R.~K.~P.~Zia，
American Journal of Physics \textbf{54}, 32 (1986)。
\end{itemize}

\section*{习题}

\exer{2.1} 计算基态氢原子气体（数密度 $10^{20}$ m$^{-3}$）的抗磁轨道磁化率，并与
100 K 下的顺磁自旋磁化率比较。

\exer{2.2} 估算一只鸭子的抗磁磁化率（假定它完全由水组成）。要多强的磁场才能在鸭子
中诱导出与一枚磁化铁屑相同的磁矩？对牛重做此题。

\exer{2.3} 计算一块 $\mathrm{CuSO_4\cdot5H_2O}$ 晶体（尺寸
$2\,\mathrm{mm}\times2\,\mathrm{mm}\times2\,\mathrm{mm}$，见表 2.1；密度
$2286\ \mathrm{kg\,m^{-3}}$，相对分子质量 249.7 g）在 10 K、1 T 磁场中的顺磁磁矩。

\exer{2.4} 利用式 2.29 与 2.30 中自旋 $\frac{1}{2}$ 粒子在磁场 B 中的配分函数与
亥姆霍兹函数表达式，证明对单位体积 $n$ 个此类无相互作用粒子，单位体积能量为
\begin{equation*}
E = -n\mu_{\mathrm{B}}B\tanh\left(\frac{\mu_{\mathrm{B}}B}{k_{\mathrm{B}}T}\right),\tag{2.74}
\end{equation*}
单位体积热容为
\begin{equation*}
C = nk_{\mathrm{B}}\left(\frac{\mu_{\mathrm{B}}B}{k_{\mathrm{B}}T}\right)^{2}\operatorname{sech}^{2}\left(\frac{\mu_{\mathrm{B}}B}{k_{\mathrm{B}}T}\right),\tag{2.75}
\end{equation*}
单位体积熵为
\begin{equation*}
S = nk_{\mathrm{B}}\left[\ln\left(2\cosh\left(\frac{\mu_{\mathrm{B}}B}{k_{\mathrm{B}}T}\right)\right) - \frac{\mu_{\mathrm{B}}B}{k_{\mathrm{B}}T}\tanh\left(\frac{\mu_{\mathrm{B}}B}{k_{\mathrm{B}}T}\right)\right].\tag{2.76}
\end{equation*}
这些结果绘于图 2.8。

\exer{2.5} 把上题结果推广到一般 $J$。写配分函数
\begin{equation*}
Z_J(y) = \frac{\sinh[(2J+1)\frac{y}{2J}]}{\sinh[\frac{y}{2J}]}\tag{2.77}
\end{equation*}
其中 $y=xJ=g_J\mu_{\mathrm{B}}JB/k_{\mathrm{B}}T$（见式 2.37），证明熵 $S(y)$ 与
恒场热容 $C(y)$ 作为 $y$ 的函数为
\begin{equation*}
S(y) = nk_{\mathrm{B}}[\ln Z_J(y) - y\,\mathrm{B}_J(y)]\tag{2.78}
\end{equation*}
\begin{equation*}
C(y) = nk_{\mathrm{B}}y^{2}\frac{\mathrm{dB}_J(y)}{\mathrm{d}y},\tag{2.79}
\end{equation*}
$\mathrm{B}_J(y)$ 是布里渊函数。进而证明 $y\ll1$ 时
\begin{equation*}
Z_J(y) = 2J+1\tag{2.80}
\end{equation*}
\begin{equation*}
S(y) = nk_{\mathrm{B}}\left[\ln(2J+1) - \frac{J+1}{3J}y^{2}\right]\tag{2.81}
\end{equation*}
\begin{equation*}
C(y) = nk_{\mathrm{B}}\frac{J+1}{3J}y^{2}\tag{2.82}
\end{equation*}
而 $y\to\infty$ 时
\begin{equation*}
S(y)\to 0\tag{2.83}
\end{equation*}
\begin{equation*}
C(y)\to 0.\tag{2.84}
\end{equation*}

\exer{2.6} 证明：对角动量为 $l$、含 $n$ 个电子的壳层，洪特规则可概括为
\begin{equation*}
\begin{aligned}
S &= \frac{2l+1-|2l+1-n|}{2}\\
L &= S\,|2l+1-n|\\
J &= S\,|2l-n|.
\end{aligned}
\end{equation*}

\exer{2.7} 求下列离子基态的谱项符号：（a）$\mathrm{Ho^{3+}}$（$4f^{10}$）；
（b）$\mathrm{Er^{3+}}$（$4f^{11}$）；（c）$\mathrm{Tm^{3+}}$（$4f^{12}$）；
（d）$\mathrm{Lu^{3+}}$（$4f^{14}$）。

\exer{2.8} 位于原点的点偶极子 $\boldsymbol{\mu}$ 在位置 $\mathbf{r}$ 处产生的磁矢势为
\begin{equation*}
\mathbf{A} = \frac{\mu_0}{4\pi r^{3}}\boldsymbol{\mu}\times\mathbf{r} = \frac{\mu_0}{4\pi}\nabla\times\left(\frac{\boldsymbol{\mu}}{r}\right).\tag{2.85}
\end{equation*}
把它用于原子中电子的哈密顿量
\begin{equation*}
\hat{\mathcal{H}} = \frac{1}{2m_{\mathrm{e}}}(\mathbf{p}+e\mathbf{A})^{2} + 2\mu_{\mathrm{B}}\mathbf{S}\cdot\mathbf{B} + V(r)\tag{2.86}
\end{equation*}
并忽略 $A^{2}$ 项，证明：核偶极子 $\boldsymbol{\mu}=g_I\mu_{\mathrm{N}}\mathbf{I}$
的存在对哈密顿量 $\hat{\mathcal{H}}_0=p^{2}/2m_{\mathrm{e}}+V(r)$ 的微扰
$\hat{\mathcal{H}}'=\hat{\mathcal{H}}-\hat{\mathcal{H}}_0$ 为
\begin{equation*}
\hat{\mathcal{H}}' = \frac{\mu_0}{4\pi}2g_I\mu_{\mathrm{B}}\mu_{\mathrm{N}}\left[(\mathbf{S}\cdot\nabla)(\mathbf{I}\cdot\nabla)\frac{1}{r} - (\mathbf{S}\cdot\mathbf{I})\nabla^{2}\frac{1}{r} + \frac{1}{r^{3}}\mathbf{L}\cdot\mathbf{I}\right].\tag{2.87}
\end{equation*}
证明时需要矢量恒等式
\begin{equation*}
\nabla\times(\nabla\times\mathbf{Q}) = \nabla(\nabla\cdot\mathbf{Q}) - \nabla^{2}\mathbf{Q}\tag{2.88}
\end{equation*}
$\mathbf{Q}$ 为矢量场。利用 $\nabla^{2}(1/r)=-4\pi\delta(\mathbf{r})$（$\delta(\mathbf{r})$
为 delta 函数），证明
\begin{equation*}
(\mathbf{S}\cdot\nabla)(\mathbf{I}\cdot\nabla)\frac{1}{r} = -\frac{1}{r^{3}}\left[\mathbf{S}\cdot\mathbf{I} - \frac{3(\mathbf{S}\cdot\mathbf{r})(\mathbf{I}\cdot\mathbf{r})}{r^{2}}\right] - \frac{4\pi}{3}\mathbf{S}\cdot\mathbf{I}\,\delta(\mathbf{r})\tag{2.89}
\end{equation*}
并由此推出
\begin{equation*}
\hat{\mathcal{H}}' = \frac{\mu_0}{4\pi}2g_I\mu_{\mathrm{B}}\mu_{\mathrm{N}}\mathbf{I}\cdot\left[\frac{\mathbf{L}}{r^{3}} - \frac{\mathbf{S}}{r^{3}} + \frac{3\mathbf{r}(\mathbf{S}\cdot\mathbf{r})}{r^{5}} + \frac{8\pi}{3}\mathbf{S}\,\delta(\mathbf{r})\right].\tag{2.90}
\end{equation*}
方括号中第一项是电子轨道磁矩与核磁矩的耦合；第二、三项是电子自旋与核磁矩的偶极
耦合；最后一项等于费米接触能（式 2.69）。

\exer{2.9} 铂的磁化率为 $2.61\times10^{-4}$，密度 $21450\ \mathrm{kg\,m^{-3}}$，
相对原子质量 $195.09\ \mathrm{g\,mol^{-1}}$。计算其摩尔磁化率（m$^3$\,mol$^{-1}$）
与质量磁化率（m$^3$\,kg$^{-1}$）。用附录 A 把这些结果换算成 cgs 单位：摩尔磁化率
（emu\,mol$^{-1}$）与质量磁化率（emu\,g$^{-1}$）。要理解金属铂的磁化率为何不随
温度变化（与居里定律相反），见第七章。

\exer{2.10} 考虑由式 2.32 的配分函数描述的一组自旋。亥姆霍兹自由能密度为
$F=U-TS=-nk_{\mathrm{B}}T\log Z$，其中 $U=-\mathbf{M}\cdot\mathbf{B}$ 是势能密度，
$S$ 是熵密度。整理可得
\begin{equation*}
TS = -\mathbf{M}\cdot\mathbf{B} + nk_{\mathrm{B}}T\log Z.\tag{2.91}
\end{equation*}
证明：温度不变而磁场改变 $\delta\mathbf{B}$ 时，
\begin{equation*}
T\,\delta S = -\mathbf{B}\cdot\delta\mathbf{M}\tag{2.92}
\end{equation*}
从而
\begin{equation*}
\delta F = \delta U - T\,\delta S = -\mathbf{M}\cdot\delta\mathbf{B}.\tag{2.93}
\end{equation*}

% 第3章 环境（Environments）
\chapter{环境}

上一章我们看到：把稀土离子当作完全自由的离子——既不彼此相互作用、也不与周围环境
相互作用——就能推断许多含稀土晶体的磁性质。然而对某些晶体中的磁性离子，通常不能
忽略这类相互作用；对许多材料，它们很大且重要。本章考虑原子与其近邻环境之间的
相互作用；下一章讨论晶体中磁性原子与相邻磁性原子之间的直接磁相互作用。

\section{晶体场}

要理解晶体造成的局域环境对原子能级的影响，首先需要回顾原子轨道的形状。s、p、d
轨道电子密度的角分布示于图 3.1。该图只画出各轨道波函数的角度部分，还有径向部分
（详见附录 C）。只有 s 轨道是球对称的，其余轨道都有显著的角依赖。这一点至关重要，
因为局域环境往往不是球对称的，不同的轨道于是会有不同的表现。

\begin{figure}
\centering
\includegraphics[width=0.8\textwidth]{fig3_01.pdf}
\caption{s、p、d 轨道的角分布。$\mathrm{d}_{z^2}$ 与 $\mathrm{d}_{x^2-y^2}$
能级合称 $e_g$ 能级；$\mathrm{d}_{xy}$、$\mathrm{d}_{xz}$、$\mathrm{d}_{yz}$
能级合称 $t_{2g}$ 能级。$\mathrm{d}_{z^2}$ 轨道有时也记作
$\mathrm{d}_{3z^2-r^2}$。}
\end{figure}

\subsection{晶体场的起源}

晶体场（crystal field）是由晶体中相邻原子产生的电场。在晶体场理论中，相邻轨道
被模型化为负的点电荷；对这一近似的改进是配体场理论（ligand field theory）——它
本质上是分子轨道理论的推广，关注中心离子的 d 轨道及其与周围离子（配体，ligands）
轨道的交叠。晶体场效应的大小与性质关键取决于局域环境的对称性。一个常见的情形是
八面体环境：在许多过渡金属化合物中，过渡金属离子位于八面体中心，每个顶点上是一个
氧之类的离子。此时晶体场主要来自氧轨道中带负电的电子的静电排斥。八面体与四面体
环境的示意图见图 3.2。

\begin{figure}
\centering
\includegraphics[width=0.5\textwidth]{fig3_02.pdf}
\caption{处于（a）八面体与（b）四面体环境中的金属原子 M。八面体环境常见于许多过渡
金属氧化物：氧阴离子位于八面体顶点，金属原子居中。四面体环境可方便地取立方体的
相间顶点来描述（如图）。}
\end{figure}

d 轨道分为两类：指向 $x$、$y$、$z$ 轴之间的 $t_{2g}$ 轨道（即 $\mathrm{d}_{xy}$、
$\mathrm{d}_{xz}$、$\mathrm{d}_{yz}$）和沿轴指向的 $e_g$ 轨道（沿 $z$ 轴有瓣的
$\mathrm{d}_{z^{2}}$ 轨道，以及沿 $x$、$y$ 轴都有瓣的 $\mathrm{d}_{x^{2}-y^{2}}$
轨道）。设想把一个含十个 d 电子的阳离子放到半径为 $r$ 的均匀带负电球的中心。
在这个球对称环境中五个 d 轨道全部简并（当然电荷的存在抬高了整个系统的能量）。
现在设想球上的电荷聚集成六个分立的点电荷，各位于一个正八面体的顶点、仍在球面上。
所有 d 轨道的总电子能量不变，但 d 轨道不再简并——这就造出了八面体环境。

为了说明环境对不同轨道的影响不同，看图 3.3：它以平面图（图 3.2(a) 在 $xy$ 面上的
投影）画出八面体环境中的两个 d 轨道。晶体场主要由相邻原子的 p 轨道产生。显然
$\mathrm{d}_{xy}$ 轨道（图 3.3(a)）与相邻 p 轨道的交叠比 $\mathrm{d}_{x^{2}-y^{2}}$
（图 3.3(b)）小，因而静电能量更低。

\begin{figure}
\centering
\includegraphics[width=0.45\textwidth]{fig3_03.pdf}
\caption{晶体场源于静电相互作用。八面体环境中 $\mathrm{d}_{xy}$ 轨道（a）的能量比
$\mathrm{d}_{x^{2}-y^{2}}$（b）低。}
\end{figure}

在八面体环境中，相邻正电荷聚于 $(\pm r,0,0)$、$(0,\pm r,0)$、$(0,0,\pm r)$，指向
$x$、$y$、$z$ 轴之间的三个轨道 $\mathrm{d}_{xy}$、$\mathrm{d}_{yz}$、$\mathrm{d}_{zx}$
能量降低，而沿 $x$、$y$、$z$ 轴指向的 $\mathrm{d}_{z^2}$ 与 $\mathrm{d}_{x^2-y^2}$
能量升高。五个能级于是按图 3.4(a) 劈裂：三重 $t_{2g}$ 能级下降，二重 $e_g$ 能级
上升。

\begin{figure}
\centering
\includegraphics[width=0.5\textwidth]{fig3_04.pdf}
\caption{（a）八面体与（b）四面体环境中的晶体场。}
\end{figure}

若局域环境不是八面体对称，晶体场甚至可能反向作用。例如四面体环境中，沿轴指向的
轨道恰好最大程度避开了位于（描写四面体的）立方体四个顶点上的原子的电荷密度
（见图 3.2(b)），故此时二重 $e_g$ 能级更低（图 3.4(b)）。

\begin{figure}
\centering
\includegraphics[width=0.45\textwidth]{fig3_05.pdf}
\caption{$3d^{6}$ 离子（如 $\mathrm{Fe^{2+}}$）（a）高自旋（弱场）与（b）低自旋
（强场）情形的电子组态。}
\end{figure}

处理 3d 电子不满的过渡金属离子时，已有的电子先填最低的（此时即 $t_{2g}$）能级再填
$e_g$ 能级。但填充的确切次序取决于晶体场能量与“把两个电子放进同一轨道”的库仑
代价之间的竞争——后者称为成对能（pairing energy），通常为正。若晶体场能量小于
成对能（弱场情形），电子加入时先逐个占据各轨道，然后才有轨道被双占据；反之若
晶体场能量大于成对能（强场情形），电子会先双重占据低能轨道，然后再去攀登高能
轨道的“险峰”。这些情形示于图 3.5。

八面体环境中，填入 1、2、3、8、9、10 个电子时方式毫无歧义；有趣的是 4、5、6、7
个电子的情形。下面以六电子为例（对应 $\mathrm{Fe^{2+}}$ 离子，见图 3.5）。

\begin{example}{3.1}
$\mathrm{Fe^{2+}}$ 离子有 $3d^{6}$ 壳层，成对能为正。弱场情形：每个轨道先填一个，
还剩一个电子，它只好不情愿地与某个 $t_{2g}$ 电子配对。结果四个未成对电子、
$S=2$（图 3.5(a)），称为高自旋组态。强场情形：六个电子全被塞进三个 $t_{2g}$
轨道，$e_g$ 轨道空着；没有未成对电子，$S=0$（图 3.5(b)），称为低自旋组态。
\end{example}

在某些含 $\mathrm{Fe^{2+}}$ 的材料中，可以用温度、压强甚至光照在低自旋与高自旋
组态之间引发自旋转变。

\subsection{轨道淬灭}

3d 离子的预期磁基态列于表 3.1。按 2.5 节的洪特规则计算 S、L、J 是轻而易举的事。
但如该节所述，预言的磁矩 $g_J[J(J+1)]^{1/2}$ 并不总与实验一致——例外是 $3d^{5}$
与 $3d^{10}$：它们的电子壳层半满或全满，$L=0$。

\begin{table}
\centering
\caption{用洪特规则得到的 3d 离子磁基态。对每个离子列出壳层组态与预言的基态
S、L、J，并给出按洪特规则计算的 $p=\mu_{\mathrm{eff}}/\mu_{\mathrm{B}}$。此值记作
$p_1=g_J[J(J+1)]^{1/2}$；下一列是由含相应离子的顺磁盐测量得到的实验值 $p_{\mathrm{exp}}$。
它与 $p_2=2[S(S+1)]^{1/2}$ 符合好得多——后者假定轨道淬灭，即 L=0、J=S、$g_J=2$。}
\small
\begin{tabular}{lcccccccc}
\toprule
离子 & 壳层 & $S$ & $L$ & $J$ & 谱项 & $p_1$ & $p_{\mathrm{exp}}$ & $p_2$\\
\midrule
$\mathrm{Ti^{3+}},\mathrm{V^{4+}}$ & $3d^{1}$ & $\frac{1}{2}$ & 2 & $\frac{3}{2}$ & ${}^{2}D_{3/2}$ & 1.55 & 1.70 & 1.73\\
$\mathrm{V^{3+}}$ & $3d^{2}$ & 1 & 3 & 2 & ${}^{3}F_{2}$ & 1.63 & 2.61 & 2.83\\
$\mathrm{Cr^{3+}},\mathrm{V^{2+}}$ & $3d^{3}$ & $\frac{3}{2}$ & 3 & $\frac{3}{2}$ & ${}^{4}F_{3/2}$ & 0.77 & 3.85 & 3.87\\
$\mathrm{Mn^{3+}},\mathrm{Cr^{2+}}$ & $3d^{4}$ & 2 & 2 & 0 & ${}^{5}D_{0}$ & 0 & 4.82 & 4.90\\
$\mathrm{Fe^{3+}},\mathrm{Mn^{2+}}$ & $3d^{5}$ & $\frac{5}{2}$ & 0 & $\frac{5}{2}$ & ${}^{6}S_{5/2}$ & 5.92 & 5.82 & 5.92\\
$\mathrm{Fe^{2+}}$ & $3d^{6}$ & 2 & 2 & 4 & ${}^{5}D_{4}$ & 6.70 & 5.36 & 4.90\\
$\mathrm{Co^{2+}}$ & $3d^{7}$ & $\frac{3}{2}$ & 3 & $\frac{9}{2}$ & ${}^{4}F_{9/2}$ & 6.63 & 4.90 & 3.87\\
$\mathrm{Ni^{2+}}$ & $3d^{8}$ & 1 & 3 & 4 & ${}^{3}F_{4}$ & 5.59 & 3.12 & 2.83\\
$\mathrm{Cu^{2+}}$ & $3d^{9}$ & $\frac{1}{2}$ & 2 & $\frac{5}{2}$ & ${}^{2}D_{5/2}$ & 3.55 & 1.83 & 1.73\\
$\mathrm{Zn^{2+}}$ & $3d^{10}$ & 0 & 0 & 0 & ${}^{1}S_{0}$ & 0 & 0 & 0\\
\bottomrule
\end{tabular}
\end{table}

偏差的原因是：对 3d 离子，晶体场相互作用远强于自旋--轨道相互作用。于是 2.5 节
表述的洪特第三条规则——它基于“自旋--轨道相互作用是仅次于静电效应的最重要能量项”
——实际上是错的。数据显示这些体系转而选择 $L=0$（从而 $J=S$、$g_J=2$）的基态，
即
\begin{equation*}
\mu_{\mathrm{eff}} = 2\mu_{\mathrm{B}}\sqrt{S(S+1)}.\tag{3.1}
\end{equation*}
如表 3.1 所示，这一式子与实验符合得好得多。这个效应称为轨道淬灭（orbital
quenching），此时称轨道矩被淬灭（quenched）。对 4f 离子，轨道在空间上离核较近、
藏在 5s 与 5p 壳层之下，晶体场项的重要性小得多，洪特第三条规则得以遵守。更高过渡金属
离子（4d、5d 系列）的情形不那么明确：较重的离子自旋--轨道劈裂更大，晶体场与自旋--
轨道相互作用的效应可能不相上下。

对自旋--轨道相互作用可以忽略的 3d 离子，让我们更细致地考察轨道淬灭。八面体环境
中的晶体场等于常数加上正比于
$x^{4}+y^{4}+z^{4}-\frac{3}{5}r^{4}+O(r^{6}/a^{6})$ 的项（见习题 3.2），因此是实
函数。晶体场哈密顿量的典型特点是可以用实函数（不含微分算符）表达，故其本征函数
全为实函数。而总角动量算符 $\hat{\mathbf{L}}$ 是厄米的，本征值为实，但算符本身是
纯虚的。\fn{1}{式 C.23 表明角动量算符 $\hat{\mathbf{L}}=-\mathrm{i}\hbar\,\mathbf{r}\times\nabla$ 是纯虚的。}于是，晶体场劈裂所实现的非简并基态 $|0\rangle$（非简并——这样才
没法与另一个基态做线性组合来“耍花招”）必须是实函数；因此 $\langle 0|\hat{\mathbf{L}}|0\rangle$
因 $\hat{\mathbf{L}}$ 纯虚而为纯虚数；但 $\hat{\mathbf{L}}$ 厄米，$\langle 0|\hat{\mathbf{L}}|0\rangle$
又必须是纯实数。既纯实又纯虚，只能是
\begin{equation*}
\langle 0|\hat{\mathbf{L}}|0\rangle = 0\tag{3.2}
\end{equation*}
即非简并态的轨道角动量各分量全部被淬灭。

这个结果还可换个角度看。$\hat{L}_z$ 的本征值为 $m_l$ 的轨道态有方位角依赖
$e^{\mathrm{i}m_l\phi}$。本征函数须为实，这要求把 $\pm m_l$ 的态做线性组合，\fn{2}{这是为了构造方位角依赖如 $\cos(m_l\phi)$ 或 $\sin(m_l\phi)$ 的态。注意 $\cos(m_l\phi)=\frac{1}{2}(e^{\mathrm{i}m_l\phi}+e^{-\mathrm{i}m_l\phi})$，$\sin(m_l\phi)=\frac{1}{2\mathrm{i}}(e^{\mathrm{i}m_l\phi}-e^{-\mathrm{i}m_l\phi})$。}所得态的 $\langle L_z\rangle=0$。轨道淬灭的一种半经典解读是：轨道角动量在晶体场中进动，
大小不变而各分量平均为零。

实际上轨道角动量可能没有被完全淬灭，因为即便在 3d 离子中自旋--轨道相互作用也非
全然可忽略。它可以作为微扰加入，混入角动量非零的态，得到淬灭的基态但 $g$ 因子
不完全等于纯自旋值 2——与 2 的差别反映混入的 $L>0$ 态。这还会使 $g$ 因子略有
各向异性：其值随磁场相对晶体轴的方向略有不同。

\begin{figure}
\centering
\includegraphics[width=0.42\textwidth]{fig3_06.pdf}
\caption{$\mathrm{Mn^{3+}}$（$3d^{4}$）的扬--特勒效应。八面体配合物（左）可以畸变
（右），使 $t_{2g}$ 与 $e_g$ 能级劈裂。畸变之所以降低能量，是因为单占据的 $e_g$ 能级被降低：$\mathrm{d}_{xz}$、$\mathrm{d}_{yz}$ 能级下降所省的能量恰被
$\mathrm{d}_{xy}$ 能级的升高所抵偿。}
\end{figure}

\subsection{扬--特勒效应}

到目前为止我们假定：只需弄清局域环境具有何种对称性，推断出电子结构，再按填充
能级的电子数推出磁性质。然而有时磁性质本身竟能反过来影响局域环境的对称性！原因
在于：如八面体这样的结构自发畸变（见图 3.6）有时在能量上有利——弹性能量的增加
恰好被畸变带来的电子能量节省所平衡。这一现象称为扬--特勒效应（Jahn--Teller
effect）。例如八面体环境中的 $\mathrm{Mn^{3+}}$ 离子（组态 $3d^{4}$）就表现出这种
行为（见图 3.6）；相反，$\mathrm{Mn^{4+}}$ 离子（$3d^{3}$）不表现出，因为畸变
没有电子能量的净降低。

为在唯象层面描述这一效应，设系统的畸变可用参数 Q 量化——沿某合适简正模坐标的
畸变距离。它带来一个二次型的能量代价：
\begin{equation*}
E(Q) = \frac{1}{2}M\omega^{2}Q^{2},\tag{3.3}
\end{equation*}
M 与 $\omega$ 分别是阴离子质量和该简正模对应的角频率。此关系绘于图 3.7(a)。显然
最小畸变能为零，出现在 Q=0（无畸变）处。

畸变还会提升某些轨道的能量、降低另一些的能量。若所有轨道或全满或全空，这无关
紧要——总能量就是式 3.3。但对部分填充的轨道，这个效应可能极其重要：系统可以实现
总能量的净降低。电子能量对 Q 的依赖可能相当复杂，但可写成对 Q 的泰勒级数；只要
畸变小，只保留线性项是合理的。设某轨道的能量含 $+AQ$ 或 $-AQ$ 项（对应电子能量的
升高或降低），A 为适当的正常数。总能量 $E(Q)$ 是电子能量与弹性能量之和：
\begin{equation*}
E(Q) = \pm AQ + \frac{1}{2}M\omega^{2}Q^{2},\tag{3.4}
\end{equation*}
$AQ$ 项符号的两种选择给出两条曲线，绘于图 3.7(b)。只看其中一条，用
$\partial E/\partial Q=0$ 求该轨道的最小能量，得
\begin{equation*}
Q_0 = \frac{A}{M\omega^{2}}\tag{3.5}
\end{equation*}
及最小能量 $E_{\min}=-A^{2}/2M\omega^{2}<0$。若只有该轨道被占据，系统便可因自发
畸变而净省能量。

\begin{figure}
\centering
\includegraphics[width=0.55\textwidth]{fig3_07.pdf}
\caption{（a）按式 3.3，八面体配合物能量对畸变 Q 的函数。（b）按式 3.4 的情形。}
\end{figure}

迄今讨论的实质上是静态扬--特勒效应：自发畸变固定在八面体的某一特定轴上。但在
较高温度下，畸变可以在不同轴之间切换，产生动态扬--特勒效应。另一类动态效应是
畸变在格点之间的快速跳迁——在例如含 $\mathrm{Mn^{3+}}$ 与 $\mathrm{Mn^{4+}}$
混合离子的材料中很重要。这些现象可以通过它们对磁共振数据的影响来探测。

还有一类效应：某些材料（如 $\mathrm{DyVO_4}$）在低于某临界温度时，各配合物的
扬--特勒畸变会在整个晶体中协同发生，称为协同扬--特勒转变（cooperative Jahn--Teller
transition）。

晶体场的讨论至此告一段落。本章余下部分描述可用于研究磁矩与其环境相互作用的
各种实验技术。

\section{磁共振技术}

磁矩的局域环境由晶体场决定，但自旋--轨道耦合以及与核的超精细相互作用对电子结构
同样重要。这些效应可以用多种磁共振实验技术研究，本章余下部分将描述这些技术。

1.1.2 节已经看到：对磁矩施加磁场可诱发该磁矩以角频率 $|\gamma B|$ 进动，$\gamma$
为旋磁比。磁场中的磁矩系统可在此频率吸收能量，于是可以观察到系统对调到正确频率的
电磁波的共振吸收。这就是磁共振（magnetic resonance）；依共振磁矩的类型不同，实验
形式多样。电磁波的适当频率取决于磁场与磁矩的大小。图 3.8 以各种单位总结了电磁辐
射的频率范围，供下文随时参照。

\begin{figure}
\centering
\includegraphics[width=0.9\textwidth]{fig3_08.pdf}
\caption{电磁波谱。光子能量以温度 $T=E/k_{\mathrm{B}}$（K）与能量 $E$（eV）标出。
相应频率 $f$ 以 Hz 标出，也以光谱学常用的 $\mathrm{cm}^{-1}$ 标出。$\mathrm{cm}^{-1}$
标尺上标了若干常见分子跃迁与激发（分子转动与振动的典型范围，以及 C--H 弯曲与伸缩
模式），典型 $\pi$ 与 $\sigma$ 键的能量亦然。光子波长 $\lambda=c/f$（$c$ 为光速）。
温度标尺上特地标出 $T_{\mathrm{CMB}}$（宇宙微波背景温度）、常压下液氦（$\mathrm{He^4}$）
与液氮（$\mathrm{N_2}$）的沸点以及室温。图中其他缩写：IR=红外，UV=紫外，R=红，
G=绿，V=紫。字母 H 标记 13.6 eV——氢 1s 电子能量的大小。频率轴上还标出电磁波谱
主要区域的名称：radio（无线电）、microwave（微波）、infrared（红外，含“近”与
“远”）、optical（光学）与 UV（紫外）。}
\end{figure}

\subsection{核磁共振}

最常用的磁共振形式是核磁共振（nuclear magnetic resonance，NMR）。这一技术在医学
成像中大量使用，为避讳“核”这个吓人的字眼而称磁共振成像（MRI）。有人得知人体
平均含有超过 $10^{27}$ 个核可能会大惊失色：按质量计，我们 $\sim$99.98\% 是核！
MRI 测量的是病人体内质子的 NMR，提供安全、无创的精细断层成像。

做任何 NMR 实验都需要一个自旋非零的核。常研究的核包括 ${}^{1}$H（质子）、
${}^{2}$H（氘核）与 ${}^{13}$C。简单 NMR 实验（见图 3.9）把样品放进一个线圈内，
线圈安装在磁体极靴之间。磁体产生沿某方向（如 $z$ 方向）的磁场 $B_0$。我们已经
知道核角动量的 $z$ 分量 $m_I$ 只能取 $-I$ 到 $I$ 之间的整数间隔值。

\begin{figure}
\centering
\includegraphics[width=0.45\textwidth]{fig3_09.pdf}
\caption{NMR 实验示意图。样品置于射频（RF）线圈内，线圈产生振荡的射频场；磁体提供
高度均匀的静磁场。静场 $\mathbf{B}_0$ 与振荡场 $\mathbf{B}_1$ 互相垂直。真实实验中
样品比图中所示小得多，从而感受到 RF 线圈的均匀场。}
\end{figure}

核的能量是磁矩\fn{3}{核带正电、电子带负电，这正是下列表达式中符号不同的原因：$\boldsymbol{\mu}=+g_{\mathrm{N}}\mu_{\mathrm{N}}\mathbf{I}$（核）与 $\boldsymbol{\mu}=-g_J\mu_{\mathrm{B}}\mathbf{J}$（电子）。}$\boldsymbol{\mu}$ 在磁场 $B_0$ 中的能量：$E=-\boldsymbol{\mu}\cdot\mathbf{B}_0$，故
\begin{equation*}
E = -g_{\mathrm{N}}\mu_{\mathrm{N}}m_IB_0.\tag{3.6}
\end{equation*}
这是 $2I+1$ 个等间距能级的阶梯，间距 $g_{\mathrm{N}}\mu_{\mathrm{N}}B_0$；Na 或 Cu
中 $I=\frac{3}{2}$ 的情形示于图 3.10(a)。\fn{4}{此处先设 J=0，从而忽略电子自旋。}用射频场激发相邻能级对之间的跃迁，正是核磁共振的
基础。RF 场 $B_1$ 沿 $x$ 方向施加，对系统的微扰正比于 $B_1\hat{I}_x$。微扰矩阵元
正比于 $\langle m_I'|\hat{I}_x|m_I\rangle$，除非 $m_I'=m_I\pm1$ 否则为零。故允许
跃迁由选择定则
\begin{equation*}
\Delta m_I = \pm 1,\tag{3.7}
\end{equation*}
描述：只有相邻态之间的跃迁可以发生。最常见的装置中，RF 线圈不仅产生激发，还自身
是高品质因数（大 Q 值）调谐回路的一部分。RF 场激发核跃迁时，能量在 RF 回路与样品
之间传递，导致回路 Q 值的微小变化。

\begin{figure}
\centering
\includegraphics[width=0.45\textwidth]{fig3_10.pdf}
\caption{（a）$I=\frac{3}{2}$ 核的四个能级。（b）较简单的 $I=\frac{1}{2}$ 核只有两个
能级，间隔 $\Delta E=g_{\mathrm{N}}\mu_{\mathrm{N}}B_0$；下（上）能级对应核磁矩平行
（反平行）于磁场 $B_0$。}
\end{figure}

实验有两种做法：固定 RF 场频率而改变磁场，或固定磁场而改变 RF 频率。通常采用
前者。关键是要有高度均匀的磁体来产生恒定场 $B_0$；否则样品不同部分处于略不同的
磁场中，在磁场扫描的不同点先后共振，使测得的共振极宽，甚至可能完全被抹平。

\begin{example}{3.2}
对质子 $I=\frac{1}{2}$，$m_I$ 只能取 $\pm\frac{1}{2}$（图 3.10(b)）。系统的两态
间隔 $\Delta E=g_{\mathrm{N}}\mu_{\mathrm{N}}B$，非常小：典型实验室磁场
$B_0\sim1$~T 中质子的能隙 $\sim10^{-7}$ eV；写成 $k_{\mathrm{B}}T$ 只相当于
$\sim1$ mK 的温度。

因此室温下、该磁场处，核平均而言只有极微弱的沿外场排齐的趋势——热随机化能量远
超微弱的排齐能量。若不采用共振技术，核磁性引起的任何效应实际上都无法探测。NMR
以恰好\fn{5}{其实不必完全精确：可以差一个线宽。线宽由多种因素决定，包括相邻核的偶极场造成的核处磁场变化。}合适的角频率 $\omega$ 的 RF 信号微扰系统：
\begin{equation*}
\hbar\omega = \Delta E\tag{3.8}
\end{equation*}
从而可以激发跨越这一能隙的跃迁。为什么需要射频？下例说明。
\end{example}

\begin{example}{3.3}
磁场 1 T 时，频率为
\begin{equation*}
\nu = \frac{\omega}{2\pi} = \frac{g_{\mathrm{N}}\mu_{\mathrm{N}}B_0}{2\pi\hbar} = 42.58\ \mathrm{MHz}\tag{3.9}
\end{equation*}
处于电磁波谱的射频区。
\end{example}

这意味着可以使用振荡器与线圈（相比之下，电子自旋共振实验须用 Gunn 二极管与微波
波导，见 3.2.2 节）。于是，在恰好合适的 RF 激发频率处系统将吸收能量——这就是核磁
共振：RF 激发频率与核能级间隔之间的共振。现代 NMR 谱仪使用 12--15 T 的磁体，RF
频率在 $\sim$500--650 MHz 范围，仍在射频区。

给定磁场下 NMR 共振频率可以从 $\gamma B/2\pi$ 略微上下偏移，取决于核的化学环境。
这些化学位移（chemical shifts，典型为百万分之几）源于绕核运动的电子对核的轻微
屏蔽。给定化学环境中核被屏蔽的程度是已知的，因此可以对特定分子“指纹化”。NMR
谱线还会因与相邻核的磁耦合而分裂，称为自旋--自旋耦合（spin--spin coupling），经由
电子传递的间接接触超精细相互作用发生。这同样给出核环境的信息。

为更细致地理解我们诱导的跃迁，考虑双能级自旋系统（$I=\frac{1}{2}$）如何从 RF
线圈吸收能量。记下能级为 $-$、上能级为 $+$。诱导两能级间跃迁的概率与跃迁方向无关，
因为振荡 RF 场对应的微扰哈密顿量 $\hat{H}\propto\hat{I}_x$ 是厄米的：
\begin{equation*}
|\langle\psi_-|\hat{\mathcal{H}}|\psi_+\rangle|^{2} = |\langle\psi_+|\hat{\mathcal{H}}|\psi_-\rangle|^{2}.\tag{3.10}
\end{equation*}
于是所谓受激跃迁（+ 与 $-$ 之间）的概率密度与方向无关，以速率 W 发生，W 正比于
所用 RF 功率。设 $t$ 时刻下能级有 $N_-(t)$ 个自旋，则每单位时间 $WN_-(t)$ 个被
激发到上能级；上能级有 $N_+(t)$ 个时，每单位时间 $WN_+(t)$ 个被激发到下能级。故
\begin{equation*}
\frac{\mathrm{d}N_+(t)}{\mathrm{d}t} = WN_-(t) - WN_+(t)\tag{3.11}
\end{equation*}
\begin{equation*}
\frac{\mathrm{d}N_-(t)}{\mathrm{d}t} = WN_+(t) - WN_-(t)\tag{3.12}
\end{equation*}
相减得
\begin{equation*}
\frac{\mathrm{d}}{\mathrm{d}t}(N_+(t)-N_-(t)) = -2W(N_+(t)-N_-(t)),\tag{3.13}
\end{equation*}
解为
\begin{equation*}
N_+(t)-N_-(t) = (N_+(0)-N_-(0))\,\mathrm{e}^{-2Wt}\tag{3.14}
\end{equation*}
初始的布居差在受激电磁跃迁驱动下指数趋于零。定义
\begin{equation*}
n(t) = N_+(t)-N_-(t),\tag{3.15}
\end{equation*}
式 3.14 成为
\begin{equation*}
n(t) = n(0)\,\mathrm{e}^{-2Wt}\tag{3.16}
\end{equation*}

自旋系统吸收或发射电磁能量的速率容易计算：每次跃迁交换能量 $\hbar\omega$。系统
在 $t$ 时刻的能量为
\begin{equation*}
E(t) = N_-E_- + N_+(E_-+\hbar\omega),\tag{3.17}
\end{equation*}
$E_-$ 是下能级能量。此式可整理成常数项加 $\frac{1}{2}\hbar\omega\,n(t)$。于是能量
吸收速率为
\begin{equation*}
\frac{\mathrm{d}E}{\mathrm{d}t} = -W\hbar\omega\,n(t)\tag{3.18}
\end{equation*}
它与 $n(t)$ 成正比，随上、下能级布居逐渐相等而以时间常数 $1/2W$ 趋于零。这也表明
吸收能量需要 $n(t)\neq0$，即布居差。布居差由系统与热模式的相互作用产生，使系统
的极化回到玻尔兹曼概率：
\begin{equation*}
\left(\frac{N_+}{N_-}\right)_0 = \mathrm{e}^{-\hbar\omega/k_{\mathrm{B}}T}\tag{3.19}
\end{equation*}
下标 0 表示系统的平衡值。吸收能量后极化可能改变；若核自旋系统与样品的热运动及
激发相互作用，自旋系统的极化可以回到平衡值。若自旋系统的极化被快速的电磁跃迁
降为零，关闭电磁跃迁后自旋系统需要一段时间“恢复”；这段时间称为 $T_1$，度量
自旋系统与其环境相互作用的时间常数。$T_1$ 是自旋--晶格弛豫时间（spin--lattice
relaxation time，也称纵向弛豫时间）。预期极化按
\begin{equation*}
n(t) = n_0\left(1-\mathrm{e}^{-t/T_1}\right).\tag{3.20}
\end{equation*}
恢复。

\begin{figure}
\centering
\includegraphics[width=0.5\textwidth]{fig3_11.pdf}
\caption{式 3.25 给出的电磁能量吸收速率随跃迁速率 W 的变化。W 小时正比于 W；W 大时
饱和于由 $1/T_1$ 决定的值。}
\end{figure}

现在把两个过程——受激跃迁与弛豫——放在一起。合并方程得
\begin{equation*}
\frac{\mathrm{d}}{\mathrm{d}t}n(t) = -2Wn(t) + \frac{n_0-n(t)}{T_1},\tag{3.21}
\end{equation*}
在
\begin{equation*}
\frac{\mathrm{d}}{\mathrm{d}t}n(t) = 0.\tag{3.22}
\end{equation*}
时有稳态解
\begin{equation*}
n(t) = \frac{n_0}{1+2WT_1},\tag{3.23}
\end{equation*}
电磁能量的稳态吸收速率为
\begin{align*}
\frac{\mathrm{d}E}{\mathrm{d}t} &= n(t)\hbar\omega W\tag{3.24}\\
&= n_0\hbar\omega\frac{W}{1+2WT_1}.\tag{3.25}
\end{align*}
此关系绘于图 3.11。低幅 RF 微扰场下，能量吸收速率正比于 W；大幅 RF 微扰场下，
它稳定在与 $T_1$ 成正比、与 W 的精确大小无关的水平。这称为饱和（saturation）。

电磁能量吸收速率正比于上、下能级的布居差，因此可用它监测布居差及其时间依赖——
当然这种观测本身必然在一定程度上扰动布居差。式 3.25 与图 3.11 表明：提高 RF 功率
只能到某个程度为止。式 3.25 还表明 $\mathrm{d}E/\mathrm{d}t$ 正比于 $\hbar\omega$：
能量更高（频率更高）的光子给出更大信号；但这当然意味着需要更大的磁场使 $g\mu_{\mathrm{N}}B$
匹配 $\hbar\omega$。这对 ESR（见下节 3.2.2）反而是优势：激发跨越更大的电子能隙
需要很强健的微波光子，从而带来更高的灵敏度。由于电子磁矩大于核磁矩，ESR 即便用
比 NMR 更小的磁场，也仍能使用频率更高的光子。
\bio{Felix Bloch\\（1905--1983）}

回到 NMR，更细致地考虑弛豫过程。磁场中自旋存在很弱的极化（弱是因为能隙
$\Delta E\ll k_{\mathrm{B}}T$）。设接通静磁场 $B_0$ 后（回顾 $B_0$ 沿 $z$ 轴）
极化取平衡值 $M_0$。射频激发的效果是逐渐破坏这一磁化强度；而如前所述，切断磁场
后，磁化强度通过与周围的弱相互作用以时间常数 $T_1$ 弛豫回平衡值。因此预期
\begin{equation*}
\frac{\mathrm{d}M_z}{\mathrm{d}t} = \frac{M_0-M_z}{T_1}.\tag{3.26}
\end{equation*}
$T_1$ 弛豫必须涉及与晶格的相互作用，因为能量必须与晶格交换：改变 $M_z$ 有能量
后果——自旋沿 $z$ 方向弛豫会改变它相对外加场的取向，从而改变其能量。

$M_x$ 与 $M_y$ 分量应为零；若不为零，将以时间常数 $T_2$ 弛豫回零：
\begin{equation*}
\frac{\mathrm{d}M_x}{\mathrm{d}t} = -\frac{M_x}{T_2},\qquad
\frac{\mathrm{d}M_y}{\mathrm{d}t} = -\frac{M_y}{T_2}.\tag{3.27}
\end{equation*}
$T_2$ 是自旋--自旋弛豫时间，是失相（dephasing）的特征时间，对应自旋系统不同部分
之间的相互作用；也可能由 $B_0$ 的不均匀性造成。它导致被观测自旋与邻近自旋相互
作用产生的进动频率差异。改变 $M_x$ 或 $M_y$ 没有能量后果，因为外加场沿 $M_z$。

外加磁场还引起自旋进动，于是 $M_x$、$M_y$、$M_z$ 的方程为
\begin{equation*}
\frac{\mathrm{d}M_x}{\mathrm{d}t} = \gamma(\mathbf{M}\times\mathbf{B})_x - \frac{M_x}{T_2},\tag{3.28}
\end{equation*}
\begin{equation*}
\frac{\mathrm{d}M_y}{\mathrm{d}t} = \gamma(\mathbf{M}\times\mathbf{B})_y - \frac{M_y}{T_2},\tag{3.29}
\end{equation*}
\begin{equation*}
\frac{\mathrm{d}M_z}{\mathrm{d}t} = \gamma(\mathbf{M}\times\mathbf{B})_z + \frac{M_0-M_z}{T_1},\tag{3.30}
\end{equation*}
称为布洛赫方程（Bloch equations）。它们于 1946 年导出，可在许多有意义的情形下
求解，常用旋转参照系法——把坐标换成在 $xy$ 平面内以共振频率旋转的坐标系。

自旋--自旋弛豫时间 $T_2$ 可用如下技术测量：在与 $B_0$ 热平衡、$B_1$ 关断时，
存在沿 $z$ 的弱磁化强度。加一个时长 $t_{\mathrm{p}}$ 的短 RF 脉冲（频率调到接近
自旋共振频率），使自旋转过角度 $\gamma B_1t_{\mathrm{p}}$（$B_1$ 为 RF 振幅）。
角度可通过改变 $t_{\mathrm{p}}$ 调节；取 $t_{\mathrm{p}}=\pi/2\gamma B_1$ 就得到
所谓的 $90^{\circ}$ 脉冲。脉冲幅度与时长选得使磁化强度恰好进动一小段时间、转入
$xy$ 平面。关断 $B_1$ 后，磁化强度在静场 $\mathbf{B}_0$ 中进动，在 $xy$ 平面内以
频率 $\gamma B_0$ 稳定旋转，在线圈中产生感应电压。若没有 $T_2$ 弛豫过程（自旋--
自旋弛豫），振荡将永远持续；正是 $T_2$ 使振荡衰减，从而给出 $T_2$ 的测量。线圈
电压的相应衰减称为自由感应衰减（free induction decay，见图 3.12）。

\begin{figure}
\centering
\includegraphics[width=0.55\textwidth]{fig3_12.pdf}
\caption{模拟的自由感应衰减信号。振荡因自旋--自旋弛豫（$T_2$ 过程）随时间衰减。
数据呈现干涉图样，来自三个不等价核——它们因在晶体中所处格位不同而感受略微不同的
磁场。}
\end{figure}

这一方法的问题在于它同时测到了磁体不均匀性导致的弛豫——样品不同部位进动频率
略异。需要一种方法把真正的自旋--自旋弛豫与磁体不均匀性弛豫分开。办法之一是自旋
回波（spin echo）技术，其原理示于图 3.13：先用 $90^{\circ}$ 脉冲（如前把自旋转入
$xy$ 平面），随后加 $180^{\circ}$ 脉冲（把自旋转过 $180^{\circ}$）。可以把这一
效应比作一场马拉松：所有选手稳步跑，但有的快有的慢。发令枪（$90^{\circ}$ 脉冲）
后选手们成群出发，但不久体力好的渐渐领先，其余的人渐渐落后。然后突然在时刻 $\tau$，
规则改变（$180^{\circ}$ 脉冲）：选手们被告知掉头跑回起点。此时慢者反而占优——
他们离起点更近。由对称性易见：无论各自速度如何，所有选手将在时刻 $2\tau$ 同时
回到起点。同理，所有自旋在 $2\tau$ 后重新对齐，不论各自的失相。这就是自旋回波。

\begin{figure}
\centering
\includegraphics[width=0.46\textwidth]{fig3_13.pdf}
\caption{自旋回波效应。（a）平衡磁化强度初始沿 $z$ 方向、平行于 $B_0$。（b）称该时刻为 $t=0$：沿 $x$ 轴的 $90^{\circ}$ 脉冲把自旋转入 $xy$ 平面。（c）由于 $z$ 轴方向的恒定场 $B_0$，自旋在 $xy$ 平面内进动，但因磁场不均匀而速率略异。（d）它们逐渐相互失相。（e）$t=\tau$ 时沿 $x$ 轴的 $180^{\circ}$ 脉冲使自旋绕 $x$ 轴转过 $180^{\circ}$；随后在 $B_0$ 中进动时次序已颠倒。（f）$2\tau$ 时它们重新聚齐，产生自旋回波信号——前提是期间没有发生其他弛豫过程。图中画的是两脉冲间隔较短的情形；若间隔长、自旋在 $180^{\circ}$ 脉冲前已转过若干圈，效应同样成立。}
\end{figure}

于是实验中可用自旋回波消除磁场不均匀导致自旋以不同速率进动的问题，化学位移造成
的场展宽也被消除。但自旋--自旋相互作用——来自相邻核的含时涨落随机磁场——的效应
无法消除：这一弛豫机制无法重聚焦，回波幅度将随 $\tau$ 衰减。回波信号的 NMR 强度
应遵循 $I(2\tau)=I(0)e^{-2\tau/T_2}$，由此可测得真正的（自旋--自旋耦合导致的）
$T_2$。

$T_1$ 可用类似方法测量，只是先加 $180^{\circ}$ 脉冲：它使磁化强度从 $+\hat{\mathbf{z}}$
转到 $-\hat{\mathbf{z}}$，随后以时间常数 $T_1$ 弛豫回 $+\hat{\mathbf{z}}$，但不
进动。于是 $180^{\circ}$ 脉冲后 $\tau$ 时刻
\begin{equation*}
\mathbf{M}(\tau) = \hat{\mathbf{z}}\,M_z(\tau) = \hat{\mathbf{z}}\,M_0(1-2\,\mathrm{e}^{-\tau/T_1}).\tag{3.31}
\end{equation*}
$180^{\circ}$ 脉冲后 $\tau$ 时刻再用 $90^{\circ}$ 脉冲把磁化强度转入 $xy$ 平面，
开始初始幅度为 $M_z(\tau)$ 的自由感应衰减。于是测量自由感应衰减初始幅度对
$180^{\circ}$--$90^{\circ}$ 脉冲间隔时间的依赖，即可得 $T_1$。其他多种脉冲序列
也可用于提取更多类型的信息。

\subsection{电子自旋共振}

可以用电子做与 NMR 类似的实验，称为电子自旋共振（electron spin resonance，ESR），
有时也称电子顺磁共振（electron paramagnetic resonance，EPR）。电子磁矩远大于核
磁矩，进动频率也高得多：典型实验室磁场下，电磁辐射已处于微波波段。实验装置示于
图 3.14：样品置于谐振腔内，微波经波导进入。通常更方便的做法是固定微波频率、扫描
磁场。谐振腔品质因数极高，用以提高微弱 ESR 信号的探测灵敏度。记录微波吸收随外加
磁场的变化。常在磁体提供的静场之外加一组调制线圈，在样品处产生小的振荡磁场，
再配合相敏检测技术进一步提高测量灵敏度。\fn{6}{锁相放大器产生驱动振荡磁场的信号；微波信号馈入锁相放大器，后者滤除任何不在振荡磁场频率处的成分，从而大幅降低噪声。}

\begin{figure}
\centering
\includegraphics[width=0.45\textwidth]{fig3_14.pdf}
\caption{ESR 实验示意图。微波经波导进入谐振腔，共振诱导的微波吸收通过监测腔的
Q 因子来测量。样品必须放在磁体中心——那里场最均匀。}
\end{figure}

ESR 跃迁的选择定则是
\begin{equation*}
\Delta m_J = \pm 1,\tag{3.32}
\end{equation*}
即熟悉的偶极选择定则——我们诱导的是磁偶极跃迁。图 3.15 展示八面体环境中
$J=1$ 离子（如 $\mathrm{Ni^{2+}}$）的这些跃迁。实验在固定频率 $h\nu$ 下扫描磁场：
每当两个相邻能级间隔为 $h\nu$ 就发生一次 ESR 跃迁。图 3.15 表明谱线的位置与数目
可以给出晶体场劈裂的信息。

\begin{figure}
\centering
\includegraphics[width=0.45\textwidth]{fig3_15.pdf}
\caption{$J=1$ 离子（如 $\mathrm{Ni^{2+}}$）的 ESR：（a）无晶体场劈裂（单条 ESR
线）；（b）有晶体场劈裂 $\Delta$（两条 ESR 线）。}
\end{figure}

图 3.16 给出自由 $\mathrm{Mn^{2+}}$ 离子（自旋 $\frac{5}{2}$）的能量随磁场的变化。
磁场 B 下不同 $m_J$ 能级间有许多可能跃迁，大小都是 $g\mu_{\mathrm{B}}B$。因此
ESR 实验或许预期在 $h\nu=g\mu_{\mathrm{B}}B$ 给出的频率 $\nu$ 处观察到单线。但对
含 $\mathrm{Mn^{2+}}$ 离子的固体做真实实验时，即便在立方环境中，晶体场也会使图像
显著复杂化。

不过，核与电子之间的超精细耦合在哈密顿量中给出 $A\mathbf{I}\cdot\mathbf{J}$ 项，
每个 $m_J$ 能级按 $m_I$ 分裂成 $2I+1$ 个超精细能级。ESR 跃迁的选择定则因此是式
3.32 加上
\begin{equation*}
\Delta m_I = 0.\tag{3.33}
\end{equation*}

\begin{figure}
\centering
\includegraphics[width=0.5\textwidth]{fig3_16.pdf}
\caption{自旋 $\frac{5}{2}$ 磁矩（如自由 $\mathrm{Mn^{2+}}$ 离子）的能量随磁场的
变化。按选择定则 $\Delta m_J=\pm1$ 的可能跃迁如图。此图忽略了晶体场。}
\end{figure}

后一定则成立是因为微波只诱导电子能级间的偶极跃迁、而非核能级间的跃迁。\fn{7}{这是因为核能级间距与微波光子能量相比极小。}由于 ESR 实验的典型场下 A 远小于
$g\mu_{\mathrm{B}}B$，只需计算 J 沿 B 的分量，磁场中能级即为
\begin{equation*}
E = g\mu_{\mathrm{B}}Bm_J + Am_Im_J\tag{3.34}
\end{equation*}
（外加磁场中核自旋的能量已忽略——它比超精细能量还小得多）。用选择定则可得共振
应发生在
\begin{align*}
h\nu &= [g\mu_{\mathrm{B}}B(m_J+1) + Am_I(m_J+1)] - [g\mu_{\mathrm{B}}Bm_J + Am_Im_J]\\
&= g\mu_{\mathrm{B}}B + Am_I\tag{3.35}
\end{align*}
处，每条 ESR 线分裂成 $2I+1$ 条超精细线。

\begin{example}{3.4}
自由 $\mathrm{Mn^{2+}}$ 离子 $I=\frac{5}{2}$，故 ESR 谱中有六条线。这一情形示于
图 3.17（晶体中的 $\mathrm{Mn^{2+}}$ 离子因晶体场而更复杂：劈裂非平凡地依赖磁场
与晶轴的夹角）。
\end{example}

\begin{figure}
\centering
\includegraphics[width=0.55\textwidth]{fig3_17.pdf}
\caption{固定磁场下自由 $\mathrm{Mn^{2+}}$ 离子的超精细劈裂。$J=\frac{5}{2}$、
$I=\frac{5}{2}$，六个 $m_J$ 能级各自分裂成六个超精细能级。$\Delta m_J=\pm1$、
$\Delta m_I=0$ 的可能跃迁如图，在 ESR 谱中给出六条不同的线。此图忽略晶体场效应——
它使劈裂强烈依赖于外加磁场相对晶轴的方向。}
\end{figure}

ESR 实验在有用的范围内探测电子自旋，因为晶体场劈裂常在 GHz 频段。顺磁盐的 ESR
研究已经提供了大量关于晶体场与配体场的细致信息。效应往往是各向异性的——磁场
相对晶轴旋转时谱线会移动。

\begin{example}{3.5}
举一个如何获得这类信息的例子：考虑轴对称环境中的 $\mathrm{Ce^{3+}}$（$4f^{1}$）
（见图 3.18）。该离子 L=3、$S=\frac{1}{2}$，自旋--轨道相互作用把它分裂成八重简并
的 $J=\frac{7}{2}$ 能级和较低的六重简并的 $J=\frac{5}{2}$ 能级。晶体场是轴对称的，
按对称性只含 $J_z$ 的偶次幂，于是六重 $J=\frac{5}{2}$ 多重态（manifold，即一组谱线的合称）分裂
成三个二重态，八重 $J=\frac{7}{2}$ 多重态分裂成四个。这些劈裂在稀土体系中远小于
自旋--轨道相互作用。二重态可被磁场进一步分裂——对典型实验室磁场而言，这几乎总是
最小的能量项，通常约 $10^{-4}$ eV。低温下只有基态被占据，观察到的正是基态跃迁。
\end{example}

许多体系的基态像 $\mathrm{Ce^{3+}}$ 一样是二重态。这是克拉默斯定理（Kramers
theorem）的结论：含奇数个电子的系统在无磁场时至少保留二重简并。涉及的态对称为
克拉默斯二重态（Kramers doublets），它们互为时间共轭\fn{8}{即互为复共轭，因而互为时间反演的版本。}\bio{Hendrik A. Kramers\\（1894--1952）}——可以被磁场分裂而不能被电场分裂。$\mathrm{Ce^{3+}}$（$4f^{1}$）的 4f 壳层有奇数个电子，因此是克拉默斯离子。

$\mathrm{Ce^{3+}}$ 的基态是克拉默斯二重态；孤立地看它像自旋 $\frac{1}{2}$ 态，
尽管量子数全然不同。事实上可以赋予它有效自旋 $\tilde{S}=\frac{1}{2}$：这是虚构的
角动量，选它使所考虑的那组能级的简并度等于 $(2\tilde{S}+1)$。这样做的动机是找到
一个有效自旋哈密顿量，用某种简单得多的语言——自旋 $\frac{1}{2}$ 的自由原子或
离子——近似描述实验情形。该哈密顿量可写为
\begin{equation*}
\tilde{H} = g\mu_{\mathrm{B}}\mathbf{B}\cdot\tilde{\mathbf{S}}\tag{3.36}
\end{equation*}
这里的 $g$ 是有效 $g$ 因子（有时称光谱分裂因子）：对真实的孤立自旋
$\frac{1}{2}$ 电子它等于 2，此处还包含了轨道贡献。常常还需要用有效 $g$ 张量
$\mathbf{g}$ 表达自旋与场的相互作用依赖于磁场取向这一事实，式 3.36 须换成
\begin{equation*}
\tilde{H} = \mu_{\mathrm{B}}(\mathbf{B}\cdot\mathbf{g}\cdot\tilde{\mathbf{S}}) = \sum_{ij}\mathbf{g}_{ij}B_i\tilde{S}_j.\tag{3.37}
\end{equation*}
晶体场的效果可再往哈密顿量里加一项——表示自旋因晶体场而偏好沿特定晶向的能量。
这称为单离子各向异性（single ion anisotropy），哈密顿量中的项 $\tilde{H}_{\mathrm{SI}}$
在单轴晶体（存在特殊轴，能量只依赖自旋与该轴的夹角）为
\begin{equation*}
\tilde{H}_{\mathrm{SI}} = -DS_z^{2},\tag{3.38}
\end{equation*}
在立方晶体为
\begin{equation*}
\tilde{H}_{\mathrm{SI}} = -D(S_x^{4}+S_y^{4}+S_z^{4})\tag{3.39}
\end{equation*}
$D$ 是各向异性常数。

\begin{figure}
\centering
\includegraphics[width=0.55\textwidth]{fig3_18.pdf}
\caption{$\mathrm{Ce^{3+}}$ 离子的能级。主导的能量劈裂是自旋--轨道相互作用（此
能量约对应 3000 K，故上多重态的热布居可忽略）。两个多重态被假设的轴对称晶体场
分裂成若干二重态组；图中晶体场劈裂被夸大以便观察。较低处的劈裂 $\sim5$ meV，约
对应 50 K。最后，外加磁场分裂所有二重态（图中只画出了下面那个的分裂）。}
\end{figure}

ESR 不仅用于研究含顺磁离子盐类的物理，还广泛用于化学。共振对原子在分子中的位置
可以极为敏感，该技术被大量用于化学反应研究，尤其是自由基（free radicals）——
带未成对电子的原子或分子碎片，往往高活性，具有重要的环境与生物学意义。ESR 不如
NMR 用得广，因为它只适用于有未成对自旋的材料。

\begin{example}{3.6}
ESR 大显身手的一个领域是研究常含过渡金属离子的生物分子。一个重要例子是血红素
（haem）——铁卟啉基团，存在于血红蛋白与肌红蛋白中。这些分子参与氧及其他小分子
（如一氧化碳）的转运。血红素中的 Fe(II) 离子在与小分子结合后从 S=2 变为 S=0
（这是自旋转变的一个例子，见 3.1.1 节）。
\end{example}

ESR 中的自旋--晶格弛豫时间常常比 NMR 短得多，因为电子磁矩与晶格振动的耦合比核
强得多。晶格振动调制晶体场，经自旋--轨道相互作用耦合到电子磁矩，常使 $T_1$ 短到
ESR 谱线宽得无法观察（线宽正比于 $T_1^{-1}$）。有利的情形是 $\mathrm{Mn^{2+}}$——
它是 S 态离子（$3d^{5}$，${}^{6}S_{5/2}$），没有可供自旋耦合的轨道矩，其 ESR 在
室温下即可轻易观察。对其他过渡金属离子，$T_1$ 的大小取决于轨道淬灭的程度。\fn{9}{轨道矩被淬灭得越彻底，电子磁矩与晶格振动的耦合越小，$T_1$ 越长，ESR 线越窄。}冷却样品可使共振变窄（$T_1$ 增大），因为晶格振动减少；这还提高了能级间相对热
布居，增强共振强度。

ESR 的一个变种称为 ENDOR（电子核双共振，Electron Nuclear Double Resonance 的
缩写），可以以极高精度测量超精细相互作用常数与核塞曼相互作用。顺磁盐中核能级间
的 NMR 跃迁往往太弱而无法直接测量。ENDOR 同时使用 RF 信号与微波信号，借力 ESR 的
高灵敏度：微波功率调到 $\Delta m_I=0$ 的跃迁，此时 ESR 信号依赖于相应能级的相对
布居；再用调到 NMR 跃迁的 RF 功率改变布居，NMR 跃迁便可经由它对 ESR 信号的影响
被探测。

\subsection{穆斯堡尔谱学}

另一类谱学基于 1957 年发现的穆斯堡尔效应（M\"ossbauer effect）。实验原理示于
图 3.19。含 ${}^{57}$Co 核的源提供现成的激发态 ${}^{57}$Fe 核；它们经级联 $\gamma$
跃迁衰变到基态，其中包含一条 14.4 keV 的 $\gamma$ 射线（对应频率
$\nu=3.5\times10^{18}$ Hz）。若这条 $\gamma$ 射线被共振吸收，就能激发所研究样品中
的跃迁——其能量须匹配样品中的能隙。让源以速度 $v$ 运动，可因多普勒效应非常轻微地
调节 $\gamma$ 射线频率。光子频率极高，多普勒频移相当可观：$v=1\ \mathrm{mm\,s^{-1}}$
导致 $\nu v/c\approx12$ MHz 的频移。于是可以探测源或吸收体核中可能由磁相互作用或
其他相互作用引起的基态劈裂。

\begin{figure}
\centering
\includegraphics[width=0.55\textwidth]{fig3_19.pdf}
\caption{穆斯堡尔技术原理示意。${}^{57}$Co 缓慢衰变到 ${}^{57}$Fe 核的激发态；
随后 91\% 的衰变迅速到达 $I=\frac{3}{2}$ 态，再衰变到 $I=\frac{1}{2}$ 基态，放出
14.4 keV 光子。实验中让 ${}^{57}$Co 源相对样品以速度 $v$ 运动（样品须含一些
${}^{57}$Fe 核）。探测器测量 $\gamma$ 射线穿过样品的透射，由此推知吸收。}
\end{figure}

倘若激发态 ${}^{57}$Fe 核发射的 $\gamma$ 光子频率不明确，这项技术将毫无用处。
这正是 ${}^{57}$Fe 核相对缓慢的衰变至关重要的原因：0.2 $\mu$s 的半衰期只对应
约 2 MHz 的频率不确定度。第二个关键特征是 ${}^{57}$Fe 原子处于固态。为守恒动量，
自由 Fe 原子会获得 $h\nu/m_{\mathrm{Fe}}c\sim80\ \mathrm{m\,s^{-1}}$ 量级的反冲
速度，实验就毁了（我们已见过 1 mm\,s$^{-1}$ 的相对速度即可产生可测效应）。然而
被刚性固定在固体中的 ${}^{57}$Fe 原子把动量传递给整个晶体：若反冲能低于某个量，
声子发射的概率趋于零，$\gamma$ 射线可以不损失任何反冲能量地发射。这种无反冲
发射与共振吸收正是穆斯堡尔效应的精髓。

上述条件意味着：低能 $\gamma$ 射线 + 强结合于晶格的核 + 低温，效应最佳。
${}^{57}$Fe 并非唯一合适的同位素（可用者还有 ${}^{119}$Sn、${}^{127}$I、
${}^{151}$Eu、${}^{197}$Au），但最常用。穆斯堡尔测量的原始结果是一张 $\gamma$
计数（或相对吸收）对源与吸收体相对速度的图。

这个效应为什么能告诉我们样品的信息？首先，共振吸收未必出现在“期望”的位置
（源静止时），而可能略有偏移（即源以某特定速度运动时）。这一同质异能移（isomer
shift）源于核体积上核电荷分布与电子电荷分布之间库仑相互作用的微小变化——与
${}^{57}$Fe 核在 $I=\frac{3}{2}$ 态体积略增有关。

此外，观察到的共振吸收线可能不止一条，而是若干条。原因可以是四极劈裂（quadrupole
splitting）或磁劈裂（magnetic splitting）。前者源于激发态 ${}^{57}$Fe 核的电
四极矩（${}^{57}$Fe 基态 $I=\frac{1}{2}$ 无电四极矩，但所关心的激发态
$I=\frac{3}{2}$，$I>\frac{1}{2}$ 的核可有非零四极矩）。若核处于电场梯度中——某些
晶体环境即是如此——核四极矩与电场梯度的相互作用把激发的 $I=\frac{3}{2}$ 态分裂成
二重态，穆斯堡尔谱中产生两条线。后者由核与局域磁场的相互作用造成：它把 $I=\frac{1}{2}$
基态分裂成二重态、把激发的 $I=\frac{3}{2}$ 态分裂成四重态，穆斯堡尔谱中可有六条线
（选择定则 $\Delta m_I=0,\pm1$）。这一效应可用来探测磁交换相互作用与局域磁场。
所有这些移动与劈裂示意于图 3.20。注意图中（也绝无法）按比例画：观察到的劈裂典型
在 $10^{7}$--$10^{8}$ Hz，比 ${}^{57}$Fe 核基态（$I=\frac{1}{2}$）与激发态
（$I=\frac{3}{2}$）之间的能隙小 11 或 12 个量级。换句话说：这是在用 14.4 keV 的
光子测量小于 1 $\mu$eV 的能隙！若非穆斯堡尔效应——这种无反冲的共振吸收与发射——
此类实验完全不可能。

\begin{figure}
\centering
\includegraphics[width=0.55\textwidth]{fig3_20.pdf}
\caption{化学位移、四极劈裂与磁劈裂对 ${}^{57}$Fe 核能级的影响。箭头表示穆斯堡尔
吸收跃迁。跃迁大小的差异被极大夸张：实际上它们彼此之差小于 $10^{11}$ 分之一。}
\end{figure}

\subsection{$\mu$ 子自旋旋转}

$\mu$ 子是自旋 $\frac{1}{2}$ 粒子（电荷 $\pm e$，质量 $250m_{\mathrm{e}}$），寿命
$2.2\ \mu$s。在自然界，$\mu$ 子是到达海平面的宇宙线的主要成分。$\mu$ 子可用于
研究样品的磁性质，但宇宙线源强不足，须用同步加速器与回旋加速器提供的更强 $\mu$
子束。

下面较详细地描述称为 $\mu$ 子自旋旋转（muon-spin rotation，常缩写为 $\mu$SR）
的技术。必须认识到：与本书后面讨论的中子与 X 射线技术截然不同，这里不涉及散射；
$\mu$ 子被注入感兴趣的样品并在其短暂余生中留在那里、永不再出来。从样品中释放
出来并给出 $\mu$ 子信息的是它们衰变成的正电子。

$\mu$ 子质量介于电子与质子之间，其磁矩与旋磁比亦然。$\mu$ 子有正负两种电荷态，
磁性实验中特别有用的是正 $\mu$ 子 $\mu^{+}$。作为小的带正电粒子，它被电子密度大
的区域吸引，停在无机材料的间隙位，或直接键连到有机分子上。相反，负 $\mu$ 子
$\mu^{-}$ 注入后停在原子核附近，对磁性质一般远不敏感。

用高能质子束轰击合适的靶产生 $\pi$ 介子，$\pi$ 介子（$\pi^{+}$）极快地（26 ns）
衰变为 $\mu$ 子：
\begin{equation*}
\pi^{+} \rightarrow \mu^{+} + \nu_{\mu},
\end{equation*}
若只选取在靶中停下来的 $\pi$ 介子所产生的 $\mu$ 子，则 $\mu$ 子束完全自旋极化。
原因在于：同时产生的中微子（$\nu_{\mu}$）自旋与其动量反平行，而 $\pi$ 介子无自旋，
故 $\mu$ 子自旋必须与其动量反平行——出射即自旋极化。这些 $\mu$ 子随后被注入样品，
但其能量不小，至少 4 MeV。注入后它们通过原子电离与电子散射，在 0.1--1 ns 内很快
损失能量到几个 keV；接着经过一连串电子俘获与丢失反应，在约一皮秒内降到几百 eV。
若最终形成 $\mu$ 子素（muonium，见图 3.21）——$\mu^{+}$ 与 $e^{-}$ 组成的类氢态——
电子俘获最终胜出，最后几个 eV 通过 $\mu$ 子素原子与宿主原子的非弹性碰撞损失。
这一切都极快，$\mu$ 子（或 $\mu$ 子素）迅速热化。而且这些效应全源于库仑作用，
不与 $\mu$ 子自旋相互作用，故 $\mu$ 子在物质中热化时几乎不退极化——这对 $\mu$SR
实验至关重要。有人可能担心 $\mu$ 子只测量到被高能入射 $\mu$ 子辐射损伤的区域。
这实际上不构成问题：空位产生有阈值能量，只有 $\mu$ 子路径的初始部分受损严重；
过了损伤点，$\mu$ 子仍有足够能量继续在样品中前行约 1 $\mu$m，远离任何诱发产生的
空位与间隙原子。

\begin{figure}
\centering
\includegraphics[width=0.4\textwidth]{fig3_21.pdf}
\caption{$\mu$ 子素：由 $\mu^{+}$ 与 $e^{-}$ 组成的类氢态。}
\end{figure}

$\mu$SR 实验中，$\mu$ 子停在感兴趣的样品里，经时间 $t$ 后以正比于
$\mathrm{e}^{-t/\tau_{\mu}}$ 的概率衰变，$\tau_{\mu}=2.2\ \mu$s 是 $\mu$ 子寿命。
$\mu$ 子衰变是三体过程
\begin{equation*}
\mu^{+} \rightarrow e^{+} + \nu_e + \bar{\nu}_{\mu}.
\end{equation*}
衰变经由弱相互作用，因而具有宇称不守恒这一不寻常特征。这一现象（也是 $\mu$ 子
中微子负螺旋度的根源）使发射出的正电子（$e^{+}$）倾向沿 $\mu$ 子衰变时自旋所指
的方向出射。

能量最高的正电子的角分布示于图 3.22。实际上各种能量的正电子都会发射，净效果
没有这么显著。实验中可施加垂直于初始 $\mu$ 子自旋方向的磁场，使 $\mu$ 子自旋
进动。对许多 $\mu$ 子重复实验，只要愿意采集足够长时间的数据，就能以任意精度跟踪
进动 $\mu$ 子集合的极化。实验示意于图 3.23(a)。考虑一个极化与动量反平行的 $\mu$
子被注入样品（反平行是因为它的产生方式，见上——故 $\mu$ 子进入样品时自旋指向它
来的方向）。若它“倒霉”地立即衰变，就没时间进动，正电子将优先进入后向探测器；
若多活一会儿，就有时间进动——活过半圈，正电子将优先进入前向探测器。于是，进动
$\mu$ 子集合发出的正电子束可以比作灯塔的光束。

\begin{figure}
\centering
\includegraphics[width=0.45\textwidth]{fig3_22.pdf}
\caption{发射正电子相对初始 $\mu$ 子自旋方向的角分布（能量最高的正电子的期望
分布）。}
\end{figure}

\begin{figure}
\centering
\includegraphics[width=0.55\textwidth]{fig3_23.pdf}
\caption{（a）$\mu$SR 实验示意。自旋极化的 $\mu$ 子束停在样品 S 中；衰变后，正电子
由前向探测器 F 或后向探测器 B 探测。若如图对样品施加横向磁场 H，$\mu$ 子将进动。
（b）前向（虚线）与后向（实线）探测器探测到的正电子数；点线是两信号的平均。
（c）不对称度函数。}
\end{figure}

前向与后向探测器探测到的正电子数的时间演化由函数 $N_{\mathrm{F}}(t)$ 与
$N_{\mathrm{B}}(t)$ 描述（图 3.23(b)）。由于 $\mu$ 子衰变是放射性过程，两项之和
呈指数衰减。$\mu$ 子极化的时间演化可由两者的归一化差——不对称度函数（asymmetry
function）$A(t)$——得到：
\begin{equation*}
A(t) = \frac{N_{\mathrm{B}}(t)-N_{\mathrm{F}}(t)}{N_{\mathrm{B}}(t)+N_{\mathrm{F}}(t)},\tag{3.40}
\end{equation*}
示于图 3.23(c)。$\mu$ 子产生时具有 100\% 自旋极化，因此与 NMR 不同，无须用脉冲
技巧就能观察自由感应衰减。$\mu$ 子会在磁场——内场或外场——中拉莫尔进动，进动可
通过测量不对称度跟踪。进动频率通过 $\omega=\gamma_{\mu}B$ 直接联系磁场，其中
$\gamma_{\mu}=ge/2m_{\mu}=2\pi\times135.5\ \mathrm{MHz\,T^{-1}}$ 是 $\mu$ 子的旋磁
比（$m_{\mu}$ 为其质量，此处 $g\approx2$）。因此，注入磁有序材料的 $\mu$ 子在内场
中进动，直接给出频率正比于内磁场的振荡信号——在这个意义上，$\mu$ 子是一台微观
磁强计。质子（NMR）、电子（ESR）与 $\mu$ 子（$\mu$SR）的拉莫尔进动频率示于
图 3.24。

\begin{figure}
\centering
\includegraphics[width=0.5\textwidth]{fig3_24.pdf}
\caption{电子、$\mu$ 子与质子的拉莫尔进动频率 $f$（及相应周期 $\tau=1/f$）随外加
磁场 B 的变化。}
\end{figure}

$\mu$ 子磁矩大，对极小的磁场（低至 $\sim10^{-5}$ T）也很敏感，因此适合研究小磁矩
磁性；对磁有序随机或极短程的材料也有价值。由于 $\mu$ 子在样品中均匀停下，每个
信号在实验谱中以正比于其体积分数的强度出现——故对多相或未完全有序的样品尤为
有用。由于（与后面讨论的衍射技术不同）不获得空间信息，单晶样品并非必需（尽管
某些情形有用），而且对某些不便做常规磁中子衍射的材料，$\mu$SR 常能提供磁有序
信息。要从 $\mu$SR 实验提取定量信息，须知道 $\mu$ 子落位（muon-site），这在某些
情形会妨碍对数据的直截了当的解释。（NMR 中核把质子定位得很好，探针位置一目了然。）
通常 $\mu$ 子可占据的间隙位是一小组，且在有利情形只有一个与观察数据一致。这一
技术已在磁性材料中得到广泛应用。

\section*{延伸阅读}

\begin{itemize}[itemsep=0.2em]
\item NMR 的优秀入门可见 P.~J.~Hore，《Nuclear magnetic resonance》，OUP 1995；
B.~Cowan，《Nuclear magnetic resonance and relaxation》，CUP 1997 也极为有用；
NMR 的经典著作是 A.~Abragam，《Principles of nuclear magnetism》，OUP 1961。
\item A.~Abragam 与 B.~Bleaney，《Electron Paramagnetic Resonance of Transition
Ions》，Dover 1986，提供了关于顺磁盐中晶体场与 ESR 实验的大量信息。
\item 晶体场可用所谓的 Stevens 算符处理，见 K.~W.~H.~Stevens，Proc.~Phys.~Soc.~A
\textbf{65}, 209 (1952) 与 M.~T.~Hutchings，Solid State Physics \textbf{16}, 227
(1966)。
\item 穆斯堡尔效应的更多信息见《M\"ossbauer spectroscopy》，D.~P.~E.~Dickson 与
F.~J.~Berry 编，CUP 1986。
\item $\mu$SR 的更多信息见 S.~J.~Blundell，Contemp.~Phys.~\textbf{40}, 175 (1999)。
\end{itemize}

\section*{习题}

\exer{3.1} $\mathrm{Sc^{2+}}$ 离子在 3d 壳层有一个电子。它处于各向异性晶体中，
晶体场对 3d 电子的作用可写为势 $A\hat{l}_z^{2}$。$A>0$ 与 $A<0$ 时 Sc 离子的最低
轨道态各是什么？自旋--轨道耦合 $\lambda\hat{\mathbf{l}}\cdot\hat{\mathbf{s}}$ 远小于
晶体场；把这一项也包括进来后，$A<0$ 与 $A>0$ 时离子的近似基态是什么？讨论沿 $z$
轴以及垂直于 $z$ 轴加小磁场对这些态的影响，并画出磁化率随温度变化的示意图。

\exer{3.2} 把等量的正点电荷放在正八面体的六个顶点上。取笛卡尔坐标系原点位于八面体
中心，证明中心附近势为
\begin{equation*}
V = \frac{q}{4\pi\epsilon_0 a}\left[6 + \frac{35}{4a^{4}}\left(x^{4}+y^{4}+z^{4}-\frac{3}{5}r^{4}\right) + O\left(\frac{r^{6}}{a^{6}}\right)\right],\tag{3.41}
\end{equation*}
$q$ 为每个电荷的大小，$a$ 为原点到每个电荷的距离。

\exer{3.3} 某化合物单位体积含 $n$ 个过渡金属离子，每个离子处于形式为
\begin{equation*}
V = D\left(x^{4}+y^{4}+z^{4}-\frac{3}{5}r^{4}\right)\tag{3.42}
\end{equation*}
的晶体场中，它对简并的 d 能级构成微扰。未微扰波函数可以有多种选取方式；先取
$\hat{L}_z$ 的本征函数并按其本征值 $m_l$ 标记：
\begin{center}
\resizebox{0.88\textwidth}{!}{$\begin{aligned}
&|\pm2\rangle = R(r)\sin^{2}\theta\,\mathrm{e}^{\pm2\mathrm{i}\phi}\\
&|\pm1\rangle = \mp2R(r)\sin\theta\cos\theta\,\mathrm{e}^{\pm\mathrm{i}\phi}\\
&|0\rangle = \sqrt{\tfrac{2}{3}}\,R(r)(3\cos^{2}\theta-1),
\end{aligned}$}\qquad(3.43)
\end{center}
$R(r)$ 是含归一化常数的波函数径向部分。系统的一般态可写为
$|\psi\rangle=\sum_{j=-2}^{2}a_j|j\rangle$，并指定为矢量：
\begin{equation*}
|\psi\rangle = \begin{pmatrix}a_2\\a_1\\a_0\\a_{-1}\\a_{-2}\end{pmatrix}.\tag{3.44}
\end{equation*}
在该基下证明
\begin{equation*}
V = \begin{pmatrix} A & 0 & 0 & 0 & 5A\\ 0 & -4A & 0 & 0 & 0\\ 0 & 0 & 6A & 0 & 0\\ 0 & 0 & 0 & -4A & 0\\ 5A & 0 & 0 & 0 & A \end{pmatrix}\tag{3.45}
\end{equation*}
$A$ 为常数，并证明本征值与本征函数为：
\begin{center}
\begin{tabular}{lccccc}
本征值： & $6A$ & $6A$ & $-4A$ & $-4A$ & $-4A$\\[2pt]
本征函数： & $|0\rangle$ & $\dfrac{|2\rangle+|-2\rangle}{\sqrt{2}}$ & $|1\rangle$ & $|-1\rangle$ & $\dfrac{|2\rangle-|-2\rangle}{\sqrt{2}}$
\end{tabular}\hfill(3.46)
\end{center}
磁场 B 诱导劈裂 $\mu_{\mathrm{B}}Bm_l$，于是 V 变为
\begin{equation*}
V = \begin{pmatrix} A+2\mu_{\mathrm{B}}B & 0 & 0 & 0 & 5A\\ 0 & -4A+\mu_{\mathrm{B}}B & 0 & 0 & 0\\ 0 & 0 & 6A & 0 & 0\\ 0 & 0 & 0 & -4A-\mu_{\mathrm{B}}B & 0\\ 5A & 0 & 0 & 0 & A-2\mu_{\mathrm{B}}B \end{pmatrix}\tag{3.47}
\end{equation*}
计算此时的本征值与本征矢。本征值的图示于图 3.25。

再证明低温下（$k_{\mathrm{B}}T\ll A$），趋于零的磁场中磁化率等于
\begin{equation*}
\chi = \begin{cases} \dfrac{2\mu_0\mu_{\mathrm{B}}^{2}n}{3k_{\mathrm{B}}T} & \text{当 } A>0\\[10pt] \dfrac{2\mu_0\mu_{\mathrm{B}}^{2}n}{5A} & \text{当 } A<0, \end{cases}\tag{3.48}
\end{equation*}
第一种情形不依赖温度且很小，第二种情形等价于 $l=1$ 自由离子的预期值。

\begin{figure}
\centering
\includegraphics[width=0.55\textwidth]{fig3_25.pdf}
\caption{晶体场把 $l=2$ 的离子分裂成一个三重态和一个激发单态；磁场 B 使这些态进一步
分裂。所列本征态指零场情形（场增大时，较高的态逐渐获得更多 $|2\rangle$ 成分，较低
的态逐渐获得更多 $|-2\rangle$ 成分）。}
\end{figure}

\exer{3.4} 某离子核自旋为零，其基态由两个简并能级组成，对应有效自旋 $S=1/2$。
施加磁通密度为 B 的磁场，产生与 B 成线性的能级分裂。观察到该离子的单条电子顺磁
共振线，频率 30 GHz、磁通密度 0.6 T。核自旋非零 I 的同位素产生超精细结构（在自旋哈密顿量
中由 $A\mathbf{I}\cdot\mathbf{S}$ 项描述），由四条近似等距的共振线组成，间隔
$10^{-2}$ T，对称分布在核自旋为零的同位素的谱线两侧。这给出了关于（a）核自旋与
（b）该同位素核磁矩的什么信息？计算参数 A 的数值。

为什么有时必须在低温下进行这些测量？

\exer{3.5} 图 3.26 给出三种顺磁盐每离子磁矩对 $B/T$ 的图。证明式 2.9 在低与高
$B/T$ 处都与数据同形。自由 $\mathrm{Cr^{3+}}$ 与 $\mathrm{Fe^{3+}}$ 离子的 3d 壳层
分别含三与五个电子，$\mathrm{Gd^{3+}}$ 的 4f 壳层含七个电子。用洪特规则计算这些
离子的 $S, L, J$ 与 $g_J$，证明所得数值能很好解释图中 $\mathrm{Gd^{3+}}$ 与
$\mathrm{Fe^{3+}}$ 的数据，但不能解释 $\mathrm{Cr^{3+}}$。从实验数据推断
$\mathrm{Cr^{3+}}$ 的实际 $g$ 因子与基态量子数，并简述数据与你最初计算不符的原因。

\begin{figure}
\centering
\includegraphics[width=0.5\textwidth]{fig3_26.pdf}
\caption{三种顺磁盐的每离子磁矩：$\mathrm{KCr(SO_4)_2\cdot12H_2O}$
（$\mathrm{Cr^{3+}}$）、$\mathrm{NH_4Fe(SO_4)_2\cdot12H_2O}$（$\mathrm{Fe^{3+}}$）
与 $\mathrm{Gd_2(SO_4)_3\cdot8H_2O}$（$\mathrm{Gd^{3+}}$）。取自 W.~Henry，
Phys.~Rev.~\textbf{88}, 559 (1952)。}
\end{figure}

\exer{3.6} 一台 NMR 谱仪工作在 60 MHz。${}^{1}$H、${}^{2}$H、${}^{13}$C 与
${}^{19}$F 核的共振预计出现在多大外加磁场处？（用表 2.3 的数据。）

\exer{3.7} 一台 ESR 谱仪工作在 9 GHz（X 波段）。要观察 DPPH（一种用于校准 ESR
谱仪的有机分子，在 $g=2$ 处给出尖锐信号）中未成对电子的信号，需要多大磁场？

\exer{3.8} 某系统的自旋哈密顿量为
\begin{equation*}
\hat{\mathcal{H}} = g_{\parallel}\mu_{\mathrm{B}}B_zS_z + g_{\perp}\mu_{\mathrm{B}}(B_xS_x+B_yS_y) + DS_z^{2}.\tag{3.49}
\end{equation*}
证明：对该系统中处于三重态（$S=1$）的两个电子，能量本征值 E 可由
\begin{equation*}
E(E-D)^{2} - E(g_{\parallel}\mu_{\mathrm{B}}B\cos\theta)^{2} - (E-D)(g_{\perp}\mu_{\mathrm{B}}B\sin\theta)^{2} = 0.\tag{3.50}
\end{equation*}
得到。求并画出 $\theta=0$ 与 $\theta=\pi/2$ 的解。示意 $\hbar\omega>D$ 时可能的
光子跃迁，进而说明如何通过测量 $\theta=0$ 与 $\theta=\pi/2$ 的 ESR 得到参数
$g_{\parallel}$、$g_{\perp}$ 与 D。

\exer{3.9} 穆斯堡尔效应中 ${}^{57}$Fe 原子发射 14.4 keV 光子。计算自由
${}^{57}$Fe 原子以及被刚性固定在 10 mg 晶体晶格中的 ${}^{57}$Fe 原子的反冲速度。
两种情形下 14.4 keV 光子的多普勒频移各是多少？某穆斯堡尔实验在源与样品相对速度
为 2 mm\,s$^{-1}$ 时得到尖锐谱线。计算相应的频移（Hz 与 cm$^{-1}$）以及能量移动
（eV）。

\exer{3.10} $\mu$ 子平均寿命 2.2 $\mu$s；假定可测比例的 $\mu$ 子能活 20 $\mu$s，
这给可测内场设置了什么下限？一次典型实验测量 $10^{7}$ 个 $\mu$ 子，其中有多少能
活 20 $\mu$s 或更久？同步加速器产生的 $\mu$ 子脉冲宽度为 50 ns。这给样品中可测
内场设置了什么上限？（提示：考虑脉冲前端与后端 $\mu$ 子产生的进动信号的差异。）
回旋加速器可产生连续 $\mu$ 子束，避免这一问题。某磁有序氧化物中被注入 $\mu$ 子
占据的间隙位内场为 0.4 T。预期测到多大的进动频率？

\exer{3.11} 晶体中 $\mathrm{Fe^{3+}}$ 离子（${}^{6}S_{5/2}$）的饱和磁矩预期为
$5\mu_{\mathrm{B}}$（见图 3.26），但由磁化率测量推得的有效磁矩预期为
$5.92\mu_{\mathrm{B}}$（见表 3.1）。为何不同？

% 第4章 磁有序与磁结构？不——第4章 相互作用（Interactions）
\chapter{相互作用}

本章讨论几类磁相互作用：它们使固体中的磁矩得以相互“沟通”，并可能产生长程序。

\section{磁偶极相互作用}

最先想到会起作用的相互作用是磁偶极相互作用。两个相距 $\mathbf{r}$ 的磁偶极子
$\boldsymbol{\mu}_1$ 与 $\boldsymbol{\mu}_2$ 的能量为
\begin{equation*}
E = \frac{\mu_0}{4\pi r^{3}}\left[\boldsymbol{\mu}_1\cdot\boldsymbol{\mu}_2 - \frac{3}{r^{2}}(\boldsymbol{\mu}_1\cdot\mathbf{r})(\boldsymbol{\mu}_2\cdot\mathbf{r})\right]\tag{4.1}
\end{equation*}
依赖其间距与相互排列的程度。容易估计其量级：两个 $\mu\approx1\mu_{\mathrm{B}}$ 的磁矩
相距 $r\approx1$ \AA{} 时约为 $\mu^{2}/4\pi r^{3}\sim10^{-23}$~J，按温度计约 1 K。
\fn{1}{能量与温度的便捷换算见图 3.8。}许多材料在远高于此的温度有序（有的高达
1000 K 左右），因此磁偶极相互作用太弱，不足以解释大多数磁性材料的有序。尽管如此，
对在毫开尔文量级有序的那些材料，它可能相当重要。

\section{交换相互作用}

交换相互作用位于长程磁有序现象的心脏。交换效应精微而不无神秘——处理的无非是
一根条形磁铁和一堆铁屑，却必须动用交换算符与全同粒子，这乍看令人诧异。但正如
磁学中常见的情形，这恰恰表明量子力学是许多日常现象的根源。交换相互作用不过是
静电相互作用：同号电荷靠近时要付出能量、分开时省下能量。

\subsection{交换的起源}

考虑只有两个电子的简单模型，其空间坐标分别为 $\mathbf{r}_1$ 与 $\mathbf{r}_2$。
联合态的波函数可以写成单电子态的乘积：若第一个电子处于 $\psi_a(\mathbf{r}_1)$、
第二个电子处于 $\psi_b(\mathbf{r}_2)$，则联合波函数为
$\psi_a(\mathbf{r}_1)\psi_b(\mathbf{r}_2)$。但这个乘积态不满足交换对称性：交换两个
电子得到 $\psi_a(\mathbf{r}_2)\psi_b(\mathbf{r}_1)$，它与初始态不成倍数。因此我们
只能构造在粒子交换下行为正确的（反）对称化乘积态（1.3.4 节已讨论）。

对电子，总波函数必须反对称，故波函数的自旋部分要么是反对称单态 $\chi_{\mathrm{S}}$
（S=0，配对称空间态），要么是对称三重态 $\chi_{\mathrm{T}}$（S=1，配反对称空间态）。
单态 $\Psi_{\mathrm{S}}$ 与三重态 $\Psi_{\mathrm{T}}$ 的波函数写为
\begin{equation*}
\begin{aligned}
\Psi_{\mathrm{S}} &= \frac{1}{\sqrt{2}}\left[\psi_a(\mathbf{r}_1)\psi_b(\mathbf{r}_2) + \psi_a(\mathbf{r}_2)\psi_b(\mathbf{r}_1)\right]\chi_{\mathrm{S}}\\
\Psi_{\mathrm{T}} &= \frac{1}{\sqrt{2}}\left[\psi_a(\mathbf{r}_1)\psi_b(\mathbf{r}_2) - \psi_a(\mathbf{r}_2)\psi_b(\mathbf{r}_1)\right]\chi_{\mathrm{T}},
\end{aligned}\tag{4.2}
\end{equation*}
空间与自旋部分都包括在内。两个可能态的能量为
\begin{equation*}
\begin{aligned}
E_{\mathrm{S}} &= \int\Psi_{\mathrm{S}}^{*}\hat{\mathcal{H}}\Psi_{\mathrm{S}}\,\mathrm{d}\mathbf{r}_1\mathrm{d}\mathbf{r}_2\\
E_{\mathrm{T}} &= \int\Psi_{\mathrm{T}}^{*}\hat{\mathcal{H}}\Psi_{\mathrm{T}}\,\mathrm{d}\mathbf{r}_1\mathrm{d}\mathbf{r}_2,
\end{aligned}
\end{equation*}
假定波函数的自旋部分 $\chi_{\mathrm{S}}$、$\chi_{\mathrm{T}}$ 已归一化。两能量之差为
\begin{equation*}
E_{\mathrm{S}} - E_{\mathrm{T}} = 2\int\psi_a^{*}(\mathbf{r}_1)\psi_b^{*}(\mathbf{r}_2)\hat{\mathcal{H}}\psi_a(\mathbf{r}_2)\psi_b(\mathbf{r}_1)\,\mathrm{d}\mathbf{r}_1\mathrm{d}\mathbf{r}_2.\tag{4.3}
\end{equation*}
式 1.70 表明单态与三重态之差可以用 $\mathbf{S}_1\cdot\mathbf{S}_2$ 参数化：单态
$\mathbf{S}_1\cdot\mathbf{S}_2=-\frac{3}{4}$，三重态 $\mathbf{S}_1\cdot\mathbf{S}_2=\frac{1}{4}$。
于是哈密顿量可以写成“有效哈密顿量”的形式
\begin{equation*}
\hat{\mathcal{H}} = \frac{1}{4}(E_{\mathrm{S}}+3E_{\mathrm{T}}) - (E_{\mathrm{S}}-E_{\mathrm{T}})\mathbf{S}_1\cdot\mathbf{S}_2.\tag{4.4}
\end{equation*}
这是常数项与依赖自旋的项之和。常数可并入其他常数能量项，第二项更有意思。交换常数
（或交换积分）$\mathsf{J}$ 定义为
\begin{equation*}
\mathsf{J} = \frac{E_{\mathrm{S}}-E_{\mathrm{T}}}{2} = \int\psi_a^{*}(\mathbf{r}_1)\psi_b^{*}(\mathbf{r}_2)\hat{\mathcal{H}}\psi_a(\mathbf{r}_2)\psi_b(\mathbf{r}_1)\,\mathrm{d}\mathbf{r}_1\mathrm{d}\mathbf{r}_2,\tag{4.5}
\end{equation*}
于是有效哈密顿量中依赖自旋的项可写作
\begin{equation*}
\hat{\mathcal{H}}^{\text{spin}} = -2\mathsf{J}\,\mathbf{S}_1\cdot\mathbf{S}_2.\tag{4.6}
\end{equation*}
若 $\mathsf{J}>0$，$E_{\mathrm{S}}>E_{\mathrm{T}}$，更倾向于三重态 $S=1$；若
$\mathsf{J}<0$，$E_{\mathrm{S}}<E_{\mathrm{T}}$，更倾向于单态 $S=0$。对两个电子导出
此式还算简单，\bio{Werner Heisenberg\\（1901--1976）}\mnote{注意有些书把 J 取为此处数值的两倍，此时式 4.6--4.8 变为
\begin{equation*}
\hat{\mathcal{H}}^{\mathrm{spin}} = -\mathsf{J}\,\mathbf{S}_1\cdot\mathbf{S}_2,\tag{4.9}
\end{equation*}
\begin{equation*}
\hat{\mathcal{H}} = -\frac{1}{2}\sum_{ij}\mathsf{J}_{ij}\,\mathbf{S}_i\cdot\mathbf{S}_j,\tag{4.10}
\end{equation*}
\begin{equation*}
\hat{\mathcal{H}} = -\sum_{i>j}\mathsf{J}_{ij}\,\mathbf{S}_i\cdot\mathbf{S}_j.\tag{4.11}
\end{equation*}}
但推广到多体系统绝非易事。尽管如此，量子力学早期人们就已认识到：式 4.6 这类相互
作用很可能存在于所有相邻原子之间。这引出海森伯模型（Heisenberg model）的哈密顿量：
\begin{equation*}
\hat{\mathcal{H}} = -\sum_{ij}\mathsf{J}_{ij}\,\mathbf{S}_i\cdot\mathbf{S}_j,\tag{4.7}
\end{equation*}
$\mathsf{J}_{ij}$ 是第 $i$ 与第 $j$ 个自旋之间的交换常数。因子 2 被省去，因为求和把
每对自旋计了两次。式 4.7 的另一种写法是
\begin{equation*}
\hat{\mathcal{H}} = -2\sum_{i>j}\mathsf{J}_{ij}\,\mathbf{S}_i\cdot\mathbf{S}_j,\tag{4.8}
\end{equation*}
$i>j$ 避免了“重复计数”，因子 2 于是回来了。常可取 $\mathsf{J}_{ij}$ 对最近邻自旋
等于常数 $\mathsf{J}$、其余为零。

交换积分的计算一般会相当复杂，这里只提几个一般性的特征。第一，若两个电子在同一原子上，
交换积分通常为正：这稳定三重态、保证空间态反对称，使两电子相互回避从而减小库仑
排斥——与洪特第一条规则一致。

当两个电子分处相邻原子时，情形截然不同。任何联合态都是一个以某原子为中心的态与
以另一原子为中心的态的组合。值得记住：长为 $L$ 的一维箱中粒子的能量正比于
$L^{-2}$——这是动能，说明被关进小盒子要付出巨大动能。因此电子可以通过成键节省
动能：成键让它们在两个原子而非一个原子的范围内游走（即住进“更大的盒子”）。
此时该考虑的不再是原子轨道而是分子轨道（见图 4.1）：成键轨道（空间对称）与
“反键”轨道（空间反对称），后者能量更高——反键轨道曲率更大，动能更大。这更倾向于
单态（反对称），因此交换积分很可能为负。

\subsection{直接交换}

如果相邻磁性原子上的电子经由交换相互作用发生耦合，这就称为直接交换（direct
exchange）：交换相互作用直接进行，无需任何中介。尽管这看似是交换相互作用最显而易
见的途径，物理实际却很少如此简单。

在很多情形下，直接交换并不是决定磁性的重要机制，因为相邻磁性轨道之间缺乏足够的直
接交叠。例如在稀土中，4f 电子高度局域、十分贴近原子核，其概率密度明显伸出原
子间距约十分之一以外的部分微乎其微，因此直接交换相互作用在稀土中不太可能十分有效。
即便在 Fe、Co、Ni 等过渡金属中——3d 轨道离核伸得更远——也很难解释直接交换何以导
致观测到的磁性质。这些材料是金属，这意味着传导电子的作用不容忽视，正确的描述必须
同时顾及电子的局域性与能带特征。

\begin{figure}
\centering
\includegraphics[width=0.55\textwidth]{fig4_01.pdf}
\caption{双原子分子的分子轨道。成键轨道（两原子轨道之和；就波函数空间部分而言在
交换下对称）能量低于反键轨道（两原子轨道之差；交换下反对称）。因此两个电子填入
成键态、反键态空着的单态基态更稳定。此图适用于氢分子 $\mathrm{H_2}$——其能量低于
两个孤立 H 原子（$E_0$）。注意氦的双原子形式 $\mathrm{He_2}$ 不能形成：两个 He 原子
的四个电子会同时填满成键与反键轨道，相比两个孤立 He 原子没有净能量节省。}
\end{figure}

因此在许多磁性材料中，必须考虑某种间接交换相互作用。

\subsection{离子固体中的间接交换：超交换}

一些离子固体（包括某些氧化物与氟化物）具有磁基态。例如 MnO（见图 4.2）与
$\mathrm{MnF_2}$ 都是反铁磁体——初看令人惊讶，因为两个体系中 $\mathrm{Mn^{2+}}$
离子的电子并无直接交叠。交换相互作用通常是短程的，此处起作用的较长程相互作用必在
某种意义上是“超”的。

这里起作用的交换机制称为超交换（superexchange）：可定义为非相邻磁性离子之间经
由置于其间的非磁性离子传递的间接交换相互作用。它的产生是因为反铁磁性在动能上有
好处。参照图 4.3 即可理解：图中两个过渡金属离子被一个氧离子隔开。为简单起见，设
过渡金属离子上的磁矩来自单个未成对电子（更复杂的情形可类似处理）。若系统完全
离子化，每个金属离子的 d 轨道上有一个未成对电子，而氧的最外层占据态有两个 p 电子。
图示表明：反铁磁耦合使体系能量降低——它让这些电子得以在整个结构上离域化，从而降低动能。

\begin{figure}
\centering
\includegraphics[width=0.4\textwidth]{fig4_02.pdf}
\caption{MnO 的晶体结构。$\mathrm{Mn^{2+}}$（锰）离子的最近邻对经由
$\mathrm{O^{2-}}$（氧）离子连接。}
\end{figure}

\begin{figure}
\centering
\includegraphics[width=0.4\textwidth]{fig4_03.pdf}
\caption{磁性氧化物中的超交换。箭头表示四个电子的自旋及其在过渡金属（M）与氧（O）
原子上的分布。设 M 有一个未成对电子因而具有磁性。若金属原子上的磁矩反铁磁耦合
（a、b、c），基态为（a），它可与（b）、（c）这样的激发组态混合：磁性电子因而在
M--O--M 单元上离域化，动能降低。若金属原子上的磁矩铁磁耦合（d、e、f），基态（d）
无法与（c）、（f）这样的激发组态混合——不相容原理禁止这些组态。铁磁组态因此付出
更多能量。}
\end{figure}

由于超交换同时涉及氧轨道与金属原子，它是二级过程，由二阶微扰论给出。二阶微扰论的
一般结论是：所涉能量近似等于跃迁矩阵元的平方除以制造激发态的能量代价。这里跃迁
矩阵元由一个称为跳跃积分（hopping integral）$t$ 的参数控制——在简单紧束缚近似中
它正比于导带的能量宽度（即带宽）。制造激发态的能量代价由库仑能 U 给出。于是
$\mathsf{J}\sim-t^{2}/U$。（四阶时可以有 $\Delta E\propto-t^{4}(\mathbf{S}_1\cdot\mathbf{S}_2)^{2}/U^{3}$
形式的相互作用，称为双二次交换（biquadratic exchange）。）

交换积分由两部分组成。第一是势交换项，代表电子排斥，更倾向于铁磁基态，但在离子
相距较远时很小。第二是动交换项，在此占主导，即上述效应。它依赖于轨道交叠程度，
因此超交换强烈依赖 M--O--M 键角。图只为一种 d 轨道而画；能与氧轨道交叠的其他 d
轨道的效应可能也需计入。

在某些情形下，超交换实际上可以是铁磁的。设想这样一个情形：通过一个氧离子，一个
磁性离子上的已占据 $e_g$ 轨道与另一磁性离子上的未占据 $e_g$ 轨道相耦合。$e_g$
电子跳到未占据轨道上有能量好处——只要到达时其自旋因洪特规则耦合而与 $t_{2g}$
电子的自旋排齐。此时超交换可以是铁磁的，但这是较弱的相互作用，也比常见的反铁磁
超交换少见。

\subsection{金属中的间接交换}

金属中磁性离子之间的交换相互作用可以由传导电子传递。局域磁矩使传导电子自旋极化，
这一极化又与距离 $r$ 之外的另一个局域磁矩耦合。这种交换是间接的——不涉及磁矩间的
直接耦合——称为 RKKY 相互作用（也称巡游交换，itinerant exchange）。RKKY 之名取自
该效应发现者 Ruderman、Kittel、Kasuya、Yosida 四人姓氏的首字母。耦合取与 $r$ 相关
的交换相互作用 $\mathsf{J}_{\mathrm{RKKY}}(r)$ 的形式：
\begin{equation*}
\mathsf{J}_{\mathrm{RKKY}}(r) \propto \frac{\cos(2k_{\mathrm{F}}r)}{r^{3}}.\tag{4.12}
\end{equation*}
（大 $r$ 时，设半径 $k_{\mathrm{F}}$ 的球形费米面。）这一相互作用是长程的，且随磁矩
间距振荡：依间距不同，可以铁磁也可以反铁磁。由于费米面十分锐利，耦合以 $\pi/k_{\mathrm{F}}$
为波长振荡。RKKY 相互作用将在第七章详细讨论。

\subsection{双交换}

在某些氧化物中可以出现铁磁交换相互作用，其原因是磁性离子可表现混合价态
（mixed valency）——即以不止一种氧化态存在。例子包括含 Mn 离子的化合物：Mn 可
处于氧化态 3 或 4，即 $\mathrm{Mn^{3+}}$ 或 $\mathrm{Mn^{4+}}$。其一即
$\mathrm{La}_{1-x}\mathrm{Sr}_x\mathrm{MnO}_3$（$0\leq x\leq1$），采取钙钛矿结构。
Sr 是二价的（以 $\mathrm{Sr^{2+}}$ 存在），La 是三价的（以 $\mathrm{La^{3+}}$ 存在）。
这意味着有 $x$ 分数的 Mn 离子是 $\mathrm{Mn^{4+}}$，$1-x$ 是 $\mathrm{Mn^{3+}}$。
系列两端（$x=0$ 与 $x=1$）都是反铁磁绝缘体——对磁性经氧由超交换传递的氧化物材料
而言这正合预期。$\mathrm{LaMnO_3}$ 只含 $\mathrm{Mn^{3+}}$，而 $\mathrm{Mn^{3+}}$
是扬--特勒离子；$\mathrm{LaMnO_3}$ 具有 A 型反铁磁有序（见 5.2 节）。然而当
$\mathrm{LaMnO_3}$ 掺 Sr 至 $x=0.175$ 时，扬--特勒畸变消失，系统变为铁磁，居里
温度在室温附近；低于该温度材料变为金属。

铁磁排列源于双交换（double exchange）机制，可借图 4.4 理解。$\mathrm{Mn^{3+}}$ 离子
上的 $e_g$ 电子只有当邻位存在同自旋空位时才能跳过去（因为跳跃不改变跳跃电子的
自旋）。若邻居是 $e_g$ 壳层无电子的 $\mathrm{Mn^{4+}}$，这本不成问题。但 $e_g$
电子与 $t_{2g}$ 能级三个电子之间存在强单心（洪特第一规则）交换相互作用，要把它们
全部排齐。因此，$e_g$ 电子跳到其 $t_{2g}$ 自旋将与自己反平行的邻居离子上，在能量上
是不利的（图 4.4(b)）。要维持给体与受体离子上的高自旋排列，相邻离子必须铁磁排列。
由于跳跃能力带来动能节省，允许图 4.4(a) 所示的跳跃过程降低总能量：系统于是铁磁
排列以省能量。而且铁磁排列进而允许 $e_g$ 电子在晶体中穿行，材料变为金属。双交换
铁磁体的导电性问题将在 8.9.5 节进一步讨论。双交换本质上是扩展体系中的铁磁超交换。
\fn{2}{铁磁超交换通常用于两个孤立离子：铁磁排列省下的动能对应跳入激发态。双交换用于扩展体系：省下的动能对应电子带宽的增益。}

\begin{figure}
\centering
\includegraphics[width=0.55\textwidth]{fig4_04.pdf}
\caption{双交换机制：参与电子转移的 $\mathrm{Mn^{3+}}$（$3\mathrm{d}^{4}$）与
$\mathrm{Mn^{4+}}$（$3\mathrm{d}^{3}$）离子之间产生铁磁耦合。单心交换相互作用
有利于跳跃，当（a）相邻离子铁磁排列时；（b）相邻离子反铁磁排列时则否。}
\end{figure}

双交换也见于磁铁矿（$\mathrm{Fe_3O_4}$）：八面体位上有等量的 $\mathrm{Fe^{2+}}$
（$3\mathrm{d}^{6}$）与 $\mathrm{Fe^{3+}}$（$3\mathrm{d}^{5}$）离子，四面体位上另有同等数目的
$\mathrm{Fe^{3+}}$ 离子。双交换相互作用使八面体位上的 $\mathrm{Fe^{2+}}$ 与
$\mathrm{Fe^{3+}}$ 离子铁磁排列；四面体位上的 $\mathrm{Fe^{3+}}$ 不参与这一相互
作用，与八面体位上的 $\mathrm{Fe^{3+}}$ 经反铁磁超交换耦合。于是两组
$\mathrm{Fe^{3+}}$ 的磁矩相消，剩下仅由 $\mathrm{Fe^{2+}}$ 贡献的净磁矩。实测的
每化学式单元磁矩非常接近仅由 $\mathrm{Fe^{2+}}$ 给出的预期值 $4\mu_{\mathrm{B}}$。

\subsection{各向异性交换相互作用}

自旋--轨道相互作用也可能以类似超交换中氧原子的方式扮演角色：此时激发态与氧无关，
而由磁性离子之一的自旋--轨道相互作用产生。于是一个离子的激发态与另一个离子的
基态之间存在交换相互作用。这称为各向异性交换相互作用（anisotropic exchange
interaction），又称 Dzyaloshinsky--Moriya 相互作用。作用于两个自旋 $\mathbf{S}_1$、
$\mathbf{S}_2$ 之间时，它在哈密顿量中给出 $\hat{\mathcal{H}}_{\mathrm{DM}}$：
\begin{equation*}
\hat{\mathcal{H}}_{\mathrm{DM}} = \mathbf{D}\cdot\mathbf{S}_1\times\mathbf{S}_2.\tag{4.13}
\end{equation*}
当晶体场对两磁离子连线中点具有反演对称时矢量 D 为零。但一般 D 可以不为零，此时
依对称性它平行或垂直于连接两自旋的连线。该相互作用的形式迫使 $S_1$ 与 $S_2$ 在
垂直于 D 的平面内接近直角，且取向使能量为负。其效果常常是使自旋倾摆（cant，即
略微转动一个小角）。它常见于反铁磁体，产生垂直于反铁磁自旋轴的小铁磁矩分量，
称为弱铁磁性（weak ferromagnetism）。例如 $\alpha$-$\mathrm{Fe_2O_3}$、
$\mathrm{MnCO_3}$ 与 $\mathrm{CoCO_3}$ 中即有此效应。

\begin{figure}
\centering
\includegraphics[width=0.5\textwidth]{fig4_05.pdf}
\caption{磁矩用相邻格点 $i$、$j$ 上的约化矩 $m_i$、$m_j$ 表示，两点相距
$r_{ij}$，磁矩间夹角为 $\phi_{ij}$。约化矩 $m_i$、$m_j$ 按定义为单位矢量。}
\end{figure}

\subsection{连续介质近似}

本节回到式 4.7 的海森伯模型。对本书后文有用的是：在忽略晶格分立性的连续介质近似
下，为这一相互作用找出一个表达式。

先设 $\mathsf{J}_{ij}$ 在 $i$、$j$ 为最近邻时等于常数 $\mathsf{J}$，否则为零。写成
\begin{equation*}
\hat{\mathcal{H}} = -\sum_{\langle ij\rangle}\mathsf{J}\,\mathbf{S}_i\cdot\mathbf{S}_j,\tag{4.14}
\end{equation*}
求和号下的 $\langle ij\rangle$ 表示只对最近邻求和。考虑经典自旋，设最近邻自旋间
夹角为 $\phi_{ij}$ 且极小，即对所有 $i$、$j$ 有 $\phi_{ij}\ll1$。我们实质上是假定系统
表现铁磁性（见下一章）但自旋并未完全排齐。

在这些假设下，系统能量可写为
\begin{equation*}
E = -\mathsf{J}S^{2}\sum_{\langle ij\rangle}\cos\phi_{ij} = \text{常数} + \frac{\mathsf{J}S^{2}}{2}\sum_{\langle ij\rangle}\phi_{ij}^{2},\tag{4.15}
\end{equation*}
最后一步用了 $\phi_{ij}\ll1$ 时的 $\cos\phi_{ij}\approx1-\phi_{ij}^{2}/2$。下面
忽略常数项（它只是完全排齐态的能量）。定义约化矩 $\mathbf{m}=\mathbf{M}/M_{\mathrm{s}}$，
$M$ 为磁化强度，$M_{\mathrm{s}}$ 为饱和磁化强度。单位矢量 $\mathbf{m}$ 跟随自旋方向，
$m_x$、$m_y$、$m_z$ 可视作格点 $\mathbf{r}_{ij}$ 处自旋的方向余弦。用图 4.5 的记号，
可写
\begin{equation*}
|\phi_{ij}| \approx |\mathbf{m}_i-\mathbf{m}_j| \approx |(\mathbf{r}_{ij}\cdot\nabla)\mathbf{m}|,\tag{4.16}
\end{equation*}
于是能量可写为
\begin{equation*}
E = \mathsf{J}S^{2}\sum_{\langle ij\rangle}[(\mathbf{r}_{ij}\cdot\nabla)\mathbf{m}]^{2}.\tag{4.17}
\end{equation*}
连续介质极限下忽略晶格分立性，写成
\begin{equation*}
E = A\int_{V}[(\nabla m_x)^{2}+(\nabla m_y)^{2}+(\nabla m_z)^{2}]\,\mathrm{d}^{3}r,\tag{4.18}
\end{equation*}
其中
\begin{equation*}
A = 2\mathsf{J}S^{2}z/a,\tag{4.19}
\end{equation*}
$a$ 是最近邻距离，$z$ 是单胞内的格点数（简单立方 $z=1$，体心立方 bcc $z=2$，面心立
方 fcc $z=4$）。

这个结果也可以从另一视角得到。若断言交换来自非均匀的磁化分布，那么当非均匀性
较为平缓时，可以仅凭对称性导出结果：表达式必须在自旋旋转下、以及在磁化分量变号下
不变。因此寻找与晶体对称性相容的、M 的导数的最低偶数阶表达式。由于正比于
$\partial M_{\alpha}/\partial x_{\beta}$ 的项在 M 反向下变号，只剩磁化梯度的二次项。
最一般的表达式是
\begin{equation*}
E = \sum_{\alpha\beta\gamma}C_{\alpha\beta}\frac{\partial M_{\gamma}}{\partial x_{\alpha}}\frac{\partial M_{\gamma}}{\partial x_{\beta}},\tag{4.20}
\end{equation*}
$C_{\alpha\beta}$ 是具有晶体对称性的张量。立方晶体中它约化为
\begin{equation*}
E = C\sum_{\alpha\gamma}\frac{\partial M_{\gamma}}{\partial x_{\alpha}}\frac{\partial M_{\gamma}}{\partial x_{\alpha}} = C\left[(\nabla M_x)^{2}+(\nabla M_y)^{2}+(\nabla M_z)^{2}\right],\tag{4.21}
\end{equation*}
与上面式 4.18 的结果等价。式 4.21 有时写作 $C|\nabla\mathbf{M}|^{2}$，但要注意对
矢量取梯度须小心。

\section*{延伸阅读}

\begin{itemize}[itemsep=0.2em]
\item 关于交换相互作用的更多信息见 C.~Herring，载于《Magnetism》，G.~Rado 与
H.~Suhl 编，vol.~2B, p.1，Academic Press, New York 1966。
\item 磁性体系中相互作用的理论论述见 K.~Yosida，《Theory of magnetism》，
Springer 1996。
\item 真实体系中各种交换相互作用的综述见 P.~A.~Cox，《Transition metal oxides》，
OUP 1995。
\item 关于交换与交换相互作用非常透彻而有益的参考书是 D.~C.~Mattis，《The theory
of magnetism I》，Springer 1981。
\end{itemize}

\section*{习题}

\exer{4.1} 证明相距 $\mathbf{r}$ 的两个磁偶极子 $\boldsymbol{\mu}_1$ 与
$\boldsymbol{\mu}_2$ 的偶极能量为
\begin{equation*}
E = \frac{\mu_0}{4\pi r^{3}}\left[\boldsymbol{\mu}_1\cdot\boldsymbol{\mu}_2 - \frac{3}{r^{2}}(\boldsymbol{\mu}_1\cdot\mathbf{r})(\boldsymbol{\mu}_2\cdot\mathbf{r})\right].\tag{4.22}
\end{equation*}

\exer{4.2} 计算距质子 1 \AA{} 与 10 \AA{} 处磁场的大小，方向分别（a）平行与
（b）垂直于质子自旋方向。

\exer{4.3} 估计金属 Fe 中相邻两个 Fe 原子的交换耦合与偶极耦合之比。（Fe 的交换
常数可粗略地取 $k_{\mathrm{B}}T_{\mathrm{C}}$，$T_{\mathrm{C}}$ 为居里温度；Fe 的
$T_{\mathrm{C}}=1043$ K。）

\exer{4.4} 用（非弹性中子散射测得的）电子带宽 0.05 eV 粗略估计由超交换耦合的磁性
氧化物中交换常数的大小。取库仑能 $\sim$1 eV。进而估计反铁磁有序温度。

\exer{4.5} 考虑两个相互作用的自旋 $\frac{1}{2}$ 电子。好量子数为 $S=0$ 与 $1$，
故有一个三重态和一个单态，二者被能隙 $\Delta$ 分开。定义 $\Delta$ 的符号使
$\Delta>0$ 时单态（$S=0$）在下、$\Delta<0$ 时三重态在下。两种情形示于图 4.6(a)、
(b)。证明该模型的磁化率为
\begin{equation*}
\chi = \frac{2Ng\mu_{\mathrm{B}}^{2}}{k_{\mathrm{B}}T(3+\mathrm{e}^{\Delta/k_{\mathrm{B}}T})},\tag{4.23}
\end{equation*}
即 Bleaney--Bowers 方程，绘于图 4.6(c)。

\begin{figure}
\centering
\includegraphics[width=0.42\textwidth]{fig4_06.pdf}
\caption{由 $\Delta$ 耦合的两个自旋产生单态与三重态：（a）$\Delta>0$ 与（b）
$\Delta<0$ 时的能级。（c）由 Bleaney--Bowers 方程给出的磁化率，分别对 $\Delta>0$
（下曲线）、$\Delta=0$（中曲线）与 $\Delta<0$（上曲线）绘制。}
\end{figure}

\exer{4.6} 两个自旋为 $\mathbf{I}_1$、$\mathbf{I}_2$ 的核相距 $\mathbf{r}$，由偶极--
偶极相互作用耦合，哈密顿量为
\begin{equation*}
\hat{\mathcal{H}} = \frac{\mu_0}{4\pi r^{3}}\left[\hat{\boldsymbol{\mu}}_1\cdot\hat{\boldsymbol{\mu}}_2 - \frac{3}{r^{2}}(\hat{\boldsymbol{\mu}}_1\cdot\mathbf{r})(\hat{\boldsymbol{\mu}}_2\cdot\mathbf{r})\right],\tag{4.24}
\end{equation*}
其中 $\boldsymbol{\mu}_i=g_I\mu_{\mathrm{N}}\mathbf{I}_i$，$i=1,2$。证明哈密顿量可表示为
\begin{equation*}
\hat{\mathcal{H}} = \frac{\mu_0 g_I^{2}\mu_{\mathrm{N}}^{2}}{4\pi r^{3}}\left[A+B+C+D+E+F\right],\tag{4.25}
\end{equation*}
其中
\begin{equation*}
\begin{aligned}
A &= (1-3\cos^{2}\theta)I_{1z}I_{2z}\\
B &= -\frac{1}{4}(1-3\cos^{2}\theta)(I_{1}^{+}I_{2}^{-}+I_{1}^{-}I_{2}^{+})\\
C &= -\frac{3}{2}\sin\theta\cos\theta\,e^{-i\phi}(I_{1z}I_{2}^{+}+I_{1}^{+}I_{2z})\\
D &= -\frac{3}{2}\sin\theta\cos\theta\,e^{i\phi}(I_{1z}I_{2}^{-}+I_{1}^{-}I_{2z})\\
E &= -\frac{3}{4}\sin\theta\,e^{-2i\phi}I_{1}^{+}I_{2}^{+}\\
F &= -\frac{3}{4}\sin\theta\,e^{2i\phi}I_{1}^{-}I_{2}^{-},
\end{aligned}\tag{4.26}
\end{equation*}
角 $\theta$ 与 $\phi$ 是矢量 $\mathbf{r}$ 的球极角，升降算符为
$I_j^{\pm}=I_{jx}\pm\mathrm{i}I_{jy}$。对三重态（即 $\mathbf{I}=\mathbf{I}_1+\mathbf{I}_2$
大小 $I=1$）的两个质子（$I_1=I_2=\frac{1}{2}$），证明上列哈密顿量中的各项诱导如下
跃迁：$A$ 与 $B$ 为 $\Delta M=0$，$C$ 为 $\Delta M=+1$，$D$ 为 $\Delta M=-1$，
$E$ 为 $\Delta M=+2$，$F$ 为 $\Delta M=-2$。

\exer{4.7} 两个原子键合形成双原子分子。分子中有两个电子，坐标为 $\mathbf{r}_1$、
$\mathbf{r}_2$；两个核分别固定在 $\mathbf{R}_1$、$\mathbf{R}_2$，电荷各为 $Ze$。
忽略核间排斥，哈密顿量可写为 $\hat{\mathcal{H}}_0+\hat{\mathcal{H}}'$，其中
\begin{equation*}
\hat{\mathcal{H}}_0 = -\frac{\hbar^{2}}{2m}(\nabla_1^{2}+\nabla_2^{2}) - \frac{Ze^{2}}{4\pi\epsilon_0|\mathbf{r}_1-\mathbf{R}_1|} - \frac{Ze^{2}}{4\pi\epsilon_0|\mathbf{r}_2-\mathbf{R}_2|} = \hat{\mathcal{H}}_1+\hat{\mathcal{H}}_2,\tag{4.27}
\end{equation*}
是单粒子哈密顿量 $\hat{\mathcal{H}}_1$ 与 $\hat{\mathcal{H}}_2$（每个电子绕自己的
核）之和，而
\begin{equation*}
\hat{\mathcal{H}}' = \frac{e^{2}}{4\pi\epsilon_0|\mathbf{r}_1-\mathbf{r}_2|} - \frac{Ze^{2}}{4\pi\epsilon_0|\mathbf{r}_1-\mathbf{R}_2|} - \frac{Ze^{2}}{4\pi\epsilon_0|\mathbf{r}_2-\mathbf{R}_1|},\tag{4.28}
\end{equation*}
包含电子排斥能与每个电子受对方核吸引的能量项。证明分子能量为
\begin{equation*}
E = E_1 + E_2 + \frac{K\pm\mathsf{J}}{1\pm S}\tag{4.29}
\end{equation*}
$E_1$、$E_2$ 分别是 $\hat{\mathcal{H}}_1$、$\hat{\mathcal{H}}_2$ 的本征值，
\begin{equation*}
K = \int\psi_a^{*}(\mathbf{r}_1)\psi_b^{*}(\mathbf{r}_2)\hat{\mathcal{H}}'\psi_a(\mathbf{r}_1)\psi_b(\mathbf{r}_2)\,\mathrm{d}\mathbf{r}_1\mathrm{d}\mathbf{r}_2 = \int|\psi_a(\mathbf{r}_1)|^{2}\hat{\mathcal{H}}'|\psi_b(\mathbf{r}_2)|^{2}\,\mathrm{d}\mathbf{r}_1\mathrm{d}\mathbf{r}_2\tag{4.30}
\end{equation*}
称为库仑积分，
\begin{equation*}
\mathsf{J} = \int\psi_a^{*}(\mathbf{r}_1)\psi_b^{*}(\mathbf{r}_2)\hat{\mathcal{H}}'\psi_a(\mathbf{r}_2)\psi_b(\mathbf{r}_1)\,\mathrm{d}\mathbf{r}_1\mathrm{d}\mathbf{r}_2\tag{4.31}
\end{equation*}
称为交换积分，
\begin{equation*}
S = \int\psi_a^{*}(\mathbf{r}_1)\psi_b^{*}(\mathbf{r}_2)\psi_a(\mathbf{r}_2)\psi_b(\mathbf{r}_1)\,\mathrm{d}\mathbf{r}_1\mathrm{d}\mathbf{r}_2\tag{4.32}
\end{equation*}
称为交叠积分，且假定本征函数为
\begin{equation*}
\Psi_{\pm} = \frac{1}{\sqrt{2}}\left[\psi_a(\mathbf{r}_1)\psi_b(\mathbf{r}_2)\pm\psi_a(\mathbf{r}_2)\psi_b(\mathbf{r}_1)\right].\tag{4.33}
\end{equation*}

\exer{4.8} 三个 $S=1$ 原子置于等边三角形的三个顶点，可用哈密顿量
\begin{equation*}
\hat{\mathcal{H}} = -2\mathsf{J}(\hat{\mathbf{S}}_1\cdot\hat{\mathbf{S}}_2+\hat{\mathbf{S}}_2\cdot\hat{\mathbf{S}}_3+\hat{\mathbf{S}}_3\cdot\hat{\mathbf{S}}_1).\tag{4.34}
\end{equation*}
描述。证明系统的能量本征值为 0、$2\mathsf{J}$ 与 $6\mathsf{J}$。

% 第5章 磁有序与磁结构（Order and magnetic structures）
\chapter{磁有序与磁结构}

上一章介绍了固体中磁矩之间起作用的各种磁相互作用。本章讨论这些相互作用可以产生
哪些不同类型的磁基态。其中一些基态示于图 5.1：包括所有磁矩平行排列的铁磁体，
相邻磁矩反平行排列的反铁磁体，磁矩方向随位置沿圆锥或圆周进动的螺形与螺旋结构，
以及磁矩冻结在随机取向上的自旋玻璃。本章的任务是粗线条地说明：上一章讨论的
相互作用如何导致这些不同的基态。下一章将在更一般的语境中考察有序现象，并看到
有序是对称性破缺的后果。

\begin{figure}
\centering
\includegraphics[width=0.75\textwidth]{fig5_01.pdf}
\caption{有序体系中各种自旋排列：（a）铁磁体，（b）反铁磁体，（c）自旋玻璃，
（d）螺形与（e）螺旋结构。}
\end{figure}
\mnote{记号提醒：本书中 $\mathsf{J}$ 表示交换常数，$J$ 表示总角动量。}

\section{铁磁性}

铁磁体即使在没有外场时也有自发磁化强度，所有磁矩沿唯一的同一方向排列。\fn{1}{事实上，许多铁磁样品由于磁畴的存在，这一点在样品内并不处处成立。每个畴内有均匀的磁化强度，但相邻畴的磁化方向不同。磁畴详见第 6.7 节。}这一效应一般来源于上一章描述的交换相互作用。处于外场
$\mathbf{B}$ 中的铁磁体，需求解的哈密顿量为
\begin{equation*}
\hat{\mathcal{H}} = -\sum_{ij}\mathsf{J}_{ij}\,\mathbf{S}_i\cdot\mathbf{S}_j + g\mu_{\mathrm{B}}\sum_j\mathbf{S}_j\cdot\mathbf{B},\tag{5.1}
\end{equation*}
此时最近邻交换常数为正，以保证铁磁排列。右边第一项是海森伯交换能（见式 4.7），
第二项是塞曼能量（见式 1.35）。为简单起见，先设\fn{2}{这一假设将在 5.1.4 节放宽。}所处理的是没有轨道角动量的系统，即 L=0、J=S。

\subsection{铁磁体的外斯模型}

为推进式 5.1 的求解，需要作一个近似。定义第 $i$ 个格点上的有效分子场（molecular
field）
\begin{equation*}
\mathbf{B}_{\mathrm{mf}} = -\frac{2}{g\mu_{\mathrm{B}}}\sum_j\mathsf{J}_{ij}\,\mathbf{S}_j.\tag{5.2}
\end{equation*}
\bio{Pierre Weiss\\（1865--1946）}

聚焦第 $i$ 个自旋。它的能量来自塞曼部分 $g\mu_{\mathrm{B}}\mathbf{S}_i\cdot\mathbf{B}$
与交换部分。第 $i$ 个自旋与邻居的总交换相互作用是
$-2\sum_j\mathsf{J}_{ij}\mathbf{S}_i\cdot\mathbf{S}_j$，因子 2 来自重复计数。\fn{3}{见第 4.2.1 节。}这一项可写作
\begin{equation*}
-2\mathbf{S}_i\cdot\sum_j\mathsf{J}_{ij}\mathbf{S}_j = -g\mu_{\mathrm{B}}\mathbf{S}_i\cdot\mathbf{B}_{\mathrm{mf}}.\tag{5.3}
\end{equation*}
于是交换相互作用被邻居自旋产生的有效分子场 $B_{\mathrm{mf}}$ 取代。有效哈密顿量
现在可写作
\begin{equation*}
\hat{\mathcal{H}} = g\mu_{\mathrm{B}}\sum_i\mathbf{S}_i\cdot(\mathbf{B}+\mathbf{B}_{\mathrm{mf}})\tag{5.4}
\end{equation*}
看上去就像处于磁场 $B+B_{\mathrm{mf}}$ 中的顺磁体的哈密顿量。这一做法背后的假设是：
所有磁性离子感受同样的分子场。这相当可疑——尤其在接近磁相变的温度——下一章
再谈。对铁磁体，分子场的作用是使邻居磁矩排齐，因为主导的交换相互作用为正
（反铁磁体中则为负）。

既然分子场量度系统有序化的效果，可以设
\begin{equation*}
\mathbf{B}_{\mathrm{mf}} = \lambda\mathbf{M}\tag{5.5}
\end{equation*}
$\lambda$ 是把分子场强度表示为磁化强度函数的常数。铁磁体 $\lambda>0$。由于交换
相互作用涉及巨大的库仑能，铁磁体中的分子场往往极为巨大。

现在可以把这个问题当作置于磁场 $B+B_{\mathrm{mf}}$ 中的简单顺磁体来处理。低温下，
即使完全不加外场，内部分子场也能把磁矩排齐。注意：磁矩的排齐产生了内部分子场，
而正是这个场导致排齐——这是一个“先有鸡还是先有蛋”的局面。低温下磁有序是
自维持的。温度升高，热涨落逐渐破坏磁化强度，在某个临界温度有序被摧毁。这一模型
就是铁磁性的外斯模型（Weiss model）。

求该模型的解需要联立求解
\begin{equation*}
\frac{M}{M_{\mathrm{s}}} = \mathrm{B}_J(y)\tag{5.6}
\end{equation*}
（见式 2.38）与
\begin{equation*}
y = \frac{g_J\mu_{\mathrm{B}}J(B+\lambda M)}{k_{\mathrm{B}}T}\tag{5.7}
\end{equation*}
（见式 2.37）。\mnote{提醒：此阶段我们假定 J=S、L=0。}若没有来自分子场的 $\lambda M$ 项，
这与 2.4.3 节对顺磁体的处理完全一样。

\begin{figure}
\centering
\includegraphics[width=0.45\textwidth]{fig5_02.pdf}
\caption{B=0 时式 5.6 与 5.7 的图解法。}
\end{figure}

这两个方程可以图解。先限于 B=0，此时 $M=k_{\mathrm{B}}Ty/g_J\mu_{\mathrm{B}}J\lambda$：
把 M 对 y 作图得到的直线斜率正比于温度 T（图 5.2）。高温时，除了原点（y=0、
$M_{\mathrm{s}}=0$）外式 5.6 与 5.7 无联立解。当直线斜率小于布里渊函数在原点的
斜率时，情形改变：低温下出现三个解——一个 $M_{\mathrm{s}}=0$，另两个 $M_{\mathrm{s}}$
为 ± 某非零值。可以证明：当直线不如原点处的布里渊函数陡时，非零解稳定而零解
不稳定（若 $T<T_{\mathrm{C}}$ 时系统处于 $M_{\mathrm{s}}=0$，任何涨落——无论多小
——都会使系统转入两个稳定态之一）。于是低于某一温度时出现非零磁化强度，随材料
冷却而增长：物质即使在无外场时也被磁化。这一自发磁化强度正是铁磁性的特征。

转变温度可由“直线 $M=k_{\mathrm{B}}Ty/g_J\mu_{\mathrm{B}}J\lambda M_{\mathrm{s}}$ 与曲线
$M=M_{\mathrm{s}}\mathrm{B}_J(y)$ 在原点斜率相等”的条件得到。$y$ 小时
$\mathrm{B}_J(y)=(J+1)y/3J+O(y^{3})$。转变温度即居里温度 $T_{\mathrm{C}}$：
\begin{equation*}
T_{\mathrm{C}} = \frac{g_J\mu_{\mathrm{B}}(J+1)\lambda M_{\mathrm{s}}}{3k_{\mathrm{B}}} = \frac{n\lambda\mu_{\mathrm{eff}}^{2}}{3k_{\mathrm{B}}}.\tag{5.8}
\end{equation*}
于是分子场 $B_{\mathrm{mf}}=\lambda M_{\mathrm{s}}=3k_{\mathrm{B}}T_{\mathrm{C}}/g_J\mu_{\mathrm{B}}(J+1)$；
对 $J=\frac{1}{2}$、$T_{\mathrm{C}}\sim10^{3}$ K 的铁磁体，
$B_{\mathrm{mf}}=k_{\mathrm{B}}T_{\mathrm{C}}/\mu_{\mathrm{B}}\sim1500$ T——巨大的
有效磁场，反映了交换相互作用的强度。

这些方程的解作为温度的函数绘于图 5.3（对一系列 J）。曲线形状随 J 略有不同，但
一般特征不变：$T\geq T_{\mathrm{C}}$ 时磁化强度为零，$T<T_{\mathrm{C}}$ 时非零；
在 $T=T_{\mathrm{C}}$ 处磁化强度连续但其梯度不连续。在此分子场模型中，非磁相与
铁磁相之间的相变因此属于二级相变（second-order phase transition）。相变的级定义为
自由能最低阶导数在转变处出现不连续的阶数：一级相变在自由能一阶导数（体积、熵、
磁化强度等量）处不连续跳变，熵的跳变给出潜热；二级相变在二阶导数（压缩率、热容
等量）处不连续。此处不连续出现在磁化强度的梯度，即自由能的二阶导数——故转变为
二级。相变与临界指数将在 6.4 节详述。一些常见铁磁体的性质列于表 5.1。

\begin{figure}
\centering
\includegraphics[width=0.45\textwidth]{fig5_03.pdf}
\caption{对不同 J 值导出的平均场磁化强度随温度的函数。}
\end{figure}

\begin{table}
\centering
\caption{一些常见铁磁体的性质。}
\begin{small}
\begin{tabular}{lcc}
\toprule
材料 & $T_{\mathrm{C}}$（K） & 磁矩（$\mu_{\mathrm{B}}$/化学式单元）\\
\midrule
Fe & 1043 & 2.22\\
Co & 1394 & 1.715\\
Ni & 631 & 0.605\\
Gd & 289 & 7.5\\
MnSb & 587 & 3.5\\
EuO & 70 & 6.9\\
EuS & 16.5 & 6.9\\
\bottomrule
\end{tabular}
\end{small}
\end{table}

\subsection{磁化率}

在 $T\geq T_{\mathrm{C}}$ 下加小场 B 给出小磁化强度，可用布里渊函数的 $y\ll1$ 近似：
\begin{equation*}
\frac{M}{M_{\mathrm{s}}} \approx \frac{g_J\mu_{\mathrm{B}}(J+1)}{3k_{\mathrm{B}}}\left(\frac{B+\lambda M}{T}\right)\tag{5.9}
\end{equation*}
于是
\begin{equation*}
\frac{M}{M_{\mathrm{s}}} \approx \frac{T_{\mathrm{C}}}{\lambda M_{\mathrm{s}}}\left(\frac{B+\lambda M}{T}\right).\tag{5.10}
\end{equation*}
整理得
\begin{equation*}
\frac{M}{M_{\mathrm{s}}}\left(1-\frac{T_{\mathrm{C}}}{T}\right) \approx \frac{T_{\mathrm{C}}B}{\lambda M_{\mathrm{s}}}\tag{5.11}
\end{equation*}
从而
\begin{equation*}
\chi = \lim_{B\to0}\frac{\mu_0M}{B} \propto \frac{1}{T-T_{\mathrm{C}}}\tag{5.12}
\end{equation*}
即居里--外斯定律（Curie--Weiss law）。

\begin{figure}
\centering
\includegraphics[width=0.45\textwidth]{fig5_04.pdf}
\caption{$B\neq0$ 时式 5.6 与 5.7 的图解法。}
\end{figure}

\subsection{磁场的效果}

加磁场的效果是把图解中的直线向右平移（见图 5.4）：一切温度下都有 $M\neq0$ 的解，
相变被消除。非零磁场中的铁磁体总有能量优势使磁化强度非零、磁矩沿场排列。相变的
消除见图 5.5——图中对一系列磁场绘出式 5.6 与 5.7 的图解。此模型无须考虑磁场施加
的方向：无论沿哪个方向加场，磁化强度都会转向跟随；模型中不存在铁磁体自身特有的
方向。真实铁磁体则不然——必须考虑材料的磁各向异性（见下一章）。

\begin{figure}
\centering
\includegraphics[width=0.45\textwidth]{fig5_05.pdf}
\caption{$J=\frac{1}{2}$ 时对不同的外加场 B 计算的平均场磁化强度随温度的函数。
只有 B=0 时才有相变。}
\end{figure}

在 $T=T_{\mathrm{C}}$ 处，磁场的效应可以解析求出：小场下 $M\propto B^{1/3}$。证明
需要取 $\mathrm{B}_J(y)$ 泰勒展开的下一项。写 $\mathrm{B}_J(y)=(J+1)y/3J-\zeta y^{3}+O(y^{5})$
（$\zeta$ 为常数），联立 $M=M_{\mathrm{s}}\mathrm{B}_J(y)$ 与
\begin{equation*}
y = \frac{g_J\mu_{\mathrm{B}}J(B+\lambda M)}{k_{\mathrm{B}}T_{\mathrm{C}}} = \frac{(B+\lambda M)}{\zeta M_{\mathrm{s}}}\tag{5.13}
\end{equation*}
得
\begin{equation*}
M = M_{\mathrm{s}}\zeta\frac{(B+\lambda M)}{\lambda} - \zeta M_{\mathrm{s}}\left(\frac{3J(B+\lambda M)}{\lambda(J+1)M_{\mathrm{s}}}\right)^{3}\tag{5.14}
\end{equation*}
从而
\begin{equation*}
B \propto (B+\lambda M)^{3}\tag{5.15}
\end{equation*}
由于 $\lambda M\gg B$，右侧由 $M^{3}$ 项主导，故 $M\propto B^{1/3}$。

\subsection{分子场的起源}

1907 年外斯提出分子场模型时，令他失望的是：为符合自然界中很大的 $T_{\mathrm{C}}$，
常数 $\lambda$ 必须取得很大。只考虑偶极场无法解释需要 $\sim10^{3}$ T 的内场
（如前所述，这是解释 Fe 居里温度所必需的）。三十年后，海森伯证明：正是涉及巨大
库仑能的交换相互作用，才是巨大分子场的来源。\fn{4}{分子场是一个方便的虚构——不要以为铁磁体中的电子真的感受 $\sim10^{3}$ T 的磁场。交换相互作用纯属静电效应。谈分子场时，我们是假装交换相互作用其实是内部磁场。要点是：如果它真的存在而海森伯交换不存在，那么分子场必须 $\sim10^{3}$ T 才能解释我们看到的现象。}

由 $\lambda$ 参数化的分子场可与由 $\mathsf{J}_{ij}$ 参数化的交换相互作用大小联系
起来。设交换相互作用只在一个离子的 z 个最近邻内起作用、取值 $\mathsf{J}$，则用
式 5.1、5.4、5.5 容易得到
\begin{equation*}
\lambda = \frac{2z\mathsf{J}}{ng^{2}\mu_{\mathrm{B}}^{2}}.\tag{5.16}
\end{equation*}
再用式 5.8，居里温度可写为
\begin{equation*}
T_{\mathrm{C}} = \frac{2z\mathsf{J}J(J+1)}{3k_{\mathrm{B}}}.\tag{5.17}
\end{equation*}

迄今我们假定 L=0、J=S——这对多数 3d 离子成立。交换发生在自旋自由度之间，因此
依赖 S；而离子的磁矩依赖总（自旋+轨道）角动量 J。对 3d 离子二者是同一回事，因为
L 被淬灭。

对 4f 离子则不然：S 不是好量子数，J 才是。于是 S 垂直于 J 的分量必平均为零，
平行于 J 的分量才守恒——必须把 S 投影到 J 上。由 $J=L+S$，$L+2S$ 等于 $g_JJ$
加一个垂直于 J 的分量，故 S 中作为好量子数的分量是 $(g_J-1)J$。各 4f 离子的
$(g_J-1)$ 列于表 5.2。可见对所谓的重稀土（Gd 至 Yb），S 与 J 平行；对所谓的轻稀土
（Ce 至 Sm），S 与 J 反平行。

\begin{table}
\centering
\caption{用洪特规则得到的 4f 离子 $g$ 因子。}
\begin{small}
\begin{tabular}{lccccccc}
\toprule
离子 & 壳层 & $S$ & $L$ & $J$ & $g_J$ & $g_J-1$ & $(g_J-1)^{2}J(J+1)$\\
\midrule
Ce$^{3+}$ & $4f^{1}$ & $\frac{1}{2}$ & 3 & $\frac{5}{2}$ & $\frac{6}{7}$ & $-\frac{1}{7}$ & 0.18\\
Pr$^{3+}$ & $4f^{2}$ & 1 & 5 & 4 & $\frac{4}{5}$ & $-\frac{1}{5}$ & 0.80\\
Nd$^{3+}$ & $4f^{3}$ & $\frac{3}{2}$ & 6 & $\frac{9}{2}$ & $\frac{72}{99}$ & $-\frac{27}{99}$ & 1.84\\
Pm$^{3+}$ & $4f^{4}$ & 2 & 6 & 4 & $\frac{3}{5}$ & $-\frac{2}{5}$ & 3.20\\
Sm$^{3+}$ & $4f^{5}$ & $\frac{5}{2}$ & 5 & $\frac{5}{2}$ & $\frac{2}{7}$ & $-\frac{5}{7}$ & 4.46\\
Eu$^{3+}$ & $4f^{6}$ & 3 & 3 & 0 & -- & -- & --\\
Gd$^{3+}$ & $4f^{7}$ & $\frac{7}{2}$ & 0 & $\frac{7}{2}$ & 2 & 1 & 15.75\\
Tb$^{3+}$ & $4f^{8}$ & 3 & 3 & 6 & $\frac{3}{2}$ & $\frac{1}{2}$ & 10.50\\
Dy$^{3+}$ & $4f^{9}$ & $\frac{5}{2}$ & 5 & $\frac{15}{2}$ & $\frac{4}{3}$ & $\frac{1}{3}$ & 7.08\\
Ho$^{3+}$ & $4f^{10}$ & 2 & 6 & 8 & $\frac{5}{4}$ & $\frac{1}{4}$ & 4.50\\
Er$^{3+}$ & $4f^{11}$ & $\frac{3}{2}$ & 6 & $\frac{15}{2}$ & $\frac{6}{5}$ & $\frac{1}{5}$ & 2.55\\
Tm$^{3+}$ & $4f^{12}$ & 1 & 5 & 6 & $\frac{7}{6}$ & $\frac{1}{6}$ & 1.17\\
Yb$^{3+}$ & $4f^{13}$ & $\frac{1}{2}$ & 3 & $\frac{7}{2}$ & $\frac{8}{7}$ & $\frac{1}{7}$ & 0.32\\
Lu$^{3+}$ & $4f^{14}$ & 0 & 0 & 0 & -- & -- & --\\
\bottomrule
\end{tabular}
\end{small}
\end{table}

用 $(g_J-1)\mathbf{J}$ 替换 $\mathbf{S}$ 中守恒的部分，表达式
$-\sum_{ij}\mathsf{J}_{ij}\mathbf{S}_i\cdot\mathbf{S}_j$ 可以换成
$-\sum_{ij}(g_J-1)^{2}\mathsf{J}_{ij}\mathbf{J}_i\cdot\mathbf{J}_j$。再用磁矩
$\boldsymbol{\mu}=-g_J\mu_{\mathrm{B}}\mathbf{J}$ 重复导出式 5.16 的计算，得
\begin{equation*}
\lambda = \frac{2z\mathsf{J}(g_J-1)^{2}}{ng_J^{2}\mu_{\mathrm{B}}^{2}}.\tag{5.18}
\end{equation*}
（对轨道淬灭的过渡金属离子特例 $g_J=2$，式 5.18 退化为式 5.16。）居里温度可写为
\begin{equation*}
T_{\mathrm{C}} = \frac{2z(g_J-1)^{2}\mathsf{J}}{3k_{\mathrm{B}}}J(J+1).\tag{5.19}
\end{equation*}
因此预期临界温度正比于 de Gennes 因子 $(g_J-1)^{2}J(J+1)$（其值亦列于表 5.2）。
de Gennes 因子最大的 Gd 是铁磁体，但稀土金属表现出各种不同的基态——包括反铁磁与
螺旋磁性，见下面几节。

\section{反铁磁性}

若交换相互作用为负（$\mathsf{J}<0$），分子场的取向使最近邻磁矩反平行排列有利——
这就是反铁磁性（antiferromagnetism）。它常出现在可以看作两套相互贯穿的子晶格的
体系中（见图 5.6）：一套子晶格上磁矩向上，另一套向下。图 5.6 中每个磁矩的最近邻
全部在另一套子晶格上。于是我们起初假定：一套子晶格上的分子场正比于另一套子晶格的
磁化强度，且无外加磁场。

\begin{figure}
\centering
\includegraphics[width=0.45\textwidth]{fig5_06.pdf}
\caption{反铁磁体可以分解为两套相互贯穿的子晶格。}
\end{figure}

\subsection{反铁磁体的外斯模型}

把“向上”子晶格标记为 +、“向下”子晶格标记为 $-$，则每套子晶格上的分子场为
\begin{equation*}
\begin{aligned}
B_{+} &= -|\lambda|M_{-}\\
B_{-} &= -|\lambda|M_{+},
\end{aligned}\tag{5.20}
\end{equation*}
$\lambda$ 是现在取负值的分子场常数。每套子晶格上
\begin{equation*}
M_{\pm} = M_{\mathrm{s}}\mathrm{B}_J\left(-\frac{g_J\mu_{\mathrm{B}}J|\lambda|M_{\mp}}{k_{\mathrm{B}}T}\right).\tag{5.21}
\end{equation*}
两套子晶格除磁矩方向外一切等价，故
\begin{equation*}
|M_{+}| = |M_{-}| \equiv M,\tag{5.22}
\end{equation*}
从而
\begin{equation*}
M = M_{\mathrm{s}}\mathrm{B}_J\left(\frac{g_J\mu_{\mathrm{B}}J|\lambda|M}{k_{\mathrm{B}}T}\right).\tag{5.23}
\end{equation*}
这与铁磁体的对应方程（式 5.6、5.7）几乎相同，每套子晶格上的分子场将严格遵循图 5.3
的形式，并在转变温度以上消失。该转变温度即奈尔温度（Néel temperature）$T_{\mathrm{N}}$：
\begin{equation*}
T_{\mathrm{N}} = \frac{g_J\mu_{\mathrm{B}}(J+1)|\lambda|M_{\mathrm{s}}}{3k_{\mathrm{B}}} = \frac{n|\lambda|\mu_{\mathrm{eff}}^{2}}{3k_{\mathrm{B}}}.\tag{5.24}
\end{equation*}
虽然每套子晶格的磁化强度都遵循图 5.3 的形式，但两者方向相反，反铁磁体的净磁化强度
$M_{+}+M_{-}$ 为零。可以定义错向磁化强度（staggered magnetization）为两套子晶格
磁化强度之差 $M_{+}-M_{-}$：它在 $T_{\mathrm{N}}$ 以下非零，因而可用作反铁磁体的
序参量（order parameter）。\fn{5}{序参量将在第 6.1 节定义。}

\subsection{磁化率}

$T>T_{\mathrm{N}}$ 时，小外场的效果可以与铁磁体同样计算——展开布里渊函数
$\mathrm{B}_J(y)=(J+1)y/3J+O(y^{3})$——得磁化率
\begin{equation*}
\chi = \lim_{B\to0}\frac{\mu_0M}{B} \propto \frac{1}{T+T_{\mathrm{N}}},\tag{5.25}
\end{equation*}
又是居里--外斯定律，只是 $-T_{\mathrm{C}}$ 换成了 $+T_{\mathrm{N}}$。

这给出解读顺磁态（即转变温度以上）磁化率数据的现成办法：把磁化率拟合为居里--外斯
形式
\begin{equation*}
\chi \propto \frac{1}{T-\theta},\tag{5.26}
\end{equation*}
$\theta$ 称为外斯温度。$\theta=0$：材料是顺磁体（见式 2.44）；$\theta>0$：材料是
铁磁体，预期 $\theta=T_{\mathrm{C}}$（见式 5.12）；$\theta<0$：材料是反铁磁体，
预期 $\theta=-T_{\mathrm{N}}$（见式 5.25）。这些可能性示于图 5.7。
\bio{Louis E. F. Néel\\（1904--2000）}

\begin{figure}
\centering
\includegraphics[width=0.32\textwidth]{fig5_07.pdf}
\caption{居里--外斯定律 $\chi\propto1/(T-\theta)$（$T>\theta$）：（a）三种情形——
$\theta=0$（顺磁体）、$\theta=\Theta>0$（铁磁体）与 $\theta=-\Theta<0$（反铁磁体）。
（b）$1/\chi$ 对 $T$ 作图为直线，温度轴截距给出 $\theta$。（c）$\chi T$ 对 $T$ 可为
常数（$\theta=0$）、随 $T$ 减小而增大（$\theta>0$）或随 $T$ 减小而减小
（$\theta<0$）。}
\end{figure}

实验测得的反铁磁体外斯温度常常离 $-T_{\mathrm{N}}$ 很远（一些常见反铁磁体的数据
见表 5.3）。这一偏差主要源于我们“一套子晶格上的分子场只依赖另一套子晶格磁化强度”
的假设。更现实的计算见习题 5.3。

\begin{table}
\centering
\caption{一些常见反铁磁体的性质。}
\begin{small}
\begin{tabular}{lccc}
\toprule
材料 & $T_{\mathrm{N}}$（K） & $\theta$（K） & $J$\\
\midrule
MnF$_2$ & 67 & $-80$ & $\frac{5}{2}$\\
MnO & 116 & $-510$ & $\frac{5}{2}$\\
CoO & 292 & $-330$ & $\frac{3}{2}$\\
FeO & 116 & $-610$ & 2\\
Cr$_2$O$_3$ & 307 & $-485$ & $\frac{3}{2}$\\
$\alpha$-Fe$_2$O$_3$ & 950 & $-2000$ & $\frac{5}{2}$\\
\bottomrule
\end{tabular}
\end{small}
\end{table}

在 $T_{\mathrm{N}}$ 以下对反铁磁体加磁场，比 $T_{\mathrm{C}}$ 以下的铁磁体复杂，
因为磁场的施加方向至关重要。磁矩沿场排列不再有能量优势：若两套子晶格磁化强度
等大反向，一套上省的能量会被另一套付出的能量抵消。

先考虑绝对零度（T=0），热扰动可以忽略：$|M_{+}|=|M_{-}|=M_{\mathrm{s}}$。若加
小磁场平行于一套子晶格的磁化方向（从而反平行于另一套），相当于给各子晶格的局域场
加上或减去一个小项；两套子晶格均已饱和，这不起作用，材料中诱导的净磁化强度为零，
故 $\chi_{\parallel}=0$。若小场垂直于磁化方向施加，则两套子晶格的磁化强度都略微
倾斜，沿外场产生一个磁化强度分量（见图 5.8），故 $\chi_{\perp}\neq0$。

\begin{figure}
\centering
\includegraphics[width=0.45\textwidth]{fig5_08.pdf}
\caption{$\chi_{\perp}$ 的起源。垂直于子晶格磁化方向加小磁场 B，两套子晶格的磁化
强度都略微倾斜，沿外场产生一个磁化强度分量。}
\end{figure}

温度升高但仍低于 $T_{\mathrm{N}}$ 时，热涨落减小每套子晶格上的分子场。这对平行
情形影响巨大——场使一套子晶格磁化增强、另一套减弱；而垂直情形下 $M_{+}$ 与
$M_{-}$ 等量减小，受小场的影响也对称，升温几乎没有影响。于是 $\chi_{\perp}$ 不
依赖温度，而 $\chi_{\parallel}$ 从 0 上升到 $T\to T_{\mathrm{N}}$ 时的
$\chi_{\perp}$（见图 5.9）。

\begin{figure}
\centering
\includegraphics[width=0.45\textwidth]{fig5_09.pdf}
\caption{温度对 $\chi_{\parallel}$ 与 $\chi_{\perp}$ 的影响。}
\end{figure}

\subsection{强磁场的效果}

先考虑 T=0 的反铁磁体在强磁场下的行为，以避开热涨落的干扰。场足够大时终将压倒
任何内部分子场，迫使所有磁矩彼此平行。但随场增大，虽然最终结局明确，通往终点的
路径却强烈依赖于外场相对子晶格磁化强度初始方向的取向。

若外场垂直于子晶格磁化，则随场增大，磁矩只是越来越弯（$\phi$ 逐渐变小，见
图 5.8），直到与外场排齐。

若外场平行于子晶格磁化，情形更有趣。小场下磁矩不转动、保持原位（图 5.10(a)）；
但达到临界场时，系统突然跳入另一种组态（图 5.10(b)），称为自旋翻转跳跃转变
（spin-flop transition）。场继续增大时角 $\theta$ 逐渐变小，最终磁矩与外场排齐。

\begin{figure}
\centering
\includegraphics[width=0.42\textwidth]{fig5_10.pdf}
\caption{磁场平行于子晶格磁化强度施加。（a）小场下无事发生，系统保持在反铁磁相。
（b）超过临界场后系统经自旋翻转跳跃转变进入自旋翻转跳跃相。}
\end{figure}

这些效应可以定量计算。设 $M_{+}$ 与场的夹角为 $\theta$（逆时针量），$M_{-}$ 与场
的夹角为 $\phi$（顺时针量），磁场沿晶体学 $z$ 轴施加。反铁磁相对应 $\theta=0$、
$\phi=\pi$（图 5.10(a)），自旋翻转跳跃相对应 $\theta=\phi$。需要判断哪个相能量
更低。

设总能量 E 是各子晶格塞曼能量与一个代表交换耦合（依赖两套子晶格磁矩相对取向）的
项之和：
\begin{equation*}
E = -MB\cos\theta - MB\cos\phi + AM^{2}\cos(\theta+\phi),\tag{5.27}
\end{equation*}
A 是与交换耦合有关的常数。为模拟磁各向异性，须加上形如
\begin{equation*}
-\frac{1}{2}\Delta(\cos^{2}\theta+\cos^{2}\phi)\tag{5.28}
\end{equation*}
的项（$\Delta$ 为小常数）——它反映磁化强度确实偏好沿某晶体学轴（此处 $z$ 轴）
排列：$\theta$、$\phi$ 偏好取 0 或 $\pi$ 而非中间值。反铁磁情形（$\theta=0$、
$\phi=\pi$）：$E=-AM^{2}-\Delta$，与场无关。自旋翻转跳跃情形（$\phi=\theta$）：
\begin{equation*}
E = -2MB\cos\theta + AM^{2}\cos2\theta - \Delta\cos^{2}\theta.\tag{5.29}
\end{equation*}
条件 $\partial E/\partial\theta=0$ 表明（忽略各向异性项）在 $\theta=\cos^{-1}[B/2AM]$
处有能量极小。代回 E 并作图得图 5.11：低于临界场 $B_{\text{spin-flop}}$ 时反铁磁相
能量最低；临界场处系统从一态切换到另一态，发生自旋翻转跳跃转变；高于此场时自旋
翻转跳跃相能量最低。

\begin{figure}
\centering
\includegraphics[width=0.45\textwidth]{fig5_11.pdf}
\caption{反铁磁相与自旋翻转跳跃相的能量随 B 的变化。}
\end{figure}

\begin{figure}
\centering
\includegraphics[width=0.45\textwidth]{fig5_12.pdf}
\caption{（a）对反铁磁体施加平行磁场时的磁化强度：起初无事发生，随后在 $B_1$ 处
发生自旋翻转跳跃转变进入自旋翻转跳跃相；磁场继续转动磁矩，直至 $B_2$ 处饱和。
（b）若自旋强烈偏好沿平行方向排列，则不发生自旋翻转跳跃，而是在 $B_3$ 处发生
自旋翻转（spin-flip）转变。两图均为绝对零度的预期曲线；有限温度会把尖角磨圆。
这也称为变磁转变（metamagnetic transition）。}
\end{figure}

大平行磁场下反铁磁体的磁化强度示于图 5.12(a)：自旋翻转跳跃转变之前毫无变化，其上
磁化强度稳步增大直至饱和。若各向异性效应很强（$\Delta$ 大），会出现另一种效应：
外场沿 z 时不再发生自旋翻转跳跃，而是发生自旋翻转（spin-flip）转变——B 达到临界
值时一套子晶格的磁化强度突然反转，系统一步进入铁磁态（图 5.12(b)）。

\subsection{反铁磁有序的类型}

反铁磁性的另一个复杂之处：把等量的上下自旋排在一个晶格上有很多种方式，且不同
排列还取决于晶格类型。图 5.13、5.14 给出若干可能的排列。

立方钙钛矿中磁性原子排在简单立方晶格上，G 型有序（图 5.13(d)）非常常见——经氧
原子的超交换相互作用迫使所有最近邻磁性原子反铁磁排列。这只对 G 型有序成立，
例如 LaFeO$_3$ 与 LaCrO$_3$。LaMnO$_3$ 也是立方钙钛矿，却表现 A 型有序
（图 5.13(a)）：铁磁排列的 (100) 面交替取向。原因是 $\mathrm{Mn^{3+}}$ 离子的
扬--特勒畸变使 (100) 面内 Mn--O 键长短交替，相邻 $\mathrm{Mn^{3+}}$ 离子的轨道
取向不同，超交换导致一个原子的占据轨道与邻居的未占据轨道之间的相互作用：面内
相互作用为铁磁，面间则因常规超交换而为反铁磁。

\begin{figure}
\centering
\includegraphics[width=0.24\textwidth]{fig5_13.pdf}
\caption{简单立方晶格上可以出现的四种反铁磁有序（A、C、E、G 型）。两种自旋态以
+、$-$ 标记。}
\end{figure}

\begin{figure}
\centering
\includegraphics[width=0.28\textwidth]{fig5_14.pdf}
\caption{体心立方晶格上可以出现的三种反铁磁有序。}
\end{figure}

\section{亚铁磁性}

以上对反铁磁性的处理假定两套子晶格等价。若因晶体学上的原因二者不等价呢？此时
两套子晶格的磁化强度未必等大反向，因而不相抵消——材料将有净磁化强度。这一现象
称为亚铁磁性（ferrimagnetism）。由于两套子晶格上的分子场不同，子晶格的自发磁化
强度对温度的依赖可以大相径庭，净磁化强度本身因此可能有复杂的温度依赖。有时低温
下一套子晶格主导磁化强度，较高温度时另一套主导；此时净磁化强度可在某个温度降至
零并变号，该温度称为补偿温度（compensation temperature）。因此亚铁磁体的磁化率
不遵循居里--外斯定律。

铁氧体（ferrites）是一族亚铁磁体：化学式 $MO\cdot\mathrm{Fe_2O_3}$ 的化合物，
M 是二价阳离子（如 $\mathrm{Zn^{2+}}$、$\mathrm{Co^{2+}}$、$\mathrm{Fe^{2+}}$、
$\mathrm{Ni^{2+}}$、$\mathrm{Cu^{2+}}$ 或 $\mathrm{Mn^{2+}}$）。晶体结构为尖晶石
结构（spinel），含两类格位：四面体位（四个氧邻居，称 A 位）与八面体位（六个氧
邻居，称 B 位）；B 位数目是 A 位两倍。两套子晶格不等价——存在两类晶体学格位，
且容纳两种不同离子。正尖晶石（normal spinel）中 $M^{2+}$ 阳离子坐 A 位，
$\mathrm{Fe^{3+}}$（${}^{6}S_{5/2}$，磁矩 $5\mu_{\mathrm{B}}$）阳离子坐 B 位。
反尖晶石（inverse spinel）中 $M^{2+}$ 阳离子坐一半 B 位，$\mathrm{Fe^{3+}}$ 阳离子
占据另一半 B 位与全部 A 位；反尖晶石中 A、B 位上 $\mathrm{Fe^{3+}}$ 阳离子的磁矩
反平行，样品总磁矩仅来自 $M^{2+}$ 离子。

\begin{table}
\centering
\caption{一些常见亚铁磁体的性质。}
\begin{small}
\begin{tabular}{lccc}
\toprule
材料 & $T_{\mathrm{C}}$（K） & 磁矩（$\mu_{\mathrm{B}}$/化学式单元） & 补偿温度（K）\\
\midrule
Fe$_3$O$_4$ & 858 & 4.1 & --\\
CoFe$_2$O$_4$ & 793 & 3.7 & --\\
NiFe$_2$O$_4$ & 858 & 2.3 & --\\
CuFe$_2$O$_4$ & 728 & 1.3 & --\\
Y$_3$Fe$_5$O$_12$ & 560 & 5.0 & --\\
Gd$_3$Fe$_5$O$_12$ & 564 & 16.0 & 290\\
Dy$_3$Fe$_5$O$_12$ & 563 & 18.2 & 220\\
Ho$_3$Fe$_5$O$_12$ & 567 & 15.2 & 137\\
\bottomrule
\end{tabular}
\end{small}
\end{table}

M=Fe 的情形即 $\mathrm{Fe_3O_4}$（半导体，与其他绝缘体铁氧体不同），已在 4.2.5 节
讨论过。

另一族亚铁磁体是石榴石（garnets），化学式 $R_3\mathrm{Fe_5O_{12}}$，R 是三价稀土
原子。晶体结构为立方，但单胞相当复杂：三个 $\mathrm{Fe^{3+}}$ 离子在四面体位，
两个在八面体位，$R^{3+}$ 离子在十二面体对称性的格位上。钇铁石榴石（YIG，
$\mathrm{Y_3Fe_5O_{12}}$）中 $Y^{3+}$ 无磁矩（$4d^{0}$），四面体位 $\mathrm{Fe^{3+}}$
的磁矩与八面体位的反平行，净磁矩为 $5\mu_{\mathrm{B}}$。

钡铁氧体（$\mathrm{BaFe_{12}O_{19}}=\mathrm{BaO}\cdot6\mathrm{Fe_2O_3}$）具有
六角结构：八个 $\mathrm{Fe^{3+}}$ 离子与其他四个反平行，净磁矩相当于四个
$\mathrm{Fe^{3+}}$ 离子，即 $20\mu_{\mathrm{B}}$。粉体形式下它用于磁记录，因为
矫顽力高（见 6.7.9 节）。一些常见亚铁磁体的性质列于表 5.4。

大多数亚铁磁体是电绝缘体，这一事实造就了它们的许多实际应用。铁磁体往往是金属，
不适合涉及振荡磁场的应用：快速变化的磁场感应电压并在导体中驱动电流（涡流，
eddy currents），电流在金属中造成电阻性发热（涡流损耗）。因此，当需要自发磁化
材料工作在高频时，常可用许多亚铁磁体——绝缘体中感应电压无法驱动显著的涡流。
固体铁氧体磁芯用于许多高频应用：需要高磁导率、低能量损耗的天线与变压器，以及
微波器件。此外，许多亚铁磁体本就是氧化物，比金属铁磁体更耐腐蚀。

\section{螺旋有序}

许多稀土金属的晶体结构使原子分层排列。先考虑与镝相关的情形：层内原子磁矩铁磁
排列，层间相互作用由最近邻交换常数 $\mathsf{J}_1$ 与次近邻交换常数 $\mathsf{J}_2$
描述。若相邻两基面（即层所在平面）中磁矩间夹角为 $\theta$（见图 5.15(a)），系统
能量可写为
\begin{equation*}
E = -2NS^{2}(\mathsf{J}_1\cos\theta + \mathsf{J}_2\cos2\theta),\tag{5.30}
\end{equation*}
N 是每层原子数。能量在 $\partial E/\partial\theta=0$ 时取极小：
\begin{equation*}
(\mathsf{J}_1 + 4\mathsf{J}_2\cos\theta)\sin\theta = 0.\tag{5.31}
\end{equation*}
解为 $\sin\theta=0$（即 $\theta=0$ 或 $\theta=\pi$，铁磁或反铁磁），或
\begin{equation*}
\cos\theta = -\frac{\mathsf{J}_1}{4\mathsf{J}_2}.\tag{5.32}
\end{equation*}
最后这个解对应螺旋有序（helical order，又称螺旋磁性 helimagnetism），当
$\mathsf{J}_2<0$ 且 $|\mathsf{J}_1|<4|\mathsf{J}_2|$ 时优于铁磁或反铁磁
（见图 5.15(b)）。螺旋的螺距一般与晶格常数不可公度，故晶体中没有两层恰好具有
相同的自旋方向。

\begin{figure}
\centering
\includegraphics[width=0.4\textwidth]{fig5_15.pdf}
\caption{（a）螺旋磁有序。（b）由最近邻交换常数 $\mathsf{J}_1$ 与次近邻交换常数
$\mathsf{J}_2$ 耦合的层模型的相图。}
\end{figure}

螺旋结构出现在许多磁性体系中，最著名的是稀土金属。许多稀土金属具有六角密堆晶体
结构；螺旋轴垂直于六角密堆面，沿通常定义的 c 轴。Tb、Dy、Ho 中自旋旋转的平面是
六角密堆面；而 Er、Tm 中易轴为 c 轴，故在某温度范围内自旋的 c 分量被正弦调制。
还可能出现螺形结构（图 5.1(d)）：自旋有沿易轴（垂直于旋转面）的分量，同时面内
分量绕螺旋进动。

稀土金属中的交换相互作用是由传导电子传递的间接 RKKY 相互作用：比超交换长程得多，
且随距离变号。上述只含最近邻与次近邻相互作用的模型因此过分简化——计算波矢依赖
的交换相互作用 $\mathsf{J}(\mathbf{q})$ 需要每个稀土金属费米面的细节；若它在某
波矢 q 取极大，则可诱发波矢为 q 的螺旋磁性。另一个因素是稀土金属中晶体场的大
影响：晶体场分裂电子态，通常不仅基态、某些激发态也有热布居——使理解这些体系
所需的分析更加复杂。

\section{自旋玻璃}

迄今考虑的材料每个单胞里都有一个或多个磁矩。若从一个非磁晶格出发，稀疏地、随机
地散布磁性原子，会怎样？直觉或许会认为：结局是彻底随机的东西，不大可能表现出从高温
无序态到低温有序态的相变。这只对了一部分：这类体系虽然内在随机，却在某个特定
温度表现出近似于相变的行为，进入一个虽非有序、却与高温无序态截然不同的态。自旋
玻璃（spin glass）可定义为：具有混合相互作用的随机磁性系统，其特征是在明确的环境
温度 $T_{\mathrm{f}}$（冻结温度）以下自旋发生随机然而协同的冻结，冻结态是亚稳的，
且没有通常的磁长程序。

\begin{example}{5.1}
研究得很好的自旋玻璃例子是 CuMn 合金，Mn 浓度为几个原子百分比。Mn 离子稀疏，
其磁矩经 Cu 中传导电子传递的 RKKY 相互作用（见 4.2.4 节与 7.7.3 节）彼此耦合。
RKKY 相互作用随符号振荡，相互作用可铁磁可反铁磁。净结果是一个内置大量阻挫
（frustration）、没有明确基态而有大量可能基态的系统。高温下 Mn 磁矩热涨落；温度
降低时磁矩变慢、局域关联的团簇逐渐建立。在 $T_{\mathrm{f}}$ 处磁矩卡进这些简并
基态之一，冻结。CuMn 中位置随机性给出自旋间距离分布，从而提供阻挫；自旋玻璃也
见于具有键随机性（bond-randomness）的体系——最近邻相互作用在 $+\mathsf{J}$ 与
$-\mathsf{J}$ 之间变化。自旋玻璃与阻挫将在 8.2 节更详细讨论。
\end{example}

\section{核有序}

迄今只考虑了电子磁矩。有些核具有磁矩——它们能有序吗？答案是能，但效应极小，
只在低温且材料没有电子磁矩时才重要。核自旋深居原子内部，彼此耦合不强，其间不
存在直接交换相互作用。核磁矩比电子磁矩小约三个量级，而偶极相互作用正比于磁矩
的平方，故任何偶极相互作用都比电子情形低六个量级。只要电子磁矩存在，任何核效应
都完全被淹没。核磁矩之小意味着核自旋体系只能在低于 $\sim1\ \mu$K 的温度显示磁
有序。金属样品中核磁矩之间也存在 RKKY 耦合的可能——交换由传导电子传递——但其
大小与核偶极相互作用同量级。

尽管如此，在若干没有电子磁矩的材料上已做出惊人的实验，核自旋有序真的被观察到：
铜在 58 nK 发生一级相变进入反铁磁有序的核基态；银中核自旋的奈尔温度为 560 pK。
这些极低温度如何获得？实验依靠一个事实：同一块样品中，晶格与核可以温度不同。
用绝热退磁（见 2.6 节）冷却核，核在由自旋--自旋弛豫时间 $T_2$（典型几毫秒）
表征的时间内达到彼此的热平衡；而核自旋要以由 $T_1$（可长达数小时）表征的时间才
弛豫回晶格温度。

还有一个巧妙的窍门：把核置于磁场中，核能级分裂并按玻尔兹曼概率布居。然后把磁场
突然反向（时间短于 $T_2$）：核自旋来不及相对能级移动，于是得到布居反转；而且
各能级之间恰具有对应负温度的正确玻尔兹曼分布。负温度态是一个有着独特磁相图的
独特热力学态。银的实验表明：低温负温度下银核在 $-1.9$ nK 发生铁磁有序。

\section{磁有序的测量}

前面几节描述了各种磁有序。但如何测量？本节讨论测量磁有序的各种技术。

\subsection{磁化强度与磁化率}

最直截了当的实验技术是常规磁化强度测量——本质上测样品的净磁矩 $\boldsymbol{\mu}$
（除以样品体积即得磁化强度）。

抽拉式磁强计（extraction magnetometer）：样品置于线圈中心，然后抽出到很远处，在线圈中
感应电压 V；样品产生的磁通 $\Phi$ 等于 $\int V\,\mathrm{d}t$。通常磁强计由两个
反绕线圈构成，使外场变化引起的电压相互抵消，只留下样品的信号。振动样品磁强计
（vibrating sample magnetometer，通称 VSM）让样品正弦上下振动，样品磁矩的运动在
固定的拾取线圈中感应电信号；信号正比于磁矩以及振动幅度与频率。交变梯度磁强计中，
样品贴在压电条上，反绕线圈提供交变场，使样品处出现交变的场梯度，样品受力
$(\boldsymbol{\mu}\cdot\nabla)\mathbf{B}$；压电条以驱动频率同频偏折，偏折量可由
压电信号测得，从而推知 $\boldsymbol{\mu}$。力矩磁强计中样品悬在扭丝上，加磁场 B
产生力矩 $\boldsymbol{\mu}\times\mathbf{B}$。最灵敏的技术之一用 SQUID（超导量子
干涉器件）：一个含“弱连接”（即约瑟夫森结）的超导环，像一个极灵敏的量子干涉仪；
让样品穿过环，感应的持续电流正比于样品磁化强度。这些技术当然只在样品磁化强度
非零时有用；反铁磁体的磁化强度在外场不大时为零。

另一策略是测磁化率——它很好地指示磁有序的类型，可拟合成居里--外斯定律。
磁化率有多种测法。它等于小外场下样品中诱导的磁化强度，故上述测磁化强度的技术
皆可用于测磁化率。此外，主要用于提取小磁化率的方法是古埃法（Gouy method）：
长圆柱样品悬于磁体极靴之间，一端在极靴间、另一端远在磁体之外零场区。通电后样品
的表观重量改变。相关的法拉第法可用更小的样品，但此时力由适当形状极靴产生的磁场
梯度提供。

\subsection{中子散射}

事实证明，中子在研究凝聚态磁性中极为有用；在研究磁结构与动力学方面，中子散射比其他
实验技术有显著优势：它给出样品中磁矩排列的最直接信息。

中子是自旋 $\frac{1}{2}$ 粒子（见表 2.3），具有非零磁矩。核反应堆燃料元件内的
裂变反应产生大量中子。反应堆出来的中子束的能谱由慢化剂温度 T 决定——通常为室温，
加热或冷却慢化剂可改变 T。因此反应堆中子具有速度分布
\begin{equation*}
n(v) \propto v^{3}\exp\left(-\frac{\frac{1}{2}m_{\mathrm{n}}v^{2}}{k_{\mathrm{B}}T}\right),\tag{5.33}
\end{equation*}
$n(v)\,\mathrm{d}v$ 是每秒通过单位面积的速度在 $v$ 到 $v+\mathrm{d}v$ 的中子数，
$m_{\mathrm{n}}$ 是中子质量。分布含玻尔兹曼因子与 $v^{3}$ 项——后者是相空间因子
$v^{2}$（同于分子运动论的麦克斯韦速度分布）与泻流过孔的因子 v 之积。此函数绘于
图 5.16(a)，其极大在
\begin{equation*}
v = \sqrt{\frac{3k_{\mathrm{B}}T}{m_{\mathrm{n}}}}\tag{5.34}
\end{equation*}
对应条件 $\frac{1}{2}m_{\mathrm{n}}v^{2}=\frac{3}{2}k_{\mathrm{B}}T$。速度为 v 的
中子的德布罗意波长 $\lambda$ 为
\begin{equation*}
\lambda = \frac{h}{m_{\mathrm{n}}v}.\tag{5.35}
\end{equation*}
分布函数按波长绘于图 5.16(b)。这表明热中子的波长与原子间距相近，从而可以进行
衍射测量。波长可调意味着衍射实验的尺度范围从直接探测小原子的波函数直到大分子的
低分辨研究。

\begin{figure}
\centering
\includegraphics[width=0.42\textwidth]{fig5_16.pdf}
\caption{（a）30 K 与 300 K 时中子分布函数随速度 v 的变化。（b）同样的分布函数按
德布罗意波长 $\lambda=h/m_{\mathrm{n}}v$ 绘出。曲线已归一化为等面积。}
\end{figure}

中子经强核力与核散射，也经电磁相互作用与晶体内磁场的变化散射。后者使中子能探测
材料磁性——中子与电子的磁矩耦合。

先考虑核散射。引起中子散射的核力程极短（$10^{-15}$--$10^{-14}$ m），远小于典型
热中子波长。因此入射中子波 $\psi_i=e^{i\mathbf{k}\cdot\mathbf{r}}$ 产生球对称的
弹性散射波 $\psi_f=-(b/r)e^{ikr}$。量 b 称为散射长度（scattering length）；表达式
中的负号保证排斥性核势对应 $b>0$。散射长度依赖具体的核以及核--中子系统的自旋态
（对自旋 I 的核，可为 $I+\frac{1}{2}$ 或 $I-\frac{1}{2}$）。

弹性中子散射实验中，弹性散射的中子在散射矢量等于倒格矢时产生强布拉格反射
（图 5.17(a, b)）。研究这些反射的实验方法包括：（i）转动晶体并测量单色束的散射，
寻找强反射；（ii）使用一系列入射波长（劳厄法，见图 5.17(c)）；（iii）测量单色
中子在粉末样品上的衍射（见图 5.17(d)）。反应堆源上一种可能的实验安排示意于
图 5.17(e)。中子也可由散裂源产生：同步加速器的高能质子轰击重金属靶（如
${}^{238}$U）产生中子；同步加速器质子束是脉冲的，中子束因此也是脉冲的。利用脉冲
中子的能量展宽可建造飞行时间谱仪：测量脉冲开始后散射中子随时间的变化，一次脉冲
即可提取一系列入射波矢的散射（见图 5.17(f)）。

中子的核散射来自与核（而非电子云）的相互作用，因此原子的散射能力与其原子序数
关系不大——这与 X 射线散射和电子散射截然相反。这带来若干优点：更容易在重原子
存在下感知轻原子（如氢）；周期表中相邻元素的散射截面一般差别显著、可以区分；
核散射的依赖性使同一元素的不同同位素可以有很不一样的中子散射长度，因此可以采用
同位素替换。中子是高度穿透的探针，能（同样是无损地）研究材料内部，而不像 X 射线
散射、电子显微术或光学方法那样仅研究表面层。

中子不带电——这正是电子云对它不可见的原因。但若原子有净磁矩，中子的磁矩可直接
与其耦合。这是中子重要的第二类散射机制：磁散射。结果是自旋向上与向下的电子磁矩
在中子“眼”里不同：样品变为磁性时，中子衍射图中可以出现新的峰。例子见图 5.18、
5.19——反铁磁体 MnO 的磁结构。$T_{\mathrm{N}}=116$ K 以上，中子衍射峰来自锰离子的
面心立方（fcc）晶格——所有锰离子等价（见图 5.19）。fcc 晶格使衍射图出现系统性
消光。$T_{\mathrm{N}}$ 以下，磁结构如图 5.18：每个 (111) 面内自旋全部平行，但相邻
两个 (111) 面的磁矩反平行。磁性单胞因而是化学单胞的两倍大，故 $T_{\mathrm{N}}$
以下中子衍射图中突然多出几个峰。磁布拉格峰的幅度可用来量度磁有序的强度，从而
可以跟踪磁有序随温度的变化。

\begin{figure}
\centering
\includegraphics[width=0.45\textwidth]{fig5_17.pdf}
\caption{（a）波矢 $\mathbf{k}$ 的入射中子被晶面以散射角 $2\theta$ 散射到
$\mathbf{k}'$；弹性散射意味着 $|\mathbf{k}|=|\mathbf{k}'|$；入射中子波长
$\lambda=2\pi/|\mathbf{k}|$。（b）当散射矢量 $\mathbf{Q}=\mathbf{k}-\mathbf{k}'$
等于倒格矢时，晶面散射的中子发生相长干涉。（c）劳厄法：入射方向固定而波长有一
范围的径束，当 $\mathbf{Q}$ 为倒格矢时给出布拉格反射。（d）粉末衍射法：单色束
入射粉末样品，散射中子位于半角 $2\theta$ 的圆锥（德拜--谢乐锥）上。（e）反应堆源
的中子衍射：单色器与旋转屏蔽选择入射中子波长，样品散射的中子由多探测器测量。
（f）散裂源：靶上散裂产生中子脉冲；探测器信号按时间记录，较快（波长较短）的
中子较早到达。}
\end{figure}

\begin{figure}
\centering
\includegraphics[width=0.42\textwidth]{fig5_18.pdf}
\caption{MnO 的磁结构。Mn 离子位于面心立方晶格上；$\mathrm{O^{2-}}$ 离子未画出
（见图 4.2）。Mn 离子按自旋态以黑白两色标记。}
\end{figure}

\begin{figure}
\centering
\includegraphics[width=0.5\textwidth]{fig5_19.pdf}
\caption{$T_{\mathrm{N}}$ 上、下 MnO 的中子衍射图。取自 C.~G.~Shull,
W.~A.~Strauser 与 E.~O.~Wollan, Phys.~Rev., \textbf{83}, 333 (1951)。}
\end{figure}

磁矩散射的幅度依赖于磁矩排列的方向，因此中子衍射可用于确定磁有序晶体中原子磁矩
的排列。

周期性体系中，散射矢量 Q 等于倒格矢 G 时出现尖锐布拉格反射。螺旋磁体中，螺旋
带来额外的周期性，可用螺旋波矢 q 描述：螺旋轴沿 q，螺距等于 $2\pi/|q|$。布拉格
反射于是也出现在 $G\pm q$ 处。由于螺旋螺距大于晶格间距，$|q|<|G|$：磁转变温度
以上观察到的布拉格峰在转变以下仍被观察到，且每个布拉格峰两侧伴随较小的卫星峰。

中子散射是确定样品磁态细节的最直接方法。但中子昂贵，且中子与物质相互作用较弱，
需要大样品（这一缺点对 6.6.4 节描述的非弹性实验最为严重）。

\subsection{其他技术}

内场可用第三章描述的技术测量。NMR、穆斯堡尔与 $\mu$SR 实验都能测铁磁体磁化强度
的温度依赖。它们是局域探针——测的不是平均磁化强度而是特定晶体学格位处的磁化
强度——因此可用于测量反铁磁体或亚铁磁体中各子晶格的磁化强度，还能给出样品整体
未磁化时磁畴内部自发磁化强度的信息。

另一种技术是 X 射线散射。磁性 X 射线散射通常是相当弱的效应（衍射实验中磁散射
名义上比非磁散射小五个量级），因此 X 射线以往主要用于结构而非磁性研究。这一局面
已因 X 射线共振交换散射（XRES）的发展而改变——当 X 射线能量扫过吸收边时磁散射
被增强。另一个重要因素是高亮度 X 射线同步辐射源的发展与优质偏振光子束的提供。
这些因素抵消了效应的内在弱点，使 X 射线得以用于磁性研究。X 射线的优点：可研究
小样品（束径可以很小）；元素分辨（把能量调到样品中特定原子的吸收边，测到的信号
就是该元素的标志）；分辨率可优于中子；还能分别测量磁化密度中自旋与轨道的贡献。

二色性（dichroism）是光的吸收对偏振的依赖（由取向长链分子构成的偏振片是光频线
二色性的著名例子）。磁二色性以类似方式工作，其效应由未成对电子造成的原子周围
磁化分布的不对称驱动。这本质上是轨道现象——光子把螺旋度 $\pm1$ 传递给吸收体——
但自旋--轨道相互作用可使二色性依赖自旋角动量。磁 X 射线圆二色（MXCD）技术测量
右旋与左旋圆偏振 X 射线衰减之差，结果可分别推知轨道与自旋磁矩。

\section*{延伸阅读}

\begin{itemize}[itemsep=0.2em]
\item 核磁有序实验的综述见 A.~S.~Oja 与 O.~V.~Lounsmaa, Reviews of Modern Physics
\textbf{67}, 1 (1997)。
\item 中子散射的描述见 M.~Dove,《Structure and dynamics》, OUP（即将出版）；
G.~L.~Squires,《Introduction to the theory of thermal neutron scattering》,
CUP 1978；W.~Marshall 与 S.~W.~Lovesey,《Theory of thermal neutron scattering》,
OUP 1971。
\item 更进阶的处理见 S.~W.~Lovesey,《Theory of neutron scattering from condensed
matter》, OUP 1984。
\item S.~W.~Lovesey 与 S.~P.~Collins,《X-ray scattering and absorption by magnetic
materials》, OUP 1996，提供了磁 X 射线散射的有用信息。
\end{itemize}

\section*{习题}

\exer{5.1} 估算铁中分子场 $B_{\mathrm{mf}}$ 的大小（以特斯拉计），并与磁化强度对
B 场的贡献 $\mu_0M$ 比较。进而说明为什么解释铁的铁磁性必须借助交换的概念。Fe 的
密度 7873 kg\,m$^{-3}$，相对原子质量 55.847，居里温度 $T_{\mathrm{C}}=1043$ K，
每个原子携带约 2.2 $\mu_{\mathrm{B}}$ 的磁矩。

\exer{5.2} 把外斯模型推广到自旋 $S>\frac{1}{2}$。证明恰低于 $T_{\mathrm{C}}$ 处
磁化强度为
\begin{equation*}
\frac{M}{M_{\mathrm{s}}} = \left(\frac{10(S+1)^{2}(T_{\mathrm{C}}-T)}{3[(S+1)^{2}+S^{2}]T_{\mathrm{C}}}\right)^{1/2}\tag{5.36}
\end{equation*}
且 $T_{\mathrm{C}}$ 处热容有大小为
\begin{equation*}
\Delta C = \frac{5}{2}nk_{\mathrm{B}}\frac{(2S+1)^{2}-1}{(2S+1)^{2}+1}.\tag{5.37}
\end{equation*}
的间断。你需要关系式
\begin{equation*}
\mathrm{B}_S(y) = \left[\frac{(2S+1)^{2}-1}{(2S)^{2}}\right]\frac{y}{3} - \left[\frac{(2S+1)^{4}-1}{(2S)^{4}}\right]\frac{y^{3}}{45} + O(y^{5}),\qquad y\ll1.\tag{5.38}
\end{equation*}

\exer{5.3} 反铁磁体每套子晶格上的分子场为
\begin{equation*}
\begin{aligned}
B_{+} &= -\Gamma M_{+} - |\lambda|M_{-}\\
B_{-} &= -\Gamma M_{-} - |\lambda|M_{+},
\end{aligned}\tag{5.39}
\end{equation*}
$\Gamma$ 是表示同一子晶格对分子场贡献的常数。推广 5.2 节的处理，证明每套子晶格上
磁化强度为
\begin{equation*}
M_{\pm} = M_{\mathrm{s}}\mathrm{B}_J\left(-\frac{g_J\mu_{\mathrm{B}}J(\Gamma M_{\pm}+|\lambda|M_{\mp})}{k_{\mathrm{B}}T}\right).\tag{5.40}
\end{equation*}
进而证明奈尔温度（每套子晶格自发磁化强度消失的温度）为
\begin{equation*}
T_{\mathrm{N}} = \frac{n(|\lambda|-\Gamma)\mu_{\mathrm{eff}}^{2}}{3k_{\mathrm{B}}}.\tag{5.41}
\end{equation*}
再证明磁化率为 $\chi\propto(T-\theta)^{-1}$，其中
\begin{equation*}
\theta = -\frac{n(|\lambda|+\Gamma)\mu_{\mathrm{eff}}^{2}}{3k_{\mathrm{B}}},\tag{5.42}
\end{equation*}
故仅当 $\Gamma=0$ 时 $|\theta|=T_{\mathrm{N}}$。

\exer{5.4} 对 5.4 节描述的螺旋磁体，能量由式 5.30 给出。证明铁磁、反铁磁与螺旋磁
情况的能量分别为 $E_{\mathrm{FM}}$、$E_{\mathrm{AFM}}$ 与 $E_{\mathrm{HM}}$：
\begin{equation*}
E_{\mathrm{FM}} = -2NS^{2}(\mathsf{J}_1+\mathsf{J}_2)\tag{5.43}
\end{equation*}
\begin{equation*}
E_{\mathrm{AFM}} = -2NS^{2}(-\mathsf{J}_1+\mathsf{J}_2)\tag{5.44}
\end{equation*}
\begin{equation*}
E_{\mathrm{HM}} = -2NS^{2}\left(\frac{-\mathsf{J}_1^{2}}{8\mathsf{J}_2} - \mathsf{J}_2\right)\tag{5.45}
\end{equation*}
故当 $\mathsf{J}_2<0$ 且 $|\mathsf{J}_1|<4|\mathsf{J}_2|$ 时螺旋磁性优于铁磁或
反铁磁。

\exer{5.5} 证明 $T\ll T_{\mathrm{C}}$ 时铁磁体的外斯模型给出磁化强度
\begin{equation*}
M = M_{\mathrm{s}}(1-2\,\mathrm{e}^{-2T_{\mathrm{C}}/T}).\tag{5.46}
\end{equation*}
注意这与实验确定的幂律行为不符——实验给出
\begin{equation*}
M = M_{\mathrm{s}}(1-AT^{3/2})\tag{5.47}
\end{equation*}
A 为常数。这一低温行为可以由自旋波效应解释（见 6.6 节）。

\exer{5.6} 亚铁磁体的磁化率 $\chi$ 可用平均场理论处理。考虑具有两套不等价子晶格的
亚铁磁体，子晶格 1、2 上的分子场 $B_1$、$B_2$ 为
\begin{equation*}
B_1 = \mu_0H - \lambda M_1\tag{5.48}
\end{equation*}
\begin{equation*}
B_2 = \mu_0H - \lambda M_2,\tag{5.49}
\end{equation*}
$M_1$、$M_2$ 是各子晶格的磁化强度，H 是外场。但两套子晶格有不同的居里常数
$C_1$、$C_2$（即不同的 $g_JJ$ 值）。证明磁化率为
\begin{equation*}
\chi = \frac{1}{T^{2}-\theta^{2}}\left[(C_1+C_2)T - \frac{2\lambda C_1C_2}{\mu_0}\right],\tag{5.50}
\end{equation*}
并求转变温度 $\theta$ 的值。证明当 $C_1=C_2=C$ 时磁化率退化为 $\theta=\lambda C/\mu_0$
的反铁磁体情形。

\exer{5.7} 证明海森伯模型的哈密顿量
\begin{equation*}
\hat{\mathcal{H}} = -\sum_{\langle ij\rangle}\mathsf{J}\,\mathbf{S}_i\cdot\mathbf{S}_j,\tag{5.51}
\end{equation*}
可改写为
\begin{equation*}
\hat{\mathcal{H}} = -\sum_{\langle ij\rangle}\mathsf{J}\left[S_i^{z}S_j^{z} + \frac{1}{2}(S_i^{+}S_j^{-}+S_i^{-}S_j^{+})\right].\tag{5.52}
\end{equation*}
进而证明：对海森伯铁磁体（$\mathsf{J}>0$），所有格点自旋都朝上（比如说）的态
$|\Phi\rangle$ 是哈密顿量的本征态，能量为
\begin{equation*}
E_0 = -NS^{2}\mathsf{J}.\tag{5.53}
\end{equation*}
考虑海森伯反铁磁体（$\mathsf{J}<0$），自旋位于两套子晶格上，每个自旋只与
另一套子晶格上的自旋相互作用。证明“想当然”的基态——每套子晶格内铁磁排列、两套
子晶格磁化方向相反——并非哈密顿量的本征态。这凸显了海森伯反铁磁体是一个复杂
而困难的问题。

\exer{5.8} $\mathrm{MnF_2}$ 具有四角晶体结构：Mn 离子位于四角单胞（a=b=0.5 nm，
c=0.3 nm）的顶角及体心位置。70 K 以下 Mn 离子自旋沿 c 轴反铁磁排列，体心离子的
自旋与顶角离子的相反。用入射波长 0.3 nm、散射角 0°--90° 测量 $\mathrm{MnF_2}$
粉末样品的中子散射。画出（a）100 K 与（b）10 K 时预期的结果。可忽略 F 离子的
散射。

\exer{5.9} 外斯模型可以找到精确解。无外场时问题归结为联立求解
\begin{equation*}
\tilde{m} = \frac{M}{M_{\mathrm{s}}} = \mathrm{B}_J(y)\tag{5.54}
\end{equation*}
与
\begin{equation*}
y = \frac{g_J\mu_{\mathrm{B}}J\lambda M}{k_{\mathrm{B}}T}\tag{5.55}
\end{equation*}
$\tilde{m}$ 是约化磁化强度，$M_{\mathrm{s}}=ng_J\mu_{\mathrm{B}}J$。定义 $x=e^{y}$
有助益。

（a）对 $J=\frac{1}{2}$、$g_J=2$，$\mathrm{B}_J(y)=\tanh y$。证明这意味着
\begin{equation*}
\tilde{m} = \frac{x-\frac{1}{2}}{x+\frac{1}{2}}\tag{5.56}
\end{equation*}
从而
\begin{equation*}
\log x = \frac{1}{2}\left[\log(1+\tilde{m})-\log(1-\tilde{m})\right].\tag{5.57}
\end{equation*}
写 $T_{\mathrm{C}}=n\mu_{\mathrm{B}}^{2}\lambda/k_{\mathrm{B}}$，证明
\begin{equation*}
\frac{T}{T_{\mathrm{C}}} = \frac{2\tilde{m}}{\log(1+\tilde{m})-\log(1-\tilde{m})}.\tag{5.58}
\end{equation*}

（b）对 $J=1$，证明
\begin{equation*}
\mathrm{B}_1(y) = \frac{2\sinh y}{2\coth y+1}\tag{5.59}
\end{equation*}
从而
\begin{equation*}
\tilde{m} = \frac{x-\frac{1}{x}}{x+1+\frac{1}{x}}\tag{5.60}
\end{equation*}
其中
\begin{equation*}
x = \exp(2\lambda M\mu_{\mathrm{B}}/k_{\mathrm{B}}T).\tag{5.61}
\end{equation*}
证明这意味着
\begin{equation*}
x = \frac{\tilde{m}\pm\sqrt{4-3\tilde{m}^{2}}}{2(1-\tilde{m})},\tag{5.62}
\end{equation*}
进而（仔细想想 $\pm$ 号）意味着
\begin{equation*}
\frac{T}{T_{\mathrm{C}}} = \frac{\frac{3}{2}\tilde{m}}{\log(\tilde{m}+\sqrt{4-3\tilde{m}^{2}})-\log(2(1-\tilde{m}))},\tag{5.63}
\end{equation*}
其中 $T_{\mathrm{C}}=8n\mu_{\mathrm{B}}^{2}\lambda/3k_{\mathrm{B}}$。本题所用方法
详述于 M.~A.~B.~Whittaker, American Journal of Physics \textbf{57}, 45 (1989)。

% 第6章 有序与对称性破缺（Order and broken symmetry）
\chapter{有序与对称性破缺}

低温下自发有序的出现是凝聚态物理的基本现象。铁磁体、反铁磁体、液晶与超导体都是
有序相，固态本身也是如此。所有这些现象共享一些基本性质与特征。例如，它们都以
这样的温度依赖为标志：某个相关物理性质在临界温度 $T_{\mathrm{c}}$ 上下表现出显著
差异。对每种相都可以定义一个序参量（order parameter）：$T>T_{\mathrm{c}}$ 时为零，
$T<T_{\mathrm{c}}$ 时非零，因此它充当系统是否有序的指示器。对铁磁性，序参量就是
磁化强度。本章将说明每种有序相都伴随着一个对称性破缺（broken symmetry）——这一
概念将在下一节解释。

\section{对称性破缺}

铁磁体中选定了一个唯一的方向，所有原子磁矩都沿它排列：它们全都选择了“向上”而非
“向下”。这其实是个相当惊人的效应，因为底层的物理方程并不区分“上”与“下”——微观
物理具有实验观察到的基态所不具备的对称性。要深入这一谜题，必须更细致地考虑
对称性。

一个有用的例子是欧拉压杆（Euler strut，见图 6.1）：一根竖直杆，底端固定于地面，
顶端均匀加载。载荷加重时杆被压缩；超过临界重量它就会弯曲。向左弯还是向右弯取决
于载荷放置得如何对称等细节；但即使在“载荷绝对居中”的理想化情形下，杆仍然必须
选择弯曲的方向。它危如累卵、如立刀锋，一个随机的热涨落也许就足以使它向某一边
弯曲。轻载杆所固有的左右对称性在弯曲时被打破——这称为对称性破缺。

\begin{figure}
\centering
\includegraphics[width=0.6\textwidth]{fig6_01.pdf}
\caption{（a）欧拉压杆。（b）压力增大时压杆被压缩但保持左右对称。（c）力超过临界值
后压杆弯曲，打破左右对称。（d）压杆也可以向相反方向打破对称。}
\end{figure}

凝聚态体系中观察到许多类似特征。驱动对称性破缺转变的参数可以是力（如外加压强），
但更常是温度。图 6.2 画出液体与固体中的原子。液体冷却时系统只有轻微收缩，仍保持
极高的对称度；但低于临界温度（熔点）后液体变为固体，对称性被打破。乍看令人惊讶——
固体画出来比液体“显得”更对称：固体中原子整齐排列而液体里原子乱作一团。关键的观察
是：液体中任何一点与其他任何点完全相同——把系统对时间平均，每个位置被原子造访的
频率一样；不存在原子排齐的独特方向或轴。总之，液体具有完全的平移与旋转对称性。
而在固体中这种高度对称几乎全部丧失：图 6.2 所画的固体仍有残余对称——它不再在任意
旋转下不变，但在四重旋转（$\pi/2$、$\pi$、$3\pi/2$、$2\pi$）下不变；不再在任意
平移下不变，但在晶格基矢的整数组合平移下不变。所以并非所有对称都丧失了，而是
液态的高对称被打破了。

\mnote{关于对称性的术语（Terminology concerning symmetry）：若描述系统的哈密顿量
在对称群诸元素所对应的变换下不变，则系统具有该对称性。群是元素的集合加上组合它
们的运算，遵循一组特定规则：群是"封闭"的，即元素的组合仍是群的成员；群组合满足
结合律；逆元与单位元均有良好定义。离散对称性指元素可数的对称群，如立方体的旋转
对称群；连续对称性则有不可数的连续多个元素，如球的旋转对称群。若把群的对称操作
施加于整个系统而系统不变，则系统具有整体对称性（global symmetry）；若对各空间
点施加不同的对称操作后哈密顿量仍不变，则为局域对称性（local symmetry）。超导体
的规范对称性就是这样一种局域对称性。}

\begin{figure}
\centering
\includegraphics[width=0.4\textwidth]{fig6_02.pdf}
\caption{液--固相变。左：高温态（统计平均）具有完全的平移与旋转对称性。右：系统在
临界温度 $T_{\mathrm{c}}$ 以下变为固体，这些对称性被打破。}
\end{figure}

铁磁体情形类似（图 6.3）。居里温度 $T_{\mathrm{C}}$ 以上的铁磁体具有完全的旋转
对称性：所有方向等价，每个磁矩可以指向任意方向。$T_{\mathrm{C}}$ 以下系统为所有
自旋“选定”了唯一方向——高温态较高的旋转对称被打破；较低温态的旋转对称降低了
（只允许绕磁化轴——图 6.3 中的向上方向——的旋转）。

\begin{figure}
\centering
\includegraphics[width=0.4\textwidth]{fig6_03.pdf}
\caption{顺磁体--铁磁体相变。左：高温态（统计平均）具有旋转对称性。右：系统在居里
温度 $T_{\mathrm{C}}$ 以下变为铁磁体，该对称性被打破。}
\end{figure}

值得注意的重要一点：对称性不可能逐渐改变——某个特定的对称性要么存在要么不存在。
因此相变是尖锐的，有序态与无序态之间有清晰的分界。低温下有序的出现可以由相当
一般的热力学考虑理解：自由能 $F=E-TS$ 联系着能量与熵。为使自由能极小，低温下
系统选择能量最低的基态——通常是有序的，即最小化 E；温度升高时，最大化 S 的态
对最小化 F 更重要，于是更倾向于无序态。

并非所有相变都涉及对称性变化。图 6.4 是水的相图：液--气两区之间的边界线终止于
一个临界点。因此可以“作弊”绕过尖锐相变——沿一条避免不连续变化的路径穿过相图
（图 6.4 中的路径 C）。临界温度（647 K）以上，气体与液体仅以密度区分；气--液
转变不涉及对称性变化。相反，固--液转变涉及对称性变化，因此熔化曲线没有临界点。

\begin{figure}
\centering
\includegraphics[width=0.5\textwidth]{fig6_04.pdf}
\caption{水的相图。路径 A 与 B 越过相边界，C 则否。路径 A 涉及对称性的变化，
B 与 C 则不。}
\end{figure}

对称性破缺转变性质中令人困惑的一点是：破缺对称的基态并不具有哈密顿量的对称性。
把铁磁体冷却通过转变温度时，系统必须选定一个特定方向让所有自旋指向——尽管底层
物理并未挑出任何特定方向。这怎么会发生？

理解这一点，考虑一个非磁的例子很有用：氨分子（$\mathrm{NH_3}$，见图 6.5）。三个
氢原子构成等边三角形的三个顶点（图 6.5(a)、(b) 为侧视），氮原子可以在三角平面
略上方或略下方。这种金字塔形状源于氮上的孤对电子与三条氮--氢键争抢空间；它给
分子一个偶极矩，同时也在空间中定义了特定方向、打破了底层哈密顿量的对称性——
因此这个金字塔不可能是系统的稳定态。

事实上可以设想在图 6.5(a) 与 (b) 两个金字塔态之间变换：让氮原子穿过平面中央，
把四面体翻转——就像伞被大风吹翻那样。若以 x 表示氮原子到平面的距离，可以设想
一条势能曲线 $V(x)$，在两个稳定组态处各有一个极小（图 6.5(c)）。该势的两个最低
能量本征态示于图 6.5(d) 与 (e)：分别为第一激发态 $\phi_1(x)$ 与基态 $\phi_0(x)$。

因此图 6.5(a)、(b) 所代表的组态不是系统的稳定组态，而是能量本征态的叠加——它们
是亚稳态。图 6.5(a) 的态由 $(\phi_0(x)-\phi_1(x))/\sqrt{2}$ 表示，图 6.5(b) 的态
由 $(\phi_0(x)+\phi_1(x))/\sqrt{2}$ 表示。若把系统制备在图 6.5(a) 的亚稳态，会
发现它在它与图 6.5(b) 的亚稳态之间振荡，频率由 $\phi_0(x)$ 与 $\phi_1(x)$ 的能量
差给出：氨中为 24 GHz，氨微波激射器（maser）用的正是这一跃迁。若把氮换成更重的
磷制成 $\mathrm{PH_3}$，频率下降约十倍。系统越重，图 6.5(a) 亚稳态的时间常数越长。

铁磁体中我们有重得多的“态”：不只是更重的原子，而是更多的原子——也许 $10^{23}$
个。所有自旋指向单一方向的系统基态其实并不是系统的真正定态，但它是亚稳态，稳定
时间超过宇宙年龄。要到达另一个态（所有自旋反转），必须同时翻转系统中每一个自旋
——极不可能。

\begin{figure}
\centering
\includegraphics[width=0.6\textwidth]{fig6_05.pdf}
\caption{（a）氨 $\mathrm{NH_3}$ 是四面体分子，可存在于两种形式之一，另一种示于
（b）。氮原子有两个稳定位置，因此可以定义该系统的势能 $V(x)$。（c）势能曲线。
（d）第一激发态波函数 $\phi_1(x)$。（e）基态波函数 $\phi_0(x)$。}
\end{figure}

\section{模型}

为了讨论对称性破缺的一些后果，方便的做法是考虑一些简单的磁性模型并尝试求解。

\subsection{铁磁性的朗道理论}

俄国物理学家朗道（Lev Landau）给出了一个简便的、直接产生相变的模型，它源于
非常一般的考虑。把磁化强度为 M 的铁磁体的自由能写成 M 的幂级数。由于“上”与
“下”没有能量差别，幂级数不能包含 M 的任何奇次幂。于是自由能 $F(M)$ 可写为
\bio{Lev D. Landau\\（1908--1968）}
\begin{equation*}
F(M) = F_0 + a(T)M^{2} + bM^{4}\tag{6.1}
\end{equation*}
$F_0$ 与 b 是常数（设 $b>0$），$a(T)$ 依赖温度。若让 $a(T)$ 在转变温度 $T_{\mathrm{C}}$
处变号，可以证明该系统给出恰当的相变。在转变附近感兴趣的区域写
$a(T)=a_0(T-T_{\mathrm{C}})$（$a_0$ 为正常数）。求系统基态需使自由能极小，即解
$\partial F/\partial M=0$：
\begin{equation*}
2M\left[a_0(T-T_{\mathrm{C}}) + 2bM^{2}\right] = 0.\tag{6.2}
\end{equation*}
左边是两项之积，任一项可为零：
\begin{equation*}
M = 0 \quad\text{或}\quad M = \pm\left[\frac{a_0(T_{\mathrm{C}}-T)}{2b}\right]^{1/2}.\tag{6.3}
\end{equation*}
第二个条件只在 $T<T_{\mathrm{C}}$ 时有效，否则等于要给负数开平方。第一个条件在
$T_{\mathrm{C}}$ 上下都适用，但在 $T_{\mathrm{C}}$ 以下它只给出一个不稳定平衡点
（可由 $\partial^{2}F/\partial M^{2}$ 判断）。于是磁化强度遵循图 6.6(b) 的曲线：
$T\geq T_{\mathrm{C}}$ 时为零，$T<T_{\mathrm{C}}$ 时非零且正比于 $(T_{\mathrm{C}}-T)^{1/2}$。

\begin{figure}
\centering
\includegraphics[width=0.4\textwidth]{fig6_06.pdf}
\caption{（a）铁磁体的自由能 $F(M)$。（b）相应的磁化强度随温度的函数。}
\end{figure}

朗道研究相变的方法称为平均场理论（mean-field theory）：它假定所有自旋“感受”到
邻居产生的同一个平均交换场，该场正比于磁化强度。这与 5.1.1 节的外斯模型完全
相同。平均场理论是能构造出来描述多种相变的最简单理论，在各种情形给出相似结果；
在不同情形有不同名字，如合金有序--无序转变的 Bragg--Williams 理论。平均场理论
不能准确解释临界区，因为“样品各处相同”的假设在临界区恰恰最不成立。

平均场理论忽略关联与涨落，而它们在 $T_{\mathrm{C}}$ 附近极为重要。非常接近临界
温度时，序参量出现大涨落；临界区的特征是所有尺度上都有涨落。烧水时，水安静而不
起眼地变热，直到接近沸点才喧哗起来、冒泡翻腾。\fn{1}{这一现象的直观演示是临界乳光（critical opalescence）：透过临界点附近的气体看物体，影像模糊混浊。原因是临界点附近密度涨落强烈，引起折射率的大幅变化。}铁磁体在 $T_{\mathrm{C}}$ 附近发生同类效应。
刻画涨落的主导长度尺度是关联长度（correlation length）$\xi$：$T\to T_{\mathrm{C}}$
时 $\xi\to\infty$，即临界点处关联长度变为无穷。涨落的重要性也可从 $F(M)$ 的形式
看出：M=0 处的曲率 $(\partial^{2}F/\partial M^{2})_{M=0}$ 在 $T=T_{\mathrm{C}}$
趋于零，M=0 的极小变得宽广平坦（见图 6.6）。因此，尽管平均场理论易于求解，其
关于临界区的预言必须谨慎对待。
\bio{Ernst Ising\\（1900--1998）}
\mnote{第一章中自旋 $z$ 分量的算符写作 $\hat{S}_z$。现在还需要标记自旋所在的格点，
故把 z 写成上标、格点标号写成下标：$\hat{S}_i^z$ 即格点 $i$ 上自旋 $z$ 分量的
算符。}

\subsection{海森伯模型与伊辛模型}

理解固体磁行为的另一条路是考虑磁相互作用的特定微观模型。一个常被研究的模型是
最近邻海森伯模型：
\begin{equation*}
\hat{\mathcal{H}} = -\sum_{\langle ij\rangle}\mathsf{J}\,\mathbf{S}_i\cdot\mathbf{S}_j,\tag{6.4}
\end{equation*}
常数 $\mathsf{J}$ 是交换积分，求和号下的 $\langle ij\rangle$ 表示对最近邻求和。
自旋 $\mathbf{S}_i$ 被当作三维矢量——允许指向三维空间中任意方向。但求和可以在 1、
2 或 3 维晶格上进行。这里要区分自旋所在晶格的维数 d 与自旋自身的维数 D（一般称
D 为序参量的维数）。海森伯模型 D=3（自旋是三维矢量）；但这些自旋可以排在 1、2、
3 维（想排 4 维也行！）的晶格上，故 d 可为 1、2、3……

相关模型是伊辛模型（Ising model）：自旋只允许朝上或朝下，即只考虑自旋的 z 分量。
其哈密顿量为
\begin{equation*}
\hat{\mathcal{H}} = -\sum_{\langle ij\rangle}\mathsf{J}\,S_i^{z}S_j^{z}.\tag{6.5}
\end{equation*}
此处序参量维数 D=1（自旋只允许沿 $\pm z$）。但同样可以把这些一维自旋排在
$d=1,2,\ldots$ 的晶格上。

\subsection{一维伊辛模型（$D=1, d=1$）}

把伊辛自旋放在一维晶格上，我们将证明：没有相变。先考虑含 $N+1$ 个自旋的链
（须考虑 N 个相邻“键”）。哈密顿量为
\begin{equation*}
\hat{\mathcal{H}} = -2\mathsf{J}\sum_{i=1}^{N}S_i^{z}S_{i+1}^{z}\tag{6.6}
\end{equation*}
设 $\mathsf{J}>0$，基态为所有相邻自旋铁磁排列（图 6.7(a)）。
基态能量为 $-N\mathsf{J}/2$，因为所有 $S_i^{z}=+\frac{1}{2}$。现在加入一个“错误”
——单个缺陷（图 6.7(b)）。
这多付出能量 $E=\mathsf{J}$：把一个有利的相互作用（省 $\mathsf{J}/2$）变成不利的
（费 $\mathsf{J}/2$），能量变化为 $\mathsf{J}$。但熵增 $S=k_{\mathrm{B}}\ln N$——
缺陷可以放在 N 个位置中的任何一个。让链变得很长（$N\to\infty$）时，缺陷的能量
代价不变（$\mathsf{J}$），熵增益却变为无穷。性质由自由能 $F=E-TS$ 决定：只要
温度非零，熵的考虑就意味着缺陷的存在使自由能骤降至 $-\infty$——缺陷会自发形成，
实际上 $T>0$ 时不存在长程序。换句话说，临界温度为零。这一论证对一维晶格上的
所有模型都成立（一维中熵总是获胜），结论：一维中不可能有长程序。

\begin{figure}
\centering
\includegraphics[width=0.6\textwidth]{fig6_07.pdf}
\caption{一维伊辛模型。（a）基态为 $N+1$ 个自旋全部铁磁排列。（b）加入单个缺陷。}
\end{figure}

\subsection{二维伊辛模型（$D=1, d=2$）}

把伊辛自旋放在二维晶格上，则会出现向磁有序态的相变，其临界温度非零。因为制造
缺陷的能量代价与熵增益都按缺陷的周长标度（见图 6.8），能量与熵可以势均力敌，
谁也没有压倒性优势。该模型的严格求解是二十世纪统计物理的杰出成就之一
（Lars Onsager 于 1944 年解决），其解超出本书范围。这说明：即便表述简单的问题
也绝非容易求解。
\bio{Lars Onsager\\（1903--1976）}

\begin{figure}
\centering
\includegraphics[width=0.5\textwidth]{fig6_08.pdf}
\caption{二维伊辛模型。自旋向上标记为 +、向下标记为 $-$。向下自旋缺陷的能量与熵
都按缺陷边界的周长标度。}
\end{figure}

\section{对称性破缺的后果}

打破对称性有种种后果：
\begin{itemize}[itemsep=0.2em]
\item \textbf{相变}：系统在温度 $T_{\mathrm{c}}$ 处发生急剧的行为变化，称为改变了
相（如液体$\to$固体、顺磁体$\to$铁磁体等）。相变附近的区域称为临界区（见 6.4 节）。
\item \textbf{刚性}：打破对称后，系统对维持该破缺对称态有强烈的能量偏好；试图
改变系统打破对称的方式会遭到抵抗。因此晶体不易弯曲、铁磁体表现永磁体行为
（见 6.5 节）。
\item \textbf{激发}：$T=0$ 时系统完全有序；有限温度下，序参量中的激发削弱有序。
晶体中这些激发称为格波，量子化为声子；铁磁体中对应的模式称为自旋波，量子化为
磁振子（见 6.6 节）。
\item \textbf{缺陷}：若宏观样品相邻两部分以不同方式打破对称，其边界将含有缺陷：
如晶体中的位错或铁磁体中的畴壁（见 6.7 节）。
\end{itemize}

表 6.1 总结了几种破缺对称相的性质。每种相都有高温无序态与低温有序态、一个序参量、
一种刚性现象、一组代表相对有序态的波动型偏离的激发，以及与“不同空间区域以不同
方式打破对称”相联系的缺陷。

\begin{table}
\centering
\caption{破缺对称相的性质。$\rho_G$ 是与倒格矢 G 对应的空间频率的电荷密度傅里叶
分量；超导体中的复波函数为 $\psi=|\psi|e^{i\phi}$；电极化强度 P 是单位体积的
电偶极矩。}
\begin{small}
\begin{tabular}{lllllll}
\toprule
现象 & 高温相 & 低温相 & 序参量 & 激发 & 刚性现象 & 缺陷\\
\midrule
晶体 & 液体 & 固体 & $\rho_G$ & 声子 & 刚性 & 位错、晶界\\
铁磁体 & 顺磁体 & 铁磁体 & M & 磁振子 & 永磁性 & 畴壁\\
反铁磁体 & 顺磁体 & 反铁磁体 & M（子晶格上） & 磁振子 & （颇为微妙） & 畴壁\\
向列相液晶 & 液体 & 取向液体 & $S=\langle\frac{1}{2}(3\cos^{2}\theta-1)\rangle$ & 指向涨落 & 各种 & 向错、点缺陷\\
铁电体 & 无极性晶体 & 极性晶体 & P & 软模 & 铁电滞回 & 畴壁\\
超导体 & 正常金属 & 超导体 & $|\psi|e^{i\phi}$ & -- & 超导电性 & 磁通线\\
\bottomrule
\end{tabular}
\end{small}
\end{table}

晶体在衍射实验中给出布拉格峰。因此可用 $\rho_G$——与倒格矢 G 对应的空间频率的
电荷密度傅里叶分量——作序参量。反铁磁体与铁磁体类似，只是序参量是反铁磁体一套
子晶格上的磁化强度。超导体有复序参量 $\psi=|\psi|e^{i\phi}$：$|\psi|^{2}$ 代表
超导电子密度（对超流或玻色凝聚体，代表凝聚分数）。波函数的相位 $\phi$ 是被打破
的对称性：正常态中它在任何点可取任意值，超导态中它在整个样品上变得均匀。对向列相
液晶（见图 6.9），序参量为 $S=\langle\frac{1}{2}(3\cos^{2}\theta-1)\rangle$——
函数 $\frac{1}{2}(3\cos^{2}\theta-1)$ 的平均值，$\theta$ 是分子长轴与指向矢
（director）$\hat{\mathbf{n}}$（沿液晶分子平均取向的单位矢量）的夹角。各向同性
（液体）态对应 S=0；完全排齐对应 S=1。

\begin{figure}
\centering
\includegraphics[width=0.5\textwidth]{fig6_09.pdf}
\caption{向列相液晶含棒状分子。（a）低温下它们如固体般有序。（c）高温下处于液态。
（b）中间温度区为向列态：分子长轴排齐，尽管分子质心的位置并无序。}
\end{figure}

\begin{table}
\centering
\caption{各种模型的临界指数。}
\begin{small}
\begin{tabular}{lcccc}
\toprule
模型 & 平均场 & Ising & Ising & Heisenberg\\
\midrule
$D$ & 任意 & 1 & 1 & 3\\
$d$ & 任意 & 2 & 3 & 3\\
$\beta$ & $\frac{1}{2}$ & $\frac{1}{8}$ & 0.326 & 0.367\\
$\gamma$ & 1 & $\frac{7}{4}$ & 1.2378(6) & 1.388(3)\\
$\delta$ & 3 & 15 & 4.78 & 4.78\\
\bottomrule
\end{tabular}
\end{small}
\end{table}

\section{相变}

6.2 节给出了铁磁性的朗道理论——一种平均场理论：所有自旋感受同样的交换场，导出
转变以下磁化强度按 $(T_{\mathrm{C}}-T)^{1/2}$ 变化。真实体系中，磁化强度在转变
附近确实按 $(T_{\mathrm{C}}-T)^{\beta}$ 变化，但指数 $\beta$ 未必等于 $\frac{1}{2}$。
因此指数给出关于相变性质的重要信息。还可以定义若干类似指数，即临界指数（critical
exponents）。实验发现相变温度 $T_{\mathrm{C}}$ 附近
\begin{equation*}
\begin{array}{cccl}
\chi & \propto & (T-T_{\mathrm{C}})^{-\gamma} & \qquad T>T_{\mathrm{C}}\\
M & \propto & (T_{\mathrm{C}}-T)^{\beta} & \qquad T<T_{\mathrm{C}}\\
M & \propto & H^{1/\delta} & \qquad T=T_{\mathrm{C}}
\end{array}\tag{6.7}
\end{equation*}
$\beta$、$\gamma$、$\delta$ 是临界指数；各模型取值见表 6.2。

如前所述，平均场理论忽略关联与涨落，这意味着它无望给出 $T_{\mathrm{C}}$ 附近——
恰恰是想计算临界指数之处——的正确描述。然而事实证明：平均场理论是四维及以上
（$d\geq4$）体系的正确描述，并给出正确的临界指数。

尽管平均场理论无法完满解释四维以下体系的临界行为，但只要接受普适性（universality）
假说，只需考虑一小组代表性的理想统计模型即可计算任何物理体系的临界指数。该假说
基于这样的观察：临界指数看起来惊人地不依赖相变的类型——无论液--气、铁磁--顺磁、
超导--正常或任何其他。按该假说，对连续相变，临界指数只取决于：
（1）系统维数 d；（2）序参量维数 D（准确地说是序参量的对称性）；（3）力是短程
还是长程。
因此只需考察特定的普适类（universality classes）——即短程力与长程力下特定的
D、d 组合——在每类中挑最简单的模型计算指数即可。（上述假设适用于静态临界指数，
而非表征时间依赖性质的动态临界指数。）

现在列出有解已知的模型。精确可解的模型包括：
（1）$d=1$ 的所有情形——如前所述，不表现连续相变；
（2）$d\geq4$ 的所有情形——给出平均场解；
（3）长程相互作用的全部情形——给出平均场解；
（4）$d=2$、$D=1$——即二维伊辛模型（见表 6.2）；
（5）任意 $d$ 的 $D=\infty$ 情形——称为球模型（spherical model）。
遗憾的是，多数真实情形对应 $d=3$ 与短程相互作用——尚未精确求解。

1970 年代初发现了一种计算临界行为的方法，即便对不能精确求解的模型也适用：考虑
体积为 $L^{d}$ 的自旋块（$L$ 远小于关联长度 $\xi$）。若把该块的线性尺寸放大
（标度）$n$ 倍，哈密顿量将相应地被标度，可以考察序参量在这一变换下如何改变
（重正化）。记该变换为 $\tau$，则哈密顿量变换为 $\mathcal{H}'=\tau(\mathcal{H})$。
临界点处 $\xi\to\infty$，故必存在极限函数 $\mathcal{H}^{*}$ 使 $\mathcal{H}^{*}$
是 $\tau$ 的不动点，即 $\tau(\mathcal{H}^{*})=\mathcal{H}^{*}$。这一寻找标度变换
不动点的方法称为重正化群（renormalization group）方法——因为标度变换的集合构成
一个群。它在凝聚态物理中极为成功，并与量子场论中的重正化有深刻的联系。不过
重正化群的更详细处理超出本书范围。

团簇模型可用于描述铁磁体的顺磁区（$T>T_{\mathrm{C}}$）。它保留平均场方法的某些
简单性（从而保留预言能力），又尝试对涨落建模：精确计算单个有限尺寸自旋团簇的
组态，再把它与样品其余部分的平均场耦合。当然关联长度在 $T_{\mathrm{C}}$ 处发散，
故越接近转变团簇尺寸需要越大。若选择团簇尺寸不受限制、可以变化的团簇模型，可以
得到与数据极好的符合。\fn{2}{见 R.~V.~Chamberlin, Nature, \textbf{408}, 337 (2000)。}

\section{刚性}

打破对称涉及基态的选择：自旋该都朝上、朝下、朝左还是朝右？宏观样品的能量在其
整个体积内以同一种方式打破对称时最低。若试图让宏观样品的不同部分以不同方式打破
对称，将出现反映额外能量代价的力——由此产生广义刚性或刚度（stiffness）。这种
刚性与有序密切相关，以下对晶体的论证可以说明。

考虑晶格常数为 $a$ 的立方晶体，设想把某一平面的键拉伸，使 $a$ 变为 $a+u$，产生
应变 $\epsilon=u/a$（见图 6.10）。这产生等于 $G\epsilon$ 的应力（G 为弹性模量），
单个键上的力为 $G\epsilon a^{2}=Gau$，键中储存的能量为 $\frac{1}{2}Gau^{2}$。高温
（远高于德拜温度\fn{3}{德拜温度 $\Theta_{\mathrm{D}}$ 是最高被占据声子模的能量除以 $k_{\mathrm{B}}$——假定声子无色散。当然总有色散，但 $\Theta_{\mathrm{D}}$ 仍刻画了系统的典型声子能量。多数德拜温度在 100--1000 K 范围。$T\gg\Theta_{\mathrm{D}}$ 时声子模可作经典处理。}）下每个键的平均储存能量为 $\frac{1}{2}k_{\mathrm{B}}T$
（能量均分定理，见附录 E），于是
\begin{equation*}
\langle u^{2}\rangle = \frac{k_{\mathrm{B}}T}{Ga}\tag{6.8}
\end{equation*}
故 $\langle u^{2}\rangle\propto1/G$：G 为零则涨落发散，反之亦然。非零弹性模量
（即刚性）的存在与涨落的有限性（即晶体的稳定）相联系。

这就是说：既然有有序的晶体，就必有刚性的晶体。晶体不易形变（与液体相反），因为
存在正比于弹性模量与 $(\nabla u)^{2}$ 的弹性能（u 是晶格位移），且晶格把力传递到
宏观样品的另一端。固体因此是刚性的。

这些想法同样适用于铁磁体。铁磁体中自旋排齐，把它们相对彼此转动要付出能量——
这就是永磁现象。磁化强度非均匀时有正比于 $(\nabla\mathbf{M})^{2}$ 的交换代价
（见 4.2.7 节，尤其式 4.18）。

最后一个例子：超导体中波函数的相位 $\phi$ 是序参量。超导体中 $\phi$ 在样品上
均匀，存在正比于 $(\nabla\phi)^{2}$ 的“相位刚度”能量。（事实上，把相位沿样品扭转
会产生正比于 $\nabla\phi$ 的超流。）

\begin{figure}
\centering
\includegraphics[width=0.5\textwidth]{fig6_10.pdf}
\caption{晶格参数为 $a$ 的立方晶体被拉伸，产生应变 $\epsilon=u/a$。}
\end{figure}

\section{激发}

固体在 T=0 有序，尽管零点涨落意味着即便此时原子也并非纯静态。非零温度下，有序
被热激发的晶格振动破坏——它们量子化为声子。声子的行为由色散关系刻画，即角频率
$\omega$ 与波矢 q 之间的关系（等价地，能量 $\hbar\omega$ 与动量 $\hbar q$ 之间的
关系）。一维单原子链的例子示于图 6.11。关键特征：声学声子在 q=0 处 $\omega=0$，
故只要波长 $\lambda=2\pi/q$ 足够长，产生一个波矢 q 的声子所需能量小到可以忽略。
只要温度非零声学声子就能被热激发，是因为从系统基态（q=0、$\omega=0$）到最低
声学声子能级之间没有需要跨越的能隙——声子色散关系在 q=0 处没有能隙（与光学声子
相反）。三维晶体中声学声子给出低温 $T^{3}$ 热容。

\begin{figure}
\centering
\includegraphics[width=0.5\textwidth]{fig6_11.pdf}
\caption{质量为 m 的原子由弹簧常数 K 的键连接的一维单原子链的声子色散。这些声子
是声学的：没有能隙（q=0 处 $\omega=0$），小 q 时 $\omega\approx v_{\mathrm{s}}q$
（$v_{\mathrm{s}}$ 为声速）。色散是无质量的，因为 q=0 处 $\omega=0$（即能量
$E=\hbar\omega=0$ 于动量 $p=\hbar q=0$）。相对论性粒子的能量满足
$E^{2}=p^{2}c^{2}+m^{2}c^{4}$，零动量能量为 $mc^{2}$——正比于质量。}
\end{figure}

每当打破一个连续的整体对称（如从液体造出固体、从顺磁体造出铁磁体），就可以以
可忽略的小能量代价在序参量中产生长波长激发。这类激发称为戈德斯通模（Goldstone
modes，有时称戈德斯通玻色子）。因为不费能量，它们是“无质量的”。粒子物理中的
例子是光子——一个戈德斯通玻色子。超导体的情形相当不同：打破的是连续的局域对称，
不产生戈德斯通模——原因相当微妙，与希格斯机制有关。

铁磁体在 T=0 完全有序；非零温度下有序被自旋波破坏，自旋波量子化为磁振子
（magnons）。关键特征：只要波长足够长，产生一个磁振子所需能量小到可以忽略。
磁振子在铁磁体中扮演声子在固体中的角色，是系统的戈德斯通模。下面将证明：对各向
同性铁磁体，磁振子色散关系在 q=0 处没有能隙。

\subsection{磁振子}

本节推导各向同性铁磁体的磁振子色散关系。这是重要问题，故给出两种推导：第一种
半经典，第二种量子力学。

（1）从自旋波色散的半经典推导开始。回忆海森伯模型的哈密顿量
\begin{equation*}
\hat{\mathcal{H}} = -\sum_{\langle ij\rangle}\mathsf{J}\,\hat{\mathbf{S}}_i\cdot\hat{\mathbf{S}}_j\tag{6.9}
\end{equation*}
（即式 6.4）。一维链中每个自旋有两个邻居，哈密顿量约化为
\begin{equation*}
\hat{\mathcal{H}} = -2\mathsf{J}\sum_i\hat{\mathbf{S}}_i\cdot\hat{\mathbf{S}}_{i+1}.\tag{6.10}
\end{equation*}
用式 C.7 计算 $\langle\hat{\mathbf{S}}_j\rangle$ 的时间依赖：
\begin{align*}
\frac{\mathrm{d}\langle\hat{\mathbf{S}}_j\rangle}{\mathrm{d}t} &= \frac{1}{\mathrm{i}\hbar}\langle[\hat{\mathbf{S}}_j, \hat{\mathcal{H}}]\rangle\tag{6.11}\\
&= -\frac{2\mathsf{J}}{\mathrm{i}\hbar}\langle[\hat{\mathbf{S}}_j, \cdots+\hat{\mathbf{S}}_{j-1}\cdot\hat{\mathbf{S}}_j+\hat{\mathbf{S}}_j\cdot\hat{\mathbf{S}}_{j+1}+\cdots]\rangle\tag{6.12}\\
&= -\frac{2\mathsf{J}}{\mathrm{i}\hbar}\langle[\hat{\mathbf{S}}_j, \hat{\mathbf{S}}_{j-1}\cdot\hat{\mathbf{S}}_j]+[\hat{\mathbf{S}}_j, \hat{\mathbf{S}}_j\cdot\hat{\mathbf{S}}_{j+1}]\rangle\tag{6.13}\\
&= \frac{2\mathsf{J}}{\hbar}\langle\hat{\mathbf{S}}_j\times(\hat{\mathbf{S}}_{j-1}+\hat{\mathbf{S}}_{j+1})\rangle.\tag{6.14}
\end{align*}
现在把每个格点上的自旋当经典矢量处理。系统基态是所有自旋排齐（设沿 $z$ 轴）：
$S_j^{z}=S$、$S_j^{x}=S_j^{y}=0$。考虑对该态的小偏离：$S_j^{z}\approx S$、
$S_j^{x},S_j^{y}\ll S$，于是
\begin{align*}
\frac{\mathrm{d}S_j^{x}}{\mathrm{d}t} &\approx \frac{2\mathsf{J}S}{\hbar}(2S_j^{y}-S_{j-1}^{y}-S_{j+1}^{y})\tag{6.15}\\
\frac{\mathrm{d}S_j^{y}}{\mathrm{d}t} &\approx -\frac{2\mathsf{J}S}{\hbar}(2S_j^{x}-S_{j-1}^{x}-S_{j+1}^{x})\tag{6.16}\\
\frac{\mathrm{d}S_j^{z}}{\mathrm{d}t} &\approx 0.\tag{6.17}
\end{align*}
寻找简正模解，令
\begin{equation*}
S_j^{x} = A\,\mathrm{e}^{\mathrm{i}(qja-\omega t)}\tag{6.18}
\end{equation*}
\begin{equation*}
S_j^{y} = B\,\mathrm{e}^{\mathrm{i}(qja-\omega t)}\tag{6.19}
\end{equation*}
q 为波矢。少许代数运算给出 $A=\mathrm{i}B$（$x$ 与 $y$ 方向运动相位差 $\pi/2$），
从而
\begin{equation*}
\hbar\omega = 4\mathsf{J}S(1-\cos qa),\tag{6.20}
\end{equation*}
这就是自旋波的色散关系，绘于图 6.12。自旋波的样子如图 6.13 所示。

\begin{figure}
\centering
\includegraphics[width=0.45\textwidth]{fig6_12.pdf}
\caption{一维自旋链的自旋波色散关系。}
\end{figure}

\begin{figure}
\centering
\includegraphics[width=0.4\textwidth]{fig6_13.pdf}
\caption{自旋链上的自旋波。（a）透视图。（b）俯视图。}
\end{figure}

（2）用量子力学方法导出式 6.20 或许更令人满意，现在重复这一推导。\fn{4}{目前取 $S=\frac{1}{2}$，但很容易推广。}考虑系统基态 $|\Phi\rangle$——所有自旋沿 $+z$：
\begin{center}
$|\Phi\rangle = \cdots\uparrow\uparrow\uparrow\uparrow\uparrow\uparrow\uparrow\uparrow\uparrow\cdots$
\end{center}
一维中海森伯模型的哈密顿量可写为
\begin{equation*}
\hat{\mathcal{H}} = -2\mathsf{J}\sum_i\left[\hat{S}_i^{z}\hat{S}_{i+1}^{z} + \frac{1}{2}(\hat{S}_i^{+}\hat{S}_{i+1}^{-}+\hat{S}_i^{-}\hat{S}_{i+1}^{+})\right],\tag{6.21}
\end{equation*}
于是 $\hat{\mathcal{H}}|\Phi\rangle=-NS^{2}\mathsf{J}|\Phi\rangle$。现在制造一个
激发：翻转格点 j 的自旋，考虑态 $|j\rangle=\hat{S}_j^{-}|\Phi\rangle$——格点 j
自旋翻转了的基态（见图 6.14）：
\begin{center}
$|j\rangle = \cdots\uparrow\uparrow\downarrow\uparrow\uparrow\uparrow\uparrow\uparrow\uparrow\cdots$
\end{center}
翻转一个自旋使系统总自旋改变了 $\frac{1}{2}-(-\frac{1}{2})=1$，因此该激发有整数
自旋、是玻色子。把哈密顿量作用到这个新态上：
\begin{equation*}
\hat{\mathcal{H}}|j\rangle = 2\left[(-NS^{2}\mathsf{J}+2S\mathsf{J})|j\rangle - S\mathsf{J}|j+1\rangle - S\mathsf{J}|j-1\rangle\right],\tag{6.22}
\end{equation*}
它不是 $|j\rangle$ 的常数倍——这个态不是哈密顿量的本征态。然而，可以通过寻找
平面波形式的解来对角化哈密顿量：
\begin{equation*}
|q\rangle = \frac{1}{\sqrt{N}}\sum_j\mathrm{e}^{\mathrm{i}qR_j}|j\rangle.\tag{6.23}
\end{equation*}
态 $|q\rangle$ 本质上是一个离域化（抹开）到所有格点的翻转自旋。由于它由“单个翻转
自旋”的态组成，$|q\rangle$ 自身的总自旋为 $NS-1$。于是容易证明
\begin{equation*}
\hat{\mathcal{H}}|q\rangle = E(q)|q\rangle,\tag{6.24}
\end{equation*}
其中
\begin{equation*}
E(q) = -2NS^{2}\mathsf{J} + 4\mathsf{J}S(1-\cos qa).\tag{6.25}
\end{equation*}
激发的能量为 $\hbar\omega=4JS(1-\cos qa)$，与上面式 6.20 的结果相同。

\begin{figure}
\centering
\includegraphics[width=0.65\textwidth]{fig6_14.pdf}
\caption{态 $|\Phi\rangle$ 由所有沿 $+z$ 方向的自旋组成；态 $|j\rangle$ 是格点 j
处自旋翻转了的基态。}
\end{figure}

\subsection{布洛赫 $T^{3/2}$ 定律}

小 q 时式 6.20 给出
\begin{equation*}
\hbar\omega \approx 2\mathsf{J}Sq^{2}a^{2},\tag{6.26}
\end{equation*}
即 $\omega\propto q^{2}$。三维中态密度为
\begin{equation*}
g(q)\,\mathrm{d}q \propto q^{2}\,\mathrm{d}q\tag{6.27}
\end{equation*}
从而在低温（只有小 q、小 $\omega$ 重要）下
\begin{equation*}
g(\omega)\,\mathrm{d}\omega \propto \omega^{1/2}\,\mathrm{d}\omega\tag{6.28}
\end{equation*}
自旋波的量子化方式与格波相同：后者称声子，前者相应称磁振子。如 6.6.1 节所示，
磁振子自旋为 1、是玻色子。

温度 T 时被激发的磁振子模数 $n_{\mathrm{magnon}}$ 由磁振子态密度对所有频率积分
得到，须乘上玻色因子 $(\exp(\hbar\omega/k_{\mathrm{B}}T)-1)^{-1}$（磁振子是玻色子）：
\begin{equation*}
n_{\mathrm{magnon}} = \int_{0}^{\infty}\frac{g(\omega)\,\mathrm{d}\omega}{\exp(\hbar\omega/k_{\mathrm{B}}T)-1},\tag{6.29}
\end{equation*}
用代换 $x=\hbar\omega/k_{\mathrm{B}}T$ 求值。低温下（三维 $g(\omega)\propto\omega^{1/2}$）
得
\begin{equation*}
n_{\mathrm{magnon}} = \left(\frac{k_{\mathrm{B}}T}{\hbar}\right)^{3/2}\int_{0}^{\infty}\frac{x^{1/2}\,\mathrm{d}x}{\mathrm{e}^{x}-1} \propto T^{3/2}.\tag{6.30}
\end{equation*}
由于每个热激发的磁振子模把总磁化强度减少 S=1，\fn{5}{因为每个磁振子模是一个离域化的单个翻转自旋，见上一节。}低温下自发磁化强度相对 T=0 值的减少为
\begin{equation*}
\frac{M(0)-M(T)}{M(0)} \propto T^{3/2}.\tag{6.31}
\end{equation*}
这一结果称为布洛赫 $T^{3/2}$ 定律（Bloch $T^{3/2}$ law），在低温区与实验数据
符合（见图 6.15）。\fn{6}{布洛赫 $T^{3/2}$ 定律其实只对畴内的自发磁化强度严格成立。图 6.15 的数据实际测量的正是它——它们是用 $\mu$SR 在零外场下测得的局域内场。}磁振子模的能量为
\begin{equation*}
E_{\mathrm{magnon}} = \int_{0}^{\infty}\frac{\hbar\omega\,g(\omega)\,\mathrm{d}\omega}{\exp(\hbar\omega/k_{\mathrm{B}}T)-1} \propto T^{5/2},\tag{6.32}
\end{equation*}
故热容 $C=\partial E_{\mathrm{magnon}}/\partial T$ 也正比于 $T^{3/2}$。

\begin{figure}
\centering
\includegraphics[width=0.5\textwidth]{fig6_15.pdf}
\caption{铁磁体中的自发磁化强度。低温下可用自旋波模型拟合、遵循布洛赫 $T^{3/2}$
律；接近临界温度时磁化强度正比于 $(T-T_{\mathrm{C}})^{\beta}$（$\beta$ 为临界
指数）。两种行为都不能在整个温度范围拟合真实数据。数据取自一种有机铁磁体，
$T_{\mathrm{C}}\approx0.67$ K，$\beta\approx0.36$——适用于三维海森伯模型。}
\end{figure}

\subsection{Mermin--Wagner--Berezinskii 定理}

然而在一维与二维，式 6.29 的积分发散，故对各向同性的一维、二维海森伯模型，
一切 $T>0$ 都有 $M\to0$。\fn{7}{对不同的维数，式 6.29 中的函数 $g(\omega)$ 对 $\omega$ 有不同的依赖。}有限温度下产生的自旋波数目发散，因此自发铁磁性不可能。这一结果由
Mermin 与 Wagner 于 1966 年首先证明，Berezinskii 独立证明。二维体系（具有连续
对称性）中长程序的缺失常称为 Mermin--Wagner--Berezinskii 定理。

该结果只适用于各向同性的海森伯铁磁体。它具有旋转对称——所有自旋方向可以整体
旋转而不费任何额外能量。这意味着长波长激发——自旋态在相当长距离上偏离基态值
——几乎不费能量。\fn{8}{若自旋波波长足够长，交换能量代价 $\propto\int(|\nabla M_x|^{2}+|\nabla M_y|^{2})\,\mathrm{d}x\,\mathrm{d}y$（二维）被最小化。见习题 6.7。}于是自旋涨落以极小的能量代价即可激发；在一维与二维，它们摧毁长程序。然而若存在
显著各向异性，把自旋从基态值转开将有能量代价。事实证明，容忍这些涨落所要付出的
各向异性能量代价随激发半径 R 的平方增长，故各向异性能量会抑制除最小者之外的所有
非线性涨落。\fn{9}{这一想法在习题 6.7 中探讨。}正是这种对称性破缺场的存在，可以稳定二维体系中的长程序。
真实体系中自旋间还存在偶极相互作用——虽远弱于交换相互作用，却是各向异性的，
能以类似方式抑制涨落的生长。

实验上观察到，一些超薄磁性薄膜可以表现自发铁磁性。有时各向异性使自旋垂直于膜面
在能量上有利（系统被称为具有垂直易磁化轴）。此时各向异性在自旋波激发谱中（基态
与激发态之间）打开能隙，长程磁有序可以在有限温度存在。偶极相互作用或膜面内自旋
能量的各向异性也能导致类似效果、稳定磁化强度。

\subsection{自旋波的测量}

用于磁结构测定的弹性中子散射已在 5.7.2 节讨论。自旋波色散可以用非弹性中子散射
测量：此时入射中子波矢的大小不再等于散射中子波矢 $k'$ 的大小，中子能量也从
\begin{equation*}
E = \frac{\hbar^{2}k^{2}}{2m_{\mathrm{n}}}\tag{6.33}
\end{equation*}
变为
\begin{equation*}
E' = \frac{\hbar^{2}k'^{2}}{2m_{\mathrm{n}}}\tag{6.34}
\end{equation*}
因为中子在样品中产生了一个能量 $\hbar\omega$、波矢 q 的激发。能量与动量守恒给出
\begin{equation*}
E = E' + \hbar\omega\tag{6.35}
\end{equation*}
\begin{equation*}
\mathbf{k} = \mathbf{k}' + \mathbf{q} + \mathbf{G},\tag{6.36}
\end{equation*}
G 是倒格矢波矢，\fn{10}{加入 G 是必要的：磁振子的色散关系与声子一样，在倒格子上是周期的。}于是测得 $\mathbf{k}$、$\mathbf{k}'$、E、$E'$ 便能确定 $\omega$ 与 q。
中子的能量与原子过程及电子过程的能量相近，即 meV 到 eV 范围。

因此可以探测从量子隧穿的 $\mu$eV，经过分子的平移、转动、振动与晶格模式，直到
材料电子结构中 eV 跃迁的能量尺度；磁振子能量典型在 $10^{-3}$--$10^{-2}$ eV，
因此可以有效地用非弹性中子散射测量。

实验通常用中子三轴谱仪（triple-axis spectrometer，见图 6.16）进行——如此命名
是因为单色器、样品与分析器晶体的角度都可以独立变化。这允许测量对应大范围可能
$\omega$ 与 $\mathbf{q}$ 的散射中子，特别地可以作固定 $\omega$ 变 q（或反之）的
扫描。用该技术获得的数据示例：图 6.17 为 $\mathrm{Co}_{0.92}\mathrm{Fe}_{0.08}$，
图 6.18 为铁磁氧化物 $\mathrm{La}_{0.7}\mathrm{Pb}_{0.3}\mathrm{MnO}_3$
（它表现庞磁电阻，见 8.9.5 节）。

\begin{figure}
\centering
\includegraphics[width=0.45\textwidth]{fig6_16.pdf}
\caption{中子三轴谱仪示意。设置旋转屏蔽与单色器晶体的角度可选定入射中子能量；
样品也可旋转；分析器晶体也可旋转，从而测量散射中子的能量。}
\end{figure}

\begin{figure}
\centering
\includegraphics[width=0.45\textwidth]{fig6_17.pdf}
\caption{$\mathrm{Co}_{0.92}\mathrm{Fe}_{0.08}$ 合金中自旋波的磁振子能谱，
室温测得（Sinclair 与 Brockhouse 1960）。图中 k 是正文中的磁振子波矢 q。}
\end{figure}

\begin{figure}
\centering
\includegraphics[width=0.45\textwidth]{fig6_18.pdf}
\caption{铁磁氧化物 $\mathrm{La}_{0.7}\mathrm{Pb}_{0.3}\mathrm{MnO}_3$ 中
10 K 非弹性中子散射测得的自旋波色散；横轴为磁振子波矢。取自 Perring 等 1996。}
\end{figure}

\section{磁畴}

若宏观体系的不同区域以不同方式打破对称，则这些区域的界面上刚性可能失效。一般
预期存在畴壁、缺陷、涡旋、位错与其他奇点。铁磁体中最重要的奇点是畴壁（domain
wall）。

外斯首先提出铁磁体包含许多称为磁畴（domains）的小区域，每个畴内局域磁化达到
饱和值；不同畴的磁化方向不必平行；畴由畴壁分隔。磁畴的存在解释了令人惊讶的
观察：某些（磁“软”的\fn{11}{即磁‘软’的材料，见第 6.7.9 节。}）铁磁样品在室温下仅用极弱的磁场（低至 $10^{-6}$ T）就能使整块样品达到饱和磁化强度（相应
$\mu_0M\sim1$ T）。如此低的外场对顺磁体几乎没有效果。\fn{12}{若某顺磁体的 $\chi\sim10^{-3}$（见表 2.1），则 $10^{-6}$ T 的外场只能产生 $\mu_0M\sim10^{-9}$ T 的磁化强度。}铁磁样品中的巨大效应在于：外场无须宏观地把磁矩排序（它们早已有序！），
只须使各磁畴排齐——这可以通过能量上无痛的畴壁移动过程实现。\fn{13}{我们将看到，某些铁磁材料中畴壁的移动实际上是耗能的；这类材料需要强外场才能从一个磁态切换到另一个。}

同一铁磁样品中，零外场下磁化强度也可能为零——这也是磁畴的表现：每个畴内磁化
饱和，但各畴磁化方向安排得使样品净磁化强度为零。

\subsection{畴壁}

相邻畴之间的边界称为畴壁。畴壁可按两畴磁化强度间的夹角分类（见图 6.19）：180°
畴壁分隔磁化相反的畴；90° 畴壁分隔磁化垂直的畴。

最常见的 180° 壁是布洛赫壁（Bloch wall，见图 6.20(a)）：磁化在平行于壁面的平面内
旋转。另一种可能组态是奈尔壁（Néel wall，见图 6.20(b)）：磁化在垂直于壁面的平面内
旋转。下面尝试计算布洛赫壁的畴壁宽度。

铁磁体中转动相邻自旋要费能量。夹角为 $\theta$ 的两个自旋 $S_1$、$S_2$
（图 6.21）能量为 $-2\mathsf{J}S_1\cdot S_2=-2\mathsf{J}S^{2}\cos\theta$；
$\theta=0$ 时能量 $-2\mathsf{J}S^{2}$。于是 $\theta\neq0$ 的能量代价在 $\theta\ll1$
时近似为\fn{14}{利用 $\theta\ll1$ 时 $\cos\theta\approx1-\theta^{2}/2$。}$\mathsf{J}S^{2}\theta^{2}$。布洛赫壁中自旋跨 N 个格点转过 $\pi$ 角（图 6.20(a)），一条自旋线的能量代价是 N 个
$\mathsf{J}S^{2}\theta^{2}$ 贡献之和，$\theta=\pi/N$，即 $\mathsf{J}S^{2}\pi^{2}/N$。
布洛赫壁中是自旋平面（见图 6.20），因此我们关心的是布洛赫壁单位面积的能量
$\sigma_{\mathrm{BW}}$。一平方米的壁上有 $1/a^{2}$ 条我们计算过的自旋线。故
\begin{equation*}
\sigma_{\mathrm{BW}} = \mathsf{J}S^{2}\frac{\pi^{2}}{Na^{2}}\tag{6.37}
\end{equation*}
它在 $N\to\infty$ 时趋于零。这一结果似乎表明：一旦形成畴壁，它就会自己解开、在
整个系统中不断长大——因为自旋相互扭曲要费能量，除非有别的相互作用阻止，它们都
会解开。这个别的相互作用就是磁晶各向异性（magnetocrystalline anisotropy）。

\begin{figure}
\centering
\includegraphics[width=0.4\textwidth]{fig6_19.pdf}
\caption{（a）180° 畴壁。（b）90° 畴壁。}
\end{figure}

\begin{figure}
\centering
\includegraphics[width=0.45\textwidth]{fig6_20.pdf}
\caption{（a）布洛赫壁，（b）奈尔壁。}
\end{figure}

\begin{figure}
\centering
\includegraphics[width=0.4\textwidth]{fig6_21.pdf}
\caption{彼此成 $\theta$ 角的两个自旋，产生能量代价 $\frac{1}{2}\mathsf{J}S^{2}\theta^{2}$。}
\end{figure}

\subsection{磁晶各向异性}

晶体具有易磁化轴与难磁化轴：沿某些晶体学方向容易磁化晶体，沿另一些则较难
（Fe、Co、Ni 单晶的情形见图 6.22）。例如在 Co 中，这一各向异性导致形如
\begin{equation*}
E = K_1\sin^{2}\theta + K_2\sin^{4}\theta\tag{6.38}
\end{equation*}
的附加能量，$K_1$、$K_2$ 是各向异性常数，$\theta$ 是磁化强度与六角密堆面堆叠
方向之间的夹角。由于这些常数为正，磁化沿堆叠方向时能量最低。式 6.38 中 E、$K_1$、
$K_2$ 是能量密度（以 J\,m$^{-3}$ 量度）。各向异性常数强烈依赖温度。

\begin{figure}
\centering
\includegraphics[width=0.55\textwidth]{fig6_22.pdf}
\caption{Fe、Co、Ni 沿不同方向加场时的磁化强度，显示各向异性。取自 Honda 与 Kaya
1926、Kaya 1928。}
\end{figure}

式 6.38 适用于单轴各向异性——能量只依赖与单一轴的夹角（Co 中即六角密堆面的堆叠
轴）。立方体系中恰当的表达是
\begin{equation*}
E = K_1(m_x^{2}m_y^{2}+m_y^{2}m_z^{2}+m_z^{2}m_x^{2}) + K_2m_x^{2}m_y^{2}m_z^{2}+\cdots,\tag{6.39}
\end{equation*}
$\mathbf{m}=(m_x,m_y,m_z)=\mathbf{M}/|\mathbf{M}|$。球坐标下为
\begin{equation*}
E = K_1\left(\frac{1}{4}\sin^{2}\theta\sin^{2}2\phi+\cos^{2}\theta\right)\sin^{2}\theta + \frac{K_2}{16}\sin^{2}2\phi\sin^{2}2\theta\sin^{2}\theta+\cdots.\tag{6.40}
\end{equation*}
各向异性能量源于自旋--轨道相互作用与角动量的部分淬灭。各向异性能量通常在
$10^{2}$--$10^{7}$ J\,m$^{-3}$，相当于每原子 $10^{-8}$--$10^{-3}$ eV。对称性低的
（磁性离子）晶格各向异性能量较大，高对称者较小。例如立方 Fe、Ni 的 $K_1$ 分别为
$4.8\times10^{4}$ 与 $-5.7\times10^{3}$ J\,m$^{-3}$，六角 Co 的 $K_1=5\times10^{5}$
J\,m$^{-3}$；低对称的永磁材料 $\mathrm{Nd_2Fe_{14}B}$ 与 $\mathrm{SmCo_5}$ 的
$K_1$ 分别为 $5\times10^{6}$ 与 $1.7\times10^{7}$ J\,m$^{-3}$。

另一项附加能量来自与样品形状相联系的退磁能，称为形状各向异性（shape anisotropy）。
薄膜中形状各向异性 $\frac{1}{2}\mu_0M^{2}\cos^{2}\theta$（$\theta$ 为膜法线与 M
的夹角）使磁化强度保持在膜面内以节省能量。

\subsection{畴壁宽度}

铁磁体的磁畴中，磁化偏好沿易方向排列；但在畴之间的畴壁里它必须转动，一部分沿
难轴、耗费能量。设各向异性能量密度取简单形式 $E=K\sin^{2}\theta$（K 为各向异性
常数），便可轻松求出各向异性能量对布洛赫壁的贡献。取 K>0，自旋偏好沿 $\theta=0$
或 $\theta=\pi$。把 N 个自旋的贡献相加，并把求和换成积分（连续极限）以简化：
能量密度贡献为
\begin{equation*}
\sum_{i=1}^{N}K\sin^{2}\theta_i \approx \frac{N}{\pi}\int_{0}^{\pi}K\sin^{2}\theta\,\mathrm{d}\theta = \frac{NK}{2}.\tag{6.41}
\end{equation*}
壁单位面积的能量贡献于是为 NKa/2。畴壁单位面积的总能量（含交换能贡献（式 6.37）
与各向异性能量贡献（式 6.41））为
\begin{equation*}
\sigma_{\mathrm{BW}} = \mathsf{J}S^{2}\frac{\pi^{2}}{Na^{2}} + \frac{NKa}{2}.\tag{6.42}
\end{equation*}
这给出所需的行为：一项正比于 1/N（趋向解开壁、使之变大），另一项正比于 N（趋向
收紧壁、使之变小）。用 $\mathrm{d}E_{\mathrm{BW}}/\mathrm{d}N=0$ 求平衡组态，得
\begin{equation*}
N = \pi S\sqrt{2\mathsf{J}/Ka^{3}}\tag{6.43}
\end{equation*}
布洛赫壁宽度为
\begin{equation*}
\delta = Na = \pi S\sqrt{\frac{2\mathsf{J}}{Ka}},\tag{6.44}
\end{equation*}
a 是晶格间距。$\mathsf{J}$ 越大壁越厚，K 越大壁越薄。畴壁单位面积的能量为
\begin{equation*}
\sigma_{\mathrm{BW}} = \pi S\sqrt{\frac{2\mathsf{J}K}{a}}.\tag{6.45}
\end{equation*}
用式 4.19 定义的参数 A（立方晶体）表示，畴壁厚度为
\begin{equation*}
\delta = \pi\sqrt{\frac{A}{K}},\tag{6.46}
\end{equation*}
单位面积能量为
\begin{equation*}
\sigma_{\mathrm{BW}} = \pi\sqrt{AK}.\tag{6.47}
\end{equation*}

\subsection{磁畴的形成}

既然形成畴壁要付出能量，人们或许奇怪：为什么一开始会形成不止一个畴？原因是磁畴
的形成节省与偶极场相联系的能量。由于 $\nabla\cdot\mathbf{H}=-\nabla\cdot\mathbf{M}$
（见附录 B），凡 M 中断或开始之处（例如样品边缘），磁场的散度不为零，产生退磁场，
充满空间，每立方米耗费 $B^{2}/2\mu_0$ 焦耳的能量。与退磁场相联系的能量称为退磁能、
静磁能或偶极能：
\begin{equation*}
-\frac{\mu_0}{2}\int_{V}\mathbf{M}\cdot\mathbf{H}_{\mathrm{d}}\,\mathrm{d}\tau.\tag{6.48}
\end{equation*}
$\mathbf{H}_{\mathrm{d}}$ 是退磁场，积分遍及样品体积（见附录 D）。对沿主轴磁化的
椭球样品，此能量化为
\begin{equation*}
\frac{\mu_0}{2}NM^{2}V\tag{6.49}
\end{equation*}
N 是退磁因子，V 是样品体积。

把样品分成磁畴可以节省这一偶极能，但每创造一个畴都因畴壁而耗能。因此，正如畴壁
尺寸是交换能与各向异性能量的平衡，磁畴的形成是退磁场代价与畴壁代价的平衡。

\begin{figure}
\centering
\includegraphics[width=0.3\textwidth]{fig6_23.pdf}
\caption{铁磁样品的三种畴结构选择：（a）均匀磁化；（b）分成两个畴；（c）简单
封闭畴结构。}
\end{figure}

图 6.23 给出铁磁样品磁畴结构的三种选择：图 6.23(a) 的单畴结构没有畴壁但偶极能
巨大；把样品分成两畴（图 6.23(b)）可以降低偶极能，代价是引入畴壁；图 6.23(c) 的
封闭畴（closure domain）结构消除偶极能但引入若干畴壁。

偶极能还能决定可形成的畴壁类型。布洛赫壁在块体中更受青睐，因为其偶极能更小；
奈尔壁则常见于薄膜（把自旋转出膜面要付偶极能代价）。

\subsection{磁化过程}

图 6.24 给出测量铁磁样品磁化强度随外加磁场变化的典型磁滞回线。若样品被外场磁化
到饱和磁化强度 $M_{\mathrm{s}}$，外场降为零时磁化强度降为剩余磁化强度
（remanent magnetization）$M_{\mathrm{r}}$；需要等于矫顽场（coercive field）
$H_{\mathrm{c}}$ 的磁场才能把磁化强度切换到相反方向。参数 $M_{\mathrm{r}}$ 与
$H_{\mathrm{c}}$ 可用来表征铁磁体。

\begin{figure}
\centering
\includegraphics[width=0.42\textwidth]{fig6_24.pdf}
\caption{磁滞回线，标出饱和磁化强度 $M_{\mathrm{s}}$、剩余磁化强度 $M_{\mathrm{r}}$
与矫顽场 $H_{\mathrm{c}}$。}
\end{figure}

\begin{figure}
\centering
\includegraphics[width=0.5\textwidth]{fig6_25.pdf}
\caption{外场对铁晶须表面磁畴图样的影响：随外场 B 从 0 增大到最大值，畴壁发生
移动。}
\end{figure}
\bio{John Kerr\\（1824--1907）\quad Heinrich Barkhausen\\（1881--1956）}

铁磁体磁化时长度会略微改变——这一效应称为磁致伸缩（magnetostriction），表明
材料的磁学与力学行为相互联系（该联系称为磁弹耦合）。晶体形变是因为可以降低其
各向异性能量——后者可依赖晶体应变。例如立方体系中，若轻微偏离严格立方对称所节省
的各向异性能量超过弹性能量的代价，则形变在能量上有利。

退磁铁磁体被磁化时发生多种过程。首先是畴壁移动：相对外场取向有利的畴吞并不利的
畴（图 6.25 示意这一过程）。更高场下可发生畴转动：各向异性能量被压倒，磁化强度
突然从原磁化方向转到最靠近场方向的晶体学易轴。最高磁场下的最后一个畴过程是各畴
的一致转动——无论易轴难轴，转到与磁场排齐。

畴壁在磁性材料中的运动细节依赖于材料的冶金学性质。由于磁弹耦合，畴壁可被材料中
的应变、表面与杂质钉扎（pinning）；畴壁钉扎因此提高矫顽力。铁磁体的磁化强度也因
畴界运动而以一系列不连续台阶变化，磁化曲线上有时能看到极小的台阶——这称为巴克
豪森效应（Barkhausen effect）。磁化过程中这些突然的不连续变化及其与系统弹性模式
的耦合，有时伴随低水平的声发射，称为磁声发射（magnetoacoustic emission）。

\subsection{磁畴的观察}

观察磁畴有多种技术。比脱粉纹法（Bitter powder technique）把细磁性颗粒的胶体
悬浮液滴在样品表面；磁性颗粒聚集在畴壁附近——那里有吸引它们的强局域场——可用
光学显微镜观察。也可以用偏振光从抛光磁化样品表面的反射观察磁畴，称为克尔显微镜
（Kerr microscopy）：入射光被偏振，其偏振面被磁化样品的反射所转动（即克尔效应，
见 8.8 节）；转角可用第二个偏振器（检偏器）测量。若把光束聚焦成一点并扫过磁性
样品表面，则在检偏器的特定设置下，一些畴显得亮、另一些暗。若样品透明，实验可在
透射下进行（法拉第效应）。

电子显微术同样有效：电子束被内场的洛伦兹力偏转，畴之间的边界（扫描经过时偏转
可以突然变号）显示得特别清楚。这称为洛伦兹显微术（Lorentz microscopy）。

磁力显微术（magnetic force microscopy\fn{15}{磁力显微镜以著名的原子力显微镜为基础。}）用一个带磁化针尖的微悬臂扫过样品表面：畴壁产生的杂散场使悬臂弯曲，
这一弯曲（例如通过悬臂共振频率的变化探测）用于成像表面。

\subsection{小磁性颗粒}

许多铁磁样品在零外场下能量最低的态是退磁态，总磁矩为零。若减小样品尺寸，表面能
（如畴壁能）相对体积能（如退磁能）变得越发昂贵。\fn{16}{这是因为表面能按（样品尺寸）$^2$ 标度，而体积能按（样品尺寸）$^3$ 标度。}于是会到达某个临界尺寸：低于它，去掉畴壁、让样品成为单个磁化
畴在能量上有利——它因此像一块小永磁体。

下面的例子估计单畴颗粒的临界尺寸。计算将假定该临界尺寸仍大于畴壁宽度。

\begin{example}{6.1}
考虑半径 r 的铁磁颗粒。单畴态（图 6.26(a)）的能量用式 D.12 并取球的
$N=\frac{1}{3}$：\mnote{球的体积 $V=\frac{4}{3}\pi r^{3}$。}
\begin{equation*}
E_{(a)} = \frac{1}{6}\mu_0M^{2}V = \frac{2}{9}\mu_0\pi M^{2}r^{3}.\tag{6.50}
\end{equation*}
图 6.26(b) 的态（适用于立方各向异性）近似消除全部退磁能，但引入面积为
$2\times\pi r^{2}$ 的 90° 畴壁，能量代价
\begin{equation*}
E_{(b)} = 2\pi r^{2}\sigma_{\mathrm{W}}^{90^{\circ}},\tag{6.51}
\end{equation*}
$\sigma_{\mathrm{W}}^{90^{\circ}}$ 是 90° 畴壁单位面积的能量代价。当
$E_{(a)}<E_{(b)}$，即
\begin{equation*}
r < \frac{9\sigma_{\mathrm{W}}^{90^{\circ}}}{\mu_0M^{2}}.\tag{6.52}
\end{equation*}
时，图 6.26(a) 的态比图 6.26(b) 更有利。

图 6.26(c) 的态（适用于单轴各向异性）较难考虑，因为两个畴都不是椭球形、仍有
一些退磁能。不过粗略地说，该组态的退磁能是图 6.26(a) 的一半（畴只有一半大小）；
畴壁代价为 $\pi r^{2}\sigma_{\mathrm{W}}^{180^{\circ}}$（$\sigma_{\mathrm{W}}^{180^{\circ}}$
是 180° 畴壁单位面积的能量）。于是
\begin{equation*}
E_{(c)} = \frac{1}{9}\mu_0\pi M^{2}r^{3} + \pi r^{2}\sigma_{\mathrm{W}}^{180^{\circ}}\tag{6.53}
\end{equation*}
故当 $E_{(a)}<E_{(c)}$，即
\begin{equation*}
r < \frac{9\sigma_{\mathrm{W}}^{180^{\circ}}}{\mu_0M^{2}}.\tag{6.54}
\end{equation*}
时，图 6.26(a) 的态比图 6.26(c) 更有利。

若 $\sigma_{\mathrm{W}}^{180^{\circ}}\sim10^{-2}$ J\,m$^{-2}$、$\mu_0M\sim1$ T，
则临界半径 $\sim10^{-7}$ m。注意：若该值小于畴壁宽度，计算便失去意义。
\end{example}

\begin{figure}
\centering
\includegraphics[width=0.3\textwidth]{fig6_26.pdf}
\caption{小球形铁磁颗粒的三种可能磁化组态。}
\end{figure}

\subsection{Stoner--Wohlfarth 模型}

可以通过 Stoner--Wohlfarth 模型计算单畴颗粒的磁化曲线。因为颗粒是单个磁畴，
没有畴壁，只须考虑一致的畴转动。

考虑磁场 H 中的单畴磁性颗粒，H 与其单轴各向异性易轴成 $\theta$ 角。若颗粒磁化
强度与磁场方向成 $\phi$ 角，系统的能量密度为
\begin{equation*}
E = K\sin^{2}(\theta-\phi) - \mu_0HM_{\mathrm{s}}\cos\phi.\tag{6.55}
\end{equation*}
对任意给定的外加磁场，能量取极小即得磁化强度的方向。

图 6.27 给出结果：外加场用无量纲参数
\begin{equation*}
h = \frac{\mu_0M_{\mathrm{s}}H}{2K}\tag{6.56}
\end{equation*}
表示，对两个外场方向绘出。该模型对 $\theta=0$ 与 $\theta=\pi/2$ 有解析解，但图中
展示的是如何数值求解。每种情形都画出在能量面（能量作为 $\phi$ 与 h 的函数）上降低
能量的磁化方向轨迹，并给出相应的磁滞回线。图 6.28 给出更大范围 $\theta$ 值的
磁滞回线，以及多晶方向平均的计算结果（适用于单畴颗粒阵列，忽略颗粒间相互作用）。
该模型表明：即使在没有畴壁钉扎这类不可逆效应的体系中，系统中的各向异性也能导致
磁滞。

\begin{figure}
\centering
\includegraphics[width=0.4\textwidth]{fig6_27.pdf}
\caption{Stoner--Wohlfarth 模型中（a）$\theta=90^{\circ}$ 与（c）$\theta=30^{\circ}$
的磁滞回线。它们可以由在能量面（作为 $h$ 与 $\phi$ 的函数，（b）$\theta=90^{\circ}$
与（d）$\theta=30^{\circ}$）上寻找能量极小点得到。}
\end{figure}

\begin{figure}
\centering
\includegraphics[width=0.4\textwidth]{fig6_28.pdf}
\caption{Stoner--Wohlfarth 模型的磁滞回线：（a）$\theta=0^{\circ}$（粗线）、
$5^{\circ}$、$15^{\circ}$、$30^{\circ}$；（b）$\theta=45^{\circ}$、$60^{\circ}$、
$75^{\circ}$ 与 $90^{\circ}$（粗线）。（c）多晶平均的计算磁滞回线。}
\end{figure}

\subsection{软磁与硬磁材料}

铁磁体沿磁滞回线绕行一周所耗散（为热）的能量正比于回线面积。面积小者称磁软
（magnetically soft）材料；面积大者称磁硬（magnetically hard）材料。磁场循环时
畴壁穿越样品；可以按照畴在样品中移动的难易区分这两大类铁磁材料。

（1）软磁材料容易磁化。软磁材料用于变压器线圈、发电机与电动机：磁化强度须每秒
反转许多次，每周期耗散的能量须最小。软材料有宽畴壁（各向异性能量 K 小）因而容易
移动，矫顽场小。低磁致伸缩也常是想要的——使内应变不诱生局域各向异性能量。例子是
坡莫合金（permalloy，一种商用 Ni/Fe 合金，另加成分），矫顽场
$B_{\mathrm{c}}\sim2\times10^{-7}$ T。

（2）硬磁材料难以磁化，因而难以退磁。硬磁体用作永磁体（如扬声器背部、电动机，
当然还有你冰箱的门上！）以及磁带记录（粉体形式）。这些应用中磁化强度须尽可能
长久保持：磁滞回线每周期的能量耗散越大越好，磁化便不会自发消失。硬磁体磁滞大、
畴壁窄（K 大），畴壁钉扎容易实现。大的离子磁矩与大的晶体场有利于硬磁性质，合适
的材料常涉及稀土，如 $\mathrm{Nd_2Fe_{14}B}$：$T_{\mathrm{C}}=585$ K，矫顽场
$B_{\mathrm{c}}=1.2$ T。

硬磁材料的一个重要应用是磁记录。磁性材料可记录的比特面密度近年来猛增，反映了
对磁性材料的理解与制备的稳步进步。大致数字：1960 年 $\sim0.1$ Mbits\,in$^{-2}$，
1980 年 $\sim10$ Mbits\,in$^{-2}$，2000 年 $\sim25$ Gbits\,in$^{-2}$（1 bit\,in$^{-2}$
即每平方英寸存储一个比特的‘0’或‘1’——这个最现代的产业却还没把单位现代化！）
随着未来面密度继续增大，磁矩的超顺磁涨落（即磁性颗粒的热涨落，见 8.3 节）将变得
重要。

传统上记录磁带与磁盘使用 $\mathrm{Fe_2O_3}$ 小颗粒。高保真录音与录像带性能的
改进基于高矫顽力颗粒：掺 Co 的 $\mathrm{Fe_2O_3}$、$\mathrm{CrO_2}$、金属颗粒
（通常为 Fe）以及钡铁氧体（BaO·6Fe$_2$O$_3$，见 5.3 节）。硬盘通常是一大片刚性
盘基，其上用真空沉积或溅射镀薄膜，通常是 Co/Cr 合金。硬盘高速旋转，读头横越盘面
（可以悬浮在盘面上方一微米的几分之一处），能访问盘面的任何部位。

\section*{延伸阅读}

\begin{itemize}[itemsep=0.2em]
\item 对称性破缺概念的更详细描述见 P.~W.~Anderson,《Basic notions of condensed
matter physics》, Addison-Wesley 1984。
\item 相变的统计力学见 J.~Yeomans,《Statistical mechanics of phase transitions》,
OUP 1992。
\item 关于临界现象也非常有帮助的是 M.~F.~Collins,《Magnetic critical scattering》,
OUP 1989。
\item 对相变与对称性破缺现象（尤其侧重软凝聚态物质）非常透彻的描述是 P.~M.~Chaikin
与 T.~C.~Lubensky,《Principles of condensed matter physics》, CUP 1995。
\item J.~J.~Binney, N.~J.~Dowrick, A.~J.~Fisher 与 M.~E.~J.~Newman,《The theory
of critical phenomena》, OUP 1992，包含对重正化群方法的精彩论述。
\item 关于磁畴与磁性的信息可见 G.~Bertotti,《Hysteresis in magnetism》, Academic
Press 1998 与 A.~Hubert 及 R.~Sch\"{a}fer,《Magnetic domains》, Springer 1998。
\end{itemize}

\section*{本章参考文献}

\begin{itemize}[itemsep=0.2em]
\item K.~Honda 与 S.~Kaya, Sci.~Rep.~Tohoku Univ., \textbf{15}, 721 (1926)；
S.~Kaya, Sci.~Rep.~Tohoku Univ., \textbf{17}, 639, 1157 (1928)。
\item 有用的一般参考书：D.~Craik,《Magnetism: principles and applications》, Wiley
1995；J.~Crangle,《Solid state magnetism》, Edward Arnold 1991。
\item T.~G.~Perring, G.~Aeppli, S.~M.~Hayden, S.~A.~Carter, J.~P.~Remeika 与
S-W.~Cheong, Phys.~Rev.~Lett., \textbf{77}, 711 (1996)。
\item R.~N.~Sinclair 与 B.~N.~Brockhouse, Phys.~Rev., \textbf{120}, 1638 (1960)。
\end{itemize}

\section*{习题}

\exer{6.1} N 个自旋的一维铁磁链由伊辛哈密顿量
\begin{equation*}
\hat{\mathcal{H}} = -2\mathsf{J}\sum_{i=1}^{N-1}\hat{S}_i^{z}\hat{S}_{i+1}^{z} - 2\mathsf{J}\hat{S}_N^{z}\hat{S}_1^{z}.
\end{equation*}
描述（最后一项给出周期性边界条件）。

证明：引入新算符
\begin{equation*}
\zeta_i = 4S_i^{z}S_{i+1}^{z}
\end{equation*}
（本征值 +1 或 $-1$）可求得该系统的配分函数，且配分函数为
$Z=(2\cosh(\mathsf{J}/2k_{\mathrm{B}}T))^{N}$。

用这些结果求 $N\to\infty$ 时链的每自旋热容表达式。推出热容的低、高温行为并画出
热容随温度变化的草图。讨论结果。

\exer{6.2} 单轴铁磁体由哈密顿量
\begin{equation*}
\hat{\mathcal{H}} = -\sum_{ij}\mathsf{J}_{ij}\hat{\mathbf{S}}_i\cdot\hat{\mathbf{S}}_j - \sum_{ij}K_{ij}\hat{S}_i^{z}\hat{S}_j^{z}.\tag{6.57}
\end{equation*}
描述。（a）证明所有自旋沿 z 轴完全排齐的态是哈密顿量的本征态。

（b）求作为波矢 q 的函数的自旋波谱表达式。

（c）把表达式简化到 $\mathsf{J}_{ij}$、$K_{ij}$ 只限最近邻（$\mathsf{J}_0$、
$K_0$）的情形，铁磁体分别为（i）一维链；（ii）二维正方晶格；（iii）三维体心立方
材料。

\exer{6.3} 用习题 6.2 在小波矢下的结果，推出海森伯模型（$K_0=0$）与伊辛模型
（$\mathsf{J}_0=0$）在习题 6.2 的结构（i）、（ii）、（iii）下自旋波数目随温度的
低温依赖。证明你关于情形（i）伊辛模型的结果与习题 6.1 一致，并证明一、二维海森伯
模型在绝对零度以上没有长程磁有序。

\exer{6.4} 把自旋波理论用于反铁磁体，证明色散关系为
\begin{equation*}
\hbar\omega = 4|\mathsf{J}S\sin(qa)|.\tag{6.58}
\end{equation*}
进而证明：反铁磁体在小 q 时 $\omega\propto|q|$，而铁磁体 $\omega\propto q^{2}$。

\exer{6.5} 磁场 H 中铁磁体的朗道理论给出自由能
\begin{equation*}
F(M) = F_0 + a_0(T-T_{\mathrm{C}})M^{2} + bM^{4} - \mu_0MH\tag{6.59}
\end{equation*}
$a_0$ 与 b 为正常数。证明这意味着
\begin{equation*}
M^{2} = u + v\frac{H}{M}\tag{6.60}
\end{equation*}
u 与 v 是应由你确定的常数。对恰好高于 $T_{\mathrm{C}}$、恰好低于 $T_{\mathrm{C}}$
与恰在 $T_{\mathrm{C}}$ 的温度，把 $M^{2}$ 对 H/M 作草图，说明该方法如何用于确定
$T_{\mathrm{C}}$。这就是阿罗特图（Arrott plot，$M^{2}$ 对 H/M 的图）的原理。

\exer{6.6} 求 Stoner--Wohlfarth 模型在 $\theta=0$ 与 $\theta=\pi/2$ 的解析解。

\exer{6.7} 为阐明 Mermin--Wagner--Berezinskii 定理背后的思想，考虑二维中的一个
激发（见图 6.29）：磁化强度作为位置 $\mathbf{r}=(x,y)$ 的函数大小恒定 $|\mathbf{M}|$
而方向变化，由矢量 $\mathbf{M}(\mathbf{r})$ 给定：
\begin{equation*}
\mathbf{M}(\mathbf{r}) = \begin{cases} |\mathbf{M}|\left(\hat{\mathbf{y}}\cos\dfrac{\pi r}{R} - \hat{\mathbf{x}}\sin\dfrac{\pi r}{R}\right) & r<R\\[8pt] -|\mathbf{M}|\hat{\mathbf{y}} & r>R \end{cases}\tag{6.61}
\end{equation*}

\begin{figure}
\centering
\includegraphics[width=0.45\textwidth]{fig6_29.pdf}
\caption{二维磁体中的一个激发。}
\end{figure}

解释为什么该激发的交换能量代价正比于
\begin{equation*}
\int\left[(\nabla M_x)^{2}+(\nabla M_y)^{2}\right]\mathrm{d}x\,\mathrm{d}y.\tag{6.62}
\end{equation*}
并证明它是常数（$\pi^{3}|\mathbf{M}|^{2}$）、与 R 无关。这意味着大尺度涨落与
小尺度涨落相比没有额外能量代价——从而摧毁一切长程序。本题基于 A.~S.~Arrott 与
B.~Heinrich 的文章中的例子，Journal of Magnetism and Magnetic Materials
\textbf{93}, 571 (1991)。

\exer{6.8} 数据以 25 Gbits\,in$^{-2}$ 的面密度写到磁记录盘上。用一些有根据的估计，
把这个面密度换算成其他单位：（a）每个氢原子截面积上的比特数；（b）每张邮票可存
《莎士比亚全集》（按其信息量计）的份数。把这一信息存储密度与人类基因组比较。

\exer{6.9} 考虑磁场中各向同性立方铁磁体（晶格参数 a），哈密顿量
\begin{equation*}
\hat{\mathcal{H}} = -\mathsf{J}\sum_{\langle ij\rangle}\hat{\mathbf{S}}_i\cdot\hat{\mathbf{S}}_j + g\mu_{\mathrm{B}}\sum_j\mathbf{B}\cdot\hat{\mathbf{S}}_j.\tag{6.63}
\end{equation*}
证明 $qa\ll1$ 时其自旋波色散为
\begin{equation*}
\hbar\omega(q) = g\mu_{\mathrm{B}}B + 2\mathsf{J}Sq^{2}a^{2}\tag{6.64}
\end{equation*}

\exer{6.10} 证明哈密顿量
\begin{equation*}
\hat{\mathcal{H}} = -\sum_{\mathbf{n},\delta}\mathsf{J}(\delta)\hat{\mathbf{S}}_{\mathbf{n}}\cdot\hat{\mathbf{S}}_{\mathbf{n}+\delta}\tag{6.65}
\end{equation*}
在 $|\mathbf{q}\cdot\delta_{\max}|\ll1$ 时（$\delta_{\max}$ 是使 $\mathsf{J}(\delta)$
有任何可观贡献的最大矢量 $\delta$）有自旋波色散
\begin{equation*}
\hbar\omega(\mathbf{q}) = \hbar\omega(0) + S\sum_{\delta}\mathsf{J}(\delta)(\mathbf{q}\cdot\delta)^{2}\tag{6.66}
\end{equation*}

\exer{6.11} 式 6.23 中的态 $|q\rangle$ 表示未翻转自旋背景中的一个离域化的单个
翻转自旋。三维中它成为
\begin{equation*}
|\mathbf{q}\rangle = \frac{1}{\sqrt{N}}\sum_{\mathbf{j}}e^{i\mathbf{q}\cdot\mathbf{R}_{\mathbf{j}}}|\mathbf{j}\rangle.\tag{6.67}
\end{equation*}
翻转自旋位于格点 i 的概率为
\begin{equation*}
P = |\langle\mathbf{q}|\mathbf{i}\rangle|^{2}.\tag{6.68}
\end{equation*}
证明 $P=1/N$，从而翻转自旋以等概率分布。

横向自旋关联函数的算符为
\begin{equation*}
\hat{S}_{\mathbf{j}}^{x}\hat{S}_{\mathbf{j}'}^{x} + \hat{S}_{\mathbf{j}}^{y}\hat{S}_{\mathbf{j}'}^{y}.\tag{6.69}
\end{equation*}
证明
\begin{equation*}
\langle\mathbf{q}|\hat{S}_{\mathbf{j}}^{x}\hat{S}_{\mathbf{j}'}^{x}+\hat{S}_{\mathbf{j}}^{y}\hat{S}_{\mathbf{j}'}^{y}|\mathbf{q}\rangle = \frac{2S}{N}\cos(\mathbf{q}\cdot(\mathbf{R}_{\mathbf{j}}-\mathbf{R}_{\mathbf{j}'}))\tag{6.70}
\end{equation*}
（$j\neq j'$），并参照图 6.13 解释这一结果。

\exer{6.12} 我们对布洛赫壁宽的计算假定：壁内 N 个自旋中相邻自旋夹角处处为
$\pi/N$。这显然是个近似！下面是改进的计算。

设自旋作为垂直于壁面的坐标 $z$ 的函数由角 $\theta(z)$ 描述：$z\to-\infty$ 时
$\theta(z)\to0$，$z\to+\infty$ 时 $\theta(z)\to\pi$（对应 180° 壁）。能量为
\begin{equation*}
E = \int_{-\infty}^{\infty}\left[A\left(\frac{\partial\theta}{\partial z}\right)^{2} + f(\theta)\right]\mathrm{d}z.\tag{6.71}
\end{equation*}
A 是连续介质交换常数，$f(\theta)$ 描述各向异性。证明能量在
\begin{equation*}
\frac{\partial f}{\partial\theta} - 2A\frac{\partial^{2}\theta}{\partial z^{2}} = 0.\tag{6.72}
\end{equation*}
时极小。把该方程乘以 $\partial\theta/\partial z$ 并对 z 积分，证明
\begin{equation*}
f(\theta) = A\left(\frac{\partial\theta}{\partial z}\right)^{2}.\tag{6.73}
\end{equation*}
现在设
\begin{equation*}
f(\theta) = K\sin^{2}\theta\tag{6.74}
\end{equation*}
（K 为各向异性常数）。证明这意味着
\begin{equation*}
\cos\theta(z) = -\tanh\left(\frac{\pi z}{\delta}\right),\tag{6.75}
\end{equation*}
其中 $\delta=\pi\sqrt{A/K}$。画出 $\theta(z)$ 的草图并论证 $\delta$ 是畴壁宽度
的合适度量。再证明此情形下
\begin{equation*}
E = 4\sqrt{AK}\tag{6.76}
\end{equation*}

\exer{6.13} 考虑铁磁薄膜中的封闭畴结构（见图 6.30）：易轴竖直、难轴水平。长 180°
壁单位面积的能量为 $\sigma_{\mathrm{W}}$，各向异性能量密度为 K。忽略 90° 壁的
贡献（即设 $L\gg D$），证明
\begin{equation*}
D \approx \sqrt{\frac{2\sigma_{\mathrm{W}}L}{K}}.\tag{6.77}
\end{equation*}
取 $\sigma_{\mathrm{W}}=2\times10^{-3}$ J\,m$^{-2}$、$K=4\times10^{4}$ J\,m$^{-3}$、
$L=4$ mm，估算 D。

\begin{figure}
\centering
\includegraphics[width=0.4\textwidth]{fig6_30.pdf}
\caption{一个封闭畴。}
\end{figure}

% 第7章 金属中的磁性（Magnetism in metals）
\chapter{金属中的磁性}

本章讨论金属的磁性质。前几章我们集中于有相互作用但局域的磁矩。金属中的传导电子是
离域的，可以在样品中自由游走——它们称为巡游电子（itinerant electrons）。有些
情形金属中的磁矩与传导电子相联系；另一些情形磁矩保持局域。两种情形下都可以出现
顺磁与抗磁行为；一定条件下铁磁也可能。本章的大部分讨论将围绕自由电子模型展开
（下一节引入）。自由电子模型对多数真实情形是粗糙的近似，但便于处理，能使讨论
走得很远。后续各节推导电子气磁性质：包括泡利顺磁性、朗道抗磁性、RKKY 相互作用
的起源、电子气的不稳定性（如自旋密度波的形成），以及局域磁矩与电子气相互作用
时发生的近藤效应。

\section{自由电子模型}

通过回顾自由电子模型开始巡游电子磁性的讨论。此模型忽略晶格的周期势，电子填充
直到费米波矢 $k_{\mathrm{F}}$ 的各个态。k 空间中相邻点相距 $2\pi/L$
（见图 7.1(a)），$V=L^{3}$ 是样品体积，故 k 与 $k+\mathrm{d}k$ 之间的态数等于
$4\pi k^{2}\mathrm{d}k$（半径 k、宽 dk 的球壳体积，见图 7.1(b)）除以 k 空间中
一个点占据的体积 $(2\pi/L)^{3}$。每个态被双占据——一个自旋向上电子与一个自旋
向下电子——故再有因子 2。于是态密度 $g(k)\mathrm{d}k$ 可写为
\begin{equation*}
g(k)\,\mathrm{d}k = \frac{2}{(2\pi/L)^{3}}\times4\pi k^{2}\,\mathrm{d}k,\tag{7.1}
\end{equation*}
其中因子 2 计入电子的两个自旋态。于是
\begin{equation*}
g(k)\,\mathrm{d}k = \frac{Vk^{2}\,\mathrm{d}k}{\pi^{2}}.\tag{7.2}
\end{equation*}
若材料有 N 个电子，绝对零度（T=0）下它们将填满直到最大波矢 $k_{\mathrm{F}}$ 的态：
\begin{equation*}
N = \int_{0}^{k_{\mathrm{F}}}g(k)\,\mathrm{d}k = \frac{Vk_{F}^{3}}{3\pi^{2}},\tag{7.3}
\end{equation*}
故
\begin{equation*}
k_{\mathrm{F}}^{3} = 3\pi^{2}n\tag{7.4}
\end{equation*}
$n=N/V$ 是单位体积电子数。费米能 $E_{\mathrm{F}}$ 定义为
\begin{equation*}
E_{\mathrm{F}} = \frac{\hbar^{2}k_{\mathrm{F}}^{2}}{2m_{\mathrm{e}}}.\tag{7.5}
\end{equation*}
态密度作为能量（$E\propto k^{2}$）的函数正比于 $E^{1/2}$：
\begin{equation*}
g(E) = \frac{\mathrm{d}n}{\mathrm{d}E} \propto E^{1/2}\tag{7.6}
\end{equation*}
于是
\begin{equation*}
n = \int_{0}^{E_{\mathrm{F}}}g(E)\,\mathrm{d}E \propto E_{\mathrm{F}}^{3/2}\tag{7.7}
\end{equation*}
从而
\begin{equation*}
\frac{\mathrm{d}n}{n} = \frac{3}{2}\frac{\mathrm{d}E_{\mathrm{F}}}{E_{\mathrm{F}}}.\tag{7.8}
\end{equation*}
费米能处的态密度为\fn{1}{式 7.9 也可以直接用式 7.2、7.4、7.5 导出。}
\begin{equation*}
g(E_{\mathrm{F}}) = \left(\frac{\mathrm{d}n}{\mathrm{d}E}\right)_{E=E_{\mathrm{F}}} = \frac{3}{2}\frac{n}{E_{\mathrm{F}}}.\tag{7.9}
\end{equation*}
顺带指出：联立式 7.4、7.5 与 7.9 可得 $g(E_{\mathrm{F}})$ 的另一个有用表达式
\begin{equation*}
g(E_{\mathrm{F}}) = \frac{m_{\mathrm{e}}k_{\mathrm{F}}}{(\pi\hbar)^{2}},\tag{7.10}
\end{equation*}
它表明 $g(E_{\mathrm{F}})\propto m_{\mathrm{e}}$。许多重要性质依赖费米能处的态密度，
因此知道它正比于电子质量是有用的。在许多感兴趣的体系中，由于能带结构或相互作用的
效应，电子质量相对其自由空间值被增强。

自由电子模型忽略晶格周期势。若把它作为微扰计入（近自由电子模型），结果是：除了
电子波矢接近倒格矢之处，它几乎没有影响；在这些 k 空间点，色散关系中出现能隙
（见图 7.2）。

\begin{figure}
\centering
\includegraphics[width=0.4\textwidth]{fig7_01.pdf}
\caption{（a）电子态在 k 空间相距 $2\pi/L$；每个态可双占据、占据体积 $(2\pi/L)^{3}$。
（b）态密度的计算：考虑 k 空间中波矢 k 与 $k+\mathrm{d}k$ 的态之间的体积，即
$4\pi k^{2}\mathrm{d}k$。}
\end{figure}

\begin{figure}
\centering
\includegraphics[width=0.45\textwidth]{fig7_02.pdf}
\caption{布里渊区边界处的能隙。}
\end{figure}

迄今一切都按 T=0 处理。$T>0$ 时，态密度 $g(E)$ 不变，但每个态的占据由费米函数
$f(E)$ 决定：
\begin{equation*}
f(E) = \frac{1}{\mathrm{e}^{(E-\mu)/k_{\mathrm{B}}T}+1},\tag{7.11}
\end{equation*}
$\mu$ 是依赖温度的化学势。此函数绘于图 7.3。T=0 时 $f(E)$ 是阶跃函数：$E<\mu$
取 1、$E>\mu$ 取 0；温度升高时台阶被抹平。当费米函数接近阶跃函数——多数金属在熔点
以下几乎所有温度的通常情形——称电子处于简并极限（degenerate limit）。费米函数
是泡利不相容原理的结果：每个电子必须有一组唯一的量子数，没有两个电子能处于同一
态。$(E-\mu)\gg k_{\mathrm{B}}T$ 时费米函数趋于麦克斯韦--玻尔兹曼形式
$\mathrm{e}^{-(E-\mu)/k_{\mathrm{B}}T}$，称为非简并极限。

费米能是 T=0 时最高被占据能级的能量，由方程
\begin{equation*}
\int_{0}^{E_{\mathrm{F}}}f(E)g(E)\,\mathrm{d}E = n.\tag{7.12}
\end{equation*}
确定。函数 $f(E)g(E)$ 示于图 7.3(c)。T=0 时容易得到：费米能精确等于化学势
$E_{\mathrm{F}}=\mu$。$T>0$ 时经冗长的计算得
\begin{equation*}
\mu = E_{\mathrm{F}}\left[1 - \frac{\pi^{2}}{12}\left(\frac{k_{\mathrm{B}}T}{E_{\mathrm{F}}}\right)^{2} + O\left(\left(\frac{k_{\mathrm{B}}T}{E_{\mathrm{F}}}\right)^{4}\right)\right]\tag{7.13}
\end{equation*}
但这意味着：对典型金属，即使在室温，把 $E_{\mathrm{F}}$ 与 $\mu$ 等同起来也精确到
0.01\%——尽管值得记住这两个量并不相同。费米面是 k 空间中能量等于化学势的点集。
若化学势位于能隙内，材料是半导体或绝缘体，没有费米面。因此，金属就是有费米面的
材料。

\begin{figure}
\centering
\includegraphics[width=0.32\textwidth]{fig7_03.pdf}
\caption{（a）式 7.11 定义的费米函数 $f(E)$：粗线为 $T=0$；温度升高时台阶被抹平。
所标温度为 $T=0$、$T=0.01\mu/k_{\mathrm{B}}$、$T=0.05\mu/k_{\mathrm{B}}$ 与
$T=0.1\mu/k_{\mathrm{B}}$。（b）自由电子气的态密度 $g(E)$ 正比于 $E^{1/2}$。
（c）与（a）相同温度下的 $f(E)g(E)$。}
\end{figure}

\section{泡利顺磁性}

金属中每个 k 态因电子有两个可能的自旋态而可双占据，因此金属中每个电子要么自旋向上、要么自旋向下。
加磁场时，电子的能量随其自旋而升降。这给出电子气的顺磁磁化率，称为泡利顺磁性
（Pauli paramagnetism）。

\subsection{初等推导}

起初我们忽略轨道贡献、取 g=2，并忽略有限温度造成的费米面模糊。如图 7.4 所示，外加
磁场中电子能带按自旋分裂成两个子带，间隔 $g\mu_{\mathrm{B}}B=2\mu_{\mathrm{B}}B$。
设 $g\mu_{\mathrm{B}}B$ 是很小的能量，能带分裂非常小。自旋向上的额外电子数密度为
$n_{\uparrow}=\frac{1}{2}g(E_{\mathrm{F}})\mu_{\mathrm{B}}B$；自旋向下的亏损
数密度亦然：$n_{\downarrow}=\frac{1}{2}g(E_{\mathrm{F}})\mu_{\mathrm{B}}B$。于是
磁化强度为
\begin{equation*}
M = \mu_{\mathrm{B}}(n_{\uparrow}-n_{\downarrow}) = g(E_{\mathrm{F}})\mu_{\mathrm{B}}^{2}B\tag{7.14}
\end{equation*}
磁化率 $\chi_{\mathrm{P}}$（下标 P 表示泡利磁化率）为
\begin{align*}
\chi_{\mathrm{P}} = \frac{M}{H}\approx\frac{\mu_0M}{B} &= \mu_0\mu_{\mathrm{B}}^{2}g(E_{\mathrm{F}})\tag{7.15}\\
&= \frac{3n\mu_0\mu_{\mathrm{B}}^{2}}{2E_{\mathrm{F}}}\tag{7.16}
\end{align*}
最后一步用式 7.9。由于 $\chi_{\mathrm{P}}\ll1$，写 $\chi_{\mathrm{P}}\approx\mu_0M/B$
是合理的（见 1.1.4 节）。

\begin{figure}
\centering
\includegraphics[width=0.45\textwidth]{fig7_04.pdf}
\caption{态密度图显示磁场 B 中能带的分裂。分裂被大大夸大。}
\end{figure}

我们的泡利顺磁性表达式不依赖温度——诚然，这是因为我们一开始就忽略了有限温度的
费米面模糊。不过若计入温度，它只是一个相当小的修正（见习题 7.1）。泡利顺磁性是弱
效应，在多数温度下远小于绝缘体中因居里定律出现的顺磁性。原因：顺磁绝缘体中每个
磁性原子上至少有一个电子作贡献；金属中只有靠近费米面的电子起作用。多数金属顺磁
磁化率之小曾是个谜，直到泡利指出它是电子服从费米--狄拉克统计（而非经典统计）的
后果。

\subsection{向局域行为的过渡}

费米--狄拉克统计的效应，以及泡利顺磁性与局域矩行为之间的过渡，可以通过重新推导
式 7.15 来说明。严格地，每种自旋态的单位体积电子数应写为
\begin{equation*}
\begin{aligned}
n_{\uparrow} &= \frac{1}{2}\int_{0}^{\infty}g(E+\mu_{\mathrm{B}}B)f(E)\,\mathrm{d}E\\
n_{\downarrow} &= \frac{1}{2}\int_{0}^{\infty}g(E-\mu_{\mathrm{B}}B)f(E)\,\mathrm{d}E.
\end{aligned}\tag{7.17}
\end{equation*}
于是小 B 时磁化强度 $M=\mu_{\mathrm{B}}(n_{\uparrow}-n_{\downarrow})$，从而
\begin{align*}
M &\approx \mu_{\mathrm{B}}^{2}B\int_{0}^{\infty}\frac{\mathrm{d}g}{\mathrm{d}E}f(E)\,\mathrm{d}E\\
&= \mu_{\mathrm{B}}^{2}B\left(\left[g(E)f(E)\right]_{0}^{\infty} - \int_{0}^{\infty}\frac{\mathrm{d}f}{\mathrm{d}E}g(E)\,\mathrm{d}E\right)\tag{7.18}
\end{align*}
第二行由第一行分部积分得到。利用 $g(0)=0$ 与 $f(\infty)=0$（见图 7.3），第一项为
零。故
\begin{equation*}
M \approx \mu_{\mathrm{B}}^{2}B\int_{0}^{\infty}\left(-\frac{\mathrm{d}f}{\mathrm{d}E}\right)g(E)\,\mathrm{d}E.\tag{7.19}
\end{equation*}
T=0 的简并极限下，$-\mathrm{d}f/\mathrm{d}E$（阶跃函数的导数）是 $E_{\mathrm{F}}$
处的 delta 函数：
\begin{equation*}
-\frac{\mathrm{d}f}{\mathrm{d}E} = \delta(E-E_{\mathrm{F}}),\tag{7.20}
\end{equation*}
于是回到 $M=\mu_{\mathrm{B}}^{2}Bg(E_{\mathrm{F}})$ 与 $\chi=\mu_0\mu_{\mathrm{B}}^{2}g(E_{\mathrm{F}})$，
与式 7.15 一致。

非简并极限下 $f(E)\approx\mathrm{e}^{-(E-\mu)/k_{\mathrm{B}}T}$，故
\begin{equation*}
-\frac{\mathrm{d}f}{\mathrm{d}E} = \frac{f}{k_{\mathrm{B}}T}\tag{7.21}
\end{equation*}
磁化强度为
\begin{equation*}
M = \frac{\mu_{\mathrm{B}}^{2}B}{k_{\mathrm{B}}T}\int_{0}^{\infty}f(E)g(E)\,\mathrm{d}E = \frac{n\mu_{\mathrm{B}}^{2}B}{k_{\mathrm{B}}T},\tag{7.22}
\end{equation*}
从而
\begin{equation*}
\chi = n\mu_0\mu_{\mathrm{B}}^{2}/k_{\mathrm{B}}T,\tag{7.23}
\end{equation*}
与式 2.28 一致。磁化率等价于单位体积 n 个局域矩（大小 $\mu_{\mathrm{B}}$）。如上
所述，多数金属 $E_{\mathrm{F}}$ 为几 eV，熔点以下一切温度简并极限都成立。然而对
载流子浓度低的材料（如掺杂半导体），$E_{\mathrm{F}}\propto n^{2/3}$ 小得多，可能
达到非简并极限，得到类居里的磁化率。

\subsection{实验技术}

我们导出的金属泡利顺磁性表达式与实验符合尚可，但可以通过修正电子--电子相互作用的
效应加以改进。金属的自旋磁化率可以从 NMR 测量提取：NMR 对传导电子自旋磁矩产生的
场远比对电子轨道运动产生的场（它导致 7.6 节考虑的抗磁效应）敏感。

传导电子自旋与核自旋之间的接触相互作用使核共振频率 $\omega$ 产生小移动
$\Delta\omega$，称为奈特移位（Knight shift）。可以这样理解：设想个别传导电子
在某个核上进进出出；核感受到的净超精细耦合是对所有电子自旋取向的平均。无外场时
电子自旋取向平均为零，净超精细耦合为零；非零静场使电子自旋极化，净超精细耦合
非零。奈特移位 $K=\Delta\omega/\omega$ 因此既正比于核处的传导电子密度（体现对耦合
强度的依赖），又正比于泡利自旋磁化率（体现外场使电子极化的程度）。

超精细相互作用的静态平均造成奈特移位；围绕平均的涨落提供 $T_1$ 弛豫机制（称为
Korringa 弛豫）。主导的 $T_1$ 过程是电子与核（或 $\mu$ 子）自旋的翻转--翻转
（flip--flop）跃迁，其中电子与核塞曼能量之差由传导电子动能的变化补偿。核与传导
电子之间的能量交换非常小，只有费米面附近 $k_{\mathrm{B}}T$ 范围内的电子能够参与——
只有它们附近有空态可供跃迁。因此对简单金属，自旋--晶格弛豫速率 $T_1^{-1}$ 正比于
温度。奈特移位（常以无量纲形式 $\Delta\omega/\omega$ 表示）与 Korringa 弛豫速率
$T_1^{-1}$ 通常由方程
\begin{equation*}
\frac{1}{T_1} \propto T\left(\frac{\Delta\omega}{\omega}\right)^{2},\tag{7.24}
\end{equation*}
联系，称为 Korringa 关系。

\section{自发自旋劈裂能带}

铁中每原子磁矩约 $2.2\mu_{\mathrm{B}}$（见表 5.1）。这个非整数值无法在原子局域矩
的基础上理解——它是能带铁磁性（band ferromagnetism，又称巡游铁磁性）的有力
证据：磁化来源于自发自旋劈裂的能带。本节探讨几个可用于理解某些材料中能带如何
自发自旋劈裂的模型。

\begin{figure}
\centering
\includegraphics[width=0.45\textwidth]{fig7_05.pdf}
\caption{态密度图显示无外加磁场时能带的自发分裂。}
\end{figure}
\bio{Edmund C. Stoner\\（1899--1968）}

在分子场理论中，我们说所有自旋“感受”邻居产生的同一平均交换场 $\lambda M$。
金属中，分子场可以经由泡利顺磁性 $\chi_{\mathrm{P}}$ 磁化电子气；电子气由此产生的
磁化强度 M 又反过来充当分子场。这是“鸡生蛋蛋生鸡”的情形（也称 bootstrap）：
这一正反馈机制能导致自发铁磁性吗？大概能——只要 $\lambda$（给定 M 能得到多少
分子场）与 $\chi_{\mathrm{P}}$（给定分子场能得到多少磁化强度）都足够大。

最好把上述启发式论证做得更严格些！要问的问题是：系统整体通过变为铁磁体能省能量
吗？

先设想：无外场时，把费米面处少量电子从自旋向下带搬到自旋向上带。具体地：取能量从
$E_{\mathrm{F}}-\delta E$ 到 $E_{\mathrm{F}}$ 的自旋向下电子，翻转其自旋放入自旋
向上带，能量从 $E_{\mathrm{F}}$ 到 $E_{\mathrm{F}}+\delta E$（图 7.5）。被搬电子
数为 $g(E_{\mathrm{F}})\delta E/2$，每个能量升高 $\delta E$，总能量变化
$g(E_{\mathrm{F}})\delta E/2\times\delta E$。总动能变化 $\Delta E_{\mathrm{K.E.}}$
为
\begin{equation*}
\Delta E_{\mathrm{K.E.}} = \frac{1}{2}g(E_{\mathrm{F}})(\delta E)^{2}.\tag{7.25}
\end{equation*}
这是能量代价，过程看似不利。然而磁化强度与分子场的相互作用给出可以压倒这一代价
的能量降低。自旋向上的数密度为 $n_{\uparrow}=\frac{1}{2}(n+g(E_{\mathrm{F}})\delta E)$，
向下的为 $n_{\downarrow}=\frac{1}{2}(n-g(E_{\mathrm{F}})\delta E)$。设每个电子磁矩
$1\mu_{\mathrm{B}}$，则磁化强度 $M=\mu_{\mathrm{B}}(n_{\uparrow}-n_{\downarrow})$。
分子场能量为
\begin{equation*}
\Delta E_{\mathrm{P.E.}} = -\int_{0}^{M}\mu_0(\lambda M')\,\mathrm{d}M' = -\frac{1}{2}\mu_0\lambda M^{2} = -\frac{1}{2}\mu_0\mu_{\mathrm{B}}^{2}\lambda(n_{\uparrow}-n_{\downarrow})^{2}.\tag{7.26}
\end{equation*}
写 $U=\mu_0\mu_{\mathrm{B}}^{2}\lambda$（U 量度库仑能\fn{2}{分子场源于交换，交换源于库仑相互作用。更高级的处理中库仑能直接以 $Un_{\uparrow}n_{\downarrow}$ 计入，给出同样结果。}），有
\begin{equation*}
\Delta E_{\mathrm{P.E.}} = -\frac{1}{2}U(g(E_{\mathrm{F}})\delta E)^{2}.\tag{7.27}
\end{equation*}
故总能变化 $\Delta E$ 为
\begin{equation*}
\Delta E = \Delta E_{\mathrm{K.E.}}+\Delta E_{\mathrm{P.E.}} = \frac{1}{2}g(E_{\mathrm{F}})(\delta E)^{2}(1-Ug(E_{\mathrm{F}})).\tag{7.28}
\end{equation*}
$\Delta E<0$ 时自发铁磁性可能，即
\begin{equation*}
Ug(E_{\mathrm{F}}) \geq 1\tag{7.29}
\end{equation*}
称为 Stoner 判据（Stoner criterion）。铁磁不稳定性的这一条件要求库仑效应强、且
费米能处态密度大。若有自发铁磁性，无外场时自旋向上、向下带将劈裂能量 $\Delta$
——即交换劈裂。

若 Stoner 判据不满足，自发铁磁性不会发生，但磁化率可能被改变。把外场与相互作用
的效应都计入即可轻松算出：能量移动 $\delta E$ 产生的磁化强度就是
$M=\mu_{\mathrm{B}}(N_{\uparrow}-N_{\downarrow})=2\mu_{\mathrm{B}}g(E_{\mathrm{F}})\delta E$。于是
\begin{align*}
\Delta E &= \frac{1}{2}g(E_{\mathrm{F}})(\delta E)^{2}(1-Ug(E_{\mathrm{F}})) - MB\tag{7.30}\\
&= \frac{M^{2}}{2\mu_{\mathrm{B}}^{2}g(E_{\mathrm{F}})}(1-Ug(E_{\mathrm{F}})) - MB\tag{7.31}
\end{align*}
它在
\begin{equation*}
\frac{M}{\mu_{\mathrm{B}}^{2}g(E_{\mathrm{F}})}(1-Ug(E_{\mathrm{F}})) - B = 0\tag{7.32}
\end{equation*}
时极小，故磁化率 $\chi$ 为
\begin{equation*}
\chi = \frac{M}{H}\approx\frac{\mu_0M}{B} = \frac{\mu_0\mu_{\mathrm{B}}^{2}g(E_{\mathrm{F}})}{1-Ug(E_{\mathrm{F}})} = \frac{\chi_{\mathrm{P}}}{1-Ug(E_{\mathrm{F}})}.\tag{7.33}
\end{equation*}
这比无库仑相互作用的 $\chi_{\mathrm{P}}$ 大一个因子 $(1-Ug(E_{\mathrm{F}}))^{-1}$，
称为 Stoner 增强（Stoner enhancement）。它负责金属 Pd 与 Pt 中测得的增强泡利
磁化率——二者都可视为处于铁磁性边缘的体系：参数 $Ug(E_{\mathrm{F}})$ 足够大使
磁化率显著增强，但还不够大（不够接近 1）以致自发铁磁。

\section{自旋密度泛函理论}

迄今的讨论用自由电子或近自由电子模型。可以用更高级的方法改进，本节讨论其中之一。
真实体系中不能忽略电子间的库仑相互作用以及交换相互作用对电子运动的影响。所有
粒子的位置与运动相互关联——粒子彼此相互作用、运动中互相施力；相互作用导致粒子
之间出现关联。这类关联理论上很难处理，但一种有用而成功的途径是密度泛函理论
（density functional theory）。

该理论认识到：多电子系统的基态能量可以写成电子密度 $n(\mathbf{r})$ 的泛函
（functional）。\fn{3}{函数（function）是把一个数映射为另一个数的规则。例如函数 $f(x)=x^{2}$ 把数 2 映为数 4。泛函（functional）是把整个函数映射为一个数的规则。例如泛函 $F[f]=\int_{-1}^{1}f(x)\,\mathrm{d}x$ 把函数 $f(x)=x^{2}$ 映为数 2/3。}泛函含三部分贡献：动能；系统中带电粒子静电相互作用的库仑能；
以及称为交换关联能（exchange-correlation energy）的项——它囊括所有多体相互
作用。密度泛函理论无须处理波函数 $\psi(\mathbf{r})$，只须考虑电子密度
$n(\mathbf{r})=|\psi(\mathbf{r})|^{2}$，带来相当大的简化。能量泛函取极小导出一个
可用于求基态能量的方程。

已经证明：极小化泛函确实给出精确的基态能量——尽管用的是电子密度而非波函数。

问题在于能量泛函并非完全已知！动能与库仑能写得出来，交换关联能却是未知的。然而，
可以采用已知的均匀电子气（密度恒定的电子系统）中多电子相互作用的结果作为交换
关联能的取值。这称为局域密度近似（local density approximation，LDA）：在电子
密度为 $n(\mathbf{r})$ 的每个空间点 r，电子对多体相互作用的响应，就好像全空间
（即位置 $\mathbf{r}'\neq\mathbf{r}$ 处）的电子密度都取 $n(\mathbf{r})$。空间不同
位置对能量的贡献全部加起来（对所有体积元积分）；各体积元的贡献不同，依赖于局域
电子密度。LDA 对完美金属（电子密度确实均匀）是精确的，但对电子密度剧烈空间变化
的体系则不那么好。\fn{4}{该理论对金属很有效，对局域体系则不然。原因：库仑能代表某个特定电子对电子密度的响应，而密度来自包括该电子在内的所有电子——电子自相互作用不可避免地、然而错误地被计入。对金属这不算太糟——任一电子只是浩瀚传导电子海洋中的一员；对局域体系则是灾难——特定态可能只被一个电子占据。}

把密度泛函理论推广到包括自旋极化效应，称为自旋密度泛函理论（spin-density
functional theory）。磁性体系中它与局域自旋密度近似（local spin-density
approximation，常称 LSDA）联用：交换关联势不仅依赖局域电子密度，还依赖局域自旋
密度（自旋向上与向下电子密度之差）。这一技术可用于对能带结构做现实计算，获得
真实体系自旋密度的定量信息。不过下一节我们退回到相对安全的简单自由电子模型！

\section{朗道能级}

泡利顺磁性是与电子自旋相联系的效应。电子的轨道贡献呢？用如下论证来评估。设自由
电子哈密顿量为
\begin{equation*}
\frac{\hat{\mathbf{p}}^{2}}{2m_{\mathrm{e}}}\psi = -\frac{\hbar^{2}}{2m_{\mathrm{e}}}\nabla^{2}\psi = E\psi,\tag{7.34}
\end{equation*}
波函数为平面波
\begin{equation*}
\psi = \frac{1}{\sqrt{V}}e^{i\mathbf{k}\cdot\mathbf{r}},\tag{7.35}
\end{equation*}
故能量 E 为
\begin{equation*}
E = \frac{\hbar^{2}k_x^{2}}{2m_{\mathrm{e}}} + \frac{\hbar^{2}k_y^{2}}{2m_{\mathrm{e}}} + \frac{\hbar^{2}k_z^{2}}{2m_{\mathrm{e}}}.\tag{7.36}
\end{equation*}
与前面一样，无磁场时电子态在 k 空间均匀分布、相邻点相距 $2\pi/L$（见图 7.6(a)）。

存在磁场 $\mathbf{B}=(0,0,B)$ 时，动量算符 $\hat{\mathbf{p}}=-\mathrm{i}\hbar\nabla$
须换成 $-\mathrm{i}\hbar\nabla+e\mathbf{A}$，磁矢势的有用选取是
$\mathbf{A}=(0,Bx,0)$。于是式 7.34 变为
\begin{equation*}
\left(\frac{\hat{p}_x^{2}}{2m_{\mathrm{e}}} + \frac{(\hat{p}_y+eB\hat{x})^{2}}{2m_{\mathrm{e}}} + \frac{\hat{p}_z^{2}}{2m_{\mathrm{e}}}\right)\psi = E\psi.\tag{7.37}
\end{equation*}
由于表达式含算符 $\hat{x}$、$\hat{p}_x$、$\hat{p}_y$、$\hat{p}_z$，波函数在 $y$
与 $z$ 方向仍是平面波：$\hat{p}_y\psi=\hbar k_y\psi$、$\hat{p}_z\psi=\hbar k_z\psi$。
于是式 7.37 可整理（用分离变量）为
\begin{equation*}
\left(\left[\frac{\hat{p}_x^{2}}{2m_{\mathrm{e}}} + \frac{1}{2}m_{\mathrm{e}}\omega_{\mathrm{c}}^{2}(x-x_0)^{2}\right] + \frac{\hbar^{2}k_z^{2}}{2m_{\mathrm{e}}}\right)\psi = E\psi,\tag{7.38}
\end{equation*}
$x_0=-\hbar k_y/eB$，且
\begin{equation*}
\omega_{\mathrm{c}} = \frac{eB}{m_{\mathrm{e}}}\tag{7.39}
\end{equation*}
是回旋频率。式 7.38 方括号内是一维谐振子的哈密顿量，故能量本征值对应
\begin{equation*}
E = \left(l+\frac{1}{2}\right)\hbar\omega_{\mathrm{c}} + \frac{\hbar^{2}k_z^{2}}{2m_{\mathrm{e}}}\tag{7.40}
\end{equation*}
l 为整数。能量本征函数是 y、z 方向平面波与 x 方向一维谐振子波函数之积。

在本处理中，$x_0$ 是 x 方向谐振动的“中心”，随 $k_y$ 取不同值。\fn{5}{若选 $\mathbf{A}=(-By,0,0)$，得到的参数 $y_0$ 正比于 $k_x$，是 y 方向谐振动的中心。可以证明对应 $x_0$ 与 $y_0$ 的算符不对易，故 $x_0$ 与 $y_0$ 不能同时取确定值。}处理这一问题的另一个有用方法是选取反映此情形固有的柱对称性的规范
（习题 7.2 这样做），得到与式 7.40 相同的能量本征值，但本征函数是 z 方向平面波与
$\sqrt{k_x^{2}+k_y^{2}}$ 的函数之积。无论用哪种规范求解，允许的态都可以看作位于
平行于 $k_z$ 的朗道管（Landau tubes）上（见图 7.7）。磁场的施加使电子形成朗道
能级（Landau levels，把态按能量看时即如此），每个朗道能级可用量子数 l 标记
（见图 7.6(b)）。朗道能级的简并度可以这样算：k 空间中相邻两个朗道能级之间的
面积（见图 7.6(b)）为
\begin{equation*}
\pi k_{l+1}^{2}-\pi k_l^{2} = \frac{2m_{\mathrm{e}}\pi}{\hbar^{2}}\left(\left(l+1+\frac{1}{2}\right)\hbar\omega_{\mathrm{c}}-\left(l+\frac{1}{2}\right)\hbar\omega_{\mathrm{c}}\right) = \frac{2m_{\mathrm{e}}\pi\omega_{\mathrm{c}}}{\hbar}\tag{7.41}
\end{equation*}
由于每个态被两个电子占据、一个态在 $k_xk_y$ 平面占面积 $(2\pi/L)^{2}$，一个朗道
能级的简并度 p 为
\begin{equation*}
p = \frac{4m_{\mathrm{e}}\pi\omega_{\mathrm{c}}}{\hbar(2\pi/L)^{2}} = \frac{m_{\mathrm{e}}L^{2}\omega_{\mathrm{c}}}{\pi\hbar}\tag{7.42}
\end{equation*}
\mnote{注意：若磁矢势选为 $\mathbf{A}=(-By,0,0)$（它给出同样的
$\mathbf{B}=\nabla\times\mathbf{A}=(0,0,B)$），会得到同样的能量本征值，但本征函数
在此情形将是 x、z 方向平面波与 y 方向一维谐振子波函数之积。“规范的选择”（即
写下 $\mathbf{A}$ 的方式）不影响能量本征值，但影响导出的波函数集合。}

\begin{figure}
\centering
\includegraphics[width=0.4\textwidth]{fig7_06.pdf}
\caption{（a）B=0 时电子态在 k 空间均匀分布，相邻点相距 $2\pi/L$。（b）磁场的施加
使允许能量只剩分立值。这些朗道能级态具有简并度——由波矢 $k_{l+1}$ 与 $k_l$ 能级
之间 k 空间的面积算得。}
\end{figure}

\begin{figure}
\centering
\includegraphics[width=0.45\textwidth]{fig7_07.pdf}
\caption{朗道管。}
\end{figure}

\section{朗道抗磁性}

用上一节的结果，现在可以计算电子气的抗磁响应。想法是：加磁场时电子分布分裂成一
系列朗道能级，系统总能量可能变化；能量随场的变化等价于系统的磁化强度。这一现象
称为朗道抗磁性（Landau diamagnetism）。

计算基于如下观察：不同 $k_z$ 的最高被占据朗道能级不同（见图 7.7——费米面在不同
$k_z$ 处切割不同的朗道管）。对特定 $k_z$，定义与垂直于场的电子运动相联系的能量
$E_{\perp}$：
\begin{equation*}
E_{\perp} = E_{\mathrm{F}} - \frac{\hbar^{2}k_z^{2}}{2m_{\mathrm{e}}}.\tag{7.43}
\end{equation*}
无磁场时，对特定 $k_z$，电子可以填满直到最大值 $E_{\perp}$ 的态；加磁场则允许分立
的朗道能级。把最高被占据朗道能级标记为 $E_l$。B=0 问题中的电子可以分解成若干
“块”，每块的平均能量等于相应的朗道能级能量（顶块例外——部分填充）。

定义参数 x：
\begin{equation*}
x = E_{\perp} - E_l\tag{7.44}
\end{equation*}
（见图 7.8）。最高被占据块的平均能量为
\begin{equation*}
\frac{1}{2}\left[E_{\perp}+E_l-\frac{\hbar\omega_{\mathrm{c}}}{2}\right] = E_{\perp}-\frac{x}{2}-\frac{\hbar\omega_{\mathrm{c}}}{4}\tag{7.45}
\end{equation*}

\begin{figure}
\centering
\includegraphics[width=0.45\textwidth]{fig7_08.pdf}
\caption{（a）B=0 时电子可以填满直到最大值 $E_{\perp}$（由式 7.43 定义）的态；
T=0 被填满的态以阴影示出。文中所述“最高被占据块”标记为 $E_l$；参数
$x=E_{\perp}-E_l$ 量度该块的中点（点划线）与最高被占据态之间的能量。（b）磁场只
允许分立的朗道能级，能量为 $\ldots E_{l-2},E_{l-1},E_l,E_{l+1}\ldots$；T=0 被占据
的朗道能级以粗线示出，最高被占据朗道能级标记为 $E_l$。}
\end{figure}

占据数为
\begin{equation*}
\frac{p}{\hbar\omega_{\mathrm{c}}}\left[E_{\perp}-\left(E_l-\frac{\hbar\omega_{\mathrm{c}}}{2}\right)\right] = \frac{p}{2}\left[1+\frac{2x}{\hbar\omega_{\mathrm{c}}}\right].\tag{7.46}
\end{equation*}
两者相乘给出顶块的能量。这是对特定 $k_z$ 求值的，但必须对费米面上所有电子（即所有
不同 $k_z$ 值）平均。这可以通过对 x 的一切可能值平均实现——设 x 在
$-\hbar\omega_{\mathrm{c}}/2$ 与 $\hbar\omega_{\mathrm{c}}/2$ 之间均匀分布：
\begin{equation*}
\langle x\rangle = \int_{-\hbar\omega_{\mathrm{c}}/2}^{\hbar\omega_{\mathrm{c}}/2}x\,\mathrm{d}x = 0,\tag{7.47}
\end{equation*}
及
\begin{equation*}
\langle x^{2}\rangle = \int_{-\hbar\omega_{\mathrm{c}}/2}^{\hbar\omega_{\mathrm{c}}/2}x^{2}\,\mathrm{d}x = \frac{\hbar^{2}\omega_{\mathrm{c}}^{2}}{12}.\tag{7.48}
\end{equation*}
由此得与垂直于场的运动相联系的顶块（标记 $E_l$ 的那块）平均能量为
\begin{equation*}
\frac{p}{2}\left[E_{\perp}-\frac{\hbar\omega_{\mathrm{c}}}{3}\right]\tag{7.49}
\end{equation*}
总平均能量为（用式 7.43）
\begin{equation*}
\frac{p}{2}\left[E_{\mathrm{F}}-\frac{\hbar\omega_{\mathrm{c}}}{3}\right].\tag{7.50}
\end{equation*}
磁场非零时：若 $0<x<\hbar\omega_{\mathrm{c}}/2$，顶块被 p 个能量 $E_l$ 的电子填满；
若 $-\hbar\omega_{\mathrm{c}}/2<x<0$，顶块全空。（注意此计算针对特定 $k_z$；态在
不同 $k_z$ 之间重新分配。）于是顶块的平均能量为
\begin{equation*}
E = \begin{cases} p(E_{\perp}-x) & \text{若 } x>0\\ 0 & \text{若 } x<0 \end{cases}\tag{7.51}
\end{equation*}
故与垂直于场的运动相联系的平均能量为
\begin{equation*}
\frac{p}{2}\left[E_{\perp}-\frac{\hbar\omega_{\mathrm{c}}}{4}\right]\tag{7.52}
\end{equation*}
总平均能量（再次用式 7.43）为
\begin{equation*}
\frac{p}{2}\left[E_{\mathrm{F}}-\frac{\hbar\omega_{\mathrm{c}}}{4}\right]\tag{7.53}
\end{equation*}
计入 $k_{\mathrm{F}}L/\pi$ 因子（$k_z$ 方向态的简并度），$B\neq0$ 与 B=0 的总能量
之差为
\begin{align*}
\Delta U &= U(B\neq0)-U(B=0)\tag{7.54}\\
&= \frac{pk_{\mathrm{F}}L}{2\pi}\left[E_{\mathrm{F}}-\frac{\hbar\omega_{\mathrm{c}}}{4}\right] - \frac{pk_{\mathrm{F}}L}{2\pi}\left[E_{\mathrm{F}}-\frac{\hbar\omega_{\mathrm{c}}}{3}\right]\tag{7.55}\\
&= \frac{pk_{\mathrm{F}}L\hbar\omega_{\mathrm{c}}}{24\pi}.\tag{7.56}
\end{align*}
用简并度 p 的表达式（式 7.42）：
\begin{equation*}
\Delta U = \frac{Vk_{\mathrm{F}}e^{2}B^{2}}{24\pi^{2}m_{\mathrm{e}}}.\tag{7.57}
\end{equation*}
故磁化率为
\begin{align*}
\chi_{\mathrm{L}} &= -\frac{\mu_0}{V}\frac{\partial^{2}\Delta U}{\partial B^{2}}\tag{7.58}\\
&= -\frac{\mu_0k_{\mathrm{F}}e^{2}}{12\pi^{2}m_{\mathrm{e}}}\tag{7.59}\\
&= -\frac{1}{3}\mu_0\left(\frac{e\hbar}{2m_{\mathrm{e}}}\right)^{2}\left[\frac{m_{\mathrm{e}}k_{\mathrm{F}}}{\pi^{2}\hbar^{2}}\right]\tag{7.60}\\
&= -\frac{1}{3}\mu_0\mu_{\mathrm{B}}^{2}g(E_{\mathrm{F}}).\tag{7.61}
\end{align*}
这就是朗道抗磁性，它与泡利顺磁性 $\chi_{\mathrm{P}}$ 的关系为
\begin{equation*}
\chi_{\mathrm{L}} = -\frac{\chi_{\mathrm{P}}}{3}.\tag{7.62}
\end{equation*}
由此或许会得出结论：一切金属都是顺磁的，因为（正的）泡利顺磁性是（负的）朗道
抗磁性的三倍。这未免轻率——我们用了自由电子模型、忽略了能带结构效应。若电子有
有效质量 $m^{*}$，费米能处态密度 $g(E_{\mathrm{F}})$ 将增强 $m^{*}/m_{\mathrm{e}}$
倍，泡利顺磁性与朗道抗磁性都增强 $m^{*}/m_{\mathrm{e}}$ 倍。但式 7.61 中最终结果
用了 $\mu_{\mathrm{B}}=e\hbar/2m_{\mathrm{e}}$；玻尔磁子是固定常数、不依赖有效
质量。这意味着朗道抗磁性还须再乘 $(m_{\mathrm{e}}/m^{*})^{2}$。净结果：
\begin{equation*}
\chi_{\mathrm{L}} = -\left(\frac{m_{\mathrm{e}}}{m^{*}}\right)^{2}\frac{\chi_{\mathrm{P}}}{3},\tag{7.63}
\end{equation*}
金属的总磁化率为
\begin{equation*}
\chi = \chi_{\mathrm{P}}\left[1-\frac{1}{3}\left(\frac{m_{\mathrm{e}}}{m^{*}}\right)^{2}\right].\tag{7.64}
\end{equation*}
多数金属 $m^{*}\sim m_{\mathrm{e}}$，故为顺磁。若 $m^{*}<m_{\mathrm{e}}/\sqrt{3}$，
净磁化率变负——这解释了铋的强抗磁性（$m^{*}\sim0.01m_{\mathrm{e}}$）。

迄今忽略的另一个效应是离子芯中束缚电子的抗磁性。原子序数增大时该效应变大（贡献
的电子更多），但仍通常远小于我们讨论的传导电子效应。

费米面上的电子运动由方程
\begin{equation*}
\hbar\dot{\mathbf{k}} = -e\mathbf{v}\times\mathbf{B},\tag{7.65}
\end{equation*}
描述，k 是费米面上的电子波矢，$\mathbf{v}=\hbar^{-1}\partial E(\mathbf{K})/\partial\mathbf{k}$
是电子速度，总是垂直于费米面。故电子运动垂直于 B 且保持在费米面上。\fn{6}{这是因为在时间 dt 内 k 的变化 $\mathrm{d}\mathbf{k}\propto(\partial E(\mathbf{K})/\partial\mathbf{k})\times\mathbf{B}$，既垂直于 B 又在费米面平面内。}磁化率的抗磁轨道贡献来自费米面上的闭合轨道。如图 7.9 所示，开放轨道也可能出现。
简单自由电子球中只可能有闭合轨道；真实材料中开、闭轨道皆可能。

\begin{figure}
\centering
\includegraphics[width=0.5\textwidth]{fig7_09.pdf}
\caption{（a）闭合轨道与（b）开放轨道的例子。轨道上的矢量表示电子的实空间速度
$\mathbf{v}=\hbar^{-1}\partial E(\mathbf{K})/\partial\mathbf{k}$——垂直于费米面的
矢量。磁场 B 如（a）取向时得到闭合轨道；B 沿 x 方向时，同一费米面上可以出现开放
轨道，如（b）所示。}
\end{figure}
\bio{P.~M. van Alphen（1906--1967）\\W.~J. de Haas（1878--1960）\\L.~V. Shubnikov（1901--1945）}

磁场增大、朗道能级突破费米面时，性质发生有趣的振荡变化。性质之所以变化，是因为
费米能处的态密度可以作为磁场的函数振荡。当朗道管扫过费米面的极值截面（对费米球
即赤道截面）时效应最大。态密度的振荡导致振荡磁化强度——这就是德哈斯--范阿尔芬
效应（de Haas--van Alphen effect），可用来测量费米面的截面积。电阻上也有类似
效应，称为舒布尼科夫--德哈斯效应（Shubnikov--de Haas effect）。两种效应都给出
性质随 1/B 的振荡。振荡周期按
\begin{equation*}
\Delta\left(\frac{1}{B}\right) = \frac{2\pi e}{\hbar A_{\mathrm{ext}}}.\tag{7.66}
\end{equation*}
给出垂直于磁场的费米面最大与最小截面积 $A_{\mathrm{ext}}$。

\section{电子气的磁性}

泡利顺磁性描述电子气对均匀外场的顺磁响应。若外场随空间变化呢？电子气将按波矢依赖
的磁化率响应这一微扰——现在就来计算（暂且忽略抗磁响应）。

\subsection{电子气的顺磁响应}

空间变化的磁场
\begin{equation*}
\mathbf{H}(\mathbf{r}) = \mathbf{H}_{\mathbf{q}}\cos\mathbf{q}\cdot\mathbf{r}\tag{7.67}
\end{equation*}
提供微扰 $\hat{\mathcal{H}}'$：
\begin{equation*}
\hat{\mathcal{H}}' = \pm\frac{g\mu_0\mu_{\mathrm{B}}}{2}|\mathbf{H}_{\mathbf{q}}|\cos\mathbf{q}\cdot\mathbf{r}\tag{7.68}
\end{equation*}
± 指电子自旋。平面波态
\begin{equation*}
\psi_{\mathbf{k}\pm}(\mathbf{r}) = \frac{1}{\sqrt{V}}e^{i\mathbf{k}\cdot\mathbf{r}}|\pm\rangle\tag{7.69}
\end{equation*}
被微弱地微扰到 $\psi_{(\mathbf{k}+\mathbf{q})\pm}(\mathbf{r})$ 与
$\psi_{(\mathbf{k}-\mathbf{q})\pm}(\mathbf{r})$ 态。这些态的幅度可用一阶微扰论计算
（见附录 C）。受扰波函数为
\begin{equation*}
\psi_{\mathbf{k}\pm}(\mathbf{r}) = \frac{1}{\sqrt{V}}\left(\mathrm{e}^{i\mathbf{k}\cdot\mathbf{r}} \pm \frac{g\mu_0\mu_{\mathrm{B}}H_{\mathbf{q}}}{4}\left[\frac{\mathrm{e}^{\mathrm{i}(\mathbf{k}+\mathbf{q})\cdot\mathbf{r}}}{E_{\mathbf{k}+\mathbf{q}}-E_{\mathbf{k}}} + \frac{\mathrm{e}^{\mathrm{i}(\mathbf{k}-\mathbf{q})\cdot\mathbf{r}}}{E_{\mathbf{k}-\mathbf{q}}-E_{\mathbf{k}}}\right]\right)|\pm\rangle\tag{7.70}
\end{equation*}
$E_k=\hbar^{2}k^{2}/2m_{\mathrm{e}}$。只保留 $H_{\mathbf{q}}$ 的领头阶：
\begin{equation*}
|\psi_{\mathbf{k}\pm}(\mathbf{r})|^{2} = \frac{1}{V}\left(1 \pm \frac{g\mu_0\mu_{\mathrm{B}}H_{\mathbf{q}}m_{\mathrm{e}}}{\hbar^{2}}\times\left[\frac{1}{(\mathbf{k}+\mathbf{q})^{2}-k^{2}} + \frac{1}{(\mathbf{k}-\mathbf{q})^{2}-k^{2}}\right]\cos\mathbf{q}\cdot\mathbf{r}\right).\tag{7.71}
\end{equation*}
受扰波函数导致的磁化强度 $M(\mathbf{r})$ 为
\begin{align*}
M(\mathbf{r}) &= \frac{g\mu_0\mu_{\mathrm{B}}}{2}(|\psi_{k+}(\mathbf{r})|^{2}-|\psi_{k-}(\mathbf{r})|^{2})\\
&= \frac{g^{2}\mu_0\mu_{\mathrm{B}}^{2}m_{\mathrm{e}}H_{\mathbf{q}}\cos\mathbf{q}\cdot\mathbf{r}}{\hbar^{2}V}\left[\frac{1}{(\mathbf{k}+\mathbf{q})^{2}-k^{2}} + \frac{1}{(\mathbf{k}-\mathbf{q})^{2}-k^{2}}\right]\tag{7.72}
\end{align*}
要计算电子气整体的磁响应，须把式 7.72 对费米球内所有电子求和。电子气的空间变化
磁化强度 $M(\mathbf{r})$ 成为
\begin{equation*}
M(\mathbf{r}) = M_{\mathbf{q}}\cos\mathbf{q}\cdot\mathbf{r}\tag{7.73}
\end{equation*}
$M_{\mathbf{q}}$ 为
\begin{align*}
M_{\mathbf{q}} &= \frac{g^{2}\mu_0\mu_{\mathrm{B}}^{2}m_{\mathrm{e}}H_{\mathbf{q}}}{\hbar^{2}V}\int_{|\mathbf{k}|<k_{\mathrm{F}}}g(\mathbf{k})\,\mathrm{d}^{3}\mathbf{k}\left[\frac{1}{(\mathbf{k}+\mathbf{q})^{2}-k^{2}} + \frac{1}{(\mathbf{k}-\mathbf{q})^{2}-k^{2}}\right]\tag{7.74}\\
&= \frac{k_{\mathrm{F}}m_{\mathrm{e}}g^{2}\mu_0\mu_{\mathrm{B}}^{2}H_{\mathbf{q}}}{\pi^{2}\hbar^{2}}\left[1+\frac{4k_{\mathrm{F}}^{2}-q^{2}}{4k_{\mathrm{F}}q}\log\left|\frac{q+2k_{\mathrm{F}}}{q-2k_{\mathrm{F}}}\right|\right]\tag{7.75}
\end{align*}
其中态密度 $g(|k|)=Vk^{2}/\pi^{2}$，并用到积分
\begin{align*}
\int_{|\mathbf{k}|<k_{\mathrm{F}}}\frac{\mathrm{d}^{3}\mathbf{k}}{(\mathbf{k}+\mathbf{q})^{2}-k^{2}} &= \int_{0}^{k_{\mathrm{F}}}2\pi k^{2}\mathrm{d}k\int_{0}^{\pi}\frac{\sin\theta\,\mathrm{d}\theta}{q(q-2k\cos\theta)}\\
&= \int_{0}^{k_{\mathrm{F}}}\frac{\pi k\,\mathrm{d}k}{q}\log\frac{q+2k}{|q-2k|}\\
&= \frac{\pi k_{\mathrm{F}}}{2}\left[1+\frac{4k_{\mathrm{F}}^{2}-q^{2}}{4k_{\mathrm{F}}q}\log\left|\frac{q+2k_{\mathrm{F}}}{q-2k_{\mathrm{F}}}\right|\right]\tag{7.76}
\end{align*}
用方程
\begin{equation*}
\chi_{\mathbf{q}} = M_{\mathbf{q}}/H_{\mathbf{q}}\tag{7.77}
\end{equation*}
（下面将证明其合理）与式 7.10 的 $g(E_{\mathrm{F}})$ 表达式，波矢依赖磁化率
$\chi_{\mathbf{q}}$ 为
\begin{equation*}
\chi_{\mathbf{q}} = \chi_{\mathrm{P}}f\left(\frac{q}{2k_{\mathrm{F}}}\right)\tag{7.78}
\end{equation*}
函数 $f(x)$ 为
\begin{equation*}
f(x) = \frac{1}{2}\left(1+\frac{1-x^{2}}{2x}\log\left|\frac{x+1}{x-1}\right|\right).\tag{7.79}
\end{equation*}
$\chi_{\mathbf{q}}$ 绘于图 7.10。q=0 时 $f(0)=1$，故 $\chi_{q=0}=\chi_{\mathrm{P}}$
——结果令人欣慰地退化为 7.2 节导出的泡利顺磁性。曲线在 $q=2k_{\mathrm{F}}$ 处的
扭折源于费米面的存在（费米球直径为 $2k_{\mathrm{F}}$）。随着维数从三维降低，这一
扭折逐渐加剧——本章后面讨论。

\begin{figure}
\centering
\includegraphics[width=0.5\textwidth]{fig7_10.pdf}
\caption{波矢依赖的顺磁磁化率。q=0 时取泡利顺磁磁化率的值。}
\end{figure}

一般地，波矢依赖磁场 $H_{\mathbf{q}}$ 可由 $\mathbf{H}(\mathbf{r})$ 的傅里叶变换
定义：
\begin{equation*}
\mathbf{H}_{\mathbf{q}} = \int\mathrm{d}^{3}\mathbf{r}\,\mathbf{H}(\mathbf{r})e^{-i\mathbf{q}\cdot\mathbf{r}},\qquad \mathbf{H}(\mathbf{r}) = \frac{1}{(2\pi)^{3}}\int\mathrm{d}^{3}\mathbf{q}\,\mathbf{H}_{\mathbf{q}}e^{i\mathbf{q}\cdot\mathbf{r}}.\tag{7.80}
\end{equation*}
同样，波矢依赖磁化强度 $M_{\mathbf{q}}$ 定义为
\begin{equation*}
\mathbf{M}_{\mathbf{q}} = \int\mathrm{d}^{3}\mathbf{r}\,\mathbf{M}(\mathbf{r})e^{-i\mathbf{q}\cdot\mathbf{r}},\qquad \mathbf{M}(\mathbf{r}) = \frac{1}{(2\pi)^{3}}\int\mathrm{d}^{3}\mathbf{q}\,\mathbf{M}_{\mathbf{q}}e^{i\mathbf{q}\cdot\mathbf{r}},\tag{7.81}
\end{equation*}
波矢依赖磁化率 $\chi_{\mathbf{q}}$ 类似地
\begin{equation*}
\chi_{\mathbf{q}} = \int\mathrm{d}^{3}\mathbf{r}\,\chi(\mathbf{r})\mathrm{e}^{-\mathrm{i}\mathbf{q}\cdot\mathbf{r}},\qquad \chi(\mathbf{r}) = \frac{1}{(2\pi)^{3}}\int\mathrm{d}^{3}\mathbf{q}\,\chi_{\mathbf{q}}\mathrm{e}^{\mathrm{i}\mathbf{q}\cdot\mathbf{r}},\tag{7.82}
\end{equation*}
式 7.80、7.81、7.82 都同时给出了逆变换。波矢依赖磁化率提供了计算电子气磁化强度的
手段。电子气中任意点 r 的磁化强度 $\mathbf{M}(\mathbf{r})$ 与磁场 $\mathbf{H}(\mathbf{r})$
的关系为
\begin{equation*}
\mathbf{M}(\mathbf{r}) = \int\mathrm{d}^{3}\mathbf{r}'\,\chi(\mathbf{r}-\mathbf{r}')\mathbf{H}(\mathbf{r}')\tag{7.83}
\end{equation*}
这是一个卷积。故 $\mathbf{M}(\mathbf{r})$ 不仅响应 r 处的场，还响应附近值的加权
平均。卷积定理意味着相应波矢依赖量之间的关系为
\begin{equation*}
\mathbf{M}_{\mathbf{q}} = \chi_{\mathbf{q}}\mathbf{H}_{\mathbf{q}},\tag{7.84}
\end{equation*}
与我们上面用的式 7.77 相同。

\subsection{电子气的抗磁响应}

波矢依赖的抗磁磁化率也可以计算，只是相当复杂。这里只引用结果：
\begin{equation*}
\chi_{\mathbf{q}} = \chi_{\mathrm{L}}\frac{3k_{\mathrm{F}}^{2}}{2q^{2}}\left[1+\frac{4k_{\mathrm{F}}^{2}}{q^{2}} - \frac{k_{\mathrm{F}}}{q}\left(1-\frac{q^{2}}{4k_{\mathrm{F}}^{2}}\right)^{2}\log\left|\frac{q+2k_{\mathrm{F}}}{q-2k_{\mathrm{F}}}\right|\right].\tag{7.85}
\end{equation*}
抗磁磁化率 $\chi_{\mathbf{q}}$ 绘于图 7.11。q=0 时表达式退化为
$\chi_{q=0}=\chi_{\mathrm{L}}=-\chi_{\mathrm{P}}/3$。式 7.85 在 $q=2k_{\mathrm{F}}$
处同样有奇异性——费米面存在之故。

\begin{figure}
\centering
\includegraphics[width=0.5\textwidth]{fig7_11.pdf}
\caption{波矢依赖的抗磁磁化率。q=0 时取朗道抗磁磁化率的值。}
\end{figure}

\subsection{RKKY 相互作用}

波矢依赖的顺磁磁化率对 delta 函数微扰系统的实空间行为有重要后果。delta 函数微扰
$\mathbf{H}(\mathbf{r})=\delta(\mathbf{r})\mathbf{H}$ 可以分解成所有可能空间频率
之和：
\begin{equation*}
\mathbf{H}(\mathbf{r}) = \delta(\mathbf{r})\mathbf{H} = \frac{1}{(2\pi)^{3}}\int\mathbf{H}_{\mathbf{q}}e^{i\mathbf{q}\cdot\mathbf{r}}\,\mathrm{d}^{3}q\tag{7.86}
\end{equation*}
$H_{\mathbf{q}}=H$。delta 函数因此可以看作一切频率之和——包括波长极长与极短的
振荡。若 $\chi_{\mathbf{q}}=\chi$ 与 q 无关，则 $M_{\mathbf{q}}=\chi H_{\mathbf{q}}$，
且
\begin{equation*}
\mathbf{M} = \frac{1}{(2\pi)^{3}}\int\mathbf{M}_{\mathbf{q}}e^{i\mathbf{q}\cdot\mathbf{r}}\,\mathrm{d}^{3}q = \chi\delta(\mathbf{r})\mathbf{H},\tag{7.87}
\end{equation*}
即磁化强度也是 delta 函数：电子气完全响应微扰、完美地跟随其行为。（此时
$\chi(\mathbf{r}-\mathbf{r}')=\chi\delta(\mathbf{r}-\mathbf{r}')$。）

然而电子气的磁化率依赖 q。由于费米面的存在，波矢高于 $q=2k_{\mathrm{F}}$ 时电子气
的响应能力大大降低。因此它虽然能跟上 delta 函数微扰的长波长分量，短波长成分
（波长短于约 $\pi/k_{\mathrm{F}}$）被衰减——电子气难以响应如此高的空间频率。
这在磁化强度中产生空间“振铃”效应，类比单缝衍射图样中的振荡。由此产生的磁化
强度具有空间振荡。

可以这样计算。实空间磁化率为
\begin{align*}
\chi(\mathbf{r}) &= \frac{1}{(2\pi)^{3}}\int\mathrm{d}^{3}\mathbf{q}\,\chi_{\mathbf{q}}\mathrm{e}^{\mathrm{i}\mathbf{q}\cdot\mathbf{r}}\\
&= \frac{1}{(2\pi)^{3}}\int\mathrm{d}^{3}\mathbf{q}\,\frac{\chi_{\mathrm{P}}}{2}\left(1+\frac{4k_{\mathrm{F}}^{2}-q^{2}}{4k_{\mathrm{F}}q}\log\left|\frac{q+2k_{\mathrm{F}}}{q-2k_{\mathrm{F}}}\right|\right)\mathrm{e}^{\mathrm{i}\mathbf{q}\cdot\mathbf{r}}\\
&= \frac{2k_{\mathrm{F}}^{3}\chi_{\mathrm{P}}}{\pi}F(2k_{\mathrm{F}}r)\tag{7.88}
\end{align*}
$r=|\mathbf{r}|$，函数 $F(x)$ 为
\begin{equation*}
F(x) = \frac{-x\cos x+\sin x}{x^{4}}\tag{7.89}
\end{equation*}
绘于图 7.12。这就是 RKKY 相互作用（4.2.4 节已介绍）。式 7.88 的积分绝非平凡，
习题 7.5 将更详细地考察。磁化率（因而 delta 函数微扰所致的磁化强度）在大距离
（$r\gg k_{\mathrm{F}}^{-1}$）正比于 $\cos(2k_{\mathrm{F}}r)/r^{3}$ 且振荡；小 r
处磁化率发散——这是我们假定的 delta 函数微扰的后果。RKKY 相互作用最初是为理解
局域核磁矩对电子气的效应而提出的：核磁矩几乎是最接近 delta 函数的微扰，但它并非
无限局域，而是摊在（尽管极小的）核体积上。RKKY 相互作用为金属中局域电子磁矩之间
的磁耦合提供重要机制：一个磁矩产生电子气的振荡磁化强度，后者可与第二个磁矩相互
作用。依两磁矩间距不同，相互作用可铁磁可反铁磁。

非磁情形也观察到类似现象——电子气对 delta 函数电荷的响应。这与金属合金中局域
电荷的研究有关：电荷密度中的振荡由类似机制产生，称为 Friedel 振荡。

\begin{figure}
\centering
\includegraphics[width=0.5\textwidth]{fig7_12.pdf}
\caption{式 7.89 的函数 $F(x)$——描述 delta 函数场产生的实空间磁化强度。}
\end{figure}

\section{电子气中的激发}

电子气还包含丰富的激发谱。优雅的处理是计算 $\chi(\mathbf{q},\omega)$——依赖
q 与 $\omega$ 的磁化率——但本书不采用此法。

在无磁场时能带就已自旋劈裂的巡游铁磁体中，有两类激发。其一是熟悉的自旋波激发，
遵循常规色散关系：我们 6.6 节的自旋波处理基于海森伯模型、适用于局域矩系统，但
事实证明金属体系中也能导出自旋波。其二是电子--空穴激发的连续谱：电子从一个自旋
劈裂带的已填满态转移到另一自旋劈裂带的空态。后者称为 Stoner 激发（Stoner
excitations）。Stoner 激发中，波矢 $\mathbf{k}+\mathbf{q}$、自旋向下的电子被
激发到波矢 $\mathbf{k}$、自旋向上的态。激发能量为
\begin{equation*}
\hbar\omega = E_{\mathbf{k}+\mathbf{q}} - E_{\mathbf{k}} + \Delta.\tag{7.90}
\end{equation*}
$E_k=\hbar^{2}k^{2}/2m_{\mathrm{e}}$，$\Delta$ 是交换劈裂——翻转一个自旋的能量
代价。结果示于图 7.13。

\begin{figure}
\centering
\includegraphics[width=0.31\textwidth]{fig7_13.pdf}
\caption{电子气中的激发。（a）自旋劈裂带间隔为交换劈裂 $\Delta$。（b）自旋向上与
自旋向下电子的费米面。（c）若存在波矢 $\mathbf{k}$ 的已填满 $\uparrow$ 态与波矢
$\mathbf{k}+\mathbf{q}$ 的空 $\downarrow$ 态，则可以产生 Stoner 激发；阴影区给出
不同 $\mathbf{q}$ 值下 $\mathbf{k}$ 的可能选取。（d）色散关系（$\omega$ 对 $q$
的图）显示自旋波支与 Stoner 激发连续谱。}
\end{figure}

顺磁金属中两带在无磁场时并不自旋劈裂。自旋波激发寿命很短、被强烈阻尼（若视作
阻尼谐振子则处于过阻尼极限）；它们称为顺磁振子（paramagnons），与常规自旋波
激发一样可以用非弹性中子散射研究。

本书迄今用的是两幅截然不同的图景：绝缘材料——磁性联系于原子上局域矩；金属——
磁性联系于完全离域的矩。许多真实材料介于两者之间：磁性联系于介于局域矩与能带铁磁
之间的自旋密度涨落。局域矩涨落在实空间局域，而弱铁磁金属中的自旋涨落可视为在
q 空间局域。

\section{自旋密度波}

Stoner 判据表明：$Ug(E_{\mathrm{F}})$ 足够大时电子气对自发铁磁性不稳定（见
7.3 节），磁化率可被增强 $(1-Ug(E_{\mathrm{F}}))^{-1}$ 倍。事实证明，波矢依赖
磁化率（无库仑相互作用时记作 $\chi_{\mathbf{q}}^{0}=\chi_{\mathrm{P}}f(q/2k_{\mathrm{F}})$）
也被库仑相互作用增强为
\begin{equation*}
\chi_{\mathbf{q}} = \frac{\chi_{\mathrm{P}}f(q/2k_{\mathrm{F}})}{1-Ug(E_{\mathrm{F}})f(q/2k_{\mathrm{F}})} = \frac{\chi_{\mathbf{q}}^{0}}{1-\alpha\chi_{\mathbf{q}}^{0}}\tag{7.91}
\end{equation*}
$\alpha=U/\mu_0\mu_{\mathrm{B}}^{2}$。可能出现：$\chi_{\mathbf{q}}^{0}$ 在某个
q≠0 处取极大。此时若 $\alpha\chi_{\mathbf{q}}^{0}$ 达到 1，参数化为 $\alpha$ 的
相互作用可使磁化率在该 q 值发散——样品中可能自发发展出振荡的静态磁化强度。
q=0 对应铁磁有序；$|q|=\pi/a$ 对应反铁磁有序。一般地，预期出现波矢 q 的螺旋结构
或自旋密度波结构。

\begin{figure}
\centering
\includegraphics[width=0.5\textwidth]{fig7_14.pdf}
\caption{一、二、三维的波矢依赖顺磁磁化率。一维情形响应函数在 q=2k$_{\mathrm{F}}$ 处
发散。所有维数下 q=0 处都取泡利顺磁磁化率的值。}
\end{figure}

我们已对三维电子气求出了波矢依赖磁化率，但对二维与一维可以重复计算（见习题
7.6）。结果示于图 7.14：降低维数时，三维曲线在 $q=2k_{\mathrm{F}}$ 处的“膝”在
二维变成扭折、在一维变成奇异性。一维的峰称为 Kohn 异常（Kohn anomaly），表明
一维电子气对形成波矢 $2k_{\mathrm{F}}$ 的自旋密度波不稳定。这一不稳定与如下事实
相关：波长 $2\pi/q$、波矢 q 的磁化强度周期性调制会在波矢 $\pm q/2$ 处打开能隙
（任何周期势都能做到）。\bio{Walter Kohn\\（1923--　）}若 $q/2=k_{\mathrm{F}}$，能隙恰在费米面处，可以降低总能量——降低电子能量的机会本身就能驱动自旋密度波的形成。

\begin{figure}
\centering
\includegraphics[width=0.45\textwidth]{fig7_15.pdf}
\caption{（a）能带中电子填充直到费米能 $E_{\mathrm{F}}$。（b）波矢 $q=2k_{\mathrm{F}}$
的周期势通过在费米面产生能隙降低总电子能量。费米面打开能隙（从（a）到（b））伴随
金属--绝缘体转变；相应的实空间畸变称为 Peierls 畸变。一维电子气对这种 Peierls
不稳定性是不稳定的。}
\end{figure}
\bio{Rudolf E. Peierls\\（1907--1995）}

波矢 $2k_{\mathrm{F}}$ 的电荷密度波的形成有类似的不稳定性。两种效应都与著名的
Peierls 不稳定性相关：一维原子链的二聚化可以自发发生，因为在费米面附近打开能隙
所节省的电子能量可以超过二聚化的弹性能量代价（见图 7.15）。

三维中自旋密度波无法在费米面所有点产生能隙。有时费米面是这样的：把费米面的一部分
平移矢量 q 后恰好落在另一部分之上——这一现象称为嵌套（nesting）。波矢 q 的自旋
密度波可以在费米面上可嵌套的区域产生能隙。真实金属中经常出现：两块费米面在 k 空间
中被固定波矢 q 近似平移相接。这可以产生磁化率峰与由此的不稳定性：密度波的形成
被称为嵌套了费米面。

若嵌套波矢 q 恰为 $\pi/a$（a 为原子间距），自旋密度波与晶格可公度，产生反铁磁
有序（图 7.16(a)）。但更常见的情形是嵌套波矢 q（等于 $2k_{\mathrm{F}}$）不是
$\pi/a$ 的简单倍数，自旋密度波不可公度（图 7.16(b)）。

\begin{figure}
\centering
\includegraphics[width=0.62\textwidth]{fig7_16.pdf}
\caption{（a）可公度的自旋密度波，波矢 $q=\pi/a$。（b）不可公度自旋密度波。}
\end{figure}

铬具有体心立方结构（立方边长 a），在 $T_{\mathrm{N}}=310$ K 以下发展出自旋密度
波结构：波矢 $\mathbf{q}\approx0.96(2\pi/a)$（故材料“几乎”是反铁磁体），随温度
缓慢变化，方向沿 [100]。极化在 $T<115$ K 时为纵向、$T>115$ K 时为横向。Cr 的
电阻率在略低于 $T_{\mathrm{N}}$ 处上升——自旋密度波在部分费米面上引入能隙、减少
载流子数目所致。若干有机金属中也发现自旋密度波。

\section{近藤效应}

非磁宿主中磁性离子的极稀合金中，若忽略 RKKY 耦合，磁性离子的磁矩可视为独立。
剩下的唯一相互作用是磁矩与传导电子自旋之间。高温下磁矩表现得像自由顺磁矩；低于
特征温度——近藤温度（Kondo temperature）$T_{\mathrm{K}}$——磁矩与传导电子的
相互作用使杂质自旋变得非磁。实际发生的是：传导电子开始在杂质自旋周围形成相反
自旋极化的云，形成准束缚态。传导电子对磁性杂质的这种磁屏蔽过程称为近藤效应
（Kondo effect），有两个深刻的实验后果。其一，磁化强度降到自由矩值以下，磁化率
因此低于居里定律的预期。其二，由于杂质矩与传导电子强烈相互作用，矩的电子散射
截面强烈增强，电阻率出现正比于 $\mathsf{J}\ln T$ 的新项（$\mathsf{J}<0$ 是局域矩
与传导电子之间的交换耦合）。

多数材料的电阻率随温度降低而减小：材料冷却时声子数目减少，而电阻率强烈依赖声子
数目。故通常高温下由正比于 T 的项主导，较低温度下由 $T^{5}$ 主导；最低温度时由
杂质浓度（不依赖温度）决定的常数项变得重要。

近藤效应的新 $\mathsf{J}\ln T\propto-\ln T$ 项（因 $\mathsf{J}<0$）在低温介入，
导致电阻率极小的出现（见图 7.17）：低于极小值温度，电阻率随降温而上升。

\begin{figure}
\centering
\includegraphics[width=0.42\textwidth]{fig7_17.pdf}
\caption{电阻率的温度依赖显示极小。电阻率是正比于 $T^{5}$ 的项与 $\mathsf{J}\ln T\propto-\ln T$ 项之和。}
\end{figure}

\section{哈伯德模型}

金属行为之所以出现，是因为电子通过在整个晶体上离域化而节省动能。然而在某些体系
中，在位库仑能（把两个电子放到同一格点的代价）强到电子无法在晶体中自由移动；
双占据的高昂代价意味着金属行为失效。电子不再能当作自由粒子处理，而电子关联——
库仑效应大时电子必须相互回避的那些特征——变得非常重要。金属行为（由紧束缚模型
中的转移积分或电子带宽的大小参数化）与库仑能（由 U 参数化，有时称 Hubbard-U）
的竞争由哈伯德模型（Hubbard model）表达。

材料中每个格点可以双占据——一个自旋向上电子与一个自旋向下电子。泡利不相容原理
禁止同自旋电子双占据同一格点。半满极限下（电子恰好每格点一个），若每个格点上
坐一个电子且不移动，在位库仑能可以极小。$U=\infty$ 时这就是基态。U 大而非无穷
时，每个格点仍单占据，但若邻居自旋反平行，电子可以短暂拜访邻居的格点——省一点
动能，因此邻居格点上电子反平行有能量好处，系统将是反铁磁体。但电子无法到次近邻
格点——那里自旋相同，泡利原理禁止同自旋双占据。因此系统还是绝缘体。这种 U 大、
每格点一个电子的反铁磁绝缘体称为莫特绝缘体（Mott insulator）。U 小于电子带宽 W
时实现金属行为。交叉发生在 $U\sim W$——存在金属--绝缘体转变。
\bio{Nevill F. Mott\\（1905--1996）}

\section{中子星}

本章以磁性凝聚态物质的一个极端例子结束：中子星。中子星形成于超新星爆发中恒星的
快速坍缩，由高密核物质组成。恒星坍缩中角动量守恒意味着前身星原有的转速大为增加
（星坍缩时转动惯量骤减、转速飙升）。中子星半径典型 $\sim10$ km，转速可达
$10^{3}\,\mathrm{s}^{-1}$。坍缩还伴随磁通压缩，因此中子星的磁场巨大，可达
$\sim10^{8}$ T。作为对比，太阳表面磁场 $\sim10^{-4}$ T，地球表面 $\sim5\times10^{-5}$
T。磁场在极区附近最强，但中子星的磁轴与地球一样相对自转轴倾斜，故随中子星旋转而
进动。变化的磁场经法拉第定律产生大电场，沿磁轴加速电子；电子沿弯曲轨迹被加速而
发出电磁辐射——形成随中子星自转而旋转的紧致辐射束。地球恰好处在束流路径上时，
我们便收到规则的辐射脉冲。因此中子星最初被称为脉冲星（pulsar）——脉动的星。

用“自由中子”模型并作非相对论处理，可以对中子星内部中子的行为做些（很粗糙的）
估计。中子与电子一样是费米子、服从泡利不相容原理。作为中子星的初步模型，可以
按 7.1 节的方法考虑中子填满中子星大小的盒子里的 k 态。中子星内部主要是中子，密度
比金属中传导电子高约 12 个量级；中子质量当然比电子质量高约 3 个量级。由于费米能
正比于 $n^{2/3}/m$，中子星中的费米能应比金属中高约 5 个量级。于是中子星中中子的
费米温度\fn{7}{作为对比，金属中的电子 $E_{\mathrm{F}}\sim1$ eV，故 $T_{\mathrm{F}}\sim10^{4}$ K。}应约 $10^{9}$ K，远高于中子星温度（可达 $\sim10^{6}$ K）——表明中子星中的中子处于简并极限。中子星表面据认为是固体，内部物质可能处于超流态。这类似于
超流氦中氦原子的配对；中子星中凝聚的可能是中子对、也有质子对。中子与质子的配对能
为 1--2 MeV，相当于 $\sim10^{10}$ K——远高于中子星温度，故可预期配对发生。

\section*{延伸阅读}

\begin{itemize}[itemsep=0.2em]
\item 朗道抗磁性的讨论见 A.~Dupr\'{e}, American Journal of Physics, \textbf{49},
34 (1981)；本章的处理改编自该文。
\item 金属能带理论与德哈斯--范阿尔芬效应的清晰描述见 J.~Singleton,《Electronic
properties of solids》, OUP 2001。
\item 朗道能级物理的有用处理见 J.~H.~Davies,《The physics of low-dimensional
semiconductors》第 6 章，CUP 1998。
\item 能带电子磁性的高级论述见 K.~Yosida,《Theory of magnetism》, Springer 1996。
\item 自旋密度泛函理论见 J.~K\"{u}bler,《Theory of itinerant electron magnetism》,
OUP 2000。
\item 金属磁体中涨落的有用信息见 T.~Moriya,《Spin fluctuations in itinerant
electron magnetism》, Springer 1985。
\item 近藤效应见 A.~C.~Hewson,《The Kondo problem to heavy fermions》, CUP 1993。
\item 中子星及相关天体的物理见 S.~L.~Shapiro 与 S.~A.~Teukolsky,《Black holes,
white dwarfs and neutron stars》, Wiley 1983。
\end{itemize}

\section*{习题}

\exer{7.1} 证明泡利顺磁性给出的磁化率可写成类居里形式
\begin{equation*}
\chi_{\mathrm{P}} = \frac{n\mu_0\mu_{\mathrm{B}}^{2}}{k_{\mathrm{B}}T_0}\tag{7.92}
\end{equation*}
常数 $T_0$ 为
\begin{equation*}
T_0 = \frac{2T_{\mathrm{F}}}{3}\tag{7.93}
\end{equation*}
$T_{\mathrm{F}}=E_{\mathrm{F}}/k_{\mathrm{B}}$ 是费米温度。再证明 $T\ll T_{\mathrm{F}}$
时 $\chi_{\mathrm{P}}$ 的温度依赖修正为
\begin{equation*}
\chi_{\mathrm{P}}(T) = \frac{\mu_0\mu_{\mathrm{B}}^{2}}{3k_{\mathrm{B}}T_0}\left(1-\frac{\pi^{2}T^{2}}{12T_{\mathrm{F}}^{2}}\right)\tag{7.94}
\end{equation*}
并估算典型金属在室温下该修正的大小。

\exer{7.2} （a）证明含 N 个电子的非简并电子气的顺磁磁化率，与 N 个轨道运动被淬灭
的独立局域电子的磁化率相同。

（b）考虑导带含 $3\times10^{22}$ 个电子每立方米、有效质量 $m^{*}=0.1m_{\mathrm{e}}$
的半导体。估算磁化率在何温度以下变得不依赖温度。低于该温度时，计算泡利顺磁性与朗道
抗磁性的大小。

\exer{7.3} 存在磁场 $\mathbf{B}=(0,0,B)$ 时，动量算符 $\hat{\mathbf{p}}=-\mathrm{i}\hbar\nabla$
须换成 $-\mathrm{i}\hbar\nabla+e\mathbf{A}$，导致式 7.40 的能量本征值。7.5 节选了
$\mathbf{A}=(0,Bx,0)$。证明改用 $\mathbf{A}=(-By,0,0)$ 可得同样结果。问题也可以
在柱坐标中求解：$A_{\phi}=\frac{1}{2}Br$、$A_z=A_r=0$。证明这一选择对 $\mathbf{B}$
与能量本征值给出相同结果。

\exer{7.4} 证明费米球的傅里叶变换与式 7.89 的 RKKY 函数相关，即
\begin{align*}
\int_{|\mathbf{k}|<k_{\mathrm{F}}}\mathrm{d}^{3}\mathbf{k}\,\mathrm{e}^{\mathrm{i}\mathbf{k}\cdot\mathbf{r}} &= \int_{0}^{k_{\mathrm{F}}}2\pi k^{2}\mathrm{d}k\int_{0}^{\pi}\mathrm{e}^{\mathrm{i}kr\cos\theta}\sin\theta\,\mathrm{d}\theta\\
&= \frac{4\pi}{r^{3}}\left[\sin k_{\mathrm{F}}r - k_{\mathrm{F}}r\cos k_{\mathrm{F}}r\right]\tag{7.95}
\end{align*}

\exer{7.5} 式 7.88 的积分求值有点复杂，本题设计为一步一步来。先证明
\begin{equation*}
\int\mathrm{d}^{3}\mathbf{q}\,\mathrm{e}^{\mathrm{i}\mathbf{q}\cdot\mathbf{r}}\,\varphi(\mathbf{q}) = \frac{2\pi}{\mathrm{i}r}\int_{-\infty}^{\infty}\mathrm{d}q\,q\,\mathrm{e}^{\mathrm{i}qr}\varphi(q),\tag{7.96}
\end{equation*}
$\varphi(q)=\varphi(-q)$ 是变量 q 的函数。进而证明
\begin{align*}
\frac{1}{(2\pi)^{3}}\int\mathrm{d}^{3}\mathbf{q}\, & \frac{\chi_{\mathrm{P}}}{2}\left(1+\frac{4k_{\mathrm{F}}^{2}-q^{2}}{4k_{\mathrm{F}}q}\log\left|\frac{q+2k_{\mathrm{F}}}{q-2k_{\mathrm{F}}}\right|\right)\mathrm{e}^{\mathrm{i}\mathbf{q}\cdot\mathbf{r}}\tag{7.97}\\
&= \frac{\chi_{\mathrm{P}}k_{\mathrm{F}}^{2}}{\mathrm{i}\pi^{2}r}\int_{-\infty}^{\infty}\mathrm{d}x\,x\,\mathrm{e}^{\mathrm{i}yx}f(x),\tag{7.98}
\end{align*}
$x=q/2k_{\mathrm{F}}$、$y=2k_{\mathrm{F}}r$，$f(x)$ 是式 7.79 定义的函数。求此
积分需要借助围道积分的技巧。

考虑两个积分
\begin{align*}
I_1 &= \frac{1}{\mathrm{i}r}\int_{-\infty}^{\infty}\mathrm{d}x\,\frac{x\,\mathrm{e}^{\mathrm{i}yx}}{2}\left(1+\frac{1-x^{2}}{2x}\log\left|\frac{x+1}{x-1}\right|\right)\tag{7.99}\\
I_2 &= \frac{1}{\mathrm{i}r}\int_{-\infty}^{\infty}\mathrm{d}x\,\frac{x\,\mathrm{e}^{\mathrm{i}yx}}{2}\left(1+\frac{1-x^{2}}{2x}\log\frac{x+1}{x-1}\right).\tag{7.100}
\end{align*}
两个被积函数在 $|x|>1$ 时相等。把围道用 x 平面上半平面的无穷大半圆补全并用
Jordan 引理，证明 $I_2$ 为零。于是利用这一结果与一条从 x=−1 到 x=1、在实轴上方
无穷小位移的围道计算 $I_1$，证明
\begin{align*}
I_1 &= \frac{1}{\mathrm{i}r}\int_{-1}^{1}\mathrm{d}x\,\frac{\mathrm{e}^{\mathrm{i}yx}(1-x^{2})}{4}\left(\log\left|\frac{x+1}{x-1}\right|-\log\frac{x+1}{x-1}\right)\\
&= \frac{\pi}{r}\int_{-1}^{1}\mathrm{d}x\,\frac{\mathrm{e}^{\mathrm{i}yx}(1-x^{2})}{4}\\
&= \frac{\pi}{r}\left[\frac{\sin y-y\cos y}{y^{3}}\right]\\
&= 2\pi k_{\mathrm{F}}F(2k_{\mathrm{F}}r),\tag{7.101}
\end{align*}
从而验证式 7.88。

\exer{7.6} （a）证明一、二、三维费米能级处的态密度分别为
\begin{equation*}
\begin{aligned}
g(E_{\mathrm{F}}) &= \frac{n}{2E_{\mathrm{F}}} = \frac{2m}{\hbar^{2}\pi k_{\mathrm{F}}}\\
g(E_{\mathrm{F}}) &= \frac{n}{E_{\mathrm{F}}} = \frac{m}{\hbar^{2}\pi}\\
g(E_{\mathrm{F}}) &= \frac{3n}{2E_{\mathrm{F}}} = \frac{mk_{\mathrm{F}}}{\hbar^{2}\pi^{2}}
\end{aligned}\tag{7.102}
\end{equation*}
（一维、二维、三维）

（b）证明一维电子气的波矢依赖磁化率为
\begin{equation*}
\chi_{\mathbf{q}} = \chi_{\mathrm{P}}\frac{k_{\mathrm{F}}}{q}\log\left|\frac{q+2k_{\mathrm{F}}}{q-2k_{\mathrm{F}}}\right|\tag{7.103}
\end{equation*}
二维为
\begin{equation*}
\chi_{\mathbf{q}} = \begin{cases} \chi_{\mathrm{P}} & q<2k_{\mathrm{F}}\\ \chi_{\mathrm{P}}\left(1-\sqrt{1-(4k_{\mathrm{F}}^{2}/q^{2})}\right) & q>2k_{\mathrm{F}}. \end{cases}\tag{7.104}
\end{equation*}
这些结果绘于图 7.14。

\exer{7.7} 证明 7.8 节描述、图 7.13 绘出的 Stoner 激发连续谱以两条抛物线
\begin{equation*}
\hbar\omega = \Delta + \frac{\hbar^{2}q}{2m_{\mathrm{e}}}(q\pm2k_{\mathrm{F}\uparrow})\tag{7.105}
\end{equation*}
为界，且 $\omega=0$ 的激发包含在
\begin{equation*}
k_{\mathrm{F}\uparrow}-k_{\mathrm{F}\downarrow} \leq q \leq k_{\mathrm{F}\uparrow}+k_{\mathrm{F}\downarrow},\tag{7.106}
\end{equation*}
的 q 范围内，如图 7.13 所示。（提示：利用 $\hbar^{2}k_{\mathrm{F}\uparrow}^{2}/2m_{\mathrm{e}}=\hbar^{2}k_{\mathrm{F}\downarrow}^{2}/2m_{\mathrm{e}}+\Delta$。）

\exer{7.8} 考虑图 7.5 所示的能带劈裂，但现在 $\delta E$ 不是无穷小量。费米能变得
依赖自旋。若电子系统每单位体积有 n 个电子，证明每单位体积自旋向上电子数 $n_{+}$
与自旋向下电子数 $n_{-}$ 为
\begin{align*}
n_{\pm} &= \int_{0}^{E_{\mathrm{F}\pm}}g_{\pm}(E)\,\mathrm{d}E\tag{7.107}\\
&= \int_{0}^{E_{\mathrm{F}\pm}}\frac{3}{4}\frac{n}{E_{\mathrm{F}}^{3/2}}E^{1/2}\,\mathrm{d}E\tag{7.108}\\
&= \frac{n}{2}\left(\frac{E_{\mathrm{F}\pm}}{E_{\mathrm{F}}}\right)^{3/2},\tag{7.109}
\end{align*}
$E_{\mathrm{F}}$ 是无交换劈裂时的费米能。

用 $n_{+}+n_{-}=n$ 并定义 $x=(n_{+}-n_{-})/n$，证明
\begin{equation*}
\frac{E_{\mathrm{F}\pm}}{E_{\mathrm{F}}} = (1\pm x)^{2/3}.\tag{7.110}
\end{equation*}
动能为
\begin{equation*}
T = \sum_{\sigma=\pm}\int_{0}^{E_{\mathrm{F}\sigma}}E\,g(E)\,\mathrm{d}E.\tag{7.111}
\end{equation*}
证明
\begin{equation*}
T = \frac{3}{10}nE_{\mathrm{F}}\left[(1+x)^{5/3}+(1-x)^{5/3}\right].\tag{7.112}
\end{equation*}
相互作用能可写为
\begin{equation*}
V = -\int_{0}^{M}\mu_0(\lambda M')\,\mathrm{d}M' = -\frac{1}{2}\mu_0\lambda M^{2} = -\frac{1}{2}Un^{2}x^{2}\tag{7.113}
\end{equation*}
其中 $\lambda$、$U=\mu_0\mu_{\mathrm{B}}^{2}\lambda$ 与
$M=\mu_{\mathrm{B}}(n_{+}-n_{-})$ 分别是分子场参数、库仑能与磁化强度。总能量为
\begin{equation*}
E = T+V.\tag{7.114}
\end{equation*}
证明条件 $\mathrm{d}E/\mathrm{d}x=0$ 在
\begin{equation*}
\frac{(1+x)^{2/3}+(1-x)^{2/3}}{x} = \frac{4Ug(E_{\mathrm{F}})}{3}.\tag{7.115}
\end{equation*}
时给出稳定解。（提示：用 $g(E_{\mathrm{F}})=3n/2E_{\mathrm{F}}$。）作图解释这一
结果。

证明当
\begin{equation*}
Ug(E_{\mathrm{F}}) > 1,\tag{7.116}
\end{equation*}
时非磁态不稳定——这就是 Stoner 判据。

解释为什么 $x=0$ 处 $\mathrm{d}^{2}E/\mathrm{d}x^{2}<0$ 时非磁态不稳定，并证明
这等价于上式。

% 第8章 竞争相互作用与低维性（Competing interactions and low dimensionality）
\chapter{竞争相互作用与低维性}

本章考察竞争相互作用与低维性使固体中出现某些极为精微、复杂、有时甚至有用的磁
行为的一些方式。竞争相互作用与低维性都天然地出现在某些体系中：例如有些晶体生长
的方式使磁矩在链中强耦合、或占据在阻挫晶格的格点上。然而这些特征也可以人为引入，
例如制备铁磁多层膜，或用电子束光刻制造铁磁线。其中许多专题目前正被世界各地的研究
小组积极研究。本章叙述这些进展中的一部分——其中许多仍笼罩着尚未揭开的谜。由于
很多专题既新又复杂，有些讨论难免过于简化；有兴趣的读者可查阅章末延伸阅读中的
文献以了解详情。尽管如此，我希望传达出若干问题的风味——在写作之时，它们正给
凝聚态物理学家带来巨大的挑战。

我们从阻挫（frustration）开始——它可以导致包括自旋玻璃在内的多种基态——然后
继续讨论低维性的效应。

\section{阻挫}

在某些晶格中，不可能同时满足体系中的所有相互作用以找到基态。这常常导致没有唯一
的基态，而是存在许多相似的低能态，体系的不满足感（指能量未被极小化）被尽可能地
分摊。此时称体系表现阻挫。例如考虑图 8.1 的晶格，其中只有最近邻反铁磁相互作用
起作用。在正方晶格（图 8.1(a)）上，让最近邻自旋反平行的要求很容易满足；但在三角
晶格上事情就不那么简单了：如图所示，若两个相邻自旋反平行排列，第三个自旋便陷入
两难——无论怎么选，两个邻居中总有一个的能量未被极小。体系因此无法达到完全满足其
微观约束的状态，但拥有大量同样不满足的态。结果，阻挫体系表现出亚稳性、磁滞效应
（依赖样品的磁历史或热历史）以及趋向平衡的时间弛豫。

\begin{figure}
\centering
\includegraphics[width=0.45\textwidth]{fig8_01.pdf}
\caption{（a）正方晶格与（b）三角晶格上的反铁磁最近邻相互作用。三角晶格表现阻挫：
第三个格点上的自旋无论怎样取向，都无法与另两个自旋同时满足反铁磁最近邻相互作用的
要求。}
\end{figure}

于是某些体系中晶格的几何可以阻挫自旋的有序化。二维中，三角晶格与共角 kagomé
晶格（见图 8.2）上的海森伯自旋表现这一效应；三维中被研究得最充分的体系具有
烧绿石（pyrochlore）结构——磁性离子占据共角四面体构成的晶格。据信这些体系不
有序，而是表现具有宏观简并的经典基态——有时称为协同顺磁性（cooperative
paramagnetism）：一切温度下自旋之间都只有短程关联。这些体系中的一些还被认为
具有无色散的自旋波支——称为零模（zero mode），它强烈影响低温热力学行为，导致
自旋涨落一直持续到零温。

\begin{figure}
\centering
\includegraphics[width=0.4\textwidth]{fig8_02.pdf}
\caption{kagomé 晶格的一部分。}
\end{figure}

\section{自旋玻璃}

自旋玻璃（spin glasses）已在 5.5 节介绍。该节把自旋玻璃定义为：随机、混合相互
作用的磁性体系，其特征是在明确的温度 $T_{\mathrm{f}}$ 处自旋发生随机然而协同的
冻结，冻结温度以下出现高度不可逆的亚稳冻结态，而没有通常的磁长程序。现在我们
更细致地考察这一定义的各个部分，先看可以用几种方式引入的随机性。

其一是位置随机性（site-randomness），可在合金中实现。常被研究的自旋玻璃是
$Cu_{1-x}Mn_x$（$x\ll1$）：少量 Mn 完全随机地替代进入 Cu 基体、没有任何短程
有序。这直接导致非磁 Cu 基体中磁性 Mn 离子间距离随机（见图 5.1(c)）。取金属间
化合物如 $GdAl_2$ 使其非晶化（熔化后从熔体极速冷却，晶格被摧毁）也可以人工设计
磁离子间的随机距离。

另一种可能是键随机性（bond-randomness）：最近邻相互作用在 $+\mathsf{J}$ 与
$-\mathsf{J}$ 之间变化。可以把磁性离子固定在规则的晶格阵列上，同时调制磁性离子
之间的间接交换相互作用来实现。这在 $Rb_2Cu_{1-x}Co_xF_4$ 中实现：Cu 与 Co 的
有效自旋都是 $\frac{1}{2}$（Co 上的晶体场把 $S=\frac{3}{2}$ 能级分裂，使低温下
只有最低的 $S=\frac{1}{2}$ 二重态被占据）。晶体场还使自旋平行或反平行于 z 轴
（单轴单离子各向异性，见 3.2.2 节）——称其具有伊辛性格。然而磁性离子间超交换
相互作用的大小与符号关键取决于耦合是 Co--Co、Co--Cu 还是 Cu--Cu，以及 Cu 上
占据的是哪个轨道。净结果是离子间的键可以取不同的交换 $\mathsf{J}$ 值，这些键
随机分布在整个样品中。

自旋玻璃固有的随机性很重要，竞争相互作用的存在同样重要。随机位置自旋玻璃中
磁矩间距离的分布导致竞争相互作用，因为相互作用是 RKKY 型的，其符号（铁磁或
反铁磁）依赖自旋间距。另一要素是磁各向异性——来自单离子各向异性或
Dzyaloshinsky--Moriya 相互作用。非晶材料中这些各向异性逐点不同，因此所谓的
非晶磁体常具有随机各向异性：磁化强度的“易轴”可以局域变化。

随机键自旋玻璃中竞争相互作用自动存在——不同的键以不同方式“拉扯”系统。这些
竞争相互作用导致阻挫：像 kagomé 晶格那样，存在多重简并的基态。自旋玻璃共享这种
多重简并基态，但还表现一种新效应——协同冻结转变，现在就来描述。

高温下一切磁性系统的行为由热涨落主导，自旋玻璃中所有自旋彼此独立。自旋玻璃从
高温冷却时，独立自旋减速并结成局域关联的单元——称为团簇（clusters）。不在团簇
里的自旋参与团簇之间的相互作用。温度降到 $T_{\mathrm{f}}$ 时，团簇的涨落逐渐
减慢。自旋间的相互作用变得更加长程，每个自旋越来越“感知”到其周围逐渐扩大的
区域内的自旋。在 $T_{\mathrm{f}}$ 处系统找到它的众多基态之一并冻结。这一过程
尚未被完全理解，但看起来是一种协同相变。然而这不是向磁有序态的相变：若体系有
长程磁有序，散射实验中应出现磁布拉格峰——实测没有。$T_{\mathrm{f}}$ 以下基态
显得“玻璃态”：具有亚稳性与缓慢的弛豫行为。$T_{\mathrm{f}}$ 以下场冷却与零场
冷却磁化率出现分岔，反映这种亚稳性（见图 8.3）。

自旋玻璃行为的一个标志是交流磁化率实部 $\chi(\omega)$ 在 $T_{\mathrm{f}}$ 附近
的尖峰。该技术用频率 $\omega$ 的很小的交变磁场测量磁化率，有时同时加恒定（直流）
磁场。峰位随 $\omega$ 略有移动。与吸收相关的 $\chi(\omega)$ 虚部在 $T_{\mathrm{f}}$
附近突然出现。冻结过程的涨落动力学可用交流磁化率、中子自旋回波（在相对长时间
尺度上测量时间关联函数的方法）与 $\mu$SR 研究。这些技术表明：$T_{\mathrm{f}}$
以上形成的团簇具有很宽的弛豫时间分布。样品冷却时，团簇长大并以依赖各自尺寸与
性质的速率涨落。团簇之外的自旋快速涨落，但它一加入正在生长的团簇（有点像“恶意
收购”！）就被迫减速、与其新东家同步。在 $T_{\mathrm{f}}$ 处随机各向异性变得
重要，系统冻结成团簇固定的随机取向。$T_{\mathrm{f}}$ 以下也观察到很宽的弛豫时间
范围，表明仍有一些自由自旋或小的超顺磁（见下文）团簇存在。自旋玻璃中的协同自旋
冻结仍未被完全理解。

\begin{figure}
\centering
\includegraphics[width=0.5\textwidth]{fig8_03.pdf}
\caption{$Cu_{1-x}Mn_x$（x=0.0108 与 x=0.0202）的静态磁化率。零场冷却后，在
0.59 mT 的小场中升温测量 $\chi$（(b) 与 (d)）。若样品在 0.59 mT 中场冷却，$\chi$
可逆（(a) 与 (c)）。改编自 Nagata 等 1979。}
\end{figure}

$Cu_{1-x}Mn_x$ 自旋玻璃通常在 x 为百分之几以下时研究。浓自旋玻璃中行为大不相同：
磁最近邻数目更多（可因化学偏聚而增强），导致出现大的、有序的磁团簇并主导磁行为。
系统于是称为团簇玻璃（cluster glass，有时称 mictomagnet）——具有自旋玻璃的某些
特征，但长程磁有序始终潜伏在背后、伺机而发（术语谓之长程磁有序是“萌芽态”
（incipient））。有时高温态实际上可以是有序铁磁体而非顺磁体，系统实际冻结到低温
非磁无序态——这称为重入自旋玻璃（re-entrant spin glass），其行为据信与温度依赖
的随机各向异性或与冻结相联系的某种各向异性有关。磁性离子浓度很高时，系统逼近
渗流极限——可以找到一条由最近邻连键构成的、贯穿整个样品的路径，路径上每个离子
都是磁性的；这使样品具有长程磁有序。这些不同态之间的转变可见于 $Au_{1-x}Fe_x$
的相图（图 8.4）。

\begin{figure}
\centering
\includegraphics[width=0.5\textwidth]{fig8_04.pdf}
\caption{$\mathrm{Au}_{1-x}\mathrm{Fe}_x$ 的相图。P=顺磁，SP=超顺磁，SG=自旋
玻璃，CG=团簇玻璃，F=铁磁。取自 Coles 等 1978。}
\end{figure}

\section{超顺磁性}

若铁磁颗粒足够小，它们将是单畴的（见 6.7.7 节），因为形成畴壁的能量代价超过
退磁能的节省。小的单畴铁磁颗粒的磁化强度常被约束在平行或反平行于特定方向。原因
可以是磁晶各向异性、形状各向异性（与退磁能及颗粒形状相联系），或其他多种因素。
下面假定颗粒的能量密度含 $K\sin^{2}\theta$ 项——$\theta$ 是磁化强度与该特定
方向的夹角，K 是量化这一各向异性能量密度的常数。于是 $\theta=0$ 或 $\pi$ 时
能量最低（见图 8.5）。体积 V 的颗粒把磁化强度从 $\theta=0$ 翻转到 $\pi$（或
反之）需要激活能 $\Delta E=KV$。对非常小的颗粒，$KV$ 与 $k_{\mathrm{B}}T$ 相比
很小，热涨落便能轻易完成这种翻转。

\begin{figure}
\centering
\includegraphics[width=0.35\textwidth]{fig8_05.pdf}
\caption{磁性颗粒的能量密度含 $K\sin^{2}\theta$ 项；$\theta=0$ 或 $\pi$ 时能量
最低。}
\end{figure}

考虑非磁基体中这些小铁磁颗粒的分布，并设颗粒相距足够远、颗粒间相互作用可忽略。
$k_{\mathrm{B}}T\gg KV$ 时系统表现得像顺磁体——只不过独立“矩”不是原子矩，而是
铁磁颗粒内的一大组矩。这样的系统称为超顺磁体（superparamagnet）。高温时颗粒上
的矩能快速涨落。颗粒上矩的弛豫时间 $\tau$ 为
\begin{equation*}
\tau = \tau_0\exp\left(\frac{KV}{k_{\mathrm{B}}T}\right),\tag{8.1}
\end{equation*}
$\tau_0$ 典型为 $10^{-9}$ s。样品冷却时涨落减慢（$\tau$ 增大，见图 8.6）；当
$\tau$ 远长于所用实验技术的测量时间 t 时，系统显得静态。若定义“远长于”为
$\tau>\alpha t$（$\alpha=100$），则低于阻塞温度（blocking temperature）
$T_{\mathrm{B}}=KV/k_{\mathrm{B}}\ln(\alpha t/\tau_0)$ 时，每个磁性颗粒看起来被
锁进它的两个极小之一。

\begin{figure}
\centering
\includegraphics[width=0.42\textwidth]{fig8_06.pdf}
\caption{按式 8.1，弛豫时间 $\tau$ 随温度 T（以 $k_{\mathrm{B}}/KV$ 标度）的依赖。
温度降低时涨落减慢（$\tau$ 增大）。}
\end{figure}

实验测量时间 t：非弹性中子散射为 $10^{-12}$--$10^{-10}$ s；穆斯堡尔为
$10^{-10}$--$10^{-7}$ s（穆斯堡尔跃迁的衰变过程 $\sim10^{-7}$ s）；$\mu$SR 为
$10^{-10}$--$10^{-5}$ s（可测比例的 $\mu$ 子能活到 $\sim10\tau_{\mu}$，
$\tau_{\mu}=2.2\ \mu$s 是 $\mu$ 子平均寿命）；交流磁化率典型探测
$10^{-1}$--$10^{-5}$ s。由于对 $\alpha t/\tau_0$ 是对数依赖，把 $\alpha$、t 或
$\tau_0$ 改变几个量级，$T_{\mathrm{B}}$ 也只有相对较小的变化。若超顺磁体系中
颗粒有一系列尺寸，它们将在不同温度“阻塞”。

超顺磁体与自旋玻璃有某些相似，但多方面截然不同：超顺磁体中磁性颗粒之间的相互
作用不重要，自旋玻璃中却至关重要；自旋玻璃表现协同相变，超顺磁体则表现超顺磁
颗粒的逐渐阻塞。超顺磁性在技术上很重要，因为许多重要记录材料是颗粒型的。这一点
至关重要：磁带被期待存储信息数年而非 $\mu$s 量级——这给颗粒的最小尺寸设了限。
研究岩石磁性时，有时需要考虑磁性在数百万年上的指数衰减。因此颗粒对超顺磁涨落的
稳定性依赖你所用的时标：同一颗粒在穆斯堡尔实验中稳定、在交流磁化率中却可能涨落；
可以稳定一个世纪，却在地质时标上衰减。

\section{一维磁体}

一维中 T>0 不可能有长程序，人们或许以为一维磁体相当乏味无趣。事实恰恰相反！
一维性意味着可能存在复杂的激发——至今远未被完全理解。

\subsection{自旋链}

一维（d=1）的自旋线称为自旋链（spin chain）。单个自旋可以被约束在平行或反平行
于特定方向（伊辛自旋），或自由指向固定平面内任意方向（XY 自旋），或自由指向
任意方向（海森伯自旋）。若晶体结构使链之间保持足够距离，自旋链可以在晶体中近似
实现。晶体场导致的单离子各向异性可使磁矩表现得像伊辛自旋（D=1）、XY 自旋
（D=2）、海森伯自旋（D=3）或介于其间。一个常被研究的体系家族基于 ABX$_3$ 型
晶体：A 是一价非磁阳离子，B 是二价磁性阳离子，X 是卤素阴离子。这给出简单六角
晶格，过渡金属离子沿 c 方向构成链。例如 CsCoCl$_3$ 中各向异性把自旋约束在特定
方向，几乎表现为一维伊辛自旋链；KCuF$_3$ 表现为一维海森伯自旋链，许多带有机
配体的 Cu 盐亦然。

这些体系常在极低温下表现三维长程序——总存在某种小的链间相互作用把链耦合起来。
正因为这种链间耦合，CsCoCl$_3$ 在 21 K 以下表现三维长程磁有序。不过在向三维
行为过渡的温度之上，仍有很宽的温度区域，其磁行为是一维的。

链上每个自旋的自旋量子数取决于原子：$\mathrm{Cu^{2+}}$（3d$^9$）链为
$S=\frac{1}{2}$，$\mathrm{Mn^{2+}}$（3d$^5$）为 $S=\frac{5}{2}$；CsCoCl$_3$ 中
$\mathrm{Co^{2+}}$（3d$^7$）的基态具有有效自旋 $S=\frac{1}{2}$（$S=\frac{3}{2}$
自由离子基态被晶体场分裂，留下基态二重态与激发态）。下一节只考虑每个格点
$S=\frac{1}{2}$ 的链。

\subsection{自旋子}

使这些链有趣的不它们的有序（因为链间相互作用不够强时它们不表现有序），而是其
激发。如 6.6.1 节所示，三维海森伯磁体的激发是磁振子——玻色子（金属磁性材料中
还可能有 Stoner 激发，见 7.8 节）。每个磁振子是 S=1 的激发，能在非弹性散射实验
中与中子作用。

伊辛自旋链的激发与畴壁的产生相联系。不存在无能隙激发（即长波极限下能量趋于零
的激发；用粒子物理的行话叫“无质量戈德斯通模”），因为哪怕产生一条畴壁也要付出
有限能量（实际上激发必须成对地产生畴壁）。激发一旦产生便可沿链自由移动。激发
能量不依赖波矢，但若链不是完全伊辛型的（真实体系常如此），色散关系会有一些调制。

海森伯自旋链的激发称为自旋子（spinons）。它们具有自旋 $\frac{1}{2}$
（与自旋 1 的磁振子相对照），是费米子。色散关系为
\begin{equation*}
\hbar\omega = \pi|\mathsf{J}\sin(qa)|,\tag{8.2}
\end{equation*}
$\mathsf{J}$ 是反铁磁交换耦合，a 与 q 是沿链方向的晶格常数与波矢。式 8.2 即
图 8.7 中的粗线。可与式 6.58 的常规自旋波色散（取 $S=\frac{1}{2}$）相比——多了
一个因子 $\pi/2$。两种情形的激发都是无能隙的：$q\to0$（长波极限）时
$\omega\to0$。中子散射实验涉及自旋改变 1：在三维系统中这意味着单个磁振子的产生
或湮灭，在一维系统中则意味着两个自旋子的产生或湮灭。因此中子实验测量的是与一对
自旋子的产生或湮灭相联系的动量 $q=q_1+q_2$ 与能量
$\hbar\omega=\hbar\omega_1+\hbar\omega_2$，实验数据在式 8.2 与
$\hbar\omega=2\pi|\mathsf{J}\sin(qa/2)|$ 之间显示连续的激发（见图 8.7）。中子
散射实验已证实在某些一维反铁磁链中这些激发确实存在（见图 8.8）。（式 8.2 的推导
见 des Cloizeaux 与 Pearson 1962。）

\begin{figure}
\centering
\includegraphics[width=0.5\textwidth]{fig8_07.pdf}
\caption{一维反铁磁海森伯自旋链中自旋子激发的色散关系（粗线）。阴影区为中子散射
实验测得的激发连续谱。}
\end{figure}

\begin{figure}
\centering
\includegraphics[width=0.5\textwidth]{fig8_08.pdf}
\caption{（a）KCuF$_3$ 中自旋子的色散关系。实验采用飞行时间技术，直线为入射能量
148.9 meV、入射动量偏离 KCuF$_3$ 的 c$^*$ 方向 8° 的中子的散射轨迹；当轨迹与
连续态相交时发生散射。（b）20 K 下观察到的散射。非磁背景以虚线标出。取自
Tennant 等 1993。}
\end{figure}

\subsection{Haldane 链}

上一节指出，自旋 $\frac{1}{2}$ 反铁磁海森伯链的激发据信是无能隙的自旋子。这一
结果据信对半整数自旋链（$S=\frac{1}{2},\frac{3}{2},\frac{5}{2},\ldots$ 的链）
同样成立。Haldane 猜想：整数自旋链（$S=1,2,3,\ldots$）会有所不同——激发谱中
会出现能隙，其起因是基态中的非线性量子涨落。整数自旋的一维链因此称为 Haldane
链（Haldane chain），激发谱中的能隙称为 Haldane 能隙。半整数与整数自旋链的这一
根本差别与费米子和玻色子在交换下的差别相关；这种不同的交换对称性具有拓扑起源，
对激发的性质有戏剧性的影响。

Haldane 猜想的多数检验在 $Ni^{2+}$（S=1）离子链材料上进行，包括 CsNiCl$_3$、
$\mathrm{Ni(C_2H_8N_2)_2NO_2ClO_4}$ 与 $Y_2BaNiO_5$。这些 S=1 自旋链的激发谱
都如预言地拥有能隙。半整数反铁磁链是无能隙的——除非磁弹耦合打开自旋 Peierls
能隙，见下一节。

\subsection{自旋 Peierls 转变}

尽管自旋 $\frac{1}{2}$ 反铁磁链无能隙，它容易遭受一种与折磨一维金属的不稳定性
（见 7.9 节）类似的、能打开能隙的不稳定性。这发生在自旋 Peierls 转变（spin-Peierls
transition）。

\begin{figure}
\centering
\includegraphics[width=0.42\textwidth]{fig8_09.pdf}
\caption{（a）均匀海森伯反铁磁链与（b）交替链（基态是 q=0 的单态，单胞加倍）中
元激发的示意。改编自 Bray 等 1983。}
\end{figure}

这种本征晶格不稳定性的驱动力是一维电子结构与三维晶格振动（声子）之间的磁弹耦合。
耦合产生的原因是链的交换能量是相邻格点间距的函数。晶格畸变影响磁能量（见图
8.9）。“自旋 Peierls”之名反映了它与 Peierls 畸变（7.9 节讨论过）的相似性。

转变温度 $T_{\mathrm{SP}}$ 以上，每条链中有均匀的反铁磁近邻交换；$T_{\mathrm{SP}}$
以下出现弹性畸变，导致二聚化，进而产生两个不相等的交替交换常数。二聚化随温度降低
渐进增大，在零温达到极大。交替链在单态基态与最低三重激发带之间有能隙。能隙大小
与二聚化程度（即晶格畸变程度）相关：均匀链（零二聚化）时为零——回到无能隙自旋子
情形。于是磁化率 $\chi(T)$ 在 $T_{\mathrm{SP}}$ 处呈现“膝”，$T_{\mathrm{SP}}$
以下 $\chi$ 相当突然地跌落，对应能隙的打开（见图 8.10）。通常的 Peierls 畸变
（自旋 Peierls 转变的电子类比，见 7.9 节）发生在 $k_{\mathrm{B}}T_{\mathrm{P}}\sim E_{\mathrm{F}}\exp(-1/\lambda_{\mathrm{el-ph}})$
量级的温度（$\lambda_{\mathrm{el-ph}}$ 为电子--声子耦合常数），而自旋 Peierls
转变发生在 $k_{\mathrm{B}}T_{\mathrm{SP}}\sim|\mathsf{J}|\exp(-1/\lambda_{\mathrm{sp-ph}})$
（$\mathsf{J}$ 为相邻自旋间的交换相互作用，$\lambda_{\mathrm{sp-ph}}$ 为自旋--
声子耦合常数）。由于 $\mathsf{J}\ll E_{\mathrm{F}}$（例如 $\mathsf{J}$ 典型为 50 K，
$E_{\mathrm{F}}$ 典型为 500--5000 K），$T_{\mathrm{SP}}$ 总是远小于 $T_{\mathrm{P}}$。

表现自旋 Peierls 转变的材料非常少，因为反铁磁链常因链间耦合在低温下变为三维有序。
只有极少数材料，自旋--声子耦合能压倒链间自旋--自旋耦合、允许自旋 Peierls 基态
形成。例子包括 CuGeO$_3$（$T_{\mathrm{SP}}=14$ K）与若干有机体系，如
MEM(TCNQ)$_2$（$T_{\mathrm{SP}}=18$ K）与 TTF-CuS$_4$C$_4$(CF$_3$)$_4$
（$T_{\mathrm{SP}}=12$ K）。\fn{1}{MEM、TCNQ 与 TTF 是具有冗长化学名称的有机分子。}

\begin{figure}
\centering
\includegraphics[width=0.5\textwidth]{fig8_10.pdf}
\caption{有机自旋 Peierls 材料 MEM(TCNQ)$_2$ 的摩尔磁化率——它由有机分子 MEM 的
堆叠与有机分子 TCNQ 的堆叠构成。高温下磁化率符合均匀海森伯反铁磁链的模型（点线）；
降温时磁化率在自旋 Peierls 转变处随激发谱能隙的打开而迅速下降。极低温（T 降低时）
的回升来自缺陷的类居里（$\sim T^{-1}$）磁化率。取自 Lovett 等 2000。}
\end{figure}

\subsection{自旋梯}

在转向二维磁体之前，可以考虑一个介于一维磁体与二维磁体之间的体系。考虑两条平行的
自旋链，链间有键，链间耦合与链内耦合强度相当。这样的系统称为两腿自旋梯（two-leg
spin ladder，见图 8.11(a)）。也可以有三腿（图 8.11(b)）、四腿乃至 n 腿自旋梯。
有前景的实验体系包括 SrCu$_2$O$_3$ 与 La$_{1-x}$Sr$_x$CuO$_{2.5}$（都是两腿
自旋 $\frac{1}{2}$），以及 Sr$_2$Cu$_3$O$_5$（三腿自旋 $\frac{1}{2}$）。事实上
存在一个通式为 Sr$_{n-1}$Cu$_{n+1}$O$_{2n}$ 的体系（n 为奇数），它是 (n+1)/2 腿自旋
$\frac{1}{2}$ 体系，因为其结构含有沿宽度方向排布 (n+1)/2 个 Cu$^{2+}$ 离子的
CuO$_2$ 正方晶格条带。

\begin{figure}
\centering
\includegraphics[width=0.3\textwidth]{fig8_11.pdf}
\caption{（a）两腿梯；（b）三腿梯。（c）两腿梯的基态由梯子每一横档上的自旋单态
组成，因此两腿梯的激发谱有能隙。（d）往梯上掺入空穴会破坏单态，（e）但若空穴
配对，这一能量代价可以被最小化。}
\end{figure}

已知自旋 $\frac{1}{2}$ 两腿梯的激发谱有有限能隙——在“强横档”极限（横档耦合
$J_{\perp}$ 大于沿腿的耦合 $\mathsf{J}$）下一眼可见：基态就是每个横档上的自旋
单态（见图 8.11(c)）。要产生激发，必须把横档单态提升为横档三重态（耗能
$J_{\perp}$），故有能隙。$J_{\perp}=0$ 时系统是两条孤立的自旋 $\frac{1}{2}$ 链，
激发谱无能隙。然而据信只要 $J_{\perp}$ 非零——无论多小——能隙就出现。对 n 腿
梯，n 为偶数时情形相同；n 为奇数时，给定横档上的自旋配成单态后总剩下一个落单。
大 $J_{\perp}$ 时系统可映射为无能隙的自旋 $\frac{1}{2}$ 链。所以：偶数腿梯的
激发谱有能隙，奇数腿梯无能隙。中子、$\mu$SR 与输运实验的结果似乎支持这一点。

往自旋梯中掺杂空穴会破坏自旋梯中的单态（见图 8.11(d)）。对两腿梯，空穴配对在
能量上有利——它们可以共享同一横档，把“受损”单态从两个减少到一个
（图 8.11(e)）。因此空穴在两腿梯上配对是有利的——这展示了在两腿自旋梯中设计
出超导电性的可能。超导电性（$T_{\mathrm{c}}\sim14$ K）实际上已在
Sr$_{14-x}$Ca$_x$Cu$_{24}$O$_{21}$（有时称 [14-24-41] 或“电话号码”化合物）
中发现（x=13.6、5 GPa）。对自旋梯的兴趣在于：它们可以有激发谱能隙、可以变为
超导，却又是简单而定义明确的体系——喜欢在一维上工作的理论家可以尝试建模。因此
这些磁性体系可能照亮高 $T_{\mathrm{c}}$ 超导电性问题（见下面的 8.5 节）。

\section{二维磁体}

二维磁性常在化学式 $A_2BX_4$ 的体系中研究：与前面一样，A 是一价非磁阳离子，
B 是二价磁性阳离子，X 是卤素阴离子。晶体结构为四角，磁性离子位于二维正方晶格上。
典型材料是 $K_2NiF_4$——这些体系常被称为具有 $K_2NiF_4$ 结构。Mermin--Wagner--%
Berezinskii 定理表明：$d\leq2$ 时，各向同性体系中热涨落禁止非零温度下长程磁
有序的存在。但它对 T=0 基态不置一词。一维中，事实证明自旋 $\frac{1}{2}$ 海森伯
反铁磁体即使在 T=0 也没有长程序。二维海森伯反铁磁体的情形不那么明确，目前有
证据表明它在 T=0 确实表现长程磁有序。这一体系之所以极受关注，是因为二维海森伯
模型看来对理解高 $T_{\mathrm{c}}$ 铜氧化物超导体很重要。这些体系包含铜氧面
（见图 8.12），铜氧面对超导性质似乎至关重要。

\begin{figure}
\centering
\includegraphics[width=0.4\textwidth]{fig8_12.pdf}
\caption{$\mathrm{La_2CuO_4}$ 中的铜氧面。}
\end{figure}

例如，纯 $La_2CuO_4$ 是反铁磁绝缘体、具有 $K_2NiF_4$ 结构。$Cu^{2+}$ 离子
（3d$^9$）在 d 带有一个空穴、自旋 $\frac{1}{2}$，位于正方晶格上；每个
$Cu^{2+}$ 离子经氧离子与邻居隔开，氧传递 $Cu^{2+}$ 自旋之间的反铁磁超交换相互
作用（见图 8.12）。每格点一个电子，你会预期系统是金属，但空穴因关联而局域化。
用额外的空穴掺杂（以 $Sr^{2+}$ 替换部分 $La^{3+}$）导致最佳掺杂样品（约 20\%
的 $La^{3+}$ 被 $Sr^{2+}$ 替换）在 40 K 左右出现高温超导。可以用 $Ni^{2+}$
（S=1）替换 $Cu^{2+}$ 制备同结构材料，但它们即使掺杂也不超导——因此 $Cu^{2+}$
有某种特殊之处。特别令人惊讶的是这些磁性离子竟有助于促成超导电性——磁性通常会
摧毁超导电性：磁性离子产生的局域磁场充当对破坏者（pair-breaker），拆散承载
超流的库珀对。因此，二维 $S=\frac{1}{2}$ 正方晶格海森伯反铁磁体的磁性质可能与
高 $T_{\mathrm{c}}$ 超导电性问题相关——尽管这一论断未被证明：写作之时，尚无
公认的高 $T_{\mathrm{c}}$ 超导理论。

二维 $S=\frac{1}{2}$ 正方晶格海森伯反铁磁体因此备受关注，吸引了大量理论与实验
研究。温度降低时，系统并不有序，但短程有序的关联区域尺寸增大。这一尺寸称为自旋--
自旋关联长度 $\xi$，低温下指数发散，导致 T=0 的真正长程序。通过把二维 $S=\frac{1}{2}$ 正方晶格海森伯反铁磁体映射到一个称为二维量子非线性
sigma 模型的连续模型，可以计算 $\xi$ 随温度的行为。这一步骤超出本书范围，但为这一方法提供了大量支持的
实验结果示于图 8.13。结果由 $Sr_2CuO_2Cl_2$ 的非弹性中子散射得到——它是理想
二维 $S=\frac{1}{2}$ 正方晶格海森伯反铁磁体的极好近似：面间距比 $La_2CuO_4$
大，且直到低温保持四角。面间相互作用导致 $T_{\mathrm{N}}=256.5$ K 的长程三维
反铁磁有序，但远高于此温度时自旋涨落由二维涨落主导。三维有序在非弹性中子散射
数据中给出尖锐的布拉格峰，而关联长度 $\xi$ 的短程有序给出宽度正比于 $\xi^{-1}$
的宽峰（长程有序时 $\xi\to\infty$，峰变得非常尖锐）。实验数据显示峰随材料冷却而
变窄；不同入射能量的数据与基于二维量子非线性 sigma 模型预言的计算曲线一同示于
图 8.13。该理论除从拉曼散射测得的 $\mathsf{J}$ 外没有可调参数——符合程度因此
极为惊人，表明这一体系开始得到细致的理解。

\begin{figure}
\centering
\includegraphics[width=0.5\textwidth]{fig8_13.pdf}
\caption{由非弹性中子散射导出的 $\mathrm{Sr}_2\mathrm{CuO}_2\mathrm{Cl}_2$ 中
的逆关联长度 $\xi^{-1}$。数据对应不同初始中子能量 $E_i$，理论曲线基于二维量子
非线性 sigma 模型的预言。取自 Greven 等 1994。}
\end{figure}

尽管如此，把这些结果应用于高 $T_{\mathrm{c}}$ 超导电性问题仍有争议。近期诱人的
结果表明：某些体系中掺杂空穴并非随机排布，而是分凝成条纹（stripes），
条纹的排布相当精微地依赖掺杂水平。这一引人入胜的协同行为看起来是谜团的一部分——
但它是拼图的关键一块还是又一个转移视线的干扰项，仍待观察。

本节迄今处理的是维数 D=3 的海森伯自旋——尽管它们位于 d=2 的晶格上，自旋可以
指向三维空间任意方向。若磁矩被约束在平面内，我们得到二维 XY 模型（D=2）。此时
$T\neq0$ 没有长程序，但仍有一个有趣的相变可以发生：低温下自旋--自旋关联函数按
代数方式衰减，而高温下按指数衰减。两个区域之间有转变温度，与涡旋（vortex）的
热稳定性相联系。

\begin{figure}
\centering
\includegraphics[width=0.4\textwidth]{fig8_14.pdf}
\caption{XY 模型中的涡旋态。}
\end{figure}

在 XY 系统中产生一个涡旋（如图 8.14 所画）的能量近似为 $\pi\mathsf{J}\ln(R/a)$，
R 是系统尺寸、a 是晶格常数。这一能量代价在 $R\to\infty$ 时发散，看起来涡旋的
产生非常昂贵。然而涡旋中心可以位于系统 $(R/a)^{2}$ 个格点中的任何一个，故涡旋
的熵为 $S=k_{\mathrm{B}}\ln[(R/a)^{2}]$。于是涡旋的自由能为
\begin{equation*}
F = (\pi\mathsf{J} - 2k_{\mathrm{B}}T)\ln(R/a)\tag{8.3}
\end{equation*}
在温度
\begin{equation*}
T_{\mathrm{KT}} = \frac{\pi\mathsf{J}}{2k_{\mathrm{B}}},\tag{8.4}
\end{equation*}
以上变为负——这就是 Kosterlitz--Thouless 转变（Kosterlitz--Thouless
transition）：涡旋可以被热自发产生。这是一种特殊的相变，因为它发生在两个完全
无序的相之间——一个在低温、一个在高温。它伴随刚性的改变：低温态有弹性刚度而
高温态没有。这一转变因此也可以看作涡旋解缚转变（vortex unbinding
transition），因为 $T_{\mathrm{KT}}$ 以下系统中的涡旋都被强烈束缚。它被称为
拓扑相变（topological phase transition），因为没有对称性破缺（与铁磁--顺磁
转变相对照）。

若自旋被约束在一条线上，我们得到二维伊辛模型（d=2、D=1）——它确实有磁相变。
$K_2CoF_4$ 与 $Rb_2CoF_4$ 之类的材料是二维伊辛模型的近似实现，已用中子散射等
技术详细研究。

\section{量子相变}

第六章考虑的相变都由温度驱动：这类相变中，样品升温经过转变时摧毁有序的是热涨落
（其能量尺度由 $k_{\mathrm{B}}T$ 控制）。然而，若转变由其他变量控制（如压强、
磁场或掺杂水平），则在该变量的某个临界值处，可以有一个原则上能发生在绝对零度的
转变。这样的零温相变称为量子相变（quantum phase transition），其发生点称为量子
临界点（quantum critical point）。相关涨落不再是热的，而是由海森伯不确定性原理
决定的量子力学涨落。量子涨落与热涨落性质颇为不同，需要完全不同的理论来描述。
量子系统由复值波函数描述，这带来经典系统所没有的特征。

温度驱动的经典相变中，关联长度与关联时间随转变的逼近而发散：序参量涨落越来越慢、
尺度越来越大。因此存在某个与热涨落相联系的频率（设为 $\omega^{*}$），在临界
温度 $T_{\mathrm{c}}$ 处趋于零。只要临界点附近 $k_{\mathrm{B}}T_{\mathrm{c}}\gg\hbar\omega^{*}$，
临界涨落将表现得经典。因此对量子相变（$T_{\mathrm{c}}=0$），量子涨落不可忽略。
预期具有量子临界点的体系表现由量子涨落控制的反常动力学，可用 NMR 与中子散射等
技术探测。

量子相变的一个例子是伊辛磁体 $LiHoF_4$：沿垂直于伊辛自旋易轴的方向加磁场，可以
在绝对零度摧毁其铁磁有序（见图 8.15）。这一反直觉的行为之所以发生，是因为磁场
便利了上下自旋态之间的量子隧穿。超过临界磁场，量子涨落足以摧毁铁磁有序。这里的
顺磁态不同于通常的高温顺磁体——后者中自旋不停地（但经典地）在上下态之间涨落；
$LiHoF_4$ 中基态有一个唯一的相位相干波函数——上下自旋态的量子叠加。

\begin{figure}
\centering
\includegraphics[width=0.45\textwidth]{fig8_15.pdf}
\caption{$LiHoF_4$ 的相图。虚线是只计入电子自旋自由度的平均场理论；实线还纳入了
核超精细相互作用。取自 Bitko 等 1996。}
\end{figure}

有若干材料表现涉及相图中磁性与超导区域相邻的量子临界点。其中包括重费米子
\fn{2}{重费米子化合物的有效质量非常大，因而得名；它们常是 Ce 或 U 的化合物。}化合物如 CePd$_2$Si$_2$（见图 8.16(a)）——常压下是反铁磁体，施加静水压\fn{3}{静水压是各向同性施加的压强，沿所有方向相同。}后可以变为超导。20 kbar、T=0 处的相变就是一个量子临界点。

第二个例子是基于有机分子 ET\fn{4}{ET 是 BEDT-TTF 的缩写，而 BEDT-TTF 又是 bis-ethylenedithiotetrathiafulvalene（双亚乙基二硫代四硫富瓦烯）的缩写。}的一族
有机超导体。可以制备很纯的 ET 电荷转移盐晶体：一对 ET 分子共同向无机阴离子
（如 $\mathrm{Cu(NCS)_2^-}$）捐出一个电子。晶体结构由 ET 分子层与无机阴离子层
交替堆叠构成。由于 ET 分子的分子交叠，有机层面内电导率高，而层间电导低。有机
材料比无机材料软得多，静水压的效果远大于在 $CePd_2Si_2$ 中。而且，通过改变
阴离子（例如使其更臃肿），可以用化学手段施加“负压”。图 8.16(b) 表明：联合
化学压强与外加压强可以得到相似的相图——同样，反铁磁相紧邻超导相。

第三个例子已经描述过：反铁磁体 $La_2CuO_4$ 掺 Sr 后变为超导。这里控制性质的
变量是 Sr 掺杂水平 x。在临界 x 值处系统开始从反铁磁切换到超导。

相图中磁性\fn{5}{在 UGe$_2$ 中，相图的超导区紧邻铁磁区（Saxena 等 2000）。}区域与超导区域的这种邻近是如下判断的证据：在这些材料的一部分中，传递超导电性的是自旋涨落而非声子。\fn{6}{在超导电性的 BCS 理论中，电子之间的吸引配对相互作用由声子传递。这对常规超导体看来正确，但对重费米子、有机与高温超导体则不是正确的描述。}然而在写作之时，这一故事仍远未收尾。

\begin{figure}
\centering
\includegraphics[width=0.4\textwidth]{fig8_16.pdf}
\caption{（a）重费米子超导体 $\mathrm{CePd}_2\mathrm{Si}_2$ 的温度--压强相图
（取自 Mathur 等 1998）：AF=反铁磁，S=超导，M=金属。（b）有机超导体
$\kappa$-(ET)$_2$Cu(NCS)$_2$、$\kappa$-(ET)$_2$Cu[N(CN)$_2$]Br、
$\kappa$-(ET)$_2$Cu(CN)[N(CN)$_2$] 与 $\kappa$-(ET)$_2$Cu[N(CN)$_2$]Cl 的
温度--压强相图：压强轴包括化学改变单胞尺寸所致的“化学压强”效应以及常规静水压
（改编自 Kanoda 1997、Lubczynski 等 1996）。（c）铜氧化物超导体
$\mathrm{La}_{2-x}\mathrm{Sr}_x\mathrm{CuO}_4$ 的相图（改编自 Kastner 等
1998）：SG=自旋玻璃。}
\end{figure}

\section{薄膜与多层膜}

迄今考虑的是出现在晶体中的二维磁体。但也可以用表面科学技术直接生长极薄的磁性
薄膜来人为构造这类体系。用分子束外延（molecular-beam epitaxy），可以在平整
衬底上生长厚度仅单个原子层的磁性薄膜。生长在超高真空条件下进行，尽量减少污染与
氧化。衬底的晶格常数与对称性选得与被生长材料紧密匹配。不仅能生长薄膜，还能
生长多层膜：可以制备三明治结构，经由非磁隔离层的磁耦合研究将在后面一节考虑。

单层膜的磁性质预计与生长它的块体材料大不相同。膜表面的原子比块体中的原子最近邻
更少。最近邻的减少收窄电子带宽、从而提高费米能级处的态密度 $g(E_{\mathrm{F}})$，
提高满足 Stoner 判据（见 7.3 节）的机会，从而增强磁有序的倾向。因此磁性薄膜
被预期具有增强的磁矩。此外，预计 Pd 或 V（通常非磁的金属）的单层会因此变得
铁磁。

这些技术能耍的另一个把戏是：恰当选择生长薄膜的衬底，生长磁性元素的亚稳相。
在晶体衬底上生长磁性薄膜时，界面处存在的力可以把膜驱入不同的晶体学相。这一替代
相可能已经是已知的高温相或高压相，也可能是在块体中完全未知的相。室温下 Ni 通常
是面心立方（fcc），Fe 通常体心立方（bcc），Co 通常六角密堆（hcp）。然而用彼此
晶格匹配的特定衬底，可以生长亚稳晶体学相：例如亚稳的 fcc Co 可以稳定在 fcc Cu
表面，bcc Co 可以稳定在 GaAs(110) 表面。由于与晶体结构改变相联系的能量
（$\sim0.1$ eV/原子）和磁结构改变（如铁磁变反铁磁）相联系的能量同量级，薄膜的
磁性质常常急剧依赖生长条件以及衬底的结构。

薄膜表面原子对称性的改变也影响磁各向异性与易磁化方向。取最低阶，铁磁层的各向
异性能量可写为
\begin{equation*}
E_{\mathrm{an}} = K\sin^{2}\theta\tag{8.5}
\end{equation*}
K 是有效各向异性常数，为三项之和：
\begin{equation*}
K = \frac{2K_{\mathrm{s}}}{t} + K_{\mathrm{v}} - \mu_0M^{2},\tag{8.6}
\end{equation*}
$\theta$ 是磁化强度与表面法线的夹角。式 8.6 第一项代表表面各向异性：$K_{\mathrm{s}}$
是表面（或界面）各向异性常数，t 是层厚；因子 2 的出现是因为每层一般有两个面。
第二项是体积各向异性，可源于晶格应变，或因为我们有一个轴垂直于膜面的单轴单晶；
$K_{\mathrm{v}}$ 是体积各向异性常数。第三项是膜的形状各向异性（见 6.7.2 节）：
它是偶极贡献，按“磁极在平面上均匀分布”的假定计算；实践中界面粗糙度可以使其
显著低于这一计算值。

厚膜中偶极项主导，磁化强度位于膜面内。薄膜中表面（界面）项可能主导（因为
$t^{-1}$ 因子），自发磁化强度可能变得垂直于膜面。这一垂直各向异性（perpendicular
anisotropy）在记录应用中特别有用。\fn{7}{记录应用中，膜上的小区域存储信息比特。这些小区域因邻居而感受的偶极场，在磁化垂直于膜面时较低，因此比特原则上可以以更高的密度存储。不过这一解释略过了大量实际问题——它们有时会使制造商偏爱磁化位于面内的体系。}$K_{\mathrm{s}}$ 的来源是表面态对称性更低——增强了表面的磁晶各向异性。它受界面
粗糙度以及每一界面处晶格失配的影响。制备具有大垂直磁化强度的体系对磁记录技术的
发展极为重要。

已有多种实验技术用于研究薄膜的磁化强度。穆斯堡尔技术已在 3.2.3 节描述；
磁光方法将在下一节介绍。极化中子反射（polarized neutron reflection，PNR）
是自旋极化中子以掠入射从磁性多层膜上反射的技术：自旋向上与向下中子的反射率
不同，可用来推断随深度变化的磁化强度（方向与大小）。磁力显微术（MFM）中，
悬臂上的磁矩扫过样品表面：只要样品上方磁场存在梯度，针尖就会受力，因此该技术对
畴壁敏感（见 6.7.6 节）。另一种显微技术是 SEMPA（扫描电子显微术与极化分析）：
常规扫描电子显微镜研究二次电子的发射（二次电子保持其自旋极化，给出表面磁化强度
的信息）。其他技术包括克尔显微镜（偏振激光束扫过表面以对磁化强度成像）与扫描
SQUID 显微术。

\section{磁光学}

磁光效应最早由 Michael Faraday 于 1845 年研究：他证明偏振光通过置于磁场中的
玻璃块后出射时偏振面被旋转。这一效应如今称为法拉第效应（Faraday effect）。
相关现象由 John Kerr 于 1877 年发现：法拉第效应发生于透射，磁光克尔效应
（magneto-optic Kerr effect）发生于反射。John Kerr 把光从电磁铁抛光的磁极上
反射，注意到偏振面被旋转。

两种效应都与自旋--轨道相互作用相关，严格的处理需要微扰论。不过一般原理可以这样
相当简单地理解：右旋与左旋圆偏振光使材料中的电荷向相反方向旋转，因此每种偏振
对轨道角动量的贡献符号相反。磁场沿磁场方向产生自旋极化，自旋--轨道相互作用于是
给两种圆偏振带来大小相同、符号相反的能量贡献。这导致右旋与左旋偏振在材料中折射率
不同。若平面偏振光入射磁性材料，应把它看作右旋与左旋圆偏振光束之和——二者以
不同速度在材料中传播；出射时两束重新组合，但它们之间的相位滞后意味着出射束的
偏振面被旋转。

各向同性介质中介电张量成为
\begin{equation*}
\epsilon = \epsilon\begin{pmatrix} 1 & \mathrm{i}Q_z & -\mathrm{i}Q_y\\ -\mathrm{i}Q_z & 1 & \mathrm{i}Q_x\\ \mathrm{i}Q_y & -\mathrm{i}Q_x & 1 \end{pmatrix},\tag{8.7}
\end{equation*}
$\mathbf{Q}=(Q_x,Q_y,Q_z)$ 称为 Voigt 矢量，沿磁场方向排列，其大小依赖材料。
介电张量中的非对角项正是磁光效应的来源。该张量给出两个圆偏振简正模，介电常数
$\epsilon_{\pm}=\epsilon(1\pm\mathbf{Q}\cdot\hat{\mathbf{k}})$，$\hat{k}$ 是光
的传播方向。圆模式以不同速度行进、以不同方式衰减。出射波组合起来给出旋转的偏振
轴以及椭圆度（不同的衰减使出射光的偏振略呈椭圆）。
\bio{Woldemar Voigt\\（1850--1919）}

金属磁体的研究用克尔效应而非法拉第效应，因为金属吸收光、不能透射研究——除非
样品非常薄。克尔效应对磁性薄膜与表面特别有用；此时该技术称为 SMOKE（表面磁光
克尔效应，surface magneto-optic Kerr effect）。它可以在超高真空（UHV）腔内的
薄膜生长过程中使用：激光束穿过腔窗入射、从样品反射、再经另一窗口出腔。光偏振的
变化可用正交偏振器监测。

克尔效应可以在多种几何下进行（见图 8.17）：极向克尔效应（polar Kerr effect）中
磁化方向 M 垂直于膜面。若 M 在膜面内，M 在光的散射面内时可测纵向克尔效应
（longitudinal Kerr effect）；M 垂直于入射面时可测横向克尔效应（transverse
Kerr effect）。由于可在生长过程中进行，SMOKE 对表面磁性、膜的磁性随膜厚的变化
以及表面各向异性的研究很有用（例如它可以用来显示磁化在什么膜厚从面内变为面外）。

磁光效应在磁记录中极为有用：可以让偏振激光束极快地扫过盘面，通过反射光偏振的
变化读出盘上编码的磁信息。虽然用磁光效应读盘，写数据却用另一束激光把盘上微小
区域局部加热超过 $T_{\mathrm{C}}$ 并对盘加磁场；只有受热区域翻转——盘上其余
信息不变。激光束然后移到下一区域写下一位。在 21 世纪之交，磁光盘已开始具有巨大
的重要性：提供了一种便宜而便携的存储较大量数据的方法。

\begin{figure}
\centering
\includegraphics[width=0.3\textwidth]{fig8_17.pdf}
\caption{（a）极向克尔效应；（b）纵向克尔效应；（c）横向克尔效应。}
\end{figure}

\section{磁电阻}

材料在磁场 H 作用下电阻 R 的变化称为磁电阻（magnetoresistance）。磁电阻
$\Delta\rho/\rho$ 通常定义为
\begin{equation*}
\frac{\Delta\rho}{\rho} = \frac{R(H)-R(0)}{R(0)}\tag{8.8}
\end{equation*}
\mnote{我们将看到，许多有趣的磁电阻是大的负效应。把磁电阻定义为 $(R(H)-R(0))/R(H)$——即除以较小的数——可以使你的新型“神奇磁性材料”的磁电阻显得大得出奇。这种把戏在研究文献中不止一次被使用过！}
这是一个技术上有用的量，因为磁电阻传感器被广泛使用（例如读取信用卡磁条的磁场）。
自由电子气不表现磁电阻——即使有效质量各向异性；磁电阻只出现在多于一种载流子的
模型中，其高场行为依赖费米面的拓扑。该效应由开尔文勋爵（Lord Kelvin，当时名
William Thomson）于 1856 年首先发现：他检验铁样品的电阻，发现纵向加磁场时铁的
电阻增加 0.2\%，横向加场时减少 0.4\%。
\bio{William Thomson（Lord Kelvin）\\（1824--1907）}

\begin{figure}
\centering
\includegraphics[width=0.45\textwidth]{fig8_18.pdf}
\caption{铁磁体中的自旋劈裂能带。}
\end{figure}

\subsection{铁磁体的磁电阻}

铁磁体中负磁电阻的观察曾是非常令人困惑的。金属承载电流时，电子向费米面不同部位
的迁移使散射最小——电子找到跨样品的耗散最小路径。因此若迫使电子走别的路径（例如
因为加了磁场），它们会走散射更多的路径。一般预期正磁电阻。低温下、厚度小于平均
自由程的样品中有时可以观察到负磁电阻：若磁场在面内且垂直于电流方向，电子路径
描绘直径更小的轨道，表面散射减少。但在铁磁体中，负磁电阻的解释必须完全不同。

Mott 对这一问题提供了非常重要的洞见：他考虑 Ni 的输运性质（自旋劈裂能带示意于
图 8.18），在 Ni 中把组态从 $(3d^84s^2)$ 变到 $(3d^94s^1)$ 或 $(3d^{10})$ 只需
几 eV。一般把 Ni 视为 $(3d^{9.4}4s^{0.6})$。d 带很窄（这是过渡金属铁磁性的必要
条件——使 $g(E_{\mathrm{F}})$ 大、Stoner 判据满足），故 $m_{\mathrm{d}}^{*}\gg m_{\mathrm{e}}$；
s 带近乎自由，$m_{\mathrm{s}}^{*}\sim m_{\mathrm{e}}$。于是电导率
\begin{equation*}
\sigma = \frac{n_{\mathrm{s}}e^{2}\tau_{\mathrm{s}}}{m_{\mathrm{s}}^{*}} + \frac{n_{\mathrm{d}}e^{2}\tau_{\mathrm{d}}}{m_{\mathrm{d}}^{*}},\tag{8.9}
\end{equation*}
由第一项主导，导电主要来自 s 电子。式 8.9 中 $n_{\mathrm{s}}$ 与 $n_{\mathrm{d}}$
分别是 s 带与 d 带中的电子数，散射时间 $\tau_{\mathrm{s}}\sim\tau_{\mathrm{d}}$。

跃迁概率主要来自 s$\to$d 跃迁。低温（$T\ll T_{\mathrm{c}}$）下所有未占据的 d 态
都是反平行的，只有一半的 s 电子能跃迁；$T>T_{\mathrm{c}}$ 时所有 s 电子都能跃迁，
散射更多。因此预期 $T_{\mathrm{c}}$ 以下电阻率下降。现在考虑对 Ni 加磁场：磁场
可能增强自旋极化、减少 s$\to$d 跃迁——因此观察到负磁电阻。

多数弹性碰撞过程中电子保持其自旋，这些碰撞由相对较短的弛豫时间 $\tau$ 刻画。
然而，即使弱的自旋--轨道耦合也允许弱的自旋翻转散射，其弛豫时间长得多，为
$\tau_{\mathrm{sf}}$。无外力时，自旋向上与向下电子平衡分布中产生的微扰先在时间
$\tau$ 内衰减成每种自旋内的均匀分布，再在时间 $\tau_{\mathrm{sf}}$ 内达到平衡
分布。$T\ll T_{\mathrm{c}}$ 时自旋翻转散射预期不占主导，铁磁体可以很好地用双电流
模型（two-current model）近似——$\uparrow$ 与 $\downarrow$ 电子电流被独立处理。
它在描述少量一种过渡金属（杂质）溶入另一种过渡金属（宿主）的合金性质上特别成功。
某些过渡金属杂质引起的散射强烈依赖自旋——这是宿主 d 带自旋劈裂、杂质 d 能级自旋
劈裂以及宿主与杂质态在自旋 $\uparrow$ 与 $\downarrow$ 方向不同杂化共同作用的
结果。例如 Fe 中的 Cr 杂质对自旋 $\uparrow$ 电子的散射强得多，导致两自旋态电阻率
之比 $\rho_{\uparrow}/\rho_{\downarrow}\sim6$。高温下与自旋波碰撞导致的自旋翻转
散射引起自旋混合（spin-mixing）——两个自旋通道之间区别的模糊。这一概念应记在
心里，它将在三明治结构巨磁电阻的讨论中回归。

\subsection{各向异性磁电阻}

铁磁体中测得的磁电阻依赖磁化强度相对材料中电流方向的取向——这一效应称为各向
异性磁电阻（anisotropic magnetoresistance）。其起源与自旋--轨道相互作用及其对
s-d 散射的影响相关。自旋--轨道相互作用降低原子波函数的对称性并混合各态；晶轴
决定 L 的方向、磁化强度决定 S 的方向，态的混合导致各向异性散射。用对称性论证，
可以预言磁电阻对磁化强度与电流密度方向依赖的一般形式（对特定晶体或多晶材料）。

\subsection{巨磁电阻}

各向异性磁电阻效应相当小，因此 1980 年代末发生了一场小小的革命：在 Fe/Cr/Fe
多层膜中发现了非常大的效应（命名为巨磁电阻，giant magnetoresistance，GMR）。
低温强场下发现了超过 50\% 的大负磁电阻（见图 8.19）。该效应与磁 Fe 层反铁磁耦合
的样品相联系。人们发现：磁性层之间经隔离层的耦合随隔离层厚度增加而振荡变号；
某些厚度下是铁磁的，厚度更大时变为反铁磁，然后又回到反铁磁。\footnote{原书如此。按耦合随隔离层厚度振荡变号的物理，此处应为“回到铁磁”——疑为原书排印之误。——译注}振荡周期看来是十个
原子间距的量级（通常在 9 到 18 \AA{} 之间），其值主要由隔离层金属决定、与铁磁
金属无关。隔离层与铁磁晶格的良好匹配有利于大的耦合。在最好的情形（如 Co/Ru
超晶格），振荡幅度随隔离层厚度 t 以 $1/t^{2}$ 衰减。

\begin{figure}
\centering
\includegraphics[width=0.5\textwidth]{fig8_19.pdf}
\caption{Fe/Cr/Fe 多层膜中的巨磁电阻（Baibich 等 1988）。}
\end{figure}

这些振荡的起源与块体金属中相距 r 的两个局域自旋之间的 RKKY 相互作用
$\mathsf{J}(r)\sim\cos(2k_{\mathrm{F}}r)/r^{3}$ 相关。对界面上所有自旋求和后，
耦合变为 $\mathsf{J}\sim\cos(2k_{\mathrm{F}}t)/t^{2}$（t 是两铁磁层的间距）。
耦合的振荡可以直接联系于隔离层费米面的拓扑。若设铁磁金属具有全满的多数自旋 d 带，
则铁磁耦合时少数自旋空穴被束缚在隔离层内，反铁磁耦合时没有束缚。粒子数守恒下，
不同耦合之间的能量差完全由隔离层（厚度 t）中少数自旋空穴能量的尺寸量子化决定：
隔离层中的能级是分立的，态密度由一系列台阶组成，每个台阶的宽度正比于 $1/t^{2}$。
t 增大时台阶变窄，总有一个台阶最终扫过费米能级。用于计算这一效应的形式理论可以
理解为德哈斯--范阿尔芬效应的一维类比：此处尺寸量子化源于垂直于层面方向上的一维
（而非二维）束缚。

三明治结构中磁耦合的类型可以直接影响观察到的磁输运行为——后者对磁层的排列极其
敏感，反铁磁耦合时 GMR 效应最大。对 GMR 的第一个解释基于双电流模型（见上），
分别考虑 $\uparrow$ 与 $\downarrow$ 电子的各自电流（$\uparrow$ 指平行于多数
自旋带）。这一讨论中我先设平均自由程 $\lambda$ 远大于 Cr 层间厚度 $t_{\mathrm{Cr}}$。
想象一个 Fe/Cr/Fe 结构：零场中两 Fe 层的磁化反平行排列，并设
$\rho_{\downarrow}\ll\rho_{\uparrow}$。有两种情形：

（1）$H>H_{\mathrm{s}}$（$H_{\mathrm{s}}$ 为饱和场）。此时 Fe 层中的磁矩平行
排列（见图 8.20(a)），电阻率 $\rho$ 为
\begin{equation*}
\frac{1}{\rho} = \frac{1}{\rho_{\uparrow}} + \frac{1}{\rho_{\downarrow}} \Rightarrow \rho\sim\rho_{\downarrow}.\tag{8.10}
\end{equation*}
散射较少的电子形成有效短路。

（2）H=0。两铁层中的磁矩现在反平行（见图 8.20(b)）。此时电子相对局域磁化在
各层中交替地为 $\uparrow$ 与 $\downarrow$，自旋 $\uparrow$ 与 $\downarrow$ 通道
被有效地“混合”（对比上文考虑的高温合金中自旋波所致的自旋混合），于是
$\rho_{\uparrow}\to\rho_{\mathrm{av}}$、$\rho_{\downarrow}\to\rho_{\mathrm{av}}$
（$\rho_{\mathrm{av}}=(\rho_{\uparrow}+\rho_{\downarrow})/2$），总电阻率 $\rho$
为
\begin{equation*}
\rho = \frac{\rho_2+\rho_1}{4} \gg \rho_1.\tag{8.11}
\end{equation*}
这再次预言了观察到的负磁电阻。巨大效应来自 $\rho_{\uparrow}$ 与 $\rho_{\downarrow}$
的不等（可以非常大），以及“自旋混合”可以简单地通过施加磁场开启与关闭的便利。

\begin{figure}
\centering
\includegraphics[width=0.6\textwidth]{fig8_20.pdf}
\caption{（a）$H>H_{\mathrm{s}}$ 时 Fe 层中的磁矩平行。（b）铬层厚度
$t_{\mathrm{Cr}}$ 选择得当时，$H=0$ 时 Fe 层中的磁矩反平行。}
\end{figure}

对更厚的 Cr 层，界面处的自旋依赖散射影响界面附近厚度约为平均自由程 $\lambda$
的层内的电子分布函数。若 $t_{\mathrm{Cr}}\gg\lambda$，预期磁电阻大致按
$\exp(-t_{\mathrm{Cr}}/\lambda)$ 下降。但若没有反铁磁耦合就不会有巨磁电阻效应；
既然耦合随 $t_{\mathrm{Cr}}$ 周期出现，磁电阻也预期随层间厚度增大而振荡，同时
指数衰减。

通过多层膜的磁耦合可以用 8.7 节描述的许多磁性薄膜实验技术测量。有用的补充是
铁磁共振（铁磁体的 NMR 或 ESR 类比）与布里渊光散射技术（用光的非弹性散射研究
激发）。二者都可用于研究磁性薄膜中的自旋波与静磁模。多层膜中的自旋经交换、偶极
与各向异性相互作用耦合在一起。作为这一磁性系统本征模的自旋波激发，其色散关系可以
相当敏感地依赖交换耦合、各向异性与磁弹效应。

制备多层膜并非实现 GMR 的唯一途径。它在非均相 Cu-Co 合金膜中也被观察到：
Cu 基体内富 Co 晶粒中磁矩的相对取向决定磁电阻，并可由外加场改变。但 GMR 不出现
在缺乏孤立的、具有合适尺寸的大型富磁晶粒的均匀合金中。合金比多层膜容易制备，因此为传感器
应用提供了激动人心的前景。

\subsection{交换各向异性}

交换各向异性（exchange anisotropy，或单向各向异性 unidirectional anisotropy）
是可以在铁磁体与反铁磁体之间观察到的界面交换。若铁磁体的居里温度 $T_{\mathrm{C}}$
高于反铁磁体的奈尔温度 $T_{\mathrm{N}}$，则把一个沉积在另一个之上（适当地选择
铁磁体与反铁磁体，且铁磁层足够厚），并在外场中冷却通过奈尔温度后，在 $T\ll T_{\mathrm{N}}$
测得的磁滞回线看起来被移动了——好像除了外加磁场之外还有另一个磁场存在。看起来
铁磁膜沿某个方向（冷却时的方向）磁化在能量上比另一方向更有利。
\begin{figure}
\centering
\includegraphics[width=0.4\textwidth]{fig8_21.pdf}
\caption{被交换各向异性移动的磁滞回线。没有交换各向异性时回线将以零磁场为中心。}
\end{figure}
有时说它被反铁磁层交换偏置（exchange biased）了。交换偏置的效果是产生一个单向
交换场 $B_{\mathrm{ex}}$，它可以与外场 B 竞争（磁层的总能量为
$-\mathbf{M}\cdot(\mathbf{B}+\mathbf{B}_{\mathrm{ex}})$）。若 B 与
$B_{\mathrm{ex}}$ 同向，效果只是移动磁滞回线（见图 8.21）；若成直角，得到难磁化
轴响应。两种情形下，每个磁场值都存在唯一使能量极小的磁化角度。

这一技术可用于给磁电阻传感头加上“交换偏置”场，把它偏置到线性区；也用于约束
三明治结构中一个软铁磁层的磁化方向，另一层的磁化则可由外场转动，从而以严格受控
的方式测量磁电阻与耦合能量。这也是自旋阀（spin valve，见图 8.22）的基础——一种
巨磁电阻传感器：两个磁层夹非磁隔离层，其中一个磁层邻接反铁磁层。三明治结构两端
的电阻于是对沿特定方向的外磁场之值敏感。

\begin{figure}
\centering
\includegraphics[width=0.42\textwidth]{fig8_22.pdf}
\caption{自旋阀。反铁磁体（AFM）把上层铁磁层（FM1）的磁矩沿特定方向“夹持”；
下层磁层（FM2）的磁矩可自由转动、可与外加磁场排齐。两铁磁层磁矩间的夹角控制器件
的电阻。零场下，层的相对排列由非磁隔离层（NM）的厚度控制。}
\end{figure}

\subsection{庞磁电阻}

锰氧化物的输运性质近来引起了极大的兴趣。LaMnO$_3$ 含有 Mn$^{3+}$ 态的 Mn——
一种扬--特勒离子；LaMnO$_3$ 表现 A 型反铁磁有序（见 5.2.4 节）。若以二价的
Sr$^{2+}$、Ca$^{2+}$ 或 Ba$^{2+}$ 离子替换 $x$ 分数的三价 La$^{3+}$ 离子，
则在 Mn 格位引入空穴。这使 $1-x$ 分数的 Mn 离子保持为 Mn$^{3+}$（3d$^4$、
t$_{2g}^{3}$e$_g^{1}$），而 $x$ 分数变为 Mn$^{4+}$（3d$^4$、t$_{2g}^{3}$e$_g^{0}$）。\footnote{原书如此。按电子计数，Mn$^{4+}$ 应为 3d$^3$（t$_{2g}^{3}$e$_g^{0}$）——疑为原书排印之误。——译注}当 $x=0.175$ 时扬--特勒畸变消失，系统变为铁磁，居里温度在室温附近；$T_{\mathrm{C}}$ 以上材料是绝缘且非磁的，$T_{\mathrm{C}}$ 以下是金属且铁磁。尤其在 $T_{\mathrm{C}}$ 附近，材料表现出极大的磁电阻效应（见图 8.23），称为庞磁电阻（colossal magnetoresistance，缩写 CMR）——“巨”这个名字已经被占用啦！CMR 的起源与金属--绝缘体转变的存在相关。

上述行为的起源部分地与双交换现象（见 4.2.5 节）相联系。Mn$^{3+}$ 离子中 t$_{2g}$
电子紧束缚于离子，而 e$_g$ 电子是巡游的。由于双交换相互作用，e$_g$ 电子在 Mn 格位
之间的跳跃只有在两个 Mn 核心自旋排齐时才被允许（实际上跳跃概率正比于
$|\cos(\theta/2)|$，$\theta$ 是两个 Mn 核心自旋的夹角）。磁场使核心自旋排齐，
因而提高电导——尤其在 $T_{\mathrm{C}}$ 附近。

实际情况其实更复杂：由于扬--特勒效应，载流子与声子相互作用。这些体系中强的
电子--声子耦合意味着：$T_{\mathrm{C}}$ 以上载流子其实是极化子（polarons）——
伴随大晶格畸变的电子。这些极化子是磁性的、在晶格中自陷。向磁态的转变可以视为
被陷极化子的解缚。电子--声子耦合还有其他标志，包括依赖磁场的结构转变与电荷有序。
电荷有序基态可以与铁磁性竞争，并在可公度掺杂值附近被增强（例如 $x=\frac{1}{2}$
处每两个 Mn 格位有一个空穴）。这些效应尚未被完全理解，仍在积极研究之中。

\begin{figure}
\centering
\includegraphics[width=0.42\textwidth]{fig8_23.pdf}
\caption{$\mathrm{La}_{1-x}\mathrm{Sr}_x\mathrm{MnO}_3$
（x=0.175）的庞磁电阻：（a）电阻率的温度依赖；（b）等温磁电阻。取自 Tokura 等
1994。}
\end{figure}

\begin{figure}
\centering
\includegraphics[width=0.42\textwidth]{fig8_24.pdf}
\caption{锰氧化物的 Ruddlesden--Popper 相。}
\end{figure}

锰钙钛矿只是表现庞磁电阻的氧化物材料之一。钙钛矿是一个称为 Ruddlesden--Popper
相的大晶体家族的成员（见图 8.24）。这些相的一般化学式为
$X_{n+1}Mn_nO_{3n+1}$（X 是镧系元素、锶或其混合），可以看作 n 层厚的钙钛矿块
\fn{8}{钙钛矿块本质上是共角 MnO$_6$ 八面体的堆叠。}的堆叠，每块之间隔着一层岩盐型的 $(\mathrm{Sr},\mathrm{Ln})_2\mathrm{O}_2$ 层——它在电与磁两方面都倾向于解耦各块。
$n=\infty$ 时实现钙钛矿化合物，而 n=1 的情形等价于高 $T_{\mathrm{c}}$ 铜氧化物
采取的 $K_2NiF_4$ 结构。锰化合物可以被制备成一系列结构（写作之时 n=1、2、3 与
$\infty$ 都已制备）。钙钛矿（$n=\infty$）中每个 MnO$_6$ 八面体被六个八面体包围；
n=2 相中这一“配位数”降为五，n=1 相中为四。最近邻数目的减少预期会造成（主要）
源于 Mn 3d 轨道的能带宽度的各向异性缩减，从而改变这些材料的电导与磁行为。因此
这些材料可以通过使用不同掺杂剂、不同掺杂水平以及调节维数而受到精细控制。

\subsection{霍尔效应}

在导电材料中沿垂直于电流方向施加磁场，会对膜中传导电子产生横向力。这一力导致
横向霍尔电压。这称为正常霍尔效应\fn{9}{通常就简单地称为霍尔效应。}——正比于外加磁场 B，因为传导电子受到的洛伦兹力是 $\mathbf{F}=e(\mathbf{E}+\mathbf{v}\times\mathbf{B})$。铁磁体中还有一个额外效应，称为反常霍尔效应
（extraordinary Hall effect），依赖磁化强度。注意它不只是“来自 M 而非 B 的
洛伦兹力”——那已包含在正常霍尔效应中，因为 $\mathbf{B}=\mu_0(\mathbf{H}+\mathbf{M})$。
经验上横向电阻率 $\rho_{\mathrm{H}}$ 为
\begin{equation*}
\rho_{\mathrm{H}} = R_0B + \mu_0R_{\mathrm{e}}M,\tag{8.12}
\end{equation*}
$R_0$ 是正常霍尔系数，$R_{\mathrm{e}}$ 是反常霍尔系数（有时称自发霍尔系数）。
正常霍尔系数 $R_0$ 相当不依赖温度，而 $R_{\mathrm{e}}$ 通常强烈依赖温度。这一
效应不仅见于铁磁体，也见于强反铁磁体或顺磁体——它只需要局域矩的存在（例如 Tb
在其铁磁态与顺磁态中都可见）。若测霍尔电阻率随场的变化，得到直线，但达到饱和场时
梯度突然改变（见图 8.25）：低场下 B 与 M 都增大，正常与反常霍尔效应都出现，梯度
由 $R_0$ 与 $R_{\mathrm{e}}$ 共同贡献；饱和场以上 M 无法再增大，梯度只来自
$R_0$。

\begin{figure}
\centering
\includegraphics[width=0.4\textwidth]{fig8_25.pdf}
\caption{铁磁体中的霍尔电阻率。}
\end{figure}

反常霍尔效应的机制与传导电子和局域矩之间的自旋--轨道相互作用相联系。载流子相对
散射中心具有轨道角动量，其符号取决于载流子从中心的哪一侧经过。由于正比于
$\mathbf{L}\cdot\mathbf{S}$ 的自旋--轨道相互作用，这一轨道角动量 L 与散射中心的
自旋角动量 S 耦合。电子于是发现从一侧经过比从另一侧经过在能量上更有利。这导致
不对称散射，从而产生横向电流——反常霍尔效应之源。

\section{有机与分子磁体}

磁性材料传统上由本质上是无机化学产物制成。然而近来基于有机化学的有机磁性材料
引起了兴趣。碳化学的丰富结构允许对每个分子作许多小调整——由此带来的可调谐性
意味着原则上可以量身定制材料以表现期望的性质。不过这只是“原则上”，因为结构--
性质关系可能证明极其复杂。凝聚态物理中研究有机材料的动力尤其来自近期三类新材料的
惊人发现：其一，导电聚合物——电导率可高达铜等常规金属，已被用于制作聚合物晶体管
与发光二极管器件；其二，超导电荷转移盐——超导转变温度可高达 15 K；其三，碳的
“足球烯”C$_{60}$——碳的一种新同素异形体，适当掺杂后可表现超导性甚至不寻常的
磁性质。

铁磁体即使在元素中也相当罕见，且全部出现在 d 或 f 区。要在有机材料上获得磁矩，
自然的做法是寻找具有未成对自旋的有机自由基。有机自由基很多，但足够稳定、能组装
成晶体结构的很少；而且即使可能，让这些自旋铁磁排列通常也不可能。因此 1990 年代
初在某些硝酮氮氧自由基（nitronyl nitroxide）分子晶体中发现铁磁性（尽管温度相当
低）特别引人瞩目。找到的第一个这类材料是对硝基苯基硝酮氮氧自由基（p-NPNN）
（见图 8.26）：仅在其一个晶相中表现直到 $T_{\mathrm{C}}\sim0.65$ K 的铁磁性。
硝酮氮氧自由基只含 C、H、N、O 元素，因此是完全有机的磁体。每个硝酮氮氧自由基
分子上有一个与两个 N--O 基团相联系的未成对自旋。分子其余部分的微小化学改变会
显著改变晶体结构，从而改变分子间交叠以及相邻分子上未成对自旋之间的磁相互作用。
因此不同化合物具有大不相同的磁基态。这些材料的转变温度大多低于 1 K。

一些在技术上最有前景的材料是分子磁体：过渡金属离子提供局域矩，有机桥连充当交换
通路。在金属离子与有机分子上都带未成对电子的材料方面也取得了进展。它们的转变
温度远高于纯有机磁体，有时高于室温。这些材料有前景的一面是光学性质：其中一些
是透明的，且变磁时变色——常与自旋转变相联系。因此它们在需要用肉眼判断磁性变化
的场合（例如电话卡或显示应用）有许多潜在应用。

\begin{figure}
\centering
\includegraphics[width=0.45\textwidth]{fig8_26.pdf}
\caption{有机铁磁体对硝基苯基硝酮氮氧自由基（p-NPNN）的分子结构。}
\end{figure}

\section{自旋电子学}

操控与放大不同自旋类型之电流的能力，是新兴的自旋电子学（spin electronics，
又称 spintronics）领域的基础——它试图用金属铁磁体、非磁金属与半导体的组合制造
器件。这一领域的基础想法都与如下事实相关：自旋向上与向下电子在铁磁体中具有不同
的迁移率。关键特征是给定铁磁体中的自旋极化程度。自旋极化 P 定义为
\begin{equation*}
P = \left|\frac{n_{\uparrow}-n_{\downarrow}}{n_{\uparrow}+n_{\downarrow}}\right|,\tag{8.13}
\end{equation*}
对元素金属铁磁体小于 1（Fe 为 0.44，Co 为 0.34，Ni 为 0.11）。但某些材料 P=1——
意味着有一个自旋劈裂带完全空着。这类材料称为半金属磁体（half-metallic magnets，
因为一种自旋的电子是金属性的、另一种自旋的是绝缘性的）；$CrO_2$、$Fe_3O_4$ 与
一些锰氧化物据信属于此类。这类材料具有很长的自旋扩散长度，可能是自旋电子学应用
中令人兴奋的材料——例如磁隧道结中用自旋极化电子隧穿过绝缘势垒。

近来出现了几种自旋晶体管的设计尝试：利用不同自旋的电子从铁磁体到半导体层的
隧穿。传统电子学只基于电子电荷、无视电子自旋。当前人们颇为乐观：在本世纪，用
磁性材料控制电子电流的自旋，可能是电子器件的有前景的新发展。

\section*{延伸阅读}

\begin{itemize}[itemsep=0.2em]
\item J.~A.~Mydosh,《Spin glasses: an experimental introduction》, Taylor and
Francis 1993 与 K.~H.~Fischer 及 J.~A.~Hertz,《Spin glasses》, CUP 1991，都包含
自旋玻璃理论与实验工作的有用论述。
\item E.~Dagotto, Reports on Progress in Physics \textbf{62}, 1525 (1999) 是
关于自旋梯的可读综述。
\item A.~M.~Tsvelik,《Quantum field theory in condensed matter physics》, CUP
1995，是包含描述低维磁性现行理论的高级教材。
\item A.~Auerbach,《Interacting electrons and quantum magnetism》,
Springer-Verlag 1994，对低维磁体的量子力学与 Haldane 能隙有生动的处理。这些专题
在 I.~Affleck, J.~Phys.: Condensed Matter, \textbf{1}, 3047 (1989) 中亦有综述。
\item M.~A.~Kastner, R.~J.~Birgeneau, G.~Shirane 与 Y.~Endoh, Rev.~Mod.~Phys.,
\textbf{70}, 897 (1998) 综述了单层铜氧化物的磁性质与高 $T_{\mathrm{C}}$ 问题。
\item E.~Manousakis, Rev.~Mod.~Phys., \textbf{63}, 1 (1991) 综述了正方晶格上
二维海森伯反铁磁体的理论。
\item P.~M.~Chaikin 与 T.~C.~Lubensky,《Principles of condensed matter
physics》, CUP 1995，是凝聚态体系统计力学的优秀信息来源。
\item 量子相变的综述见 S.~L.~Sondhi, S.~M.~Girvin, J.~P.~Carini 与 D.~Shahar,
Rev.~Mod.~Phys., \textbf{69}, 315 (1997)。
\item 侧重自旋波激发的磁性多层膜综述是 R.~E.~Camley 与 R.~L.~Stamps, J.~Phys.:
Condensed Matter, \textbf{5}, 3727 (1993)。
\item 有机磁性综述见 O.~Kahn,《Molecular magnetism》, VCH (1993)。
\item K.~De'Bell, A.~B.~MacIsaac 与 J.~P.~Whitehead, Rev.~Mod.~Phys., \textbf{72},
225 (2000) 论述了薄膜磁体中的偶极效应。
\item 磁学研究各专题的当前状况可从重要磁学会议的会议录中了解——它们不时发表在
Journal of Applied Physics、Journal of Magnetism and Magnetic Materials 与
Physica B 上。磁学研究各方面的最新观点发表在 Journal of Magnetism and Magnetic
Materials 的第 $100\times n$ 卷（n 为整数）。
\end{itemize}

\section*{本章参考文献}

\begin{itemize}[itemsep=0.2em]
\item M.~N.~Baibich, J.~M.~Broto, A.~Fert, F.~Nguyen van Dau, F.~Petroff,
P.~Etienne, G.~Creuzet, A.~Friederich 与 J.~Chazelas, Phys.~Rev.~Lett., \textbf{61},
2472 (1988)。
\item D.~Bitko, T.~F.~Rosenbaum 与 G.~Aeppli, Phys.~Rev.~Lett. \textbf{77}, 940
(1996)。
\item J.~W.~Bray, L.~V.~Interrante, I.~S.~Jacobs, J.~C.~Bonner, 载于 Extended
linear chain compounds 3, p.353, J.~S.~Miller 编, New York Plenum Press (1983)。
\item B.~R.~Coles, B.~V.~B.~Sarkissian 与 R.~H.~Taylor, Phil.~Mag. \textbf{37},
489 (1978)。
\item J.~des Cloizeaux 与 J.~J.~Pearson, Phys.~Rev. \textbf{128}, 2131 (1962)。
\item M.~Greven, R.~J.~Birgeneau, Y.~Endoh, M.~A.~Kastner, B.~Keimer, M.~Matsuda,
G.~Shirane 与 T.~R.~Thurston, Phys.~Rev.~Lett. \textbf{72}, 1096 (1994)。
\item K.~Kanoda, Hyp.~Int., \textbf{104}, 235 (1997)。
\item M.~A.~Kastner, R.~J.~Birgeneau, G.~Shirane 与 Y.~Endoh, Rev.~Mod.~Phys.
\textbf{70}, 897 (1998)。
\item B.~W.~Lovett, S.~J.~Blundell, F.~L.~Pratt, Th.~Jest\"{a}dt, W.~Hayes,
S.~Tagaki 与 M.~Kurmoo, Phys.~Rev.~B, \textbf{61}, 12241 (2000)。
\item W.~Lubczynski, S.~V.~Demishev, J.~Singleton, J.~M.~Caulfield, L.~du Croo de
Jongh, C.~J.~Kepert, S.~J.~Blundell, W.~Hayes, M.~Kurmoo 与 P.~Day, J.~Phys.:
Condens.~Matter, \textbf{8}, 6005 (1996)。
\item N.~D.~Mathur, F.~M.~Grosche, S.~R.~Julian, I.~R.~Walker, D.~M.~Freye,
R.~K.~W.~Haselwinner 与 G.~G.~Lonzarich, Nature, \textbf{394}, 39 (1998)。
\item S.~Nagata, P.~H.~Keesom 与 H.~R.~Harrison, Phys.~Rev.~B \textbf{19}, 1633
(1979)。
\item S.~S.~Saxena, P.~Agarwal, K.~Ahilan, F.~M.~Grosche, R.~K.~W.~Haselwimmer,
M.~J.~Steiner, E.~Pugh, I.~R.~Walker, S.~R.~Julian, P.~Monthoux, G.~G.~Lonzarich,
A.~Huxley, I.~Sheikin, D.~Braithwaite 与 J.~Flouquet, Nature, \textbf{406}, 587
(2000)。
\item D.~A.~Tennant, T.~G.~Perring, R.~A.~Cowley 与 S.~E.~Nagler, Phys.~Rev.~Lett.
\textbf{70}, 4003 (1993)。
\item Y.~Tokura, A.~Urushibara, Y.~Moritomo, T.~Arima, A.~Asamitsu, G.~Kido 与
N.~Furukawa, J.~Phys.~Soc.~Jpn., \textbf{63}, 3391 (1994)。
\end{itemize}


% ---------------- 附录 ----------------
\appendix
\ctexset{chapter/name={附录,}, chapter/number=\Alph{chapter}}
\renewcommand{\chaptermark}[1]{\markboth{附录\Alph{chapter}\quad #1}{}}
%% 附录 A 电磁学中的单位（Units in electromagnetism）
\chapter{电磁学中的单位}

电磁学中的单位是一片潜在的地雷阵。本书通篇使用 SI 单位制，但二十世纪的许多文献
使用非 SI 记号。因此本附录试图解开旧单位制的若干谜团，供在新单位制下长大、需要
解读较老文献的读者使用。

从十九世纪末到第二次世界大战结束后不久，cgs 单位制是关心此类事务的国际委员会
偏爱的单位制。在 cgs 制中，距离以厘米量度、质量以克、时间以秒——cgs 之名由此
而来。SI 制中我们用米（100 倍于厘米）、千克（1000 倍于克）与秒。距离与质量单位
的这一选择差异看似无关痛痒，最终却在其他单位中引起巨大变化。例如，力的量纲是
质量×长度/时间$^{2}$，故 cgs 制中力的单位达因比 SI 力的单位牛顿小
$1000\times100=10^{5}$ 倍。cgs 的能量单位是尔格（1 达因·厘米），比 SI 能量单位
焦耳小 $10^{5}\times10^{2}=10^{7}$ 倍。

不幸的是事情还没完：一谈到磁场的定义，情况就急剧变糟。原因在于两个单位制对场
B、H、M 的定义选择不同：
\begin{equation*}
\mathbf{B} = \mu_0(\mathbf{H}+\mathbf{M}) \qquad\text{（SI）}
\end{equation*}
\begin{equation*}
\mathbf{B} = \mathbf{H} + 4\pi\mathbf{M} \qquad\text{（cgs）}\tag{A.1}
\end{equation*}
于是在 cgs 制中 B 与 H 量纲相同，但单位却令人困惑地不同（前者以高斯（G）计，后者
以奥斯特（Oe）计）。H 与 M 关系的差异源于两种单位制对“距磁单极 q 为 r 处磁场
强度 H 的大小”的定义方程作了不同选择：
\begin{equation*}
H = \frac{q}{4\pi r^{2}} \quad\text{（SI）}\qquad H = \frac{q}{r^{2}} \quad\text{（cgs）}\tag{A.2}
\end{equation*}
省去 $4\pi$ 因子使这一方程以及磁偶极产生磁场的类似方程在 cgs 制中相当简洁。然而
$4\pi$ 因子并未从理论中完全消失，而是冒了出来，出现在式 A.1 中一个更不受欢迎的
地方。cgs 制按式 A.2 用并不存在的磁单极定义 H：一个单位磁单极在 1 厘米远处产生
1 奥斯特的场（若那里恰好潜伏着另一个单位磁单极，它将受到 1 达因的力）。SI 制中
H 的定义更有用——用电流来定义（因为人们在实验室里通常用电流源产生磁场，而不是
费心安排单位磁单极）。于是每米 n 匝、载流 I 安培的长螺线管在 SI 制中产生的场为
$H=nI$。

SI 制的逻辑是：让 $4\pi$ 出现在式 A.2 这类方程中，提醒我们点磁单极的场具有球对称
性（$4\pi$ 与球的表面积 $4\pi r^{2}$ 有关）。距点电荷 r 处的电场 E 也是如此：
\begin{equation*}
E = \frac{q}{4\pi\epsilon_0 r^{2}} \quad\text{（SI）}\qquad E = \frac{q}{r^{2}} \quad\text{（cgs）},\tag{A.3}
\end{equation*}
其中适用类似的论证。采用 SI，就必须忍受几个由对称性论证自然产生的额外因子出现在
点电荷与偶极子的场定义方程中；其余方程则干净得多、人为因子更少。

cgs 制唯一可说的好处是电场与磁场具有相同的量纲，某些相对论性方程因此看起来更
顺眼；且自由空间中 B=H 也有某种对称性（尽管这会带来混淆的机会）。但 cgs 制中
额外的 c 因子撒得到处都是：例如，电荷在电磁场中以速度 v 运动时的洛伦兹力方程
变为
\begin{equation*}
\mathbf{F} = q(\boldsymbol{\varepsilon} + \mathbf{v}\times\mathbf{B}) \quad\text{（SI）}
\end{equation*}
\begin{equation*}
\mathbf{F} = q\left(\boldsymbol{\varepsilon} + \frac{\mathbf{v}\times\mathbf{B}}{c}\right) \quad\text{（cgs）}.\tag{A.4}
\end{equation*}

更糟糕的还在后面。事实上有两种在 cgs 中建立电磁学的途径：静电单位（esu）或电磁
单位（emu），两者之间的换算需要转换因子。磁学中通常采用 emu 制。emu 的磁矩单位
通常就缩写为“emu”，是一个量纲为体积的单位（cgs 中为 cm$^{3}$）。磁化率
$\chi=M/H$ 在两种单位制中都无量纲，但在 cgs 中常写作 emu\,cm$^{-3}$——即使乍看
不像，它仍然是无量纲的。在 cgs 磁化率与 SI 磁化率之间换算时总有一个恼人的
$4\pi$ 因子要记住——它源于 cgs 制中的 $\mathbf{B}=\mathbf{H}+4\pi\mathbf{M}$
（见上）。例如退磁因子 N 无量纲，在 SI 中 $0<N<1$，在 cgs 中 $0<N<4\pi$。
遗憾的是 emu 在磁学研究的某些领域仍被广泛使用。表 A.1 总结了 SI 与 cgs 之间
相互换算最有用的结果。

玻尔磁子虽不是 SI 单位，$\mu_{\mathrm{B}}=9.274\times10^{-24}$ J\,T$^{-1}$ 却是
磁矩的有用量度——它对应氢原子 1s 电子的磁矩。对顺磁体，摩尔磁化率 $\chi_{\mathrm{m}}$
由居里定律给出（把式 2.44 稍作改写），SI 单位下为
\begin{equation*}
\chi_{\mathrm{m}} = \frac{\mu_0N_{\mathrm{A}}\mu_{\mathrm{eff}}^{2}\mu_{\mathrm{B}}^{2}}{3k_{\mathrm{B}}T}\tag{A.5}
\end{equation*}
$N_{\mathrm{A}}$ 是阿伏伽德罗常数。于是 $\chi_m T$ 不依赖温度（见图 2.10），可以把
它与有效磁矩联系起来。整理式 A.5 得
\begin{equation*}
\mu_{\mathrm{eff}} = \left[3k_{\mathrm{B}}/\mu_0N_{\mathrm{A}}\mu_{\mathrm{B}}^{2}\right]^{1/2}\sqrt{\chi_{\mathrm{m}}T},\tag{A.6}
\end{equation*}

\begin{table}[!ht]
\centering
\caption{SI 单位制与 cgs 单位制中的单位。缩写：m=米，g=克，N=牛顿，J=焦耳，
T=特斯拉，G=高斯，A=安培，Oe=奥斯特，Wb=韦伯，Mx=麦克斯韦。emu 是电磁单位
（electromagnetic unit）的缩写。注意磁化率在 SI 单位中无量纲。}
\begin{small}
\begin{tabular}{llcll}
\toprule
量 & 符号 & \multicolumn{3}{c}{换算} \\
\midrule
长度 & $x$ & $10^{-2}$ m & $=1$ & cm\\
质量 & $m$ & $10^{-3}$ kg & $=1$ & g\\
力 & $F$ & $10^{-5}$ N & $=1$ & dyne\\
能量 & $E$ & $10^{-7}$ J & $=1$ & erg\\
磁感应强度 & $\mathbf{B}$ & $10^{-4}$ T & $=1$ & G\\
磁场强度 & $\mathbf{H}$ & $10^{3}/4\pi$ A\,m$^{-1}$ & $=1$ & Oe\\
磁矩 & $\mu$ & $10^{-3}$ J\,T$^{-1}$ 或 A\,m$^{2}$ & $=1$ & erg\,G$^{-1}$\\
磁化强度（单位体积的矩） & $\mathbf{M}$ & $10^{3}$ A\,m$^{-1}$ 或 J\,T$^{-1}$\,m$^{-3}$ & $=1$ & Oe\\
磁化率 & $\chi$ & $4\pi$ & $=1$ & emu\,cm$^{-3}$\\
摩尔磁化率 & $\chi_m$ & $4\pi\times10^{-6}$ m$^{3}$\,mol$^{-1}$ & $=1$ & emu\,mol$^{-1}$\\
质量磁化率 & $\chi_g$ & $4\pi\times10^{-3}$ m$^{3}$\,kg$^{-1}$ & $=1$ & emu\,g$^{-1}$\\
磁通量 & $\phi$ & $10^{-8}$ T\,m$^{2}$ 或 Wb & $=1$ & G\,cm$^{-2}$ 或 Mx\\
退磁因子 & N & \multicolumn{2}{c}{$0<N<1$} & $0<N<4\pi$\\
\bottomrule
\end{tabular}
\end{small}
\end{table}

从而
\begin{equation*}
\mu_{\mathrm{eff}} = 797.8\sqrt{\chi_{\mathrm{m}}^{\mathrm{SI}}T}\tag{A.7}
\end{equation*}
\begin{equation*}
\mu_{\mathrm{eff}} = 2.827\sqrt{\chi_{\mathrm{m}}^{\mathrm{cgs}}T}\tag{A.8}
\end{equation*}
其中 $\mu_{\mathrm{eff}}$ 以每化学式单元玻尔磁子计，$\chi_{\mathrm{m}}^{\mathrm{SI}}$
以 m$^{3}$\,mol$^{-1}$ 计，$\chi_{\mathrm{m}}^{\mathrm{cgs}}$ 以 emu\,mol$^{-1}$
计。这些数值关系对从 $\chi_{\mathrm{m}}T$ 对 T 的图中提取有效磁矩很有用。

\section*{延伸阅读}

\begin{itemize}[itemsep=0.2em]
\item B.~I.~Bleaney 与 B.~Bleaney,《Electricity and Magnetism》, OUP 1989。
\item J.~D.~Jackson,《Classical Electrodynamics》, Wiley 1962。
\end{itemize}

%% 附录 B 电磁学（Electromagnetism）
\chapter{电磁学}

本附录简要回顾电磁学中的若干关键结果。这些结果的详细推导请查阅延伸阅读所列
文献。

\section{磁矩}

磁通密度为 $\mathbf{B}$ 的磁场中，载流 I、长 $\mathrm{d}\mathbf{r}$ 的直导线段
所受洛伦兹力 $\mathrm{d}\mathbf{F}$ 为 $\mathrm{d}\mathbf{F}=I\,\mathrm{d}\mathbf{r}\times\mathbf{B}$
（见图 B.1）。

\begin{figure}
\centering
\includegraphics[width=0.42\textwidth]{figB_01.pdf}
\caption{电流元所受的力为 $\mathrm{d}\mathbf{F}=I\,\mathrm{d}\mathbf{r}\times\mathbf{B}$。}
\end{figure}

\begin{figure}
\centering
\includegraphics[width=0.42\textwidth]{figB_02.pdf}
\caption{取向垂直于 z 方向的正方形元电流环（边长 $\mathrm{d}x$ 与 $\mathrm{d}y$）。}
\end{figure}

这一结果可用于计算图 B.2 中元电流环所受的力偶：
\begin{equation*}
\begin{aligned}
\mathrm{d}G_x &= -I\,\mathrm{d}x\,\mathrm{d}y\,B_y\\
\mathrm{d}G_y &= I\,\mathrm{d}x\,\mathrm{d}y\,B_x,
\end{aligned}\tag{B.1}
\end{equation*}
于是定义环的磁矩 $\mathrm{d}\mu_z=I\,\mathrm{d}x\,\mathrm{d}y$，并对环的任意取向
重复该论证，力偶 $\mathrm{d}\mathbf{G}$ 为
\begin{equation*}
\mathrm{d}\mathbf{G} = \mathrm{d}\boldsymbol{\mu}\times\mathbf{B}.\tag{B.2}
\end{equation*}
对有限大小的环，积分得
\begin{equation*}
\mathbf{G} = \boldsymbol{\mu}\times\mathbf{B}.\tag{B.3}
\end{equation*}
只有磁场非均匀时环上才有净力。图 B.2 中元环在 x 方向所受的力为
\begin{equation*}
\mathrm{d}F_x = I\,\mathrm{d}y\left(B_z+\frac{\partial B_z}{\partial x}\mathrm{d}x\right) - I\,\mathrm{d}y\,B_z = \mathrm{d}\mu_z\frac{\partial B_z}{\partial x},\tag{B.4}
\end{equation*}
故
\begin{equation*}
\mathrm{d}\mathbf{F} = (\mathrm{d}\boldsymbol{\mu}\cdot\nabla)\mathbf{B}\tag{B.5}
\end{equation*}
对有限大小的磁矩积分得
\begin{equation*}
\mathbf{F} = (\boldsymbol{\mu}\cdot\nabla)\mathbf{B}.\tag{B.6}
\end{equation*}
因此只有非均匀磁场才能对磁矩施加力。

与磁场 B 排齐的磁矩 $\boldsymbol{\mu}$ 所受力矩为零（由式 B.3，$\boldsymbol{\mu}$
与 B 平行时 G=0）。因此只能通过施加力矩使 $\boldsymbol{\mu}$ 相对 B 转开。力矩
把 $\boldsymbol{\mu}$ 转到与磁场成 $\theta$ 角所做的功 W 为
\begin{equation*}
W = \int_{0}^{\theta}G\,\mathrm{d}\theta' = \int_{0}^{\theta}\mu B\sin\theta'\,\mathrm{d}\theta' = \mu B(1-\cos\theta).\tag{B.7}
\end{equation*}
磁矩的势能 U 因此可写为 $-\mu B\cos\theta$（略去常数项）。故可写
\begin{equation*}
U = -\boldsymbol{\mu}\cdot\mathbf{B}.\tag{B.8}
\end{equation*}
非均匀场 B 中磁矩 $\boldsymbol{\mu}$ 沿 x 方向所受的力按式 B.6 为
$\mu\,\partial B/\partial x$。要阻止它沿 x 运动，须施加恢复力
$-\mu\,\partial B/\partial x$。若允许它移动 $\delta x$ 进入更强的场区，该力做的功为
\begin{equation*}
\mathrm{d}W = -\mu\frac{\partial B}{\partial x}\delta x = -\mu\,\delta B,\tag{B.9}
\end{equation*}
这是样品自由能的增量（与习题 2.10 的结果一致）。

原点处磁矩产生的磁场为
\begin{equation*}
\mathbf{B}(\mathbf{r}) = -\mu_0\nabla\left(\frac{\boldsymbol{\mu}\cdot\mathbf{r}}{4\pi r^{3}}\right)\tag{B.10}
\end{equation*}
磁矢势（取最方便的规范）为
\begin{equation*}
\mathbf{A}(\mathbf{r}) = \frac{\mu_0}{4\pi}\frac{\boldsymbol{\mu}\times\mathbf{r}}{r^{3}}.\tag{B.11}
\end{equation*}

\section{自由空间中的麦克斯韦方程组}

自由空间中的麦克斯韦方程组为
\begin{align*}
\nabla\cdot\boldsymbol{\varepsilon} &= \rho/\epsilon_0\tag{B.12}\\
\nabla\cdot\mathbf{B} &= 0\tag{B.13}\\
\nabla\times\boldsymbol{\varepsilon} &= -\frac{\partial\mathbf{B}}{\partial t}\tag{B.14}\\
\nabla\times\mathbf{B} &= \mu_0\mathbf{J} + \epsilon_0\mu_0\frac{\partial\boldsymbol{\varepsilon}}{\partial t},\tag{B.15}
\end{align*}
描述电场 $\boldsymbol{\varepsilon}$、磁感应强度 B、电荷密度 $\rho$ 与电流密度
$\mathbf{J}$ 之间的相互关系。式 B.12 表明电场从正电荷发散、向负电荷汇聚；电荷
密度因此是电场的源或汇（见图 B.3）。式 B.13 表明磁场没有这种发散：不存在磁荷
（磁单极），B 场线只能构成回路，永不在任何地方起始或终止。式 B.14 表明：只有
在磁场随时间变化的空间区域周围才有电场回路——这引出电磁感应的法拉第定律。
不存在变化电场时，式 B.15 表明电流周围存在磁感应的回路。注意本节一切都是对
自由空间中的电磁场而言。
\bio{James Clerk Maxwell\\（1831--1879）}

\begin{figure}
\centering
\includegraphics[width=0.42\textwidth]{figB_03.pdf}
\caption{电场从正电荷发散、向负电荷汇聚。}
\end{figure}

\section{自由电流与束缚电流}

把物质纳入图景之后，事情变得稍微复杂，值得弄清原委。这里只讨论对磁性的效应，
尽管相应的电场推导背后是类似的思路。真实物质中含有一种与我们熟悉的、沿导线流动
的电流不同类型的电流：原子中含有电子绕核运动形成的微观束缚电流（bound
currents）。于是式 B.15 中的总电流 J 可以分解为两个分量
\begin{equation*}
\mathbf{J} = \mathbf{J}_{\text{free}} + \mathbf{J}_{\text{bound}},\tag{B.16}
\end{equation*}
分别来自自由电流（沿铜线流动的 10 安培）与束缚电流（绕原子旋转的电子）。

若 M 在样品中均匀，流过的唯一净束缚电流绕材料边缘。若 M 非均匀，材料内部会
产生束缚电流。我们需要束缚电流与磁化强度之间的一般关系，即
\begin{equation*}
\mathbf{J}_{\text{bound}} = \nabla\times\mathbf{M}.\tag{B.17}
\end{equation*}
现在概述式 B.17 的推导（更多细节见延伸阅读）。原点处单个点偶极子的矢势由式 B.11
给出。考虑体积 V 内的磁化样品，位置 $\mathbf{r}'$ 处磁化强度为
$\mathbf{M}(\mathbf{r}')$。位置 $\mathbf{r}$ 处的矢势为
\begin{equation*}
\mathbf{A}(\mathbf{r}) = \frac{\mu_0}{4\pi}\int_V\mathrm{d}^{3}r'\,\frac{\mathbf{M}(\mathbf{r}')\times(\mathbf{r}-\mathbf{r}')}{|\mathbf{r}-\mathbf{r}'|^{3}}.\tag{B.18}
\end{equation*}
利用 $\nabla_r(1/r)=-\mathbf{r}/r^{3}$、矢量恒等式
\begin{equation*}
\nabla\times(f\mathbf{g}) = f\nabla\times\mathbf{g} - \mathbf{g}\times\nabla f,\tag{B.19}
\end{equation*}
（f 与 g 分别是标量与矢量函数）并把对 $\mathbf{r}'$ 的积分变换为对 $\mathbf{r}$
的积分，此方程可以简化。结果是位置 $\mathbf{r}$ 处的矢势可写为
\begin{align*}
\mathbf{A}(\mathbf{r}) &= \frac{\mu_0}{4\pi}\int_V\frac{\nabla\times\mathbf{M}\,\mathrm{d}\tau}{r} + \frac{\mu_0}{4\pi}\int_S\nabla\times\left(\frac{\mathbf{M}}{r}\right)\mathrm{d}\tau\tag{B.20}\\
&= \frac{\mu_0}{4\pi}\int_V\frac{\nabla\times\mathbf{M}\,\mathrm{d}\tau}{r} + \frac{\mu_0}{4\pi}\int_V\frac{\mathbf{M}\times\mathrm{d}\mathbf{S}}{r}\tag{B.21}\\
&= \frac{\mu_0}{4\pi}\int_V\frac{\mathbf{J}_{\text{bound}}\,\mathrm{d}\tau}{r} + \frac{\mu_0}{4\pi}\int_S\frac{\mathbf{K}_{\text{bound}}\,\mathrm{d}S}{r}\tag{B.22}
\end{align*}
第一个积分遍及样品体积，第二个积分遍及样品表面。式 B.22 中的束缚电流与式 B.17
一致。面电流为 $\mathbf{K}_{\text{bound}}=\mathbf{M}\times\hat{\mathbf{n}}$，
矢量 $\hat{\mathbf{n}}$ 垂直于表面。

式 B.15 的问题在于：B 的旋度不仅联系自由电流密度，还联系更棘手的束缚电流密度。
这促使我们定义一个只“看见”自由电流密度的新“修正”场——尽管如下所示这只是
部分成功。利用束缚体电流的结果，定义新场 H（称为磁场强度，或干脆叫 H 场）：
\begin{equation*}
\mathbf{B} = \mu_0(\mathbf{H}+\mathbf{M})\tag{B.23}
\end{equation*}
于是可以写出只含 $\mathbf{J}_{\text{free}}$ 的式 B.15 的新版本：不存在变化电场时
\begin{equation*}
\nabla\times\mathbf{H} = \mathbf{J}_{\text{free}}.\tag{B.24}
\end{equation*}
这很有用——它把 H 场联系到能在安培计上读出的量（$\mathbf{J}_{\text{bound}}$
可不能）。遗憾的是，这一做法简化了一条麦克斯韦方程，却复杂化了另一条：式 B.13
从简单的 $\nabla\cdot\mathbf{B}=0$ 变成略微更绕弯的
\begin{equation*}
\nabla\cdot\mathbf{H} = -\nabla\cdot\mathbf{M}.\tag{B.25}
\end{equation*}
于是 H 与 B 不同，并非无散的（用数学行话说不是螺线的）。H 的表现就好像磁单极
（即电荷的磁对应物）存在一样。

\section{物质中的麦克斯韦方程组}

存在物质时麦克斯韦方程组于是成为
\begin{align*}
\nabla\cdot\mathbf{D} &= \rho_{\mathrm{free}}\tag{B.26}\\
\nabla\cdot\mathbf{B} &= 0\tag{B.27}\\
\nabla\times\boldsymbol{\varepsilon} &= -\frac{\partial\mathbf{B}}{\partial t}\tag{B.28}\\
\nabla\times\mathbf{H} &= \mathbf{J}_{\text{free}} + \frac{\partial\mathbf{D}}{\partial t},\tag{B.29}
\end{align*}
自由电荷与束缚电荷之间作了类似的区分，电位移 D 由
\begin{equation*}
\mathbf{D} = \epsilon_0\mathbf{E} + \mathbf{P}\tag{B.30}
\end{equation*}
给出，P 是电极化强度。

\section{边界条件}

\fn{1}{在物理学文献中它们传统上被称为“药盒”（pill-boxes）——尽管药丸早已不再
用那种形状的盒子出售了！}
B 场与 H 场的一个重要差别在它们遇到两个不同区域的边界时显现。区域 1 与 2 之间的
边界条件（见图 B.4）可推导如下。考虑跨在表面上的鞋油盒状的盒子。\fn{2}{我们还要确保这一点：让盒子的厚度在半径趋于零之前先趋于零。}$\nabla\cdot\mathbf{B}=0$
结合散度定理意味着 $\int_S\mathbf{B}\cdot\mathrm{d}\mathbf{S}=0$。由于盒子的面积
主要由其顶面与底面代表，B 的垂直分量连续，即
\begin{equation*}
\mathbf{B}_{\perp 1} = \mathbf{B}_{\perp 2}.\tag{B.31}
\end{equation*}
这一论证对 $\mathbf{H}$ 不适用，因为 $\nabla\cdot\mathbf{H}=-\nabla\cdot\mathbf{M}$
可能非零：若盒内磁化强度有散度（区域 1 磁化而区域 2 未磁化时完全可能如此），
垂直分量可能不连续。平行分量的论证来自方程 $\nabla\times\mathbf{H}=\mathbf{J}_{\mathrm{free}}$，
它意味着 $\oint\mathbf{H}\cdot\mathrm{d}\mathbf{l}=\int\mathbf{J}_{\mathrm{free}}\cdot\mathrm{d}\mathbf{S}$。
当图 B.4 右侧所示的安培回路收缩到零时，回路包围的自由电流也趋于零，故 $\mathbf{H}$
的平行分量连续，即
\begin{equation*}
H_{\|1} = H_{\|2}.\tag{B.32}
\end{equation*}
这一论证对 $\mathbf{B}$ 不适用：因为 $\nabla\times\mathbf{B}=\mu_0(\mathbf{J}_{\mathrm{free}}+\mathbf{J}_{\mathrm{bound}})$，
安培回路在收缩到零的过程中虽不包围自由电流，却可能在极限下包围一些束缚面电流。

\begin{figure}
\centering
\includegraphics[width=0.45\textwidth]{figB_04.pdf}
\caption{区域 1 与 2 之间的电磁边界条件。}
\end{figure}

\section*{延伸阅读}

\begin{itemize}[itemsep=0.2em]
\item B.~I.~Bleaney 与 B.~Bleaney,《Electricity and magnetism》, OUP 1989。
\item D.~J.~Griffiths,《Introduction to electrodynamics》, Prentice-Hall 1989。
\item J.~D.~Jackson,《Classical electrodynamics》, John-Wiley 1962。
\end{itemize}

%% 附录 C 量子物理与原子物理（Quantum and atomic physics）
\chapter{量子物理与原子物理}

\section{量子力学}

本节简要回顾量子力学的若干基本特征。量子力学中的基本对象是波函数（又称态函数）
$\psi$——它包含从观察获得的关于系统的全部知识。波函数 $\psi$ 是复函数，
$|\psi|^{2}=\psi\psi^{*}$ 是概率密度。\fn{1}{此处符号 $*$ 表示复共轭。}每个可观察
量都对应一个厄米算符。进行测量时，得到的结果是该算符的本征值之一。厄米算符具有
实本征值与彼此正交的本征函数。若算符 $\hat{A}$（算符都戴“帽子”）的第 $i$ 个
本征函数为 $\phi_i$（设已归一化）、本征值为 $a_i$，则
\begin{equation*}
\hat{A}\phi_i = a_i\phi_i.\tag{C.1}
\end{equation*}
测量算符 $\hat{A}$ 之后得到的期望值由
\begin{equation*}
\langle\hat{A}\rangle = \int\mathrm{d}\tau\,\psi^{*}\hat{A}\psi,\tag{C.2}
\end{equation*}
给出，$\mathrm{d}\tau$ 是体积元。波函数 $\psi$ 可以用 $\hat{A}$ 的本征函数展开：
\begin{equation*}
\psi = \sum_i c_i\phi_i,\tag{C.3}
\end{equation*}
利用 $\phi_i$ 的正交归一性，$\langle\hat{A}\rangle$ 于是为
\begin{equation*}
\langle\hat{A}\rangle = \sum_i|c_i|^{2}a_i,\tag{C.4}
\end{equation*}
得到第 $i$ 个本征值 $a_i$ 的概率是 $|c_i|^{2}$。测量是一种剧烈的过程：取波函数
$\psi$，测量与算符 $\hat{A}$ 相联系的物理量并得到结果 $a_n$（$\hat{A}$ 的第 $n$
个本征值），系统就被迫进入态函数 $\phi_n$——至少哥本哈根诠释如此。

两个算符 $\hat{A}$ 与 $\hat{B}$ 的对易子定义为
$[\hat{A},\hat{B}]=\hat{A}\hat{B}-\hat{B}\hat{A}$。定义对易子的表达式称为对易
关系。

\begin{example}{C.1}
对易子 $[\hat{x},\hat{p}]=\mathrm{i}\hbar$，其中 $\hat{x}$ 与 $\hat{p}$ 分别是
位置与动量算符。
\end{example}

若两个算符 $\hat{A}$ 与 $\hat{B}$ 对易（$[\hat{A},\hat{B}]=0$），称它们相容
（compatible）：测量其一不以任何方式影响另一者的取值。若不对易，则二者之间存在
不确定关系：
\begin{equation*}
\Delta A\,\Delta B \geq \frac{1}{2}|\langle\mathrm{i}[\hat{A},\hat{B}]\rangle|,\tag{C.5}
\end{equation*}
$\Delta A\equiv\sqrt{\langle(\hat{A}-\langle\hat{A}\rangle)^{2}\rangle}$。

不做测量时 $\psi$ 的时间依赖由薛定谔方程给出：
\begin{equation*}
\hat{\mathcal{H}}\psi = \mathrm{i}\hbar\frac{\mathrm{d}\psi}{\mathrm{d}t}\tag{C.6}
\end{equation*}
$\hat{\mathcal{H}}$ 是哈密顿量。算符期望值的时间依赖为
\begin{equation*}
\frac{\mathrm{d}\langle\hat{A}\rangle}{\mathrm{d}t} = \frac{1}{\mathrm{i}\hbar}\langle[\hat{A},\hat{\mathcal{H}}]\rangle,\tag{C.7}
\end{equation*}
前提是算符 $\hat{A}$ 自身不依赖时间。与哈密顿量对易的可观察量是守恒量，称为运动
恒量（constant of the motion）或好量子数。

\section{狄拉克左右矢记号}

本节回顾狄拉克（Paul Dirac）发明、本书偶有使用而磁学文献中常见的一套记号。矢量
$\mathbf{a}$ 与 $\mathbf{b}$ 的标量积为 $\mathbf{a}\cdot\mathbf{b}=\sum_ia_ib_i$。
$\mathbf{a}$ 的长度为 $\sqrt{\mathbf{a}\cdot\mathbf{a}}$，必须为正实数。这一记号
使内积在 $\mathbf{a}$ 与 $\mathbf{b}$ 之间显得完全对称，但其中藏有微妙之处。若
矢量是复的，标量积应写为
\begin{equation*}
\mathbf{a}^{\dagger}\cdot\mathbf{b} = \sum_ia_i^{*}b_i = (a_1^{*}\quad a_2^{*}\quad a_3^{*})\begin{pmatrix}b_1\\b_2\\b_3\end{pmatrix}.\tag{C.8}
\end{equation*}
这一定义保证 $a^{\dagger}\cdot a=\sum_i|a_i|^{2}$ 为正实数，复矢量的长度才有良好
定义。式 C.8 强调：标量积中的两个矢量其实不应等量齐观——一个是行矢量、另一个是
列矢量；一个取了复共轭、另一个没有。数学家说这两个矢量其实“住”在不同的空间里：
b 住在矢量空间，a 住在它的对偶空间（dual space）。要把矢量变成其对偶，须取伴随
（adjoint），以 † 号标记（意思是“取复共轭并转置”）：
\begin{equation*}
\begin{pmatrix}a_1\\a_2\\a_3\end{pmatrix}^{\dagger} = (a_1^{*}\quad a_2^{*}\quad a_3^{*}).\tag{C.9}
\end{equation*}
然而在量子力学中，态函数有时是矢量（如 $\mathbf{a}=\begin{pmatrix}a_1\\a_2\end{pmatrix}$），
有时是连续函数 $a(x)$。若是函数，则 $a(x)$ 与 $b(x)$ 的标量积为
\begin{equation*}
a^{\dagger}\cdot b = \int\mathrm{d}x\,a^{*}(x)b(x).\tag{C.10}
\end{equation*}
狄拉克用把态函数写成右矢（ket）$|a\rangle$ 的办法绕开了这一记号难题——无论处理
的是有限分量矢量还是函数（后者可视为无穷维矢量）。这只是记号：$|a\rangle+|b\rangle=|c\rangle$
就是 $\mathbf{a}+\mathbf{b}=\mathbf{c}$（或等价地 $\overrightarrow{a}+\overrightarrow{b}=\overrightarrow{c}$
或 $a(x)+b(x)=c(x)$）的另一种写法。可以用伴随把右矢 $|a\rangle$ 变成左矢（bra）
$\langle a|$：
\begin{equation*}
|a\rangle^{\dagger} = \langle a|.\tag{C.11}
\end{equation*}
$\langle a|$ 与 $|b\rangle$ 的标量积写作
\begin{equation*}
\langle a|b\rangle = \sum_ia_i^{*}b_i = \int\mathrm{d}x\,a^{*}(x)b(x),\tag{C.12}
\end{equation*}
即 bra-(c)-ket——“括号”（狄拉克在此展现出幽默感的证据）。算符期望值的表达式
可写为
\begin{equation*}
\langle\hat{A}\rangle = \langle\psi|\hat{A}|\psi\rangle.\tag{C.13}
\end{equation*}
这一记号传达了标量积的不对称性，被广泛使用并大大简化了许多计算。它在标记特定
态时尤其有用，例如：
\begin{equation*}
|\text{薛定谔的猫}\rangle = \frac{|\text{活}\rangle+|\text{死}\rangle}{\sqrt{2}}.\tag{C.14}
\end{equation*}

\section{玻尔模型}

玻尔模型是描述氢原子中电子运动的半经典模型：它给出正确的结果，却是出于错误的
理由。尽管如此，它足够简单，可用于快速推导诸如原子轨道半径对核电荷 Z 与有效
质量的依赖之类结果。

考虑电子绕电荷 +Ze 的原子核运动（见图 C.1）。令静电力与向心力相等：
\begin{equation*}
\frac{Ze^{2}}{4\pi\epsilon_0 r^{2}} = \frac{m_{\mathrm{e}}v^{2}}{r}.\tag{C.15}
\end{equation*}
电子角动量 $m_{\mathrm{e}}vr$ 以 $\hbar$ 为单位量子化：
\begin{equation*}
m_{\mathrm{e}}vr = n\hbar,\tag{C.16}
\end{equation*}
n 是整数（主量子数）。于是轨道半径可写为
\begin{equation*}
r = \frac{4\pi\epsilon_0\hbar^{2}n^{2}}{Ze^{2}m_{\mathrm{e}}} = \left(\frac{4\pi\epsilon_0\hbar^{2}}{m_{\mathrm{e}}e^{2}}\right)\frac{n^{2}}{Z} = \frac{a_0n^{2}}{Z},\tag{C.17}
\end{equation*}
\mnote{这些结果可写作 $a_0=\frac{h}{m\alpha c}$ 与 $E=-\frac{1}{2}m(\alpha c)^{2}\frac{Z^{2}}{n^{2}}$，$\alpha$ 为精细结构常数。}
$a_0$ 是玻尔半径：
\begin{equation*}
a_0 = \frac{4\pi\epsilon_0\hbar^{2}}{m_{\mathrm{e}}e^{2}}.\tag{C.18}
\end{equation*}
能量 E 为
\begin{align*}
E &= \frac{1}{2}mv^{2} - \frac{Ze^{2}}{4\pi\epsilon_0 r}\tag{C.19}\\
&= -\frac{1}{2}\left(\frac{Ze^{2}}{4\pi\epsilon_0 r}\right)\tag{C.20}\\
&= -\frac{\hbar^{2}}{2m_{\mathrm{e}}}\frac{1}{a_0^{2}}\frac{Z^{2}}{n^{2}}\tag{C.21}\\
&= -\frac{m_{\mathrm{e}}e^{4}}{(4\pi\epsilon_0\hbar)^{2}}\frac{Z^{2}}{n^{2}}.\tag{C.22}
\end{align*}

\begin{figure}
\centering
\includegraphics[width=0.35\textwidth]{figC_01.pdf}
\caption{玻尔模型。}
\end{figure}

\section{轨道角动量}

角动量算符为
\begin{equation*}
\hbar\hat{\mathbf{L}} = \hat{\mathbf{r}}\times\hat{\mathbf{p}} = -\mathrm{i}\hbar\,\hat{\mathbf{r}}\times\nabla.\tag{C.23}
\end{equation*}
角动量的 z 分量的算符为
\begin{align*}
\hbar\hat{L}_z = \mathrm{i}\hbar\left[y\frac{\partial}{\partial x} - x\frac{\partial}{\partial y}\right]\tag{C.24}\\
= -\mathrm{i}\hbar\frac{\partial}{\partial\phi}\tag{C.25}
\end{align*}
其本征函数为 $e^{\mathrm{i}m_l\phi}$、本征值 $m_l\hbar$，$m_l$ 是磁量子数。下列
对易关系容易证明：
\begin{equation*}
[\hat{L}_i, \hat{L}^{2}] = 0\tag{C.26}
\end{equation*}
及
\begin{equation*}
[\hat{L}_i, \hat{L}_j] = \mathrm{i}\epsilon_{ijk}\hat{L}_k\tag{C.27}
\end{equation*}
交错张量（alternating tensor）$\epsilon_{ijk}$ 定义为
\begin{equation*}
\epsilon_{ijk} = \begin{cases} 1 & \text{若 } ijk \text{ 是 123 的偶置换}\\ -1 & \text{若 } ijk \text{ 是 123 的奇置换}\\ 0 & \text{若 } i, j, k \text{ 中任何两个相等} \end{cases}\tag{C.28}
\end{equation*}
式 C.27 是
\begin{equation*}
[\hat{L}_x, \hat{L}_y] = \mathrm{i}\hat{L}_z
\end{equation*}
及其循环置换的简写。该式使用了爱因斯坦求和约定：任何重复两次的指标都默认求和。
例如 $a_ib_i$ 是 $\sum_ia_ib_i$ 的简写。交错张量在如下表达式中很有用：
\begin{equation*}
(\mathbf{a}\times\mathbf{b})_i = \epsilon_{ijk}a_jb_k.
\end{equation*}
此式中 j 与 k 重复出现、默认求和。例如
\begin{equation*}
(\mathbf{a}\times\mathbf{b})_1 = \epsilon_{123}a_2b_3+\epsilon_{132}a_3b_2 = a_2b_3-a_3b_2.
\end{equation*}

算符 $\hat{L}^{2}$ 的本征函数
\begin{equation*}
|l, m_l\rangle = Y_{lm_l}(\theta,\phi) \propto P_l^{m_l}(\cos\theta)e^{\mathrm{i}m_l\phi},\tag{C.29}
\end{equation*}
称为球谐函数，本征值 $l(l+1)$；l 是角动量量子数，$P_l^{m_l}(\cos\theta)$ 是连带
勒让德多项式（见表 C.1）。

\begin{table}
\centering
\caption{$P_{l}^{|m_l|}(\cos\theta)$。}
\begin{small}
\begin{tabular}{lcccc}
\toprule
 & $m_l=0$ & $m_l=1$ & $m_l=2$ & $m_l=3$\\
\midrule
$l=0$ & 1 & -- & -- & --\\
$l=1$ & $\cos\theta$ & $\sin\theta$ & -- & --\\
$l=2$ & $\frac{1}{2}(3\cos^{2}\theta-1)$ & $\frac{1}{3}\sin\theta\cos\theta$ & $\frac{1}{3}\sin^{2}\theta$ & --\\
$l=3$ & $\frac{1}{2}(5\cos^{3}\theta-3\cos\theta)$ & $\frac{2}{3}\sin\theta(5\cos^{2}\theta-1)$ & $\frac{1}{15}\sin^{2}\theta\cos\theta$ & $\frac{1}{15}\sin^{3}\theta$\\
\bottomrule
\end{tabular}
\end{small}
\end{table}

于是
\begin{equation*}
\hat{L}^{2}|l, m_l\rangle = l(l+1)|l, m_l\rangle.\tag{C.30}
\end{equation*}
类似地
\begin{equation*}
\hat{L}_z|l, m_l\rangle = m_l|l, m_l\rangle.\tag{C.31}
\end{equation*}
升降算符 $\hat{L}_{\pm}$ 定义为
\begin{equation*}
\hat{L}_{\pm} = \hat{L}_x \pm \mathrm{i}\hat{L}_y.\tag{C.32}
\end{equation*}
由此可证
\begin{equation*}
\hat{L}_{\pm}|l, m_l\rangle = \sqrt{l(l+1)-m_l(m_l\pm1)}\,|l, m_l\pm1\rangle.\tag{C.33}
\end{equation*}

某些态没有轨道角动量。单态（S=0）就是一个好例子，原因在于它们具有实波函数——
这可以直截了当地证明。无磁场时哈密顿量为实，故若 $\psi$ 是哈密顿量 $\hat{\mathcal{H}}$
能量为 E 的本征函数，则 $\psi^{*}$ 亦然（$\hat{\mathcal{H}}\psi=E\psi$ 蕴含
$\hat{\mathcal{H}}\psi^{*}=E\psi^{*}$）。既然该态按定义是单态，$\psi=\psi^{*}$，
$\psi$ 为实。$\hat{\mathbf{L}}$ 各分量的算符都含 i，故若 $\psi$ 为实，期望值
$\langle\psi|\hat{L}_{\alpha}|\psi\rangle$（$\alpha=x,y,z$）都是纯虚数；但期望值
是可测量、必须是纯实数——因此全为零：该态没有轨道角动量。

\section{氢原子}

本书考虑的大多数原子不是氢，但薛定谔方程只能对氢原子精确求解。本节把氢原子的
一些结果列表备查。

球对称势的薛定谔方程可以用分离变量求解，本征函数 $\Psi$ 可写为径向与角向部分
之积：
\begin{equation*}
\Psi = R(r)\Theta(\theta)\Phi(\phi)\tag{C.34}
\end{equation*}
$\Phi(\phi)$ 的方程是
\begin{equation*}
\frac{\mathrm{d}^{2}\Phi}{\mathrm{d}\phi^{2}} + m_l^{2}\Phi = 0\tag{C.35}
\end{equation*}
容易解得
\begin{equation*}
\Phi = A\,\mathrm{e}^{\mathrm{i}m_l\phi}.\tag{C.36}
\end{equation*}
波函数径向部分的解可以求出，形式为
\begin{equation*}
R_{nl}(r) = -2\left(\frac{Z}{na_0}\right)^{\frac{3}{2}}\left[\frac{(n-l-1)!}{n[(n+l)!]^{3}}\right]^{\frac{1}{2}}\rho^{l}\,\mathrm{e}^{-\rho/2}L_{n+l}^{2l+1}(\rho)\tag{C.37}
\end{equation*}
n 是主量子数，
\begin{equation*}
\rho = \frac{2Zr}{na_0},\tag{C.38}
\end{equation*}
连带拉盖尔多项式为
\begin{equation*}
L_{\alpha}^{\beta}(x) = \frac{\mathrm{d}^{\beta}}{\mathrm{d}x^{\beta}}\left(\mathrm{e}^{x}\frac{\mathrm{d}^{\alpha}}{\mathrm{d}x^{\alpha}}[\mathrm{e}^{-x}x^{\alpha}]\right) = \sum_{j=0}^{\alpha-\beta}\frac{\alpha!^{2}(-1)^{j+\beta}}{(j+\beta)!j!(\alpha-\beta-j)!}x^{j}.\tag{C.39}
\end{equation*}
原子态的量子数范围是 $n\geq1$、$0\leq l\leq n-1$、$-l\leq m_l\leq l$。表 C.2 列出
n=1、2、3 的一些径向波函数。氢原子的若干径向波函数 $R(r)$ 绘于图 C.2，同时画出
径向概率密度 $r^{2}R(r)^{2}$。

\begin{table}
\centering
\caption{径向波函数。}
\begin{small}
\begin{tabular}{cccl}
\toprule
$n$ & $l$ & $L_{n+l}^{2l+1}(x)$ & $R_{nl}(\rho)$\\
\midrule
1 & 0 & $L_1^1(x)=-1$ & $2(Z/a_0)^{3/2}e^{-\rho/2}$\\
2 & 0 & $L_2^1(x)=2x-4$ & $(1/2\sqrt{2})(Z/a_0)^{3/2}(2-\rho)e^{-\rho/2}$\\
2 & 1 & $L_3^3(x)=-6$ & $(1/2\sqrt{6})(Z/a_0)^{3/2}\rho e^{-\rho/2}$\\
3 & 0 & $L_3^1(x)=-3x^{2}+18x-18$ & $(1/9\sqrt{3})(Z/a_0)^{3/2}(6-6\rho+\rho^{2})e^{-\rho/2}$\\
3 & 1 & $L_4^3(x)=24x-96$ & $(1/9\sqrt{6})(Z/a_0)^{3/2}(4-\rho)\rho e^{-\rho/2}$\\
3 & 2 & $L_5^5(x)=-120$ & $(1/9\sqrt{30})(Z/a_0)^{3/2}\rho^{2}e^{-\rho/2}$\\
\bottomrule
\end{tabular}
\end{small}
\end{table}

\begin{figure}
\centering
\includegraphics[width=0.24\textwidth]{figC_02.pdf}
\caption{n=1、2、3、4 的径向波函数 $R(r)$（左）与相应概率密度 $r^{2}R(r)^{2}$
（右）。}
\end{figure}

\section{g 因子}

磁场 B 中电子的能量为
\begin{equation*}
E = g\mu_{\mathrm{B}}m_sB\tag{C.40}
\end{equation*}
g 称为 g 因子。能级因此分裂 $g\mu_{\mathrm{B}}B$。狄拉克电子理论（超出本书范围）
的自然结论是 g 恰好等于 2。实际上 g 因子并不完全等于 2，而取值
\begin{equation*}
g = 2\left(1+\frac{\alpha}{2\pi}+\cdots\right) = 2.0023\cdots\tag{C.41}
\end{equation*}
$\alpha$ 是精细结构常数——无量纲量：
\begin{equation*}
\alpha = \frac{e^{2}}{4\pi\epsilon_0\hbar c} = \frac{1}{137.04}.\tag{C.42}
\end{equation*}
这一由量子电动力学（QED）得到的理论值与实验符合到惊人的精度。与 g=2 的偏差可以
这样解释：电磁相互作用源于虚光子的作用。电子实际上可以发射一个虚光子、随后再
吸收它。若在虚光子存续期间测量电子的磁矩，实际测到的是电子--光子对的磁矩——
包含一个额外的轨道分量。也可能在多于一个虚光子被产生的时段测量系统，不过这类
过程随虚光子数目增多而概率变小。平均而言，系统的磁矩比最初预期的略高——这解释了
$g-2$ 的非零值。该表达式是 $\alpha$ 的幂级数，逐项反映越来越曲折的虚光子产生与
吸收过程的贡献。

\section{d 轨道}

d 波函数（$l=2$）的角向部分——以凝聚态物理最常考虑的形式——可以由形如
$Y_{2m}e^{\mathrm{i}m_l\phi}$ 的函数的线性组合构造，其中 $Y_{2m}$ 是表 C.1 中
l=2 行所列的函数：
\begin{align*}
Y_{xy} &= \frac{Y_{22}-Y_{2-2}}{\mathrm{i}\sqrt{2}}\tag{C.43}\\
Y_{x^{2}-y^{2}} &= \frac{Y_{22}+Y_{2-2}}{\sqrt{2}}\tag{C.44}\\
Y_{yz} &= \frac{-Y_{21}-Y_{2-1}}{\mathrm{i}\sqrt{2}}\tag{C.45}\\
Y_{zx} &= \frac{-Y_{21}+Y_{2-1}}{\sqrt{2}}\tag{C.46}\\
Y_{z^{2}} &= Y_{20}\tag{C.47}
\end{align*}
给出如下结果（含归一化前置因子）：
\begin{align*}
Y_{xy} &= \sqrt{\frac{15}{16\pi}}\sin^{2}\theta\sin2\phi = \sqrt{\frac{15}{4\pi}}\frac{xy}{r^{2}}\tag{C.48}\\
Y_{x^{2}-y^{2}} &= \sqrt{\frac{15}{16\pi}}\sin^{2}\theta\cos2\phi = \sqrt{\frac{15}{16\pi}}\frac{x^{2}-y^{2}}{r^{2}}\tag{C.49}\\
Y_{yz} &= \sqrt{\frac{15}{4\pi}}\sin\theta\cos\theta\sin\phi = \sqrt{\frac{15}{4\pi}}\frac{yz}{r^{2}}\tag{C.50}\\
Y_{zx} &= \sqrt{\frac{15}{4\pi}}\sin\theta\cos\theta\cos\phi = \sqrt{\frac{15}{4\pi}}\frac{zx}{r^{2}}\tag{C.51}\\
Y_{z^{2}} &= \sqrt{\frac{15}{16\pi}}(3\cos^{2}\theta-1) = \sqrt{\frac{15}{16\pi}}\frac{2z^{2}-x^{2}-y^{2}}{r^{2}}.\tag{C.52}
\end{align*}

\begin{figure}
\centering
\includegraphics[width=0.35\textwidth]{figC_03.pdf}
\caption{内禀自旋--轨道相互作用。在核的参考系中电子绕核运动；但在与电子共动的
惯性系中则是核在绕行。}
\end{figure}

这些轨道连同 s、p 轨道示于图 3.1。

\section{自旋--轨道相互作用}

原子中的自旋--轨道相互作用来源如下。考虑绕原子运动的电子：图 C.3 上图描绘了核
静止参考系中的情形；下图中原子被画在与电子共动的惯性系里——核看起来在绕电子
运动。绕行的核构成电流，在原点处产生磁场
\begin{equation*}
\mathbf{B} = \frac{\boldsymbol{\varepsilon}\times\mathbf{v}}{c^{2}},\tag{C.53}
\end{equation*}
其中
\begin{equation*}
\boldsymbol{\varepsilon} = -\nabla V(r) = -\frac{\mathbf{r}}{r}\frac{\mathrm{d}V(r)}{\mathrm{d}r}\tag{C.54}
\end{equation*}
是核在电子处产生的电场，$V(r)$ 是相应的势能。式 C.53 来自狭义相对论中电磁场的
变换。这一磁场与电子自旋作用，在哈密顿量中给出
\begin{align*}
H_{\mathrm{so}} &= -\frac{1}{2}\mathbf{m}\cdot\mathbf{B}\tag{C.55}\\
&= \frac{e\hbar^{2}}{2m_{\mathrm{e}}c^{2}r}\frac{\mathrm{d}V(r)}{\mathrm{d}r}\,\mathbf{S}\cdot\mathbf{L}\tag{C.56}
\end{align*}
轨道角动量为 $\hbar\mathbf{L}=m_{\mathrm{e}}\mathbf{r}\times\mathbf{v}$，磁矩
$\mathbf{m}=(ge\hbar/2m)\mathbf{S}$；式 C.55 中的因子 $\frac{1}{2}$ 是相对论的
托马斯因子。用狄拉克方程可以极其优雅地得到这一结果——相对论修正自动包含在内。
\fn{2}{对应电子轨道上各个瞬时静止系，存在无穷多组洛伦兹变换。速度方向改变时洛伦兹变换不对易——我们的推导忽略了这一点。净洛伦兹变换包含一个转动，这一托马斯进动恰好给出二分之一的因子。}\bio{L.~H. Thomas\\（1903--1992）}这一效应称为内禀自旋--轨道相互作用，是原子中电子波函数的自旋部分与轨道部分之间的相互作用。对类氢原子中的库仑场
\begin{equation*}
\frac{1}{r}\frac{\mathrm{d}V(r)}{\mathrm{d}r} = \frac{Ze}{4\pi\epsilon_0 r^{3}}\tag{C.57}
\end{equation*}
对量子数为 l 与 n 的电子态有
\begin{equation*}
\langle r^{-3}\rangle = \frac{Z^{3}}{a_0^{3}n^{3}l(l+\frac{1}{2})(l+1)}\tag{C.58}
\end{equation*}
故自旋--轨道劈裂为
\begin{equation*}
\frac{Z^{4}e^{2}\hbar^{2}\langle\mathbf{S}\cdot\mathbf{L}\rangle}{4\pi\epsilon_0 a_0^{3}n^{3}l(l+\frac{1}{2})(l+1)}.\tag{C.59}
\end{equation*}

\section{Landé g 因子}

Landé g 因子 $g_J$ 由
\begin{equation*}
g_J = \frac{3}{2} + \frac{S(S+1)-L(L+1)}{2J(J+1)}\tag{C.60}
\end{equation*}
给出。式 C.60 中 $g_J$ 的表达式可以从磁矩算符的表达式出发得到：
\begin{equation*}
\hat{\boldsymbol{\mu}} = \mu_{\mathrm{B}}(g_L\hat{\mathbf{L}} + g_S\hat{\mathbf{S}}).\tag{C.61}
\end{equation*}
其中 $g_L=1$ 与 $g_S=2$ 分别是轨道角动量与自旋角动量的 g 因子。在许多原子中 S
与 L 可能不是好量子数，但 J 是。因此平行于 J 的磁矩分量是守恒量，垂直于 J 的
分量则不是。于是写
\begin{equation*}
\hat{\boldsymbol{\mu}} = g_J\mu_{\mathrm{B}}\hat{\mathbf{J}},\tag{C.62}
\end{equation*}
$g_J$ 是待定常数。式 C.62 意味着 $g_J$ 是 $\hat{\mathbf{L}}+2\hat{\mathbf{S}}$
在 $\hat{\mathbf{J}}$ 上的投影。把式 C.61 两边乘 $\hat{\mathbf{J}}$：
\begin{equation*}
\hat{\boldsymbol{\mu}}\cdot\mathbf{J} = \mu_{\mathrm{B}}(g_L\hat{\mathbf{L}}\cdot\hat{\mathbf{J}} + g_S\hat{\mathbf{S}}\cdot\hat{\mathbf{J}}).\tag{C.63}
\end{equation*}
把式 C.62 两边乘 $\hat{\mathbf{J}}$ 得\fn{3}{在此及后续步骤中，我们利用 $S^{2}$ 的本征值为 $S(S+1)$、$L^{2}$ 的本征值为 $L(L+1)$、$J^{2}$ 的本征值为 $J(J+1)$ 这一事实。}
\begin{equation*}
\hat{\boldsymbol{\mu}}\cdot\mathbf{J} = g_J\mu_{\mathrm{B}}\hat{\mathbf{J}}^{2} = g_J\mu_{\mathrm{B}}J(J+1).\tag{C.64}
\end{equation*}
又 $\mathbf{L}^{2}=(\mathbf{J}-\mathbf{S})^{2}=\mathbf{J}^{2}+\mathbf{S}^{2}-2\mathbf{S}\cdot\mathbf{J}$，故
\begin{equation*}
\mathbf{S}\cdot\mathbf{J} = \frac{1}{2}(\mathbf{J}^{2}-\mathbf{L}^{2}+\mathbf{S}^{2})\tag{C.65}
\end{equation*}
且 $\mathbf{S}^{2}=(\mathbf{J}-\mathbf{L})^{2}=\mathbf{J}^{2}+\mathbf{L}^{2}-2\mathbf{L}\cdot\mathbf{J}$，故
\begin{equation*}
\mathbf{L}\cdot\mathbf{J} = \frac{1}{2}(\mathbf{J}^{2}+\mathbf{L}^{2}-\mathbf{S}^{2}).\tag{C.66}
\end{equation*}
联立式 C.63 与 C.64，并代入式 C.65、C.66 的结果：
\begin{equation*}
\begin{aligned}
g_J &= g_L\left(\frac{J(J+1)+L(L+1)-S(S+1)}{2J(J+1)}\right)\\
&\quad + g_S\left(\frac{J(J+1)-L(L+1)+S(S+1)}{2J(J+1)}\right)
\end{aligned}\tag{C.67}
\end{equation*}
$g_L=1$、$g_S=2$ 时即退化为式 C.60。

\section{微扰论}

考虑哈密顿量 $\hat{\mathcal{H}}_0$，其本征函数 $|\phi_i\rangle$ 与本征值 $E_i$
已知：
\begin{equation*}
\hat{\mathcal{H}}_0|\phi_i\rangle = E_i|\phi_i\rangle.\tag{C.68}
\end{equation*}
本征函数 $|\phi_i\rangle$ 全部正交：
\begin{equation*}
\langle\phi_i|\phi_j\rangle = \int\mathrm{d}\tau\,\phi_i^{*}\phi_j = \delta_{ij}\tag{C.69}
\end{equation*}
$\mathrm{d}\tau$ 是体积元。现设在 $\hat{\mathcal{H}}_0$ 上加微扰 $\hat{V}$，新
哈密顿量为 $\hat{\mathcal{H}}$：
\begin{equation*}
\hat{\mathcal{H}} = \hat{\mathcal{H}}_0 + \hat{V}.\tag{C.70}
\end{equation*}
加微扰前系统处于能量 $E_k$ 的态 $|\phi_k\rangle$。系统的新本征函数将为 $|\psi\rangle$、
本征值为 E：
\begin{equation*}
\hat{\mathcal{H}}|\psi\rangle = E|\psi\rangle,\tag{C.71}
\end{equation*}
但 $|\psi\rangle$ 与 E 目前都未知。$|\psi\rangle$ 可以用旧本征函数展开：
\begin{equation*}
|\psi\rangle = \sum_ja_j|\phi_j\rangle.\tag{C.72}
\end{equation*}
代入式 C.71，左乘 $\langle\phi_i|$ 并在适当体积上积分：
\begin{equation*}
\sum_ja_j\int\mathrm{d}\tau\,\phi_i^{*}(\hat{\mathcal{H}}_0+\hat{V})\phi_j = E\sum_ja_j\int\mathrm{d}\tau\,\phi_i^{*}\phi_j,\tag{C.73}
\end{equation*}
用原本征函数的正交性（式 C.69）得
\begin{equation*}
\sum_jV_{ij}a_j = (E-E_i)a_i,\tag{C.74}
\end{equation*}
称为久期方程（secular equation），其中
\begin{equation*}
V_{ij} = \langle\phi_i|\hat{V}|\phi_j\rangle = \int\mathrm{d}\tau\,\phi_i^{*}\hat{V}\phi_j\tag{C.75}
\end{equation*}
是微扰的矩阵元。从现在起我们假定非简并情形，即假定没有态是简并的。

若微扰很小，$|\psi\rangle$ 将非常接近初态 $|\phi_k\rangle$：$a_k\approx1$，且
$i\neq k$ 时 $|a_i|\ll1$。既然波函数只有微小变化，可以设想新能量为
\begin{equation*}
E \approx E_k,\tag{C.76}
\end{equation*}
即完全不变。利用 $a_k\approx1$ 与 $i\neq k$ 时 $a_i\ll1$ 可以改进这一零阶近似。
把它们代入久期方程可得能量与 $a_i$ 的近似表达式。对 i=k，久期方程中的求和由
j=k 的项主导：
\begin{equation*}
E \approx E_k + V_{kk},\tag{C.77}
\end{equation*}
态能量的移动简单地就是 $V_{kk}$。这一包含“一阶微扰论”能量修正的结果在本书
正文中用得非常多。久期方程也可对 $i\neq k$ 求值：
\begin{equation*}
a_i = \frac{V_{ik}}{E_k-E_i}\tag{C.78}
\end{equation*}
（其中 E 已换成 $E_k$）。把式 C.78 代回式 C.74 得下一级能量近似：
\begin{equation*}
E \approx E_k + V_{kk} + \sum_{i\neq k}\frac{|V_{ik}|^{2}}{E_k-E_i}\tag{C.79}
\end{equation*}
这就是“二阶微扰论”结果。这一过程可以继续下去，给出 E 的幂级数表达式，级数的
逐项包含微扰矩阵元的更高次幂。

\section*{延伸阅读}

\begin{itemize}[itemsep=0.2em]
\item A.~I.~Rae,《Introduction to quantum mechanics》第 3 版, IOP Publishing
1992。
\item J.~J.~Sakurai,《Modern quantum mechanics》第 2 版, Addison-Wesley 1994。
\item P.~A.~M.~Dirac,《The principles of quantum mechanics》第 4 版, OUP 1958。
\item B.~H.~Bransden 与 C.~J.~Joachain,《Physics of atoms and molecules》,
Longman 1983。
\item P.~W.~Atkins,《Molecular quantum mechanics》, OUP 1983。
\item G.~K.~Woodgate,《Elementary atomic structure》, OUP 1980。
\end{itemize}

%% 附录 D 磁学中的能量与退磁场（Energy in magnetism and demagnetizing fields）
\chapter{磁学中的能量与退磁场}

\section{能量}

计算磁化介质在磁场中的能量是件出人意料精微的事。可以得到各种表达式，但它们所指
的对象可能不同。在磁学中，这取决于你谈的是磁矩单独的能量，还是磁矩加上提供磁场
的东西二者的能量，以及这一能量究竟如何划分。

例如，把一把螺丝刀靠近大磁体的极靴，螺丝刀会被强烈吸向场最强的区域。用被磁化的
螺丝刀替换极靴间隙中的一部分空气似乎是能量上有利的事——这暗示可磁化介质（螺丝刀）
的存在降低了能量。然而，在带铁芯的磁体中建立电流比在无铁芯的磁体中要费更多能量
——这暗示可磁化介质（此处为铁芯）的存在提高了能量。那么能量究竟是降低还是升高？
答案取决于你考虑的是哪个能量。第一个例子里，我们只量度了握螺丝刀的手臂感受到的
力所省下的能量，却忽略了磁体电源为把电流维持在我们挥舞螺丝刀之前的水平而必须
多做的功。这提醒我们：处理磁性材料的能量学必须极其小心（详细处理见延伸阅读中
Heine 的文章）。

附录 B 中我们已经看到，$\delta W=-\boldsymbol{\mu}\cdot\delta\mathbf{B}$ 是磁化
介质应考虑的恰当自由能。这一表达式在统计力学中表述磁性时将有用武之地（见附录 E）。
本附录集中于：铁磁材料的存在如何因退磁场的引入而同时改变 H 与 B——退磁场自身
是要耗费能量的。

\section{退磁因子}

铁磁体内部的磁化强度 M 到达表面时必须突然停止，因此 M 有散度。利用方程
\begin{equation*}
\nabla\cdot\mathbf{H} = -\nabla\cdot\mathbf{M},\tag{D.1}
\end{equation*}
（它由 $\nabla\cdot\mathbf{B}=0$ 得来，见附录 B）我们发现 $\mathbf{H}$ 有等大反向
的散度。情形就好像铁磁体表面留下了磁单极，而这些单极充当 $\mathbf{H}$ 的源。
由此产生的 H 场称为退磁场（demagnetizing field）。图 D.1 以简单情形展示这一局面：
无限大平板或薄膜铁磁材料的磁化。若磁化位于板面内，M 的散度只出现在（设为无穷远
的）两端，因此没有退磁场。若磁化垂直于板面，顶、底表面产生磁极，板内出现退磁场
$\mathbf{H}_{\mathrm{d}}=-\mathbf{M}$。

\begin{figure}
\centering
\includegraphics[width=0.6\textwidth]{figD_01.pdf}
\caption{无限大平板的退磁（截面视图）。（a）若磁化位于板面内，板表面不产生磁极
（除两端极小的磁极外）。（b）若磁化如图垂直于板面，顶（底）面上 M 的负（正）散度
产生 H 的正（负）散度，从而如图在顶（底）面产生正（负）磁极。（c）这导致一个从
正磁极流向负磁极的退磁场。}
\end{figure}

对任意形状的铁磁体，退磁场可以是位置极其复杂的函数。但对椭球铁磁体它取相对简单
的形式：此时它在铁磁体内部是均匀的，值为
\begin{equation*}
\mathbf{H}_{\mathrm{d}} = -\mathbf{N}\mathbf{M}\tag{D.2}
\end{equation*}
N 是退磁张量。一般地可写
\begin{equation*}
(H_{\mathrm{d}})_i = -\sum_jN_{ij}M_j.\tag{D.3}
\end{equation*}
若 $\mathbf{M}$ 沿椭球的一个主轴，$\mathsf{N}_{ij}$ 可以对角化为
\begin{equation*}
\mathsf{N} = \begin{pmatrix}N_x & 0 & 0\\ 0 & N_y & 0\\ 0 & 0 & N_z\end{pmatrix}.\tag{D.4}
\end{equation*}
退磁张量满足
\begin{equation*}
\mathrm{Tr}\,\mathsf{N} = N_x+N_y+N_z = 1.\tag{D.5}
\end{equation*}

\begin{example}{D.1}
现在考察几个特例。

（1）球：$N_x=N_y=N_z=\frac{1}{3}$，故
\begin{equation*}
\mathbf{H}_{\mathrm{d}} = -\frac{\mathbf{M}}{3}\tag{D.6}
\end{equation*}

（2）沿 z 的细长圆柱棒：$N_x=N_y=\frac{1}{2}$、$N_z=0$。因为若磁化严格沿棒方向，
产生的磁极在棒的两端——可以假定远得不值得一提。

（3）垂直于 z 的平板：$N_x=N_y=0$、$N_z=1$。即上面考虑过的情形。
\end{example}

\section{任意形状的铁磁体}

本节说明如何计算外加磁场中任意形状铁磁体的退磁能。对此问题，磁场
$\mathbf{H}(\mathbf{r})$ 与 $\mathbf{B}(\mathbf{r})$ 可以分解为两个分量：
\begin{equation*}
\mathbf{H}(\mathbf{r}) = \mathbf{H}_{\mathrm{a}}(\mathbf{r}) + \mathbf{H}_{\mathrm{d}}(\mathbf{r})\tag{D.7}
\end{equation*}
\begin{equation*}
\mathbf{B}(\mathbf{r}) = \mathbf{B}_{\mathrm{a}}(\mathbf{r}) + \mathbf{B}_{\mathrm{d}}(\mathbf{r}),\tag{D.8}
\end{equation*}
外场为 $\mathbf{H}_{\mathrm{a}}(\mathbf{r})=\mathbf{B}_{\mathrm{a}}(\mathbf{r})/\mu_0$，
退磁场为 $\mathbf{H}_{\mathrm{d}}(\mathbf{r})$；把 $\mathbf{B}_{\mathrm{d}}(\mathbf{r})$
写成退磁场产生的磁通密度 $\mu_0\mathbf{H}_{\mathrm{d}}(\mathbf{r})$ 与材料内部
$\mu_0\mathbf{M}(\mathbf{r})$ 之和，即
$\mathbf{B}_{\mathrm{d}}(\mathbf{r})=\mu_0(\mathbf{H}_{\mathrm{d}}(\mathbf{r})+\mathbf{M}(\mathbf{r}))$。
退磁场源于磁体表面上凡 $\nabla\cdot\mathbf{M}(\mathbf{r})\neq0$ 处产生的磁极。

现在陈述一个重要结果。退磁场 $\mathbf{B}_{\mathrm{d}}(\mathbf{r})$ 与
$\mathbf{H}_{\mathrm{d}}(\mathbf{r})$ 满足
\begin{equation*}
\int_{\text{全空间}}\mathbf{B}_{\mathrm{d}}\cdot\mathbf{H}_{\mathrm{d}}\,\mathrm{d}\tau = 0.\tag{D.9}
\end{equation*}
证明如下：由于 $\nabla\times\mathbf{H}_{\mathrm{d}}=0$，退磁场 $\mathbf{H}_{\mathrm{d}}$
可以写成标量函数 $\phi$ 的梯度（如 $\mathbf{H}_{\mathrm{d}}=-\nabla\phi$）。于是
\begin{equation*}
\nabla\cdot(\phi\mathbf{B}_{\mathrm{d}}) \equiv \phi\nabla\cdot\mathbf{B}_{\mathrm{d}} + (\nabla\phi)\cdot\mathbf{B}_{\mathrm{d}} = (\nabla\phi)\cdot\mathbf{B}_{\mathrm{d}}\tag{D.10}
\end{equation*}
最后一步用了 $\nabla\cdot\mathbf{B}_{\mathrm{d}}=0$。从而
\begin{equation*}
\int_{\text{全空间}}\mathbf{B}_{\mathrm{d}}\cdot\mathbf{H}_{\mathrm{d}}\,\mathrm{d}\tau = -\int_{\text{全空间}}\nabla\cdot(\phi\mathbf{B}_{\mathrm{d}})\,\mathrm{d}\tau = \lim_{R\to\infty}\int_{\text{面}}\phi\,\mathbf{B}_{\mathrm{d}}\cdot\mathrm{d}\mathbf{S} = 0,\tag{D.11}
\end{equation*}
面积分取在半径 R 趋于无穷的球面上。若铁磁体尺度有限，积分必为零：$R\to\infty$ 时
至多 $B_d\sim R^{-2}$、$\phi\sim R^{-1}$。

若不存在外场，退磁场的能量 E 就是能量密度 $\frac{1}{2}\mu_0H_d^{2}$ 对全空间的
积分。利用式 D.9 可以把它化为铁磁体体积 V 内的积分，使能量表示为
\begin{align*}
E = \frac{\mu_0}{2}\int_{\text{全空间}}\mathbf{H}_{\mathrm{d}}^{2}\,\mathrm{d}\tau &= -\frac{1}{2}\int_{\text{全空间}}(\mathbf{B}_{\mathrm{d}}-\mu_0\mathbf{M})\cdot\mathbf{H}_{\mathrm{d}}\,\mathrm{d}\tau\\
&= -\frac{\mu_0}{2}\int_{\text{全空间}}\mathbf{M}\cdot\mathbf{H}_{\mathrm{d}}\,\mathrm{d}\tau\\
&= -\frac{\mu_0}{2}\int_V\mathbf{M}\cdot\mathbf{H}_{\mathrm{d}}\,\mathrm{d}\tau,\tag{D.12}
\end{align*}
最后一行来自 M=0（体外）。应当注意：尽管最后的积分只遍及体内，它表达的却是全部
杂散场（包括体外杂散场）的能量。

若存在外场，铁磁体的能量可以写成总场能量密度 $\frac{1}{2}\mu_0H^{2}$ 与外场能量
密度 $\frac{1}{2}\mu_0H_a^{2}$ 之差——这去掉了均匀外场对全空间积分给出的无穷大
能量贡献。于是能量可以有效地写为
\begin{equation*}
E = \frac{\mu_0}{2}\int_{\text{全空间}}(\mathbf{H}^{2}-\mathbf{H}_{\mathrm{a}}^{2})\,\mathrm{d}\tau.\tag{D.13}
\end{equation*}
利用 $\nabla\times\mathbf{H}_a=0$，有
\begin{equation*}
\int_{\text{全空间}}\mathbf{H}_{\mathrm{a}}\cdot\mathbf{B}_{\mathrm{d}}\,\mathrm{d}\tau = 0.\tag{D.14}
\end{equation*}
再利用下列方程：
\begin{equation*}
\mathbf{H}^{2} = \mathbf{H}_{\mathrm{a}}^{2} + 2\mathbf{H}_{\mathrm{d}}\cdot\mathbf{H}_{\mathrm{a}} + \mathbf{H}_{\mathrm{d}}^{2},\tag{D.15}
\end{equation*}
\begin{equation*}
\mathbf{H}_{\mathrm{d}}\cdot\mathbf{H}_{\mathrm{a}} = \mathbf{H}_{\mathrm{a}}\cdot\left(\frac{\mathbf{B}_{\mathrm{d}}}{\mu_0}-\mathbf{M}\right).\tag{D.16}
\end{equation*}
并再用式 D.9，能量 E 可以写成熟悉的形式
\begin{equation*}
E = -\mu_0\int_V\mathbf{M}\cdot\mathbf{H}_{\mathrm{a}}\,\mathrm{d}\tau - \frac{\mu_0}{2}\int_V\mathbf{M}\cdot\mathbf{H}_{\mathrm{d}}\,\mathrm{d}\tau,\tag{D.17}
\end{equation*}
即塞曼项与退磁项之和，每个积分只遍及铁磁体体积 V。式 D.17 的能量尽管如此仍隐含
地包括体外退磁场相联系的能量。

\begin{example}{D.2}
退磁项中二分之一因子可以这样理解：避免重复计数。塞曼项是磁化强度与外场的相互
作用；退磁项是自能——产生场的矩与自己的场相互作用。二分之一因子保证这一能量
不被计两次。

半径 a、磁化强度均匀为 M 的铁磁球是导出式 D.12 诸想法的详细示例（见图 D.2）。
退磁场为
\begin{equation*}
\mathbf{H}_{\mathrm{d}} = -\frac{\mathbf{M}}{3}\tag{D.18}
\end{equation*}
\begin{equation*}
\mathbf{B}_{\mathrm{d}} = \frac{2\mu_0\mathbf{M}}{3}.\tag{D.19}
\end{equation*}
\begin{figure}
\centering
\includegraphics[width=0.4\textwidth]{figD_02.pdf}
\caption{磁化强度均匀为 M 的铁磁球。球内区域（体积 $V_{\text{sphere}}=\frac{4}{3}\pi a^{3}$）
记为 V，球外区域记为 $\bar{V}$。}
\end{figure}
于是球内（体积 $V_{\text{sphere}}=\frac{4}{3}\pi a^{3}$ 的空间，记为 V）简单积分
给出：
\begin{equation*}
-\frac{\mu_0}{2}\int_V\mathbf{M}\cdot\mathbf{H}_{\mathrm{d}}\,\mathrm{d}\tau = \frac{1}{6}\mu_0M^{2}V_{\text{sphere}}\tag{D.20}
\end{equation*}
\begin{equation*}
\frac{\mu_0}{2}\int_V\mathbf{H}_{\mathrm{d}}^{2}\,\mathrm{d}\tau = \frac{1}{18}\mu_0M^{2}V_{\text{sphere}}\tag{D.21}
\end{equation*}
\begin{equation*}
\frac{1}{2}\int_V\mathbf{B}_{\mathrm{d}}\cdot\mathbf{H}_{\mathrm{d}}\,\mathrm{d}\tau = -\frac{1}{9}\mu_0M^{2}V_{\text{sphere}}\tag{D.22}
\end{equation*}
球外（记为 $\bar{V}$ 的空间）$B_d=\mu_0H_d$（球外 M=0），退磁场来自球的磁偶极矩，
其径向与极角分量为
\begin{equation*}
\begin{aligned}
(H_{\mathrm{d}})_r &= 2Ma^{3}\cos\theta/3r^{3}\\
(H_{\mathrm{d}})_{\theta} &= Ma^{3}\sin\theta/3r^{3},
\end{aligned}\tag{D.23}
\end{equation*}
对球外区域积分得：
\begin{equation*}
\frac{\mu_0}{2}\int_{\bar{V}}\mathbf{H}_{\mathrm{d}}^{2}\,\mathrm{d}\tau = \frac{1}{2}\int_{\bar{V}}\mathbf{B}_{\mathrm{d}}\cdot\mathbf{H}_{\mathrm{d}}\,\mathrm{d}\tau = \frac{1}{2\mu_0}\int_{\bar{V}}\mathbf{B}_{\mathrm{d}}^{2}\,\mathrm{d}\tau = \frac{1}{9}\mu_0M^{2}V_{\text{sphere}}.\tag{D.24}
\end{equation*}
对全空间积分于是给出
\begin{equation*}
-\frac{\mu_0}{2}\int_{\text{全空间}}\mathbf{M}\cdot\mathbf{H}_{\mathrm{d}}\,\mathrm{d}\tau = \frac{\mu_0}{2}\int_{\text{全空间}}\mathbf{H}_{\mathrm{d}}^{2}\,\mathrm{d}\tau = \frac{1}{6}\mu_0M^{2}V_{\text{sphere}}\tag{D.25}
\end{equation*}
与式 D.12 一致。联立式 D.22 与 D.24 得
\begin{equation*}
\frac{1}{2}\int_{\text{全空间}}\mathbf{B}_{\mathrm{d}}\cdot\mathbf{H}_{\mathrm{d}}\,\mathrm{d}\tau = 0\tag{D.26}
\end{equation*}
与式 D.9 一致。
\end{example}

\section*{延伸阅读}

\begin{itemize}[itemsep=0.2em]
\item V.~Heine, Proc.~Camb.~Phil.~Soc. \textbf{52}, 546 (1956)，对电磁场中物体的
热力学有深刻的处理。
\item 相关论述还可见 C.~J.~Adkins,《Equilibrium thermodynamics》, CUP 1983 与
J.~R.~Waldram,《The theory of thermodynamics》, CUP 1985。
\item 静磁能与退磁因子的深入论述见 A.~Aharoni,《Introduction to the theory of
ferromagnetism》, OUP 1996。
\item I.~Brevik, J.~Phys.: Condensed Matter \textbf{7}, 8065 (1995) 对计算浸入
铁磁液体中的铁磁体能量的问题给出了有益的分析。
\end{itemize}

%% 附录 E 统计力学（Statistical mechanics）
\chapter{统计力学}

本附录简要勾画主正文将用到的统计力学中的若干结果。

\section{配分函数与热力学函数}

温度 T 下占据能量为 $E_i$ 的态 i 的概率 $p_i$ 为
\begin{equation*}
p_i = \frac{\mathrm{e}^{-E_i/k_{\mathrm{B}}T}}{Z}\tag{E.1}
\end{equation*}
Z 是单粒子配分函数，用于归一化概率：
\begin{equation*}
\sum_ip_i = 1.\tag{E.2}
\end{equation*}
这意味着 Z 由
\begin{equation*}
Z = \sum_i\mathrm{e}^{-E_i/k_{\mathrm{B}}T}\tag{E.3}
\end{equation*}
给出。定义 $\beta=1/k_{\mathrm{B}}T$，可以求出 $\partial Z/\partial\beta$：
\begin{equation*}
\frac{\partial Z}{\partial\beta} = -\sum_iE_i\,\mathrm{e}^{-\beta E_i}.\tag{E.4}
\end{equation*}
既然已经找到概率分布，就可以用它计算各量的期望值。能量 $\langle E\rangle$ 为
\begin{equation*}
\langle E\rangle = \sum_iE_ip_i = \frac{\sum_iE_i\,\mathrm{e}^{-\beta E_i}}{Z} = -\frac{1}{Z}\frac{\partial Z}{\partial\beta} = -\frac{\partial\ln Z}{\partial\beta}.\tag{E.5}
\end{equation*}
每粒子的熵为
\begin{equation*}
S = -k_{\mathrm{B}}\sum_ip_i\ln p_i\tag{E.6}
\end{equation*}
从而
\begin{equation*}
S = -k_{\mathrm{B}}\sum_i\frac{\mathrm{e}^{-\beta E_i}}{Z}[-\beta E_i-\ln Z] = k_{\mathrm{B}}\ln Z + \frac{\langle E\rangle}{T}.\tag{E.7}
\end{equation*}
在式 E.8 中我们用了附录 B 的结果：可利用的自由能为 $-\boldsymbol{\mu}\cdot\delta\mathbf{B}$。

对 N 个独立可分辨的粒子，恰当的配分函数变为 $Z^{N}$，总能量为
$E=N\langle E\rangle$。热力学第一定律可以表示为
\begin{equation*}
\mathrm{d}E = T\,\mathrm{d}S - p\,\mathrm{d}V - M\,\mathrm{d}B.\tag{E.8}
\end{equation*}
亥姆霍兹自由能 F（又称亥姆霍兹函数）定义为
\begin{equation*}
F = E - TS,\tag{E.9}
\end{equation*}
与配分函数 Z 的关系为
\begin{equation*}
F = -Nk_{\mathrm{B}}T\ln Z.\tag{E.10}
\end{equation*}
用式 E.8 与 E.9，$\mathrm{d}F$ 可写为
\begin{equation*}
\mathrm{d}F = -S\,\mathrm{d}T - p\,\mathrm{d}V - M\,\mathrm{d}B\tag{E.11}
\end{equation*}
故
\begin{equation*}
M = -\left(\frac{\partial F}{\partial B}\right)_{T,V} = Nk_{\mathrm{B}}T\left(\frac{\partial\ln Z}{\partial B}\right)_{T,V}.\tag{E.12}
\end{equation*}
熵为\footnote{原书式 (E.13) 左端印作 $M$；按上文“熵由下式给出”，应为 $S$——疑为原书排印之误。——译注}
\begin{equation*}
M = -\left(\frac{\partial F}{\partial T}\right)_{V,B}.\tag{E.13}
\end{equation*}

\section{能量均分定理}

体系的能量常可写成一些二次项之和。例如弹簧常数 K 的弹簧上的质量 m，其能量
$E=\frac{1}{2}mv^{2}+\frac{1}{2}Kx^{2}$ 就含两个这样的二次项。低温下量子效应主导，
我们预期谐振子解与能级阶梯。但高温下（$k_{\mathrm{B}}T$ 远大于能级间距设定的能量
尺度），系统可以经典地处理。

能量均分定理（equipartition theorem）断言：在经典极限下——$k_{\mathrm{B}}T$ 远
大于量子能级间距、且能量含若干独立二次项时——能量的期望值等于 $\frac{1}{2}k_{\mathrm{B}}T$
乘以能量中独立二次项的数目。

证明很直接：若 $E(x_1,x_2,\ldots,x_n)=\sum_{i=1}^{n}\alpha_ix_i^{2}$（$\alpha_i$
为常数），则
\begin{equation*}
\langle E\rangle = \frac{\int_{-\infty}^{\infty}\int_{-\infty}^{\infty}\cdots\int_{-\infty}^{\infty}\mathrm{d}x_1\mathrm{d}x_2\cdots\mathrm{d}x_n\,e^{-\beta E(x_1,x_2,\ldots x_n)}E}{\int_{-\infty}^{\infty}\int_{-\infty}^{\infty}\cdots\int_{-\infty}^{\infty}\mathrm{d}x_1\mathrm{d}x_2\cdots\mathrm{d}x_n\,e^{-\beta E(x_1,x_2,\ldots x_n)}}.\tag{E.14}
\end{equation*}
看起来吓人，但对每个变量的积分可以轻松分离。于是
\begin{align*}
\langle E\rangle &= \sum_{i=1}^{n}\alpha_i\left(\frac{\int_{-\infty}^{\infty}\mathrm{d}x_i\,x_i^{2}\,\mathrm{e}^{-\beta\alpha_ix_i^{2}}\prod_{j\neq i}\int_{-\infty}^{\infty}\mathrm{d}x_j\,\mathrm{e}^{-\beta\alpha_ix_j^{2}}}{\int_{-\infty}^{\infty}\mathrm{d}x_i\,\mathrm{e}^{-\beta\alpha_ix_i^{2}}\prod_{j\neq i}\int_{-\infty}^{\infty}\mathrm{d}x_j\,\mathrm{e}^{-\beta\alpha_ix_j^{2}}}\right)\\
&= \sum_{i=1}^{n}\left(\frac{\frac{\alpha_i}{2}\sqrt{\frac{\pi}{\alpha_i^{3}\beta^{3}}}}{\sqrt{\frac{\pi}{\alpha_i\beta}}}\right)\\
&= \sum_{i=1}^{n}\frac{1}{2\beta}\\
&= \frac{n}{2}k_{\mathrm{B}}T,\tag{E.15}
\end{align*}
其中用了标准积分\footnote{原书乘积积分的指数印作 $\alpha_i x_j^{2}$；为使积分分解成立应为 $\alpha_j x_j^{2}$——疑为原书排印之误，上式照录。——译注}
\begin{equation*}
\int_{-\infty}^{\infty}\mathrm{d}x\,e^{-\gamma x^{2}} = \sqrt{\frac{\pi}{\gamma}} \qquad\text{及}\qquad \int_{-\infty}^{\infty}\mathrm{d}x\,x^{2}e^{-\gamma x^{2}} = \frac{1}{2}\sqrt{\frac{\pi}{\gamma^{3}}}.\tag{E.16}
\end{equation*}

\begin{example}{E.1}
这一结果的三例：
（1）弹簧常数 K 的弹簧上质量 m：$E=\frac{1}{2}mv^{2}+\frac{1}{2}Kx^{2}$（两个二次
项），故 $\langle E\rangle=k_{\mathrm{B}}T$。
（2）气体中质量 m 的分子：$E=\frac{1}{2}mv_x^{2}+\frac{1}{2}mv_y^{2}+\frac{1}{2}mv_z^{2}$
（三个二次项），故 $\langle E\rangle=\frac{3}{2}k_{\mathrm{B}}T$。
（3）含 N 个原子的固体（因而 3N 根弹簧、能量含 6N 个二次项）：
$\langle E\rangle=3Nk_{\mathrm{B}}T$。
\end{example}

除了给能量加上 $\frac{1}{2}k_{\mathrm{B}}T$，每个二次自由度还给热容加上
$\frac{1}{2}k_{\mathrm{B}}$。

\section*{延伸阅读}

\begin{itemize}[itemsep=0.2em]
\item F.~Reif,《Fundamentals of statistical and thermal physics》, McGraw-Hill,
1965。
\item J.~R.~Waldram,《The theory of thermodynamics》, CUP 1985。
\item P.~M.~Chaikin 与 T.~C.~Lubensky,《Principles of condensed matter physics》,
CUP 1995。
\item A.~M.~Glazer 与 J.~S.~Wark,《Statistical mechanics》, OUP 2001。
\end{itemize}

%% 附录 F 部分习题答案与提示（Answers and hints to selected problems）
\chapter{部分习题答案与提示}

\exer{1.1} $9.27\times10^{-24}$ A\,m$^{2}$；8.40 GHz；34.7 $\mu$eV；8.40 GHz。这里我们只考虑
自旋角动量。

\exer{1.4} 这是标量算符 $(\hat{\mathbf{S}}\cdot\mathbf{X})$ 与矢量算符
$(\hat{\mathbf{S}})$ 的对易子。算出它的一个分量，即可求值
\begin{align*}
[\hat{\mathbf{S}}\cdot\mathbf{X}, \hat{\mathbf{S}}\cdot\mathbf{Y}] &= \frac{1}{4}\left((\boldsymbol{\sigma}\cdot\mathbf{X})(\boldsymbol{\sigma}\cdot\mathbf{Y}) - (\boldsymbol{\sigma}\cdot\mathbf{y}) (\boldsymbol{\sigma}\cdot\mathbf{X})\right)\\
&= \frac{\mathrm{i}}{4}\left(\boldsymbol{\sigma}\cdot(\mathbf{X}\times\mathbf{Y}) - \boldsymbol{\sigma}\cdot(\mathbf{Y}\times\mathbf{X})\right)\\
&= \frac{\mathrm{i}}{2}\boldsymbol{\sigma}\cdot(\mathbf{X}\times\mathbf{Y})\\
&= \mathrm{i}\mathbf{S}\times\mathbf{X}\cdot\mathbf{Y},\tag{F.1}
\end{align*}
结果得证。\footnote{首行 $(\boldsymbol{\sigma}\cdot\mathbf{y})$ 原书如此，应为
$(\boldsymbol{\sigma}\cdot\mathbf{Y})$——疑为原书排印之误。末行 $\mathbf{S}$ 上的
帽子原书即未加，亦照录。——译注}

\exer{1.5} 用式 1.59 可以证明
\begin{equation*}
\hat{S}_z(\hat{S}_{\pm}|S, S_z\rangle) = (\hat{S}_{\pm}\hat{S}_z \pm \hat{S}_{\pm})|S, S_z\rangle = (S_z\pm1)(\hat{S}_{\pm}|S, S_z\rangle),\tag{F.2}
\end{equation*}
即自旋的 z 分量被升高或降低。式 1.58 与 1.61 可以写作
\begin{equation*}
\hat{S}_{+}\hat{S}_{-} - \hat{S}_{-}\hat{S}_{+} = 2\hat{S}_z\tag{F.3}
\end{equation*}
\begin{equation*}
\hat{S}_{+}\hat{S}_{-} + \hat{S}_{-}\hat{S}_{+} = 2(\hat{S}_x^{2}+\hat{S}_y^{2}) = \hat{\mathbf{S}}^{2}-\hat{S}_z^{2},\tag{F.4}
\end{equation*}
相加或相减得
\begin{align*}
\hat{S}_{+}\hat{S}_{-} &= \hat{\mathbf{S}}^{2}-\hat{S}_z^{2}+\hat{S}_z\tag{F.5}\\
\hat{S}_{-}\hat{S}_{+} &= \hat{\mathbf{S}}^{2}-\hat{S}_z^{2}-\hat{S}_z.\tag{F.6}
\end{align*}
故
\begin{align*}
\langle S, S_z|\hat{S}_{-}\hat{S}_{+}|S, S_z\rangle &= S(S+1)-S_z^{2}-S_z\\
&= S(S+1)-S_z(S_z+1)\tag{F.7}
\end{align*}
及
\begin{align*}
\langle S, S_z|\hat{S}_{+}\hat{S}_{-}|S, S_z\rangle &= S(S+1)-S_z^{2}+S_z\\
&= S(S+1)-S_z(S_z-1)\tag{F.8}
\end{align*}
所需的归一化于是得证。

\exer{1.6} 不失一般性，取 $\mathbf{B}=(0,0,B)$，则 $\mathbf{B}\times\mathbf{r}=B(-y,-x,0)$，
$\frac{1}{2}\nabla\times(\mathbf{B}\times\mathbf{r})=(0,0,B)$。其他选取也可行，例如
$\mathbf{A}=-By(1,0,0)$ 或 $\mathbf{A}=Bx(0,1,0)$。事实上
$\mathbf{A}=B((\alpha-1)y,\alpha x,0)$ 都可行，$\alpha$ 为任意实数。

\exer{1.7} 提示：计算 $[\boldsymbol{\sigma}\cdot(\mathbf{p}+e\mathbf{A})]^{2}/2m_{\mathrm{e}}$。

\exer{1.8} 不失一般性，新轴可取在 $xz$ 平面内、与 $z$ 轴成 $\theta$ 角。算符为
\begin{equation*}
\hat{S}_{\theta} = \cos\theta\hat{S}_z + \sin\theta\hat{S}_x.\tag{F.9}
\end{equation*}
本征值为 1 的本征矢为
\begin{equation*}
(2(1-\cos\theta))^{-1/2}\begin{pmatrix}\sin\theta\\1-\cos\theta\end{pmatrix},\tag{F.10}
\end{equation*}
故所求概率为
\begin{equation*}
\frac{\sin^{2}\theta}{2(1-\cos\theta)} = \cos^{2}\frac{\theta}{2}.\tag{F.11}
\end{equation*}

\exer{1.10}
\begin{equation*}
H_{\mathrm{d}}/H_{\mathrm{a}}:\ \text{(a)}\ 3\times10^{-5}\quad \text{(b)}\ -8.8\times10^{-4}
\end{equation*}
\begin{equation*}
B_{\mathrm{d}}/B_{\mathrm{a}}:\ \text{(a)}\ -6\times10^{-5}\quad \text{(b)}\ 1.76\times10^{-3}
\end{equation*}

\exer{1.12} 提示：$\sigma_m^{2}=I$，故
\begin{equation*}
\sigma^{n} = \begin{cases} I & n \text{ 为偶}\\ \sigma_m & n \text{ 为奇}. \end{cases}\tag{F.12}
\end{equation*}
于是
\begin{align*}
\mathrm{e}^{\mathrm{i}\alpha\sigma_m} &= I\left(1-\frac{\alpha^{2}}{2!}+\frac{\alpha^{4}}{4!}-\cdots\right)\\
&\quad + \mathrm{i}\sigma_m\left(\alpha-\frac{\alpha^{3}}{3!}+\frac{\alpha^{5}}{5!}-\cdots\right)\\
&= I\cos\alpha + \mathrm{i}\sigma_m\sin\alpha.\tag{F.13}
\end{align*}

\exer{1.14} （a）力矩 $=I\ddot{\theta}=-pE\sin\theta$，故
$\ddot{\theta}+pE\sin\theta/I=0$；$\theta\ll1$ 时给出简谐运动，$\omega^{2}=pE/I$。
（b）哈密顿量为 $H=\hat{L}^{2}/2I-pE\cos\theta$。此计算中需求出若干对易子：
\begin{equation*}
[\hat{\theta}, \hat{L}] = \mathrm{i}\hbar\tag{F.14}
\end{equation*}
\begin{equation*}
[\hat{\theta}, \hat{L}^{2}] = 2\mathrm{i}\hbar\hat{L}\tag{F.15}
\end{equation*}
\begin{equation*}
[\hat{L}, \cos\theta] = \mathrm{i}\hbar\sin\theta,\tag{F.16}
\end{equation*}
结果随之而来。电偶极子不同于磁矩之处在于它不与角动量相联系（电偶极子只是空间
分开的两个相反电荷；磁偶极子则是相互绕行的两个相反电荷，因而与角动量相联系）。
注意哈密顿量不含耗散，故电偶极子永远振荡、磁偶极子永远进动。只有允许它们与热浴
交换能量（耗散过程）而改变自身能量时，两种偶极子才会沿外场排列。

\exer{2.1} $-1.65\times10^{-15}$；$7.8\times10^{-12}$。

\exer{2.2} 鸭子主要是水，质量 2--3 kg。一枚铁屑约 $0.1\,\mathrm{mm^{3}}=10^{-10}\,\mathrm{m^{3}}$
（估计），每个铁原子携带 $2.2\mu_{\mathrm{B}}$。铁的密度为 $7873\,\mathrm{kg\,m^{-3}}$、
相对原子质量 55.847 g。由此可以算出所需磁场为 1 mT（地磁场约 0.05 mT）。一头牛
约 400 kg，设它也主要由水组成，则质量（从而体积）约大 200 倍，产生同样磁化强度
所需磁场小得多（$\sim2\ \mu$T）。当然我们忽略了血红蛋白的效应。

\exer{2.3} $2.74\times10^{-5}\ \mathrm{J\,T^{-1}}$。

\exer{2.4} 提示：用 $Z=2\cosh(\mu_{\mathrm{B}}B/k_{\mathrm{B}}T)$ 与
$F=-nk_{\mathrm{B}}T\log Z$。则 $E=-n(\mathrm{d}\log Z/\mathrm{d}\beta)$，
$C=\partial E/\partial T$。熵可以用 $S=-(\partial F/\partial T)_B$ 或
$S=(E-F)/T$ 得到。

\exer{2.5} 提示：一个有用的中间结果是
\begin{equation*}
\frac{\partial\log Z_J}{\partial y} = \mathrm{B}_J(y).\tag{F.17}
\end{equation*}

\exer{2.7} （a）${}^{5}\mathrm{I}_8$（b）${}^{4}\mathrm{I}_{15/2}$
（c）${}^{3}\mathrm{H}_6$（d）${}^{1}\mathrm{S}_0$。

\exer{2.8} 提示：（1）你需要计算含 $\mathbf{p}\cdot\mathbf{A}+\mathbf{A}\cdot\mathbf{p}$
的项。它正比于
\begin{equation*}
\mathbf{p}\cdot\boldsymbol{\mu}\times\mathbf{r} + \boldsymbol{\mu}\times\mathbf{r}\cdot\mathbf{p} = 2\mathbf{r}\times\mathbf{p}\cdot\boldsymbol{\mu},\tag{F.18}
\end{equation*}
最后的等号成立是因为 $\boldsymbol{\mu}$ 不依赖位置、与 r 和 p 对易；而 r 与 p
可以互易是因为它们出现在标量三重积中——你从不把 r 与 p 的同一分量相乘。

（2）$\nabla^{2}(\frac{1}{r})=\nabla\cdot\nabla(\frac{1}{r})=4\pi\delta(\mathbf{r})$。
一个类似的结果是 $(\mathbf{S}\cdot\nabla)\nabla(\frac{1}{r})=\mathbf{S}\frac{4\pi}{3}\delta(\mathbf{r})$，
因子 3 是因为 $\mathbf{S}$ 从奇点的三个分量中挑出一个。

\exer{2.9} $1.89\times10^{-4}\ \mathrm{emu\,mol^{-1}}$；$9.68\times10^{-7}\ \mathrm{emu\,g^{-1}}$。

\exer{2.10} 固定 T 时，
\begin{equation*}
T\delta S = -\mathbf{B}\cdot\delta\mathbf{M} - \mathbf{M}\cdot\delta\mathbf{B} + \frac{\partial(nk_{\mathrm{B}}T\log Z)}{\partial\mathbf{B}}\cdot\delta\mathbf{B},\tag{F.19}
\end{equation*}
用 $M=nk_{\mathrm{B}}T(\partial\log Z/\partial B)$ 即得结果。

\exer{3.1} $\mathrm{Sc^{2+}}$ 有一个 3d 电子（$l=2$、$s=\frac{1}{2}$），有 5 个由
$l_z=-2,-1,0,1,2$ 刻画的轨道态。能级为 $0, A, 4A$（计自旋后简并度分别为 2、4、4）；
$A>0$ 时基态二重简并（$E=0$），$A<0$ 时四重简并（$E=4A$）。用记号 $\{l_z, s_z\}$
表示态。$A>0$ 时基态能级为 $\{0,\frac{1}{2}\}$ 与 $\{0,-\frac{1}{2}\}$，不被自旋--
轨道相互作用分裂（因 $\lambda l_zs_z=0$）。任意方向的磁场都会分裂这些能级，在
$\mu_{\mathrm{B}}B\ll k_{\mathrm{B}}T\ll A$ 时给出类居里磁化率。这一条件保证不布居
激发态（$k_{\mathrm{B}}T\ll A$）且留在布里渊函数的线性段（$\mu_{\mathrm{B}}B\ll k_{\mathrm{B}}T$）。

$A<0$ 时基态能级为 $|2,\frac{1}{2}\rangle$、$|2,-\frac{1}{2}\rangle$、
$|-2,\frac{1}{2}\rangle$ 与 $|-2,-\frac{1}{2}\rangle$；自旋--轨道相互作用使
$|2,-\frac{1}{2}\rangle$ 与 $|-2,\frac{1}{2}\rangle$ 能量降低 $\lambda$，
$|2,\frac{1}{2}\rangle$ 与 $|-2,-\frac{1}{2}\rangle$ 升高 $\lambda$。此时只有平行
于 z 的磁场才能分裂较低的能级，故 $\mu_{\mathrm{B}}B\ll k_{\mathrm{B}}T\ll\lambda$
时磁化率是类居里的。磁场垂直于 z 时磁化率不依赖温度。

\exer{3.2} 用
\begin{align*}
(1+x)^{-1/2} &= 1-\frac{x}{2}+\frac{(-\frac{1}{2})(-\frac{3}{2})}{2!}x^{2}+\frac{(-\frac{1}{2})(-\frac{3}{2})(-\frac{5}{2})}{3!}x^{3}\\
&\quad + \frac{(-\frac{1}{2})(-\frac{3}{2})(-\frac{5}{2})(-\frac{7}{2})}{4!}x^{4}+\cdots\\
&= 1-\frac{x}{2}+\frac{3x^{2}}{8}-\frac{5x^{3}}{16}+\frac{35x^{4}}{128}+\cdots.\tag{F.20}
\end{align*}

\exer{3.3} 把 V 化为球坐标：
\begin{equation*}
V = \frac{Dr^{4}}{8}\left[\sin^{4}\theta\left(\mathrm{e}^{4\mathrm{i}\phi}+\mathrm{e}^{-4\mathrm{i}\phi}+6\right) + 8\cos^{2}\theta - \frac{24}{5}\right].\tag{F.21}
\end{equation*}
对矩阵元，除 $V_{2,-2}$ 与 $V_{-2,2}$ 外所有非对角元为零。非零矩阵元的计算直截了当
但繁琐。举一例：
\begin{align*}
V_{1,1} &= \int_{0}^{\infty}r^{6}R^{2}(r)\,\mathrm{d}r\,\frac{D}{8}\int_{0}^{\pi}\left(6\sin^{2}\theta+8\cos^{4}\theta-\frac{24}{5}\right)\\
&\quad\times4\sin^{3}\theta\cos^{2}\theta\,\mathrm{d}\theta\int_{0}^{2\pi}\mathrm{d}\phi\\
&= -4A.\tag{F.22}
\end{align*}
本征值与本征矢为
\begin{center}
\begin{tabular}{ll}
本征值 & 本征矢\\
\midrule
$A+5A\left(1+\frac{4\mu_{\mathrm{B}}^{2}B^{2}}{25A^{2}}\right)^{1/2}$ & $\dfrac{(|2\rangle+x_{+}|-2\rangle)}{(1+x_{+}^{2})^{1/2}}$\\[10pt]
$-4A+\mu_{\mathrm{B}}B$ & $|1\rangle$\\
$6A$ & $|0\rangle$\\
$-4A-\mu_{\mathrm{B}}B$ & $|-1\rangle$\\
$A-5A\left(1+\frac{4\mu_{\mathrm{B}}^{2}B^{2}}{25A^{2}}\right)^{1/2}$ & $\dfrac{(|2\rangle+x_{-}|-2\rangle)}{(1+x_{-}^{2})^{1/2}}$\\
\end{tabular}
\end{center}
其中
\begin{equation*}
x_{\pm} = \pm\left(1+\left(\frac{2\mu_{\mathrm{B}}B}{5A}\right)^{2}\right)^{1/2} - \frac{2\mu_{\mathrm{B}}B}{5A}.\tag{F.23}
\end{equation*}
低场极限下本征值与本征矢为
\begin{center}
\begin{tabular}{ll}
本征值 & 本征矢\\
\midrule
$6A+\dfrac{2\mu_{\mathrm{B}}^{2}B^{2}}{5A}$ & $\dfrac{(|2\rangle+|-2\rangle)}{\sqrt{2}}$\\[10pt]
$-4A+\mu_{\mathrm{B}}B$ & $|1\rangle$\\
$6A$ & $|0\rangle$\\
$-4A-\mu_{\mathrm{B}}B$ & $|-1\rangle$\\
$-4A-\dfrac{2\mu_{\mathrm{B}}^{2}B^{2}}{5A}$ & $\dfrac{(|2\rangle-|-2\rangle)}{\sqrt{2}}$\\
\end{tabular}
\end{center}
故若 $k_{\mathrm{B}}T\ll A$，只布居最低能级：$A>0$ 时是三重态、$A<0$ 时是二重态。
只考虑这些最低能级，配分函数可写为
\begin{equation*}
Z = \begin{cases} 1+e^{-\mu_{\mathrm{B}}B/k_{\mathrm{B}}T}+e^{\mu_{\mathrm{B}}B/k_{\mathrm{B}}T} & A>0\\ 1+e^{-2\mu_{\mathrm{B}}^{2}B^{2}/5Ak_{\mathrm{B}}T} & A<0, \end{cases}\tag{F.24}
\end{equation*}
再用 $F=-Nk_{\mathrm{B}}T\log Z$ 与 $M=-\partial F/\partial B$ 即得结果。

\exer{3.4} $I=\frac{3}{2}$、$A=2\times10^{-6}$ eV。关于核磁矩说不了太多。声子引起
的谱线展宽在高温很大。分辨精细结构常需要低温。

\exer{3.5} 测得 $g_JJ=7$（$\mathrm{Gd^{3+}}$）、5（$\mathrm{Fe^{3+}}$）与 3
（$\mathrm{Cr^{3+}}$）。洪特规则的预言对三者都符合，唯独 $\mathrm{Cr^{3+}}$
例外——按洪特规则应预期 $J=\frac{3}{2}$、$g_J=\frac{2}{5}$。这是因为发生轨道
淬灭、$g_J=2$。（$\mathrm{Fe^{3+}}$ 也发生轨道淬灭，但察觉不到——洪特规则本来
就给出 $L=0$。）

\exer{3.6} 60 MHz 对应的磁场为
\begin{equation*}
\begin{array}{ll}
{}^{1}\mathrm{H} & 1.41\ \mathrm{T}\\
{}^{2}\mathrm{H} & 9.18\ \mathrm{T}\\
{}^{13}\mathrm{C} & 5.60\ \mathrm{T}\\
{}^{19}\mathrm{F} & 1.50\ \mathrm{T}
\end{array}
\end{equation*}

\exer{3.7} 0.322 T。

\exer{3.8} 先导出 $S=1$ 粒子的 $\hat{S}_z$ 与 $\hat{S}_x$ 矩阵表示：
\begin{equation*}
\hat{S}_z = \begin{pmatrix}1 & 0 & 0\\ 0 & 0 & 0\\ 0 & 0 & -1\end{pmatrix}\tag{F.25}
\end{equation*}
\begin{equation*}
\hat{S}_x = \frac{1}{\sqrt{2}}\begin{pmatrix}0 & 1 & 0\\ 1 & 0 & 1\\ 0 & 1 & 0\end{pmatrix}.\tag{F.26}
\end{equation*}
于是哈密顿量为
\begin{equation*}
\hat{\mathcal{H}} = \begin{pmatrix} D+g_{\parallel}\mu_{\mathrm{B}}B\cos\theta & \frac{g_{\perp}\mu_{\mathrm{B}}B\sin\theta}{\sqrt{2}} & 0\\ \frac{g_{\perp}\mu_{\mathrm{B}}B\sin\theta}{\sqrt{2}} & 0 & \frac{g_{\perp}\mu_{\mathrm{B}}B\sin\theta}{\sqrt{2}}\\ 0 & \frac{g_{\perp}\mu_{\mathrm{B}}B\sin\theta}{\sqrt{2}} & D-g_{\parallel}\mu_{\mathrm{B}}B\cos\theta \end{pmatrix}\tag{F.27}
\end{equation*}
求本征值即得所需结果。$\theta=0$ 时本征值简化为
\begin{equation*}
E = 0, D\pm g_{\parallel}\mu_{\mathrm{B}}B.\tag{F.28}
\end{equation*}
$\theta=\pi/2$ 时本征值简化为
\begin{equation*}
E = D, \frac{D}{2}\pm\sqrt{\left(\frac{D}{2}\right)^{2}+(g_{\perp}\mu_{\mathrm{B}}B)^{2}}.\tag{F.29}
\end{equation*}

\exer{3.9} $82.4\,\mathrm{m\,s^{-1}}$；$7.7\times10^{-19}\,\mathrm{m\,s^{-1}}$；
$9.4\times10^{14}\,\mathrm{Hz}$；$9.1\times10^{-9}\,\mathrm{Hz}$；
$2.3\times10^{7}\,\mathrm{Hz}$；$0.1\,\mu\mathrm{eV}$；$7.7\times10^{-4}\,\mathrm{cm^{-1}}$。

\exer{3.10} 若你能轻松测量的最低频率周期为 20 $\mu$s，即频率 50 kHz，对应约
$4\times10^{-4}$ T 的场。活过 20 $\mu$s 的 $\mu$ 子比例 $e^{-20/2.2}=1.1\times10^{-4}$，
即 $10^{7}$ 个注入 $\mu$ 子中约一千出头。

这一结果来自平均寿命 $\tau_{\mu}$ 的 $\mu$ 子指数衰减。活得比 T 久的比分为
\begin{equation*}
\frac{\int_T^{\infty}\mathrm{e}^{-t/\tau_{\mu}}\,\mathrm{d}t}{\int_0^{\infty}\mathrm{e}^{-t/\tau_{\mu}}\,\mathrm{d}t} = \mathrm{e}^{-T/\tau_{\mu}}.\tag{F.30}
\end{equation*}
若脉冲宽度为 50 ns，当自旋进动周期为 100 ns 时会发生相消干涉，对应频率 10 MHz，
给出场上限 $\sim$0.07 T。

若 $B=0.4$ T，$f=54.2$ MHz。

\exer{3.11} 对 $\mathrm{Fe^{3+}}$，$J=\frac{5}{2}$、$g_J=2$。饱和磁矩为
$g_JJ\mu_{\mathrm{B}}=5\mu_{\mathrm{B}}$。磁化率量的是
$\mu_{\mathrm{eff}}^{2}=g_J^{2}J(J+1)\mu_{\mathrm{B}}^{2}$，故
\begin{equation*}
g_J\sqrt{J(J+1)}\mu_{\mathrm{B}} = 2\times\sqrt{\frac{5}{2}\times\frac{7}{2}}\,\mu_{\mathrm{B}} = \sqrt{35}\,\mu_{\mathrm{B}} = 5.92\mu_{\mathrm{B}}.\tag{F.31}
\end{equation*}
饱和磁矩涉及 $\hat{J}_z$——它是沿特定方向的饱和磁矩测量；相比之下，磁化率涉及
$\hat{J}^{2}$。

\exer{4.1} 推导此结果有多种途径。其一从偶极子 1 的矢势出发：
\begin{equation*}
\mathbf{A}_1 = \frac{\mu_0}{4\pi}\nabla\times\frac{\boldsymbol{\mu}_1}{r},\tag{F.32}
\end{equation*}
它产生磁场
\begin{equation*}
\mathbf{B}_1 = \nabla\times\mathbf{A}_1 = \frac{\mu_0}{4\pi}\left[\nabla\left(\nabla\cdot\frac{\mu_1}{r}\right) - \nabla^{2}\left(\frac{\mu_1}{r}\right)\right],\tag{F.33}
\end{equation*}
给出能量
\begin{align*}
E &= -\boldsymbol{\mu}_2\cdot\mathbf{B}_1\\
&= -\frac{\mu_0}{4\pi}\left[(\boldsymbol{\mu}_1\cdot\nabla)(\boldsymbol{\mu}_1\cdot\nabla)\frac{1}{r} - (\boldsymbol{\mu}_1\cdot\boldsymbol{\mu}_2)\nabla^{2}\frac{1}{r}\right]\\
&= \frac{\mu_0}{4\pi r^{3}}\left[\boldsymbol{\mu}_1\cdot\boldsymbol{\mu}_2 - \frac{3(\boldsymbol{\mu}_1\cdot\mathbf{r})(\boldsymbol{\mu}_1\cdot\mathbf{r})}{r^{2}}\right].\tag{F.34}
\end{align*}
\footnote{原书如此。按推导，式 F.34 括号中第二、三项里的 $\boldsymbol{\mu}_1$
应为 $\boldsymbol{\mu}_2$（即 $3(\boldsymbol{\mu}_1\cdot\mathbf{r})(\boldsymbol{\mu}_2\cdot\mathbf{r})/r^{2}$）。——译注}
也可以用磁标势 $\phi_1=\mu_1\cdot r/4\pi r^{3}$、$B_1=-\mu_0\nabla\phi_1$ 推出
结果。

\exer{4.2} 1 \AA：（a）1.855 T（b）1.855 mT；10 \AA：（a）0.927 T（b）0.927 mT。

\exer{4.3} Fe（bcc）的晶格常数 $a=2.87$ \AA。最近邻距离为 $a/\sqrt{2}$，偶极--
偶极能量 $\sim(\mu_0\times(2.2\mu_{\text{B}})^{2}/4\pi r^{3})=30\ \mu$eV。
$\mathsf{J}\approx k_{\text{B}}T_{\text{C}}\sim0.09$ eV——大约大 3000 倍。

\exer{4.4} 带宽 $\sim0.05$ eV，库仑能 $\sim1$ eV。
$\mathsf{J}\sim-\frac{t^{2}}{U}=-2.5$ meV。故若 $|\mathsf{J}|=k_{\mathrm{B}}T_{\mathrm{N}}$，
则 $T_{\mathrm{N}}=29$ K。

\exer{4.5} 配分函数为
\begin{equation*}
Z = 1+\mathrm{e}^{-\beta\Delta}(\mathrm{e}^{\beta g\mu_{\mathrm{B}}B}+1+\mathrm{e}^{-\beta g\mu_{\mathrm{B}}B}).\tag{F.35}
\end{equation*}
由此可导出 F，进而 M：
\begin{equation*}
M = \frac{2n\beta(g\mu_{\mathrm{B}})^{2}B}{3+e^{\beta\Delta}},\tag{F.36}
\end{equation*}
$\chi\ll1$ 的结果随之而来。注意高温 $T\gg\Delta/k_{\mathrm{B}}$ 时
$\chi\to n\mu_0g^{2}\mu_{\mathrm{B}}^{2}/2k_{\mathrm{B}}T$——居里定律：所有能级
都被占据，磁化率是一个 $S=1$ 系统（$S(S+1)=2$）的 $\frac{3}{4}$ 与一个 $S=1$
系统（$S(S+1)=0$）\footnote{原书如此。由 $S(S+1)=0$ 可知应为 $S=0$——疑为原书
排印之误。——译注}的 $\frac{1}{4}$ 之和，即经典居里定律的 $\frac{3}{2}$。

低温 $T\ll\Delta/k_{\mathrm{B}}$ 时：$\Delta>0$ 则
$\chi\propto e^{-\Delta/k_{\mathrm{B}}T}\to0$（非磁——只有单态被占据）；$\Delta<0$
则 $\chi\to2n\mu_0g^{2}\mu_{\mathrm{B}}^{2}/3k_{\mathrm{B}}T$（对应被占据的三重态
$S(S+1)=2$，即经典居里定律的两倍）。

\exer{4.6} 提示：先证明
\begin{equation*}
\boldsymbol{\mu}_{I1}\cdot\boldsymbol{\mu}_{I2} \propto \hat{I}_{1z}\hat{I}_{2z} + \frac{1}{2}(\hat{I}_{1+}\hat{I}_{2-}+\hat{I}_{1-}\hat{I}_{2+})\tag{F.37}
\end{equation*}
并用 $\mathbf{r}=r(\sin\theta\cos\phi,\sin\theta\sin\phi,\cos\theta)$ 证明
\begin{equation*}
\boldsymbol{\mu}_I\cdot\hat{\mathbf{r}} \propto \cos\theta\hat{I}_z + \frac{1}{2}\sin\theta[e^{-i\phi}\hat{I}_{+}+e^{i\phi}\hat{I}_{-}].\tag{F.38}
\end{equation*}
结果随之很快得出。

\exer{4.8} 注意若 $\mathbf{S}=\mathbf{S}_1+\mathbf{S}_2+\mathbf{S}_3$，则
\begin{equation*}
\mathbf{S}^{2} = \mathbf{S}_1^{2}+\mathbf{S}_2^{2}+\mathbf{S}_3^{2} + 2(\mathbf{S}_1\cdot\mathbf{S}_2+\mathbf{S}_2\cdot\mathbf{S}_3+\mathbf{S}_3\cdot\mathbf{S}_1).\tag{F.39}
\end{equation*}
于是哈密顿量可以用 $\mathbf{S}^{2}$、$\mathbf{S}_1^{2}$、$\mathbf{S}_2^{2}$、
$\mathbf{S}_3^{2}$ 改写，给出 $E=\mathsf{J}(6-S(S+1))$。对 S 的各可能取值即得
所需能量。

\exer{5.1} $B_{\mathrm{mf}}=2.1\times10^{3}$ T。Fe 中单位体积原子数
$n=8.49\times10^{28}\,\mathrm{m^{-3}}$；用 $M=n\cdot2.2\mu_{\mathrm{B}}$ 得
$\mu_0M=2.2$ T。故 $B_{\mathrm{mf}}$ 比 $\mu_0M$ 大 $10^{3}$ 倍。

\exer{5.2} 磁化强度由
\begin{equation*}
\frac{M}{M_{\mathrm{s}}} = \mathrm{B}_S(y) = \alpha y - \beta y^{3}+\cdots,\tag{F.40}
\end{equation*}
给出，$\alpha$ 与 $\beta$ 是常数，而
\begin{equation*}
y = \frac{\gamma M}{TM_{\mathrm{s}}},\tag{F.41}
\end{equation*}
$\gamma$ 是由 $T_{\mathrm{C}}/\alpha$ 给出的正常数。$T_{\mathrm{C}}$ 附近这些方程
给出
\begin{equation*}
\left(\frac{M}{M_{\mathrm{s}}}\right)^{2} \approx \frac{\alpha^{3}}{\beta}\frac{(T_{\mathrm{C}}-T)}{T_{\mathrm{C}}}\tag{F.42}
\end{equation*}
用可由式 5.38 得到的 $\alpha$、$\beta$ 值即导出所需结果。

能量为 $E=-\frac{1}{2}\lambda M^{2}$；用式 F.42 可以证明在 $T_{\mathrm{C}}$ 处
\begin{equation*}
\frac{\mathrm{d}(M^{2})}{\mathrm{d}T} = -\frac{M_{\mathrm{s}}^{2}\alpha^{3}}{\beta T_{\mathrm{C}}},\tag{F.43}
\end{equation*}
由此可得比热的结果。

一般地容易证明
\begin{equation*}
C = \begin{cases} \frac{5}{2}nk_{\mathrm{B}}\left[\frac{(2S+1)^{2}-1}{(2S+1)^{2}+1}\right]\left(3\left(\frac{T}{T_{\mathrm{C}}}\right)^{2}-2\left(\frac{T}{T_{\mathrm{C}}}\right)\right) & T\leq T_{\mathrm{C}}\\ 0 & T>T_{\mathrm{C}}. \end{cases}\tag{F.44}
\end{equation*}

\exer{5.4} 能量为
\begin{equation*}
E = -2NS^{2}(\mathsf{J}_1\cos\theta + \mathsf{J}_2\cos2\theta),\tag{F.45}
\end{equation*}
能量极小 $\partial E/\partial\theta=0$ 给出 $\sin\theta=0$（从而 $\theta=0$ 或
$\pi$）或 $\cos\theta=-\mathsf{J}_1/4\mathsf{J}_2$。这三个解分别是铁磁（$\theta=0$）、
反铁磁（$\theta=\pi$）与螺旋磁（$\cos\theta=-\mathsf{J}_1/4\mathsf{J}_2$）；
代回能量方程即复现式 5.45。螺旋磁解只在 $|\mathsf{J}_1|<4|\mathsf{J}_2|$ 时可能。

哪个解能量最低可以通过比较各解的能量评估；或者通过计算 $\partial^{2}E/\partial\theta^{2}$
考察各解的稳定性。

\exer{5.5} $y\gg1$ 时
\begin{equation*}
\tanh y = \frac{e^{y}-e^{-y}}{e^{y}+e^{-y}} = \frac{1-e^{-2y}}{1+e^{-2y}} \approx 1-2e^{-2y}.\tag{F.46}
\end{equation*}
把它与
\begin{equation*}
\frac{M}{M_{\mathrm{s}}} = \tanh y,\tag{F.47}
\end{equation*}
以及 $T_{\mathrm{C}}$ 附近
\begin{equation*}
\frac{M}{M_{\mathrm{s}}} = \frac{T}{T_{\mathrm{C}}}y,\tag{F.48}
\end{equation*}
的事实联立，即可导出最终结果。

\exer{5.6} 回忆铁磁体的平均场理论：磁化强度由
\begin{equation*}
M = \frac{C}{\mu_0T}(\lambda M + \mu_0H),\tag{F.49}
\end{equation*}
给出，C 是居里常数。整理得
\begin{equation*}
M\left(1-\frac{\lambda C}{\mu_0T}\right) = \frac{HC}{T},\tag{F.50}
\end{equation*}
故磁化率为
\begin{equation*}
\chi = \frac{M}{H} = \frac{C}{T-\theta}\tag{F.51}
\end{equation*}
外斯温度 $\theta=\lambda C/\mu_0$。

对本题的铁磁体，每套子晶格的磁化强度为
\begin{equation*}
M_1 = \frac{C_1}{\mu_0T}(\mu_0H-\lambda M_2)\tag{F.52}
\end{equation*}
\begin{equation*}
M_2 = \frac{C_2}{\mu_0T}(\mu_0H-\lambda M_1)\tag{F.53}
\end{equation*}
即
\begin{equation*}
\begin{pmatrix} T & \lambda C_1/\mu_0\\ \lambda C_2/\mu_0 & T \end{pmatrix}\begin{pmatrix}M_1\\M_2\end{pmatrix} = H\begin{pmatrix}C_1\\C_2\end{pmatrix}.\tag{F.54}
\end{equation*}
此方程矩阵的行列式在 $T^{2}-\lambda^{2}C_1C_2/\mu_0=0$ 时为零，故可推断顺磁区为
$T>\theta$，$\theta=\lambda(C_1C_2)^{1/2}/\mu_0$。磁化率为
$\chi=(M_1+M_2)/H$；求逆式 F.54 得
\begin{equation*}
\begin{pmatrix}M_1\\M_2\end{pmatrix} = \frac{H}{T^{2}-\theta^{2}}\begin{pmatrix} T & -\lambda C_1/\mu_0\\ -\lambda C_2/\mu_0 & T \end{pmatrix}\begin{pmatrix}C_1\\C_2\end{pmatrix}\tag{F.55}
\end{equation*}
从而得到最终答案。

\exer{5.7} 第一部分直接由 $\hat{S}^{\pm}=\hat{S}^{x}\pm\mathrm{i}\hat{S}^{y}$ 的
应用得到。

对海森伯铁磁体，$\hat{S}_i^z\hat{S}_j^z$ 作用在基态上产生 $S^{2}$，基态是
$\hat{S}_i^z\hat{S}_j^z$ 的本征态。这样的项有 N 个，故 $E=-NS^{2}$。
$\hat{S}_i^{+}\hat{S}_j^{-}$ 与 $\hat{S}_i^{-}\hat{S}_j^{+}$ 不产生任何结果——
两种情形里升算符都湮灭基态。于是基态是二者的本征态、本征值为零。

对海森伯反铁磁体，$\hat{S}_i^{+}\hat{S}_j^{-}$ 与 $\hat{S}_i^{-}\hat{S}_j^{+}$
两项产生“想当然”基态之外的其他态，因此这一“想当然”的基态不是系统的本征态。

\exer{5.8} 可能的动量转移范围是 $0\,\mathrm{nm^{-1}}$ 到
$\sqrt{k_{\mathrm{in}}}=2\sqrt{2}\pi/\lambda=29.6\,\mathrm{nm^{-1}}$。

（a）100 K 时只看到化学（晶格）布拉格反射：
\begin{equation*}
Q = \left[\left(\frac{2\pi}{a}\right)^{2}(h^{2}+k^{2}) + \left(\frac{2\pi}{c}\right)^{2}l^{2}\right]^{1/2}.\tag{F.56}
\end{equation*}
对称性是 bcc：$h+k+l$ 为偶数时出现反射。

（b）10 K 时还看到磁布拉格反射。

两种情形观察到的反射列于下表。列出散射矢量 Q 的大小与入射束被散射的角 $2\theta$；
磁反射以星号标出，最后一列给出各反射的多重性 M。

\begin{table}
\centering
\begin{small}
\begin{tabular}{cccclcc}
\toprule
$h$ & $k$ & $l$ & Q（nm$^{-1}$） & $2\theta$（°） & & M\\
\midrule
0 & 1 & 0 & 12.566 & 34.915 & $*$ & 4\\
1 & 0 & 0 & 12.566 & 34.915 & $*$ & \\
1 & 1 & 0 & 17.772 & 50.208 & & 4\\
0 & 0 & 1 & 20.944 & 60.000 & $*$ & 2\\
0 & 1 & 1 & 24.425 & 71.337 & & 8\\
1 & 0 & 1 & 24.425 & 71.337 & & \\
0 & 2 & 0 & 25.133 & 73.740 & & 4\\
2 & 0 & 0 & 25.133 & 73.740 & & \\
1 & 1 & 1 & 27.468 & 81.952 & $*$ & 8\\
1 & 2 & 0 & 28.099 & 84.261 & $*$ & 8\\
2 & 1 & 0 & 28.099 & 84.261 & $*$ & \\
\bottomrule
\end{tabular}
\end{small}
\end{table}

\exer{5.9} （a）$\tilde{m}$ 的第一个表达式来自 $\tanh y$ 的定义。整理得
\begin{equation*}
x^{2} = \frac{1+\tilde{m}}{1-\tilde{m}},\tag{F.57}
\end{equation*}
取对数即式 5.57；用式 F.48 即得式 5.58。

（b）把式 5.60 整理成 x 的显式：
\begin{equation*}
(1-\tilde{m})x^{2} - \tilde{m}x - (1+\tilde{m}) = 0,\tag{F.58}
\end{equation*}
这是二次方程。$4-3\tilde{m}^{2}\geq0$ 时有实解，易得式 5.62。由
$0\leq\tilde{m}\leq1$ 知 $1\leq4-3\tilde{m}^{2}\leq2$，故 $x>0$ 需要式 5.62 中的
正根。于是
\begin{equation*}
y = \log(\tilde{m}+\sqrt{4-3\tilde{m}^{2}}) - \log(2(1-\tilde{m})),\tag{F.59}
\end{equation*}
且此情形 $y=3\tilde{m}T_{\mathrm{C}}/2T$，最终结果随之而来。

\exer{6.1} 平均能量为
\begin{equation*}
\langle E\rangle = -\frac{N\mathsf{J}}{4}\tanh\left(\frac{\mathsf{J}}{2k_{\mathrm{B}}T}\right),\tag{F.60}
\end{equation*}
每自旋热容为
\begin{equation*}
C = \frac{\mathsf{J}^{2}}{4k_{\mathrm{B}}T^{2}\cosh^{2}(\mathsf{J}/2k_{\mathrm{B}}T)}.\tag{F.61}
\end{equation*}
高温下它表现为 $\mathsf{J}^{2}/4k_{\mathrm{B}}T^{2}$，低温下如
$\mathsf{J}^{2}e^{-\mathsf{J}/k_{\mathrm{B}}T}/4k_{\mathrm{B}}T^{2}$。因此在
T=0 与 $T=\infty$ 处都为零、中间有一个宽极大——没有尖峰，这印证了纯一维模型没有
相变的论点。

\exer{6.2}
\begin{equation*}
\text{(b)}\quad \hbar\omega(\mathbf{q}) = S[\mathsf{J}(\mathbf{0})-\mathsf{J}(\mathbf{q})+K(\mathbf{0})].
\end{equation*}
\begin{equation*}
\text{(c)(i)}\quad \hbar\omega(\mathbf{q}) = 2S\mathsf{J}_0(1-\cos qa) + 2SK_0.
\end{equation*}
\begin{equation*}
\text{(ii)}\quad \hbar\omega(\mathbf{q}) = 2S\mathsf{J}_0(2-\cos q_xa-\cos q_ya) + 4SK_0.
\end{equation*}
\begin{equation*}
\text{(iii)}\quad \hbar\omega(\mathbf{q}) = 8S\mathsf{J}_0\left(1-\cos\frac{q_xa}{2}\cos\frac{q_ya}{2}\cos\frac{q_za}{2}\right) + 8SK_0.
\end{equation*}
\exer{6.3} 小 q 时用 $\cos qa\approx1-(qa)^{2}/2$ 得
（i）$\hbar\omega = S(\mathsf{J}_0a^{2}q^{2}+2K_0)$
（ii）$\hbar\omega = S(\mathsf{J}_0a^{2}q^{2}+4K_0)$
（iii）$\hbar\omega = S(\mathsf{J}_0a^{2}q^{2}+8K_0)$

伊辛情形（$\mathsf{J}_0=0$）：$\hbar\omega=2dSK_0\equiv\Delta$，d 对（i）（ii）
（iii）分别为 1、2、4。此模型中自旋波无论波矢为何都“花费”同样的能量 $\Delta$。
自旋波数目正比于 $e^{-\Delta/k_{\mathrm{B}}T}$，其能量正比于
$\Delta e^{-\Delta/k_{\mathrm{B}}T}$。比热容易算得正比于
\begin{equation*}
\frac{\Delta^{2}}{k_{\mathrm{B}}T^{2}}\,\mathrm{e}^{-\Delta/k_{\mathrm{B}}T},
\end{equation*}
与习题 6.1 一致（本题的 $\Delta$ 即该题的 $\mathsf{J}$）。

海森伯情形（$K_0=0$）：$\hbar\omega=Dq^{2}$，自旋刚度 $D=S\mathsf{J}_0a^{2}$。
自旋波数目 $N_{\mathrm{s}}$ 为
\begin{equation*}
N_{\mathrm{s}} \propto \int\frac{\mathrm{d}^{n}q}{\mathrm{e}^{Dq^{2}/k_{\mathrm{B}}T}-1},\tag{F.62}
\end{equation*}
小 q 时
\begin{equation*}
\frac{1}{\mathrm{e}^{Dq^{2}/k_{\mathrm{B}}T}-1} \approx \frac{k_{\mathrm{B}}T}{Dq^{2}}.\tag{F.63}
\end{equation*}
积分在一、二维发散（$\mathrm{d}^{n}q$ 分别为 $\mathrm{d}q$ 与 $2\pi q\,\mathrm{d}q$），
三维收敛（$\mathrm{d}^{n}q=4\pi q^{2}\mathrm{d}q$）。因此纯（各向同性）海森伯
模型在一、二维不可能有长程序。

\exer{6.4} 反铁磁情形可以循铁磁情形的推导，但须记住自己位于哪套子晶格。得到
\begin{equation*}
\begin{aligned}
\hbar\dot{S}_j^{+} &= 2\mathrm{i}\mathsf{J}S[2S_j^{+}+S_{j-1}^{+}+S_{j+1}^{+}] \quad j \text{ 为奇}\\
\hbar\dot{S}_j^{+} &= -2\mathrm{i}\mathsf{J}S[2S_j^{+}+S_{j-1}^{+}+S_{j+1}^{+}] \quad j \text{ 为偶}.
\end{aligned}\tag{F.64}
\end{equation*}
代入
\begin{equation*}
S_j^{+} = \begin{cases} u\,\mathrm{e}^{\mathrm{i}(qja-\omega t)} & j \text{ 为奇}\\ v\,\mathrm{e}^{\mathrm{i}(qja-\omega t)} & j \text{ 为偶}, \end{cases}\tag{F.65}
\end{equation*}
得
\begin{equation*}
\mathrm{i}\hbar\omega\begin{pmatrix}u\\v\end{pmatrix} = 4\mathrm{i}\mathsf{J}S\begin{pmatrix} 1 & \cos qa\\ -\cos qa & -1 \end{pmatrix}\begin{pmatrix}u\\v\end{pmatrix},\tag{F.66}
\end{equation*}
从而
\begin{equation*}
\begin{vmatrix} 1-\frac{\hbar\omega}{4\mathsf{J}S} & \cos qa\\ \cos qa & 1+\frac{\hbar\omega}{4\mathsf{J}S} \end{vmatrix} = 0\tag{F.67}
\end{equation*}
故 $\hbar\omega=4\mathsf{J}S\sin qa$。

\exer{6.5} 用 $\partial F/\partial M=0$，少许整理即得
\begin{equation*}
M^{2} = -\frac{2a_0(T-T_{\mathrm{C}})}{4b} + \frac{\mu_0H}{4bM},\tag{F.68}
\end{equation*}
正是所需形式。把 $M^{2}$ 对 H/M 作图得直线，$M^{2}$ 轴上的截距在 $T=T_{\mathrm{C}}$
处变号。

\exer{6.6} 合适单位下模型为
\begin{equation*}
\epsilon = \frac{1}{2}\sin^{2}(\theta-\phi) - h\cos\phi.\tag{F.69}
\end{equation*}
$\theta=0$ 时能量极小给出
\begin{equation*}
\sin\phi(\cos\phi+h) = 0\tag{F.70}
\end{equation*}
解为 $\phi=0, \pi, \cos^{-1}(-h)$。计算 $\partial^{2}E/\partial\theta^{2}$\footnote{原书
如此。按题意能量 $\epsilon$ 为 $\phi$ 的函数，此处应为 $\partial^{2}E/\partial\phi^{2}$
——疑为原书排印之误。——译注}可考察
这些解的稳定性。$\theta=\pi/2$ 时能量极小给出
\begin{equation*}
\sin\phi(h-\cos\phi) = 0\tag{F.71}
\end{equation*}
解为 $\phi=0, \pi, \cos^{-1}(h)$。

\exer{6.7} 磁化强度的分量为
\begin{equation*}
M_x = |\mathbf{M}|\sin\frac{\pi r}{R}\tag{F.72}
\end{equation*}
\begin{equation*}
M_y = |\mathbf{M}|\cos\frac{\pi r}{R},\tag{F.73}
\end{equation*}
用
\begin{equation*}
\frac{\mathrm{d}r}{\mathrm{d}x} = \frac{\mathrm{d}}{\mathrm{d}x}\sqrt{x^{2}+y^{2}+z^{2}} = \frac{x}{r}\tag{F.74}
\end{equation*}
之类的结果可以求出 $M_x$ 与 $M_y$ 的所有偏导数。于是
\begin{equation*}
(\nabla M_x)^{2}+(\nabla M_y)^{2} = |\mathbf{M}|^{2}\frac{\pi^{2}}{R^{2}},\tag{F.75}
\end{equation*}
最终结果由此而来。

\exer{6.8} 用合理的单位，25 Gbits\,in$^{-2}$ 即 $3.9\times10^{13}$ bits\,m$^{-2}$。

（a）取 $\pi a_0^{2}$ 为一个氢原子的面积，得 $3.4\times10^{-7}$ bits（氢原子
面积）$^{-1}$。

（b）莎士比亚的剧本多在两万到四万词之间；每词平均约 5 个字母，共约 40 部戏，
大约 4 Mbytes 即 32 Mbits（每字节 8 比特）。拿一张英国邮票算，我把每张邮票算出
约 600 份全集——用数据压缩还能翻倍。“我对这些数目字全不在行”（《哈姆雷特》
第二幕第二场）。

DNA 每碱基对存 1 比特，约每立方纳米 1 比特。若做成一张薄片，相当于约
$6\times10^{5}$ Gbits\,in$^{-2}$；但 DNA 的巨大优势在于它能卷曲，是体积存储技术。

\exer{6.9} 最好的做法是把磁场取在 $-z$ 方向——不失一般性。（为何 $-z$？这样
基态所有自旋平行于 z。电子带负电、自旋与其磁矩反平行，因而将与外场反平行排列。）
你会发现每个磁振子添加能量 $+g\mu_{\mathrm{B}}B$。

\exer{6.11} 方程
\begin{equation*}
\langle q|i\rangle = \frac{1}{\sqrt{N}}\sum_j e^{-i\mathbf{q}\cdot\mathbf{R}_j}\langle j|i\rangle\tag{F.76}
\end{equation*}
可以用 $\langle j|i\rangle=\delta_{ij}$ 简化，结果随之而来。

下一部分的提示：$S^{+}$ 会湮灭它所作用的态——除非它作用的自旋小于 S。

本题表明每个自旋都有一个垂直于磁化方向的小横向分量；这些横向分量随空间的变化
正如图 6.13 所画。

\exer{6.12} 需要求解
\begin{equation*}
\Delta E = 0 = \int_{-\infty}^{\infty}\left[2A\frac{\partial\theta}{\partial z}\frac{\partial}{\partial z}(\delta\theta) + \frac{\partial f(\theta)}{\partial\theta}\delta\theta\right]\mathrm{d}z\tag{F.77}
\end{equation*}
积分的第一部分分部积分后给出
\begin{equation*}
\frac{\partial f}{\partial\theta} - 2A\frac{\partial^{2}f}{\partial z^{2}} = 0.\tag{F.78}
\end{equation*}
题目本身细心地引导读者完成其余部分。

\exer{6.13} 由结构中壁的长度带来的壁能量密度约为 $\sigma_{\mathrm{W}}L/D$。
小三角形内部的自旋耗费各向异性能量。每个三角形面积 $D^{2}/4$，每个水平长度 D
中有两个，总代价为 $2\times K\times(D^{2}/4)\times(1/D)=KD/2$。把这些项相加并
求导即得所需结果。$D=20\ \mu$m。

\exer{7.1} 用 $g(E_{\mathrm{F}})=3n/2E_{\mathrm{F}}$，第一个结果很快得到。要考虑
温度依赖，值得回顾关于费米函数的一些数学细节。设
$h(E)=\int_0^E g(E')\mathrm{d}E'$（$g(E)=\mathrm{d}h/\mathrm{d}E$），考虑积分
\begin{align*}
I &= \int_{0}^{\infty}\frac{\mathrm{d}h}{\mathrm{d}E}f(E)\,\mathrm{d}E\\
&= [f(E)g(E)]_{0}^{\infty} - \int_{0}^{\infty}h(E)\frac{\mathrm{d}f}{\mathrm{d}E}\,\mathrm{d}E\\
&= -\int_{0}^{\infty}h(E)\frac{\mathrm{d}f}{\mathrm{d}E}\,\mathrm{d}E.\tag{F.79}
\end{align*}
令 $x=(E-\mu)/k_{\mathrm{B}}T$，则
\begin{equation*}
\frac{\mathrm{d}f}{\mathrm{d}E} = -\frac{1}{k_{\mathrm{B}}T}\frac{e^{x}}{(e^{x}+1)^{2}}.\tag{F.80}
\end{equation*}
写
\begin{equation*}
h(E) = \sum_{s=0}^{\infty}\frac{x^{s}}{s!}\left(\frac{\mathrm{d}^{s}h}{\mathrm{d}x^{s}}\right)_{x=0},\tag{F.81}
\end{equation*}
有
\begin{equation*}
I = \sum_{s=0}^{\infty}\frac{1}{s!}\left(\frac{\mathrm{d}^{s}h}{\mathrm{d}x^{s}}\right)_{x=0}\int_{-E_{\mathrm{F}}/k_{\mathrm{B}}T}^{\infty}\frac{x^{s}e^{x}\,\mathrm{d}x}{(e^{x}+1)^{2}}.\tag{F.82}
\end{equation*}
把下限换成 $-\infty$ 可简化其中的积分。奇 s 时积分为零；偶 s 时
\begin{align*}
\int_{-\infty}^{\infty}\frac{x^{s}\mathrm{e}^{x}\,\mathrm{d}x}{(\mathrm{e}^{x}+1)^{2}} &= 2\int_{0}^{\infty}\frac{x^{s}\mathrm{e}^{x}\,\mathrm{d}x}{(\mathrm{e}^{x}+1)^{2}}\\
&= 2\int_{0}^{\infty}\mathrm{d}x\sum_{n=0}^{\infty}\mathrm{e}^{x}x^{s}\left[(n+1)(-1)^{n+1}\mathrm{e}^{-nx}\right]\\
&= 2\sum_{n=1}^{\infty}(-1)^{n+1}n\int_{0}^{\infty}x^{s}\mathrm{e}^{-nx}\,\mathrm{d}x\\
&= 2(s!)\sum_{n=1}^{\infty}\frac{(-1)^{n+1}}{n^{s}}\\
&= 2(s!)(1-2^{1-s})\zeta(s),\tag{F.83}
\end{align*}
$\zeta(s)$ 是黎曼 zeta 函数（$\zeta(0)=-\frac{1}{2}$、$\zeta(2)=\frac{\pi^{2}}{6}$、
$\zeta(4)=\frac{\pi^{4}}{90}$）。

于是积分为
\begin{align*}
I &= \sum_{s=0,\,s\text{ 偶}}^{\infty}2\left(\frac{\mathrm{d}^{s}h}{\mathrm{d}x^{s}}\right)_{x=0}(1-2^{1-s})\zeta(s)\\
&= h + \frac{\pi^{2}}{6}\left(\frac{\mathrm{d}^{2}h}{\mathrm{d}x^{2}}\right)_{x=0} + \frac{7\pi^{4}}{360}\left(\frac{\mathrm{d}^{4}h}{\mathrm{d}x^{4}}\right)_{x=0} + \cdots\\
&= \int_{-\infty}^{\mu}g(E)\,\mathrm{d}E + \frac{\pi^{2}}{6}(k_{\mathrm{B}}T)^{2}\left(\frac{\mathrm{d}g}{\mathrm{d}E}\right)_{E=\mu}\\
&\quad + \frac{7\pi^{4}}{360}(k_{\mathrm{B}}T)^{4}\left(\frac{\mathrm{d}^{3}g}{\mathrm{d}E^{3}}\right)_{E=\mu} + \cdots.\tag{F.84}
\end{align*}
这意味着数密度为
\begin{equation*}
n = \int_{0}^{\mu}g(E)\,\mathrm{d}E + \frac{\pi^{2}}{6}(k_{\mathrm{B}}T)^{2}\left(\frac{\mathrm{d}g}{\mathrm{d}E}\right)_{E=\mu} + \cdots.\tag{F.85}
\end{equation*}
一阶近似下
\begin{equation*}
\int_{0}^{\mu}g(E)\,\mathrm{d}E = \int_{0}^{E_{\mathrm{F}}}g(E)\,\mathrm{d}E + (\mu-E_{\mathrm{F}})g(E_{\mathrm{F}}),\tag{F.86}
\end{equation*}
故对恒定的电子数，n 在一切温度取同值，从而
\begin{equation*}
\mu = E_{\mathrm{F}} - \frac{\pi^{2}}{6}(k_{\mathrm{B}}T)^{2}\frac{g'(E_{\mathrm{F}})}{g(E_{\mathrm{F}})}.\tag{F.87}
\end{equation*}
磁场的作用是把能级移动 $\pm\mu_{\mathrm{B}}B$。每个自旋态的数密度为
\begin{align*}
n_{\pm} &= \frac{1}{2}\int_{0}^{E_{\mathrm{F}}}g(E)\,\mathrm{d}E + (\mu-E_{\mathrm{F}}\mp\mu_{\mathrm{B}}B)g(E_{\mathrm{F}}\mp\mu_{\mathrm{B}}B)\\
&\quad + \frac{\pi^{2}}{6}(k_{\mathrm{B}}T)^{2}\left(\frac{\mathrm{d}g}{\mathrm{d}E}\right)_{E=\mu\mp\mu_{\mathrm{B}}B} + \cdots\tag{F.88}
\end{align*}
用近似关系
\begin{equation*}
g(E_{\mathrm{F}}\mp\mu_{\mathrm{B}}B) = g(E_{\mathrm{F}}) \mp \mu_{\mathrm{B}}B\left(\frac{\mathrm{d}g}{\mathrm{d}E}\right)_{E=\mu}\tag{F.89}
\end{equation*}
与
\begin{equation*}
\left(\frac{\mathrm{d}g}{\mathrm{d}E}\right)_{E=\mu\mp\mu_{\mathrm{B}}B} = \left(\frac{\mathrm{d}g}{\mathrm{d}E}\right)_{E=\mu} \mp \mu_{\mathrm{B}}B\left(\frac{\mathrm{d}^{2}g}{\mathrm{d}E^{2}}\right)_{E=\mu}\tag{F.90}
\end{equation*}
得
\begin{equation*}
n_{+}-n_{-} = -\frac{\pi^{2}}{6}(k_{\mathrm{B}}T)^{2}\mu_{\mathrm{B}}B\left[\left(\frac{g'}{g}\right)^{2}-\frac{g''}{g}\right] + \mu_{\mathrm{B}}Bg(E_{\mathrm{F}}).\tag{F.91}
\end{equation*}
故磁化率为
\begin{equation*}
\chi = \mu_0\mu_{\mathrm{B}}g(E_{\mathrm{F}})\left[1-\frac{\pi^{2}}{6}(k_{\mathrm{B}}T)^{2}\left[\left(\frac{g'}{g}\right)^{2}-\frac{g''}{g}\right]\right].\tag{F.92}
\end{equation*}
$g(E)\propto E^{1/2}$ 时易证
\begin{equation*}
\frac{g'}{g} = \frac{1}{2E_{\mathrm{F}}},\qquad \left(\frac{g'}{g}\right)^{2}-\frac{g''}{g} = \frac{1}{2E_{\mathrm{F}}^{2}}\tag{F.93}
\end{equation*}
（在费米能处取值）。最终结果随之而来。对 $E_{\mathrm{F}}\sim10$ eV 的金属
（如 Sn），$T_{\mathrm{F}}\sim10^{5}$ K，修正约为 $10^{-5}$。

\exer{7.2} （a）这在式 7.23 的推导中已经示明。

（b）交叉温度为 $2T_{\mathrm{F}}/3=27$ K。泡利顺磁磁化率为 $1.3\times10^{-8}$。
朗道抗磁磁化率为
\begin{equation*}
\chi_{\mathrm{L}} = -\left(\frac{m_{\mathrm{e}}}{m^{*}}\right)^{2}\frac{\chi_{\mathrm{P}}}{3} = -4.3\times10^{-7}.\tag{F.94}
\end{equation*}

\exer{7.3} 用 $A_{\phi}=\frac{1}{2}Br$，薛定谔方程成为
\begin{align*}
\Bigg[\frac{-\hbar^{2}}{2m_{\mathrm{e}}}\left(\frac{\partial^{2}}{\partial r^{2}}+\frac{1}{r}\frac{\partial}{\partial r}+\frac{1}{r^{2}}\frac{\partial^{2}}{\partial\phi^{2}}+\frac{\partial^{2}}{\partial z^{2}}\right) & - \frac{\mathrm{i}\hbar eB}{2m_{\mathrm{e}}}\frac{\partial}{\partial\phi}\\
& + \left.\frac{e^{2}B^{2}r^{2}}{8m_{\mathrm{e}}}\right]\psi(r,\phi,z) = E\psi(r,\phi,z).\tag{F.95}
\end{align*}
寻找如下形式的解
\begin{equation*}
\psi(r,\phi,z) = R(r)e^{\mathrm{i}l\phi}e^{\mathrm{i}k_z z}\tag{F.96}
\end{equation*}
导致
\begin{align*}
\Bigg(-\frac{\hbar^{2}}{2m_{\mathrm{e}}}\left(\frac{\partial^{2}}{\partial r^{2}}+\frac{1}{r}\frac{\partial}{\partial r}\right) - \frac{\hbar^{2}l^{2}}{2m_{\mathrm{e}}} + \frac{1}{2}m_{\mathrm{e}}\omega_{\mathrm{c}}^{2}r^{2} + \frac{\hbar k_z^{2}}{2m_{\mathrm{e}}}\Bigg)R(r) &= \left(E-\frac{l}{2}\hbar\omega_{\mathrm{c}}\right)R(r).\tag{F.97}
\end{align*}
解为
\begin{equation*}
\psi(r,\theta,z) = \exp\left[\mathrm{i}(l\phi+k_z z)-\frac{eBr^{2}}{4\hbar}\right]r^{|l|}L_{n-1}^{|l|}\left(\frac{eBr^{2}}{2\hbar}\right),\tag{F.98}
\end{equation*}
$L_{\alpha}^{\beta}(x)$ 是连带拉盖尔多项式（附录 C.5），能量为
\begin{equation*}
E = \left(n+\frac{l}{2}+\left|\frac{l}{2}\right|-\frac{1}{2}\right)\hbar\omega_{\mathrm{c}} + \frac{\hbar k_z^{2}}{2m_{\mathrm{e}}}.\tag{F.99}
\end{equation*}
在这一所谓对称规范中，所有态都绕空间中同一点旋转——以牺牲平移对称换取了旋转
对称的显现。

\exer{7.5} 提示：
\begin{equation*}
\frac{z+1}{z-1} = 1+\frac{2}{z-1}\tag{F.100}
\end{equation*}
由此可以证明
\begin{equation*}
1+\frac{1-z^{2}}{2z}\log\left(\frac{z+1}{z-1}\right) = \frac{2}{3z^{2}}+O\left(\frac{1}{z^{3}}\right).\tag{F.101}
\end{equation*}
另外 $\log(-1)=\mathrm{i}\pi$，故
\begin{equation*}
\log\left|\frac{x+1}{x-1}\right| - \log\left(\frac{x+1}{x-1}\right) = \mathrm{i}\pi.\tag{F.102}
\end{equation*}

\exer{7.6} （a）这在式 7.23 的推导中已经示明。（见习题 7.2(a)。）

（b）从下式出发很有用：
\begin{equation*}
\chi_{\mathbf{q}} = \mu_0\mu_{\mathrm{B}}^{2}\frac{2m_{\mathrm{e}}}{\hbar^{2}}I
\end{equation*}
I 是积分
\begin{equation*}
I = \int_{|\mathbf{k}|<k_{\mathrm{F}}}g(\mathbf{k})\,\mathrm{d}^{n}k\left[\frac{1}{(k+q)^{2}-k^{2}} + \frac{1}{(k-q)^{2}-k^{2}}\right].\tag{F.103}
\end{equation*}
一维中
\begin{align*}
\int_{|\mathbf{k}|<k_{\mathrm{F}}}\frac{\mathrm{d}k}{(k+q)^{2}-k^{2}} &= \int_{-k_{\mathrm{F}}}^{k_{\mathrm{F}}}\frac{\mathrm{d}k}{q(q+2k)}\\
&= \frac{1}{2q}\int_{q-2k_{\mathrm{F}}}^{q+2k_{\mathrm{F}}}\frac{\mathrm{d}u}{u}\\
&= \frac{1}{2q}\int_{|q-2k_{\mathrm{F}}|}^{q+2k_{\mathrm{F}}}\frac{\mathrm{d}u}{u}\\
&= \frac{1}{2q}\log\left|\frac{q+2k_{\mathrm{F}}}{q-2k_{\mathrm{F}}}\right|.\tag{F.104}
\end{align*}
结果随之而来。

二维中要考虑的积分是
\begin{align*}
I &= \int_{0}^{k_{\mathrm{F}}}k\,\mathrm{d}k\int_{0}^{2\pi}\frac{\mathrm{d}\theta}{q(q-2k\cos\theta)} + \frac{\mathrm{d}\theta}{q(q+2k\cos\theta)}\\
&= 2\int_{0}^{k_{\mathrm{F}}}k\,\mathrm{d}k\int_{0}^{2\pi}\frac{\mathrm{d}\theta}{q^{2}-4k^{2}\cos\theta}\tag{F.105}
\end{align*}
而
\begin{equation*}
\int_{0}^{2\pi}\frac{\mathrm{d}\theta}{q^{2}-4k^{2}\cos\theta} = \begin{cases} \frac{2\pi}{q(q^{2}-4k^{2})^{1/2}} & q>2k_{\mathrm{F}}\\ 0 & q<2k_{\mathrm{F}} \end{cases}\tag{F.106}
\end{equation*}
于是积分成为
\begin{equation*}
I = \begin{cases} \pi\left[1-\sqrt{1-\left(\frac{2k_{\mathrm{F}}}{q}\right)}\right] & q>2k_{\mathrm{F}}\\ \pi & q<2k_{\mathrm{F}} \end{cases}\tag{F.107}
\end{equation*}
最终结果随之而来。

%% 附录 G 符号、常数与有用公式（Symbols, constants and useful equations）
\chapter{符号、常数与有用公式}

\section{符号表}


\begin{center}
\begin{small}
\begin{tabular}{ll@{\hspace{3em}}ll}
\toprule
$a$ & 单胞尺寸 & $\mathbf{H}$ & 磁场强度\\
$a_0$ & 玻尔半径 & $H_{\mathrm{a}}$ & 外加场\\
$A$ & 面积 & $\mathbf{H}_{\mathrm{d}}$ & 退磁场\\
$A$ & 超精细耦合常数 & $H_{\mathrm{i}}$ & 内场\\
$A$ & 连续介质交换常数 & $\mathcal{H}$ & 哈密顿量\\
$A$ & 质量数 & $\mathrm{i}$ & $\sqrt{-1}$\\
$\mathbf{A}$ & 磁矢势 & $I$ & 电流\\
$b$ & 散射长度 & $I$ & 核自旋量子数\\
$\mathbf{B}$ & 磁通密度（磁场） & $J$ & 总角动量量子数\\
$B_{\mathrm{a}}$ & 外加场 & $\mathsf{J}$ & 交换常数/交换积分\\
$\mathbf{B}_{\mathrm{d}}$ & 退磁场 & $\mathbf{k}$ & 波矢\\
$B_{\mathrm{i}}$ & 内场 & $k_{\mathrm{B}}$ & 玻尔兹曼常数（$1.3807\times10^{-23}$ J\,K$^{-1}$）\\
$\mathbf{B}_{\mathrm{mf}}$ & 分子场 & $k_{\mathrm{F}}$ & 费米波矢\\
$\mathrm{B}_J$ & 布里渊函数 & $K$ & 库仑积分\\
$c$ & 自由空间光速（$2.9979\times10^{8}$ m\,s$^{-1}$） & $K$ & 各向异性常数\\
$\mathbf{D}$ & 电位移 & $K$ & 奈特移位\\
$D$ & 轴各向异性常数 & $K$ & 弹簧常数\\
$\mathrm{d}\mathbf{S}$ & 面元 & $K_{\mathrm{s}}$ & 表面各向异性常数\\
$\mathrm{d}S$ & 熵变 & $K_{\mathrm{v}}$ & 体积各向异性常数\\
$\mathrm{e}$ & $\exp(1)=2.718281828$ & $l$ & 轨道角动量量子数\\
$e$ & 电子电荷大小（$1.6022\times10^{-19}$ C） & $l$ & 朗道能级指标\\
$E$ & 能量 & $m_{\mathrm{e}}$ & 电子质量（$9.109\times10^{-31}$ kg）\\
$E_{\mathrm{F}}$ & 费米能 & $m_l$ & 轨道磁量子数\\
$E_{\mathrm{S}}$ & 单态能量 & $m_I$ & 核磁量子数\\
$E_{\mathrm{T}}$ & 三重态能量 & $m_J$ & 磁量子数\\
$\boldsymbol{\mathcal{E}}$ & 电场 & $m_{\mathrm{n}}$ & 中子质量\\
$F$ & 亥姆霍兹自由能（亥姆霍兹函数） & $m_s$ & 自旋磁量子数\\
$g(E)$ & 能量空间的态密度 & $\mathbf{M}$ & 磁化强度（单位体积磁矩）\\
$g(k)$ & 波矢空间的态密度 & $M_{\pm}$ & 反铁磁体子晶格上的磁化强度\\
$g$ & g 因子 & $M_{\mathrm{s}}$ & 饱和磁化强度\\
$g_{\mathrm{e}}$ & 电子 g 因子 & $n$ & 主量子数\\
$g_I$ & 核 g 因子 & $n$ & 单位体积原子/磁矩数\\
$g_L$ & 轨道角动量 g 因子 & $\mathsf{N}$ & 退磁张量\\
$g_S$ & 自旋角动量 g 因子 & $N$ & 退磁因子\\
$\mathbf{g}$ & 有效 g 张量 & $N$ & 原子核中的中子数\\
$\mathbf{G}$ & 力矩 & $N_{\mathrm{A}}$ & 阿伏伽德罗常数\\
$\mathbf{G}$ & 倒格矢 & $\mathbf{p}$ & 偶极矩\\
$h$ & 普朗克常数（$6.626\times10^{-34}$ J\,s） & $\mathbf{p}$ & 正则动量\\
$\hbar$ & 普朗克常数/$2\pi$（$1.0546\times10^{-34}$ J\,s） & $\hat{\mathbf{p}}$ & 动量算符\\
\bottomrule
\end{tabular}
\end{small}
\end{center}


\begin{center}
\begin{small}
\begin{tabular}{ll@{\hspace{3em}}ll}
\toprule
$\mathbf{q}$ & 磁振子波矢 & $\mu$ & 磁矩\\
$\mathbf{q}$ & 螺旋波矢 & $\mu_{\mathrm{r}}$ & 相对磁导率\\
$\mathbf{Q}$ & 散射矢量 & $\mu_{\mathrm{B}}$ & 玻尔磁子\\
$\mathbf{r}$ & 位置矢量 & $\mu_{\mathrm{N}}$ & 核磁子\\
$\hat{\mathbf{r}}$ & 位置算符 & $\mu_{\mathrm{eff}}$ & 有效磁矩（每化学式单元的 $\mu_{\mathrm{B}}$ 数）\\
$R$ & 电阻 & $\mu_0$ & 自由空间磁导率（$4\pi\times10^{-7}$ H\,m$^{-1}$）\\
$R_{\mathrm{e}}$ & 反常霍尔系数 & $\nu$ & 频率\\
$R_{\mathrm{o}}$ & 正常霍尔系数 & $\xi$ & 关联长度\\
$S$ & 熵 & $\rho$ & 电阻率\\
$s, S$ & 自旋量子数 & $\rho_{\mathrm{H}}$ & 霍尔电阻率\\
$\hat{\mathbf{S}}$ & 自旋角动量算符 & $\boldsymbol{\sigma}$ & 电导率\\
$\hat{S}_{+}$ & 升算符 & $\hat{\boldsymbol{\sigma}}$ & 泡利自旋矩阵\\
$\hat{S}_{-}$ & 降算符 & $\tau$ & 周期\\
$\tilde{S}$ & 有效自旋 & $\tau$ & 超顺磁弛豫时间\\
$t$ & 时间 & $\phi$ & 角度\\
$t$ & 跳跃积分 & $\Phi$ & 磁通量\\
$t_{\mathrm{p}}$ & 脉冲长度 & $\chi$ & 磁化率\\
$T$ & 温度 & $\chi_g$ & 质量磁化率\\
$T_{\mathrm{c}}$ & 临界温度 & $\chi_{\mathrm{L}}$ & 朗道抗磁磁化率\\
$T_{\mathrm{C}}$ & 居里温度 & $\chi_{\mathrm{m}}$ & 摩尔磁化率\\
$T_{\mathrm{N}}$ & 奈尔温度 & $\chi_{\mathrm{P}}$ & 泡利自旋磁化率\\
$T_{\mathrm{P}}$ & Peierls 转变温度 & $\chi_{\mathbf{q}}$ & 波矢依赖磁化率\\
$T_{\mathrm{SP}}$ & 自旋 Peierls 转变温度 & $\chi$ & 自旋波函数\\
$T_1$ & 自旋--晶格弛豫时间 & $\psi$ & 波函数\\
$T_2$ & 自旋--自旋弛豫时间 & $\Psi_{\mathrm{S}}$ & 单态波函数\\
$T_{\mathrm{F}}$ & 费米温度 & $\Psi_{\mathrm{T}}$ & 三重态波函数\\
$U$ & 库仑能 & $\omega$ & 角频率\\
$\mathbf{v}$ & 速度 & $\omega_{\mathrm{c}}$ & 回旋频率\\
$v$ & 速率 & $\omega_{\mathrm{L}}$ & 拉莫尔进动频率\\
$V$ & 势或势能 & $\epsilon_0$ & 自由空间电容率（$8.854\times10^{-12}$ F\,m$^{-1}$）\\
$W$ & 跃迁速率 & $\epsilon_{\mathrm{r}}$ & 相对介电常数\\
$Y_{lm_l}(\theta,\phi)$ & 球谐函数 & $\theta$ & 外斯温度\\
$Z$ & 配分函数 & $\lambda$ & 波长\\
$Z$ & 原子序数 & $\lambda$ & 自旋--轨道常数\\
$\alpha$ & 精细结构常数 & $\lambda_{\text{el-ph}}$ & 电子--声子耦合常数\\
$\beta$ & $=1/k_{\mathrm{B}}T$ & $\lambda_{\text{spin-ph}}$ & 自旋--声子耦合常数\\
$\gamma$ & 旋磁比 & & \\
$\gamma_{\mathrm{e}}$ & 电子旋磁比 & & \\
$\gamma_{\mu}$ & $\mu$ 子旋磁比 & & \\
$\gamma_{\mathrm{N}}$ & 核旋磁比 & & \\
$\Delta$ & 交换劈裂 & & \\
\bottomrule
\end{tabular}
\end{small}
\end{center}

\section{基本常数}

\begin{center}
\smallskip
{基本常数。\\
{\footnotesize\itshape Fundamental constants.}}
\begin{small}
\begin{tabular}{lcl}
\toprule
玻尔半径 & & $a_0 = 5.292\times10^{-11}$ m\\
自由空间光速 & & $c = 2.9979\times10^{8}$ m\,s$^{-1}$\\
电子电荷 & & $e = 1.6022\times10^{-19}$ C\\
普朗克常数 & & $h = 6.626\times10^{-34}$ J\,s\\
 & $h/2\pi=$ & $\hbar = 1.0546\times10^{-34}$ J\,s\\
玻尔兹曼常数 & & $k_{\mathrm{B}} = 1.3807\times10^{-23}$ J\,K$^{-1}$\\
电子静止质量 & & $m_{\mathrm{e}} = 9.109\times10^{-31}$ kg\\
质子静止质量 & & $m_{\mathrm{p}} = 1.6726\times10^{-27}$ kg\\
阿伏伽德罗常数 & & $N_{\mathrm{A}} = 6.022\times10^{23}$ mol$^{-1}$\\
标准摩尔体积 & & $22.414\times10^{-3}$ m$^{3}$\,mol$^{-1}$\\
摩尔气体常数 & & $R = 8.315$ J\,mol$^{-1}$\,K$^{-1}$\\
精细结构常数 & $\dfrac{e^{2}}{4\pi\varepsilon_0\hbar c}=$ & $\alpha = (137.04)^{-1}$\\[8pt]
自由空间电容率 & & $\epsilon_0 = 8.854\times10^{-12}$ F\,m$^{-1}$\\
自由空间磁导率 & & $\mu_0 = 4\pi\times10^{-7}$ H\,m$^{-1}$\\
玻尔磁子 & & $\mu_{\mathrm{B}} = 9.274\times10^{-24}$ A\,m$^{2}$ 或 J\,T$^{-1}$\\
核磁子 & & $\mu_{\mathrm{N}} = 5.051\times10^{-27}$ A\,m$^{2}$ 或 J\,T$^{-1}$\\
中子磁矩 & & $\mu_{\mathrm{n}} = -1.9130\mu_{\mathrm{N}}$\\
质子磁矩 & & $\mu_{\mathrm{p}} = 2.7928\mu_{\mathrm{N}}$\\
\bottomrule
\end{tabular}
\end{small}
\end{center}

\section{有用公式}

（1）恒等式：
\begin{gather*}
\mathrm{e}^{\mathrm{i}\theta} = \cos\theta + \mathrm{i}\sin\theta\\
\sin\theta = \frac{\mathrm{e}^{\mathrm{i}\theta}-\mathrm{e}^{-\mathrm{i}\theta}}{2\mathrm{i}},\qquad
\cos\theta = \frac{\mathrm{e}^{\mathrm{i}\theta}+\mathrm{e}^{-\mathrm{i}\theta}}{2}\\
\sin(\theta+\phi) = \sin\theta\cos\phi+\cos\theta\sin\phi,\qquad
\cos(\theta+\phi) = \cos\theta\cos\phi-\sin\theta\sin\phi\\
\cos^{2}\theta+\sin^{2}\theta = 1\\
\cos2\theta = \cos^{2}\theta-\sin^{2}\theta,\qquad
\sin2\theta = 2\cos\theta\sin\theta\\
\sinh x = \frac{\mathrm{e}^{x}-\mathrm{e}^{-x}}{2},\qquad
\cosh x = \frac{\mathrm{e}^{x}+\mathrm{e}^{-x}}{2}
\end{gather*}

（2）级数展开（$|x|<1$ 时有效）：
\begin{gather*}
(1+x)^{n} = 1+nx+\frac{n(n-1)}{2!}x^{2}+\frac{n(n-1)(n-2)}{3!}x^{3}+\cdots\\
\mathrm{e}^{x} = 1+x+\frac{x^{2}}{2!}+\frac{x^{3}}{3!}+\frac{x^{4}}{4!}+\cdots\\
\sin x = x-\frac{x^{3}}{3!}+\frac{x^{5}}{5!}-\cdots,\qquad
\cos x = 1-\frac{x^{2}}{2!}+\frac{x^{4}}{4!}-\cdots\\
\log(1+x) = x-\frac{x^{2}}{2}+\frac{x^{3}}{3}-\cdots
\end{gather*}

（3）积分：
\begin{gather*}
\int_{0}^{\infty}x^{n}\mathrm{e}^{-x}\,\mathrm{d}x = n!\\
\int_{-\infty}^{\infty}\mathrm{e}^{-\alpha x^{2}}\,\mathrm{d}x = \sqrt{\frac{\pi}{\alpha}}\\
\int_{-\infty}^{\infty}x^{2n}\mathrm{e}^{-\alpha^{2}x^{2}}\,\mathrm{d}x = \frac{(2n)!}{n!2^{2n}\alpha^{n}}\sqrt{\frac{\pi}{\alpha}}
\end{gather*}
不定积分（$a>0$）：
\begin{equation*}
\begin{aligned}
\int\frac{\mathrm{d}x}{x^{2}+a^{2}} &= \frac{1}{a}\tan^{-1}\frac{x}{a}\\
\int\frac{\mathrm{d}x}{x^{2}-a^{2}} &= \frac{1}{2a}\ln\left|\frac{x-a}{x+a}\right|\\
\int\frac{\mathrm{d}x}{\sqrt{x^{2}+a^{2}}} &= \sinh^{-1}\frac{x}{a}\\
\int\frac{\mathrm{d}x}{\sqrt{x^{2}-a^{2}}} &= \begin{cases} \cosh^{-1}\frac{x}{a} & x>a\\ -\cosh^{-1}\frac{x}{a} & x<-a \end{cases}\\
\int\frac{\mathrm{d}x}{\sqrt{a^{2}-x^{2}}} &= \sin^{-1}\frac{x}{a}
\end{aligned}\tag{G.1}
\end{equation*}

（4）矢量算符：
\begin{itemize}[itemsep=0.2em]
\item grad（梯度）作用于标量场产生矢量场：
$\operatorname{grad}\psi=\nabla\psi=\left(\frac{\partial\psi}{\partial x},\frac{\partial\psi}{\partial y},\frac{\partial\psi}{\partial z}\right)$
\item div（散度）作用于矢量场产生标量场：
$\operatorname{div}\mathbf{A}=\nabla\cdot\mathbf{A}=\frac{\partial A_x}{\partial x}+\frac{\partial A_y}{\partial y}+\frac{\partial A_z}{\partial z}$
\item curl（旋度）作用于矢量场产生另一个矢量场：
$\operatorname{curl}\mathbf{A}=\nabla\times\mathbf{A}=\begin{vmatrix}\mathbf{i} & \mathbf{j} & \mathbf{k}\\ \partial/\partial x & \partial/\partial y & \partial/\partial z\\ A_x & A_y & A_z\end{vmatrix}$
\end{itemize}
其中 $\psi(\mathbf{r})$ 与 $\mathbf{A}(\mathbf{r})$ 分别是任意给定的标量场与矢量场。

（5）矢量恒等式：
\begin{align*}
&\nabla\cdot(\nabla\psi) = \nabla^{2}\psi\\
&\nabla\times(\nabla\psi) = 0\\
&\nabla\cdot(\nabla\times\mathbf{A}) = 0\\
&\nabla\cdot(\psi\mathbf{A}) = \mathbf{A}\cdot\nabla\psi+\psi\nabla\cdot\mathbf{A}\\
&\nabla\times(\psi\mathbf{A}) = \psi\nabla\times\mathbf{A}-\mathbf{A}\times\nabla\psi\\
&\nabla\times(\nabla\times\mathbf{A}) = \nabla(\nabla\cdot\mathbf{A})-\nabla^{2}\mathbf{A}\\
&\nabla\cdot(\mathbf{A}\times\mathbf{B}) = \mathbf{B}\cdot\nabla\times\mathbf{A}-\mathbf{A}\cdot\nabla\times\mathbf{B}\\
&\nabla(\mathbf{A}\cdot\mathbf{B}) = (\mathbf{A}\cdot\nabla)\mathbf{B}+(\mathbf{B}\cdot\nabla)\mathbf{A}+\mathbf{A}\times(\nabla\times\mathbf{B})+\mathbf{B}\times(\nabla\times\mathbf{A})\\
&\nabla\times(\mathbf{A}\times\mathbf{B}) = (\mathbf{B}\cdot\nabla)\mathbf{A}-(\mathbf{A}\cdot\nabla)\mathbf{B}+\mathbf{A}(\nabla\cdot\mathbf{B})-\mathbf{B}(\nabla\cdot\mathbf{A})
\end{align*}
这些恒等式用交错张量与求和约定即可轻松证明。交错张量 $\epsilon_{ijk}$ 定义为
\begin{equation*}
\epsilon_{ijk} = \begin{cases} 1 & \text{若 } ijk \text{ 是 } 123 \text{ 的偶置换}\\ -1 & \text{若 } ijk \text{ 是 } 123 \text{ 的奇置换}\\ 0 & \text{若 } i, j, k \text{ 中任何两个相等} \end{cases}
\end{equation*}
于是矢量积可以写作
\begin{equation*}
(\mathbf{A}\times\mathbf{B})_i = \epsilon_{ijk}A_jB_k.
\end{equation*}
这里使用求和约定：重复两次的指标默认求和。标量积于是为
\begin{equation*}
\mathbf{A}\cdot\mathbf{B} = A_iB_i.
\end{equation*}
可以利用恒等式
\begin{equation*}
\epsilon_{ijk}\epsilon_{ilm} = \delta_{jl}\delta_{km}-\delta_{jm}\delta_{kl}
\end{equation*}
其中 $\delta_{ij}$ 是克罗内克 delta：
\begin{equation*}
\delta_{ij} = \begin{cases} 1 & i=j\\ 0 & i\neq j \end{cases}
\end{equation*}
矢量三重积为
\begin{equation*}
\mathbf{A}\times(\mathbf{B}\times\mathbf{C}) = (\mathbf{A}\cdot\mathbf{C})\mathbf{B}-(\mathbf{A}\cdot\mathbf{B})\mathbf{C}.
\end{equation*}

（6）柱坐标：
\begin{gather*}
\nabla^{2}\psi = \frac{1}{r}\frac{\partial}{\partial r}\left(r\frac{\partial\psi}{\partial r}\right) + \frac{1}{r^{2}}\frac{\partial^{2}\psi}{\partial\phi^{2}} + \frac{\partial^{2}\psi}{\partial z^{2}}\\
\nabla\psi = \left(\frac{\partial\psi}{\partial r}, \frac{1}{r}\frac{\partial\psi}{\partial\phi}, \frac{\partial\psi}{\partial z}\right)
\end{gather*}

（7）球极坐标：
\begin{align*}
\nabla^{2}\psi &= \frac{1}{r^{2}}\frac{\partial}{\partial r}\left(r^{2}\frac{\partial\psi}{\partial r}\right) + \frac{1}{r^{2}\sin\theta}\frac{\partial}{\partial\theta}\left(\sin\theta\frac{\partial\psi}{\partial\theta}\right) + \frac{1}{r^{2}\sin^{2}\theta}\frac{\partial^{2}\psi}{\partial\phi^{2}}\\
\nabla\psi &= \left(\frac{\partial\psi}{\partial r}, \frac{1}{r}\frac{\partial\psi}{\partial\theta}, \frac{1}{r\sin\theta}\frac{\partial\psi}{\partial\phi}\right)
\end{align*}

\vspace{2em}
\begin{center}
\small\kaishu ——全书完——
\end{center}


\end{document}
